% !TEX program = xelatex
\documentclass[UTF8,a4paper,11pt,fontset=windows]{ctexart}
\usepackage[margin=2.2cm]{geometry}
\usepackage{amsmath,amssymb,amsthm,mathrsfs,enumitem,booktabs,longtable,array}
\usepackage[unicode,hidelinks]{hyperref}
\usepackage{xurl,fancyhdr,microtype}
\setlength{\headheight}{15pt}
\pagestyle{fancy}\fancyhf{}
\fancyhead[L]{概率论期中复习}\fancyhead[R]{原笔记例题与历次作业}\fancyfoot[C]{\thepage}
\setlength{\parindent}{2em}\setlength{\parskip}{0.35em}
\setlength{\emergencystretch}{2em}
\setlist[enumerate]{itemsep=0.35em,topsep=0.4em}
\setcounter{tocdepth}{2}
\renewcommand{\proofname}{证明}
\newcommand{\R}{\mathbb R}\newcommand{\N}{\mathbb N}\newcommand{\Z}{\mathbb Z}
\newcommand{\C}{\mathbb C}\newcommand{\Q}{\mathbb Q}\newcommand{\F}{\mathbb F}
\renewcommand{\P}{P}\newcommand{\E}{\mathbb E}
\newcommand{\Var}{\operatorname{Var}}\newcommand{\Cov}{\operatorname{Cov}}
\newcommand{\dd}{\mathrm d}\newcommand{\dif}{\mathop{}\!\mathrm d}
\newcommand{\PP}{P}\newcommand{\EE}{\mathbb E}\newcommand{\Exp}{\operatorname{Exp}}
\newcommand{\indep}{\perp\!\!\!\perp}
\newcommand{\basis}[1]{\par\begingroup\small\raggedright\noindent\textbf{所用结论与原文依据}\quad #1\par\endgroup\smallskip}
\newcommand{\origin}[1]{\par\noindent{\small #1}\par\medskip}
\newcounter{statement}[section]\renewcommand{\thestatement}{\thesection.\arabic{statement}}
\newcommand{\statementstart}[4]{\par\medskip\refstepcounter{statement}\label{#3:#4}\noindent\textbf{#1~\thestatement\quad #2}\par\smallskip}
\newenvironment{mydefinition}[2]{\statementstart{定义}{#1}{def}{#2}}{\par\medskip}
\newenvironment{mytheorem}[2]{\statementstart{定理}{#1}{th}{#2}}{\par\medskip}
\newenvironment{proposition}[2]{\statementstart{命题}{#1}{prop}{#2}}{\par\medskip}
\newenvironment{mycorollary}[2]{\statementstart{推论}{#1}{cor}{#2}}{\par\medskip}
\newenvironment{mylemma}[2]{\statementstart{引理}{#1}{lem}{#2}}{\par\medskip}
\title{概率论期中复习\\\large 原笔记前三章例题与第一至十二周作业详解}
\author{}\date{2026年11月23日}
\begin{document}
\maketitle
\noindent 本册分为概率空间与随机事件、一维随机变量、随机向量三部分。每部分先列原笔记中的定义与结论，再列该章全部例题与相应阶段的作业。每题先完整陈述已知条件与所求，再给出解答或证明；多问例题先集中列出各问，解答按顺序对应。原笔记定理不重复证明；作业要求证明的内容仍给出完整证明。第四章不在本册范围内。

原笔记共收录78个例题。作业中的教材题已按指定页码补全题面；选做题保留，重复布置的题目保留各周出处。仅由讲义编号引用且未找到可核对题面的第四周第2题按要求略去。第一周第4项是未指定题目的开放阅读，保留该项说明。

题面存在条件不足、单位不一致或措辞冲突时，原文保留，解答中分别说明字面结论和所采用的明确解释。
\tableofcontents
\clearpage
\part{概率空间与随机事件}
% !TEX root = 概率论期中复习.tex
% 自动抽取自 chapters/chapter1.tex；原例题及其解答保留。

\clearpage\section{随机试验和事件}\label{k:events}
事件的全集为$\Omega$；$AB=A\cap B$，$A^c=\Omega\setminus A$。并表示至少一个发生，交表示同时发生，差表示前者发生而后者不发生。

\begin{enumerate}
\item \textbf{互斥（互不相容）事件}

若 $A\cap B=\varnothing$，则称事件 $A$ 与 $B$ 互斥或互不相容，其概率意义是二者不能同时发生。若对一列事件 $A_n\ (n\geq1)$，任意 $i\ne j$ 都有 $A_i\cap A_j=\varnothing$，则称该列事件两两互斥。

\item \textbf{补事件}

事件 $A$ 的补集
\[
A^c=\left\{\omega\in\Omega\mathrel{\big|}\omega\notin A\right\}
\]
称为 $A$ 的补事件，其概率意义是 $A$ 不发生。并且
\[
A\setminus B=A\cap B^c.
\]
事件运算约定先进行补运算，再进行交运算，最后进行并运算或差运算。

\item \textbf{完备事件组}

若事件 $A_1,A_2,\ldots$ 两两互斥，且
\[
\bigcup_{i=1}^{\infty}A_i=\Omega,
\]
则称 $A_1,A_2,\ldots$ 构成一个\textbf{完备事件组}，或称为 $\Omega$ 的一个可数划分。
\end{enumerate}

事件的运算性质就是集合的运算性质。对任意事件 $A,B,C$，有
\[
A\cup B=B\cup A,\qquad A\cap B=B\cap A,
\]
\[
(A\cup B)\cup C=A\cup(B\cup C),\qquad
(A\cap B)\cap C=A\cap(B\cap C),
\]
\[
A\cup(B\cap C)=(A\cup B)\cap(A\cup C),\qquad
A\cap(B\cup C)=(A\cap B)\cup(A\cap C).
\]
此外，德摩根律为
\[
\left(\bigcup_{i=1}^{n}A_i\right)^c=\bigcap_{i=1}^{n}A_i^c,
\qquad
\left(\bigcap_{i=1}^{n}A_i\right)^c=\bigcup_{i=1}^{n}A_i^c,
\]
以及
\[
\left(\bigcup_{i=1}^{\infty}A_i\right)^c=\bigcap_{i=1}^{\infty}A_i^c,
\qquad
\left(\bigcap_{i=1}^{\infty}A_i\right)^c=\bigcup_{i=1}^{\infty}A_i^c.
\]

\begin{mydefinition}{随机试验}{random-experiment}
在概率论中，将满足下列条件的试验称为随机试验，简称试验：
\begin{enumerate}
\item 试验进行前，无法确切知道一定会出现哪个结果；
\item 已知试验所有可能的结果；
\item 试验结束后，能够明确知道哪个试验结果出现。
\end{enumerate}
\end{mydefinition}
\origin{原笔记第1.1节，定义“随机试验”；源码第12行。}

\begin{mydefinition}{样本点与样本空间}{sample-space}
\begin{enumerate}
\item 一个随机试验的任意一个可能结果称为一个\textbf{样本点}（sample），记为 $\omega$；
\item 该随机试验所有可能结果组成的集合称为该试验的\textbf{样本空间}（sample space），记为 $\Omega$。
\end{enumerate}
\end{mydefinition}
\origin{原笔记第1.1节，定义“样本点与样本空间”；源码第21行。}

\begin{mydefinition}{事件}{event}
\begin{enumerate}
\item 样本空间 $\Omega$ 的一个子集称为一个\textbf{事件}，通常用大写字母 $A,B,C,\ldots$ 表示；
\item 对事件 $A$，若试验结果 $\omega\in A$，则称事件 $A$ 发生；否则称事件 $A$ 不发生。
\end{enumerate}
\end{mydefinition}
\origin{原笔记第1.1节，定义“事件”；源码第34行。}

\begin{mydefinition}{古典概率模型}{classical-probability-model}
若某随机试验的样本空间 $\Omega$ 只含有限个样本点，且每个样本点出现的概率相等，则称该随机试验为\textbf{古典概率模型}，简称\textbf{古典概型}。对任意事件 $A\subseteq\Omega$，记 $A$ 中元素个数为 $n(A)$，则
\[
P(A)=\dfrac{n(A)}{n(\Omega)}.
\]
\end{mydefinition}
\origin{原笔记第1.1节，定义“古典概率模型”；源码第142行。}

\begin{proposition}{计数的乘法原理}{multiplication-principle}
若一个过程由 $r$ 个试验构成，第 $i$ 个试验有 $n_i$ 种可能结果，则该过程共有
\[
n_1n_2\cdots n_r
\]
种可能结果。
\end{proposition}
\origin{原笔记第1.1节，命题“计数的乘法原理”；源码第151行。}

\subsection{例题：抽取方式与排列计数}\label{note:c1:e1}
\paragraph{题目}
设 $n,m$ 为正整数。求下列计数问题中的不同结果数：
\begin{enumerate}
\item 从 $n$ 个不同的球中有放回地每次取一个，依次取 $m$ 次，按取出顺序排成一列。
\item 从 $n$ 个不同的球中无放回地取出 $m$ 个球，按顺序排成一列，其中 $m\leq n$。
\item $n$ 个不同的人围圆桌坐一圈，将旋转后重合的坐法视为同一种坐法。
\item 从 $n$ 个不同的球中无放回地取出 $m$ 个球，不计顺序，其中 $m\leq n$。
\end{enumerate}

\basis{原笔记第1.1节，例“抽取方式与排列计数”；源码\texttt{chapters/chapter1.tex}第158行。}
\begin{proof}[解答]
\begin{enumerate}
\item 每次都有 $n$ 种选择，由乘法原理，结果数为
\[n^m.\]
\item 依次有 $n,n-1,\ldots,n-m+1$ 种选择，因此排列数为
\[A_n^m=n(n-1)\cdots(n-m+1)=\frac{n!}{(n-m)!}.\]
\item 固定其中一人的位置，其余 $n-1$ 人按顺序安排，恰好消除旋转造成的重复，故圆排列数为
\[(n-1)!.\]
\item 每个不计顺序的组合对应 $m!$ 个有序排列，因此组合数为
\[C_n^m=\binom nm=\frac{A_n^m}{m!}=\frac{n!}{(n-m)!m!}.\]
\end{enumerate}
\end{proof}

\begin{proposition}{多组分配}{multinomial-count}
将 $n$ 个不同元素分成 $r$ 个组，其中第 $i$ 组有 $n_i$ 个元素，组内不计顺序，且
\[
n_1+n_2+\cdots+n_r=n,
\]
则不同分法共有
\[
\dfrac{n!}{n_1!n_2!\cdots n_r!}
\]
种。
\end{proposition}
\origin{原笔记第1.1节，命题“多组分配”；源码第175行。}

\begin{mycorollary}{多项式定理}{multinomial-theorem}
\[
(x_1+x_2+\cdots+x_r)^n
=\sum_{n_1+n_2+\cdots+n_r=n}
\dfrac{n!}{n_1!n_2!\cdots n_r!}
x_1^{n_1}x_2^{n_2}\cdots x_r^{n_r}.
\]
其中
\[
\binom{n}{n_1,n_2,\ldots,n_r}
:=\dfrac{n!}{n_1!n_2!\cdots n_r!}
\]
称为多项式系数。
\end{mycorollary}
\origin{原笔记第1.1节，推论“多项式定理”；源码第202行。}

\begin{proposition}{正整数解的个数}{positive-integer-solutions}
设 $n\geq1$，方程
\[
x_1+x_2+\cdots+x_r=n,\qquad x_i>0\quad(i=1,2,\ldots,r)
\]
有
\[
\binom{n-1}{r-1}
\]
个不同的正整数解。
\end{proposition}
\origin{原笔记第1.1节，命题“正整数解的个数”；源码第217行。}

\begin{proposition}{非负整数解的个数}{nonnegative-integer-solutions}
设 $n\geq1$，方程
\[
x_1+x_2+\cdots+x_r=n,\qquad x_i\geq0\quad(i=1,2,\ldots,r)
\]
有
\[
\binom{n+r-1}{r-1}
\]
个不同的非负整数解。
\end{proposition}
\origin{原笔记第1.1节，命题“非负整数解的个数”；源码第238行。}

\begin{mydefinition}{几何概率模型}{geometric-probability-model}
设随机试验的样本空间 $\Omega$ 是 $\mathbb{R}^d$ 中一个具有体积 $m(\Omega)$ 的子集，且
\[
0<m(\Omega)<+\infty.
\]
若事件 $A\subseteq\Omega$ 具有体积 $m(A)$，并规定
\[
P(A)=\dfrac{m(A)}{m(\Omega)},
\]
则称该随机试验为\textbf{几何概率模型}，简称\textbf{几何概型}。
\end{mydefinition}
\origin{原笔记第1.1节，定义“几何概率模型”；源码第263行。}

\subsection{例题：贝特朗奇论}\label{note:c1:e2}
\paragraph{题目}
随机地从半径为 $r$ 的圆中取一根弦，求该弦长度大于圆内接正三角形边长 $\sqrt{3}r$ 的概率。设事件 $A$ 表示“弦长大于 $\sqrt{3}r$”。

\begin{enumerate}
\item 以弦中点 $M$ 确定弦，规定 $M$ 在圆内均匀取值，求 $P(A)$。
\item 固定弦的一个端点，规定弦与该端点处切线的夹角 $\theta$ 在 $[0,\pi]$ 上均匀取值，求 $P(A)$。
\item 规定弦到圆心的距离 $D$ 在 $[0,r]$ 上均匀取值，求 $P(A)$。
\end{enumerate}
说明三种结果为何不同。

\basis{原笔记第1.1节，例“贝特朗奇论”；源码\texttt{chapters/chapter1.tex}第279行。}
\begin{proof}[解答]
“随机取弦”的说法本身并未说明等可能的参照标准，因此可以得到不同的概率。
\begin{enumerate}
\item 若以弦的中点 $M$ 确定一根弦，并规定 $M$ 在大圆内均匀取值，则弦长大于 $\sqrt{3}r$ 当且仅当 $M$ 落在半径为 $r/2$ 的同心小圆内。因此
\[
P(A)=\dfrac{\pi(r/2)^2}{\pi r^2}=\dfrac14.
\]

\item 若固定弦的一个端点，并规定弦与该端点处切线的夹角 $\theta$ 在 $[0,\pi]$ 上均匀取值，则
\[
A=\left(\dfrac{\pi}{3},\dfrac{2\pi}{3}\right),
\qquad
P(A)=\dfrac{\pi/3}{\pi}=\dfrac13.
\]

\item 若以弦到圆心的距离 $D$ 确定一根弦，并规定 $D$ 在 $[0,r]$ 上均匀取值，则 $A=\{0\leq D<r/2\}$，从而
\[
P(A)=\dfrac{r/2}{r}=\dfrac12.
\]
\end{enumerate}

三种理解得到不同答案的原因在于：它们分别选择了不同的样本空间和等可能标准。求概率前，必须先明确随机试验的样本空间以及概率模型。
\end{proof}

\subsection{例题：布丰投针实验}\label{note:c1:e3}
\paragraph{题目}
在平面上画出间距为 $d$ 的一组平行直线。随机投掷一根长度为 $l$ 的针，设
\[
0<l\leq d.
\]
记针与任一直线相交为事件 $A$。令 $X$ 表示针的中点到最近平行线的距离，$\Theta$ 表示针与平行线的夹角。按照几何概型，取
\[
0\leq X\leq \dfrac d2,
\qquad
0\leq \Theta\leq \dfrac{\pi}{2},
\]
并把 $(X,\Theta)$ 看作在该矩形内等可能取值。

\begin{enumerate}
\item 求针与任一平行线相交的概率 $P(A)$。
\item 若独立投针 $N$ 次，其中 $H>0$ 次与直线相交，给出 $N$ 较大时 $\pi$ 的实验估计式。
\end{enumerate}

\basis{原笔记第1.1节，例“布丰投针实验”；源码\texttt{chapters/chapter1.tex}第305行。}
\begin{proof}[解答]
\begin{enumerate}
\item 针与直线相交当且仅当
\[
X\leq \dfrac l2\sin\Theta.
\]
因此，事件 $A$ 对应的面积为
\[
\int_0^{\pi/2}\dfrac l2\sin\theta\,\mathrm{d}\theta
=\dfrac l2,
\]
而样本空间的面积为 $\dfrac d2\cdot\dfrac{\pi}{2}=\dfrac{d\pi}{4}$。故
\[
P(A)=\dfrac{l/2}{d\pi/4}=\dfrac{2l}{\pi d}.
\]
\item 若独立投针 $N$ 次，其中有 $H$ 次与直线相交，则当 $N$ 较大时，频率 $H/N$ 可近似看作 $P(A)$，从而得到 $\pi$ 的实验估计式
\[
\pi\approx\dfrac{2lN}{dH}.
\]
\end{enumerate}
\end{proof}

\clearpage\section{可测空间和概率空间}\label{k:probability}
事件列$A_n\uparrow A$表示$A_n\subseteq A_{n+1}$且$A=\bigcup_nA_n$；$A_n\downarrow A$表示$A_{n+1}\subseteq A_n$且$A=\bigcap_nA_n$。

\begin{mydefinition}{频率}{frequency}
在 $N$ 次重复试验中，若事件 $A$ 出现 $N(A)$ 次，则比值
\[
f_N(A)=\dfrac{N(A)}{N}
\]
称为事件 $A$ 发生的\textbf{频率}。
\end{mydefinition}
\origin{原笔记第1.2节，定义“频率”；源码第363行。}

\begin{proposition}{频率的性质}{frequency-properties}
对任意事件 $A$，有
\[
0\leq f_N(A)\leq1,
\qquad
f_N(\Omega)=1,
\qquad
f_N(\varnothing)=0.
\]
并且
\[
f_N(A^c)=1-f_N(A).
\]
若 $A_1,A_2,\ldots,A_r$ 两两互斥，则
\[
f_N\left(\bigcup_{i=1}^{r}A_i\right)
=\sum_{i=1}^{r}f_N(A_i).
\]
\end{proposition}
\origin{原笔记第1.2节，命题“频率的性质”；源码第371行。}

\begin{mydefinition}{$\sigma$-域与可测空间}{sigma-algebra}
设 $\Omega$ 是非空集合，$\mathscr{F}$ 是由 $\Omega$ 的一些子集构成的集合类。若 $\mathscr{F}$ 满足：
\begin{enumerate}
\item $\Omega\in\mathscr{F}$；
\item 若 $A\in\mathscr{F}$，则 $A^c\in\mathscr{F}$；
\item 若 $A_n\in\mathscr{F}$（$n\geq1$），则
\[
\bigcup_{n=1}^{\infty}A_n\in\mathscr{F},
\]
\end{enumerate}
则称 $\mathscr{F}$ 为 $\Omega$ 的一个\textbf{$\sigma$-域}或\textbf{事件域}，在分析学中我们称之为\textbf{$\sigma$-代数}，$\mathscr{F}$ 中的元素称为事件或可测集，二元组 $(\Omega,\mathscr{F})$ 称为\textbf{可测空间}。
\end{mydefinition}
\origin{原笔记第1.2节，定义“$\sigma$-域与可测空间”；源码第412行。}

\begin{proposition}{$\sigma$-域的基本封闭性}{sigma-algebra-closure}
若 $\mathscr{F}$ 是 $\Omega$ 的一个 $\sigma$-域，则：
\begin{enumerate}
\item $\varnothing\in\mathscr{F}$；
\item 若 $A_n\in\mathscr{F}$（$n\geq1$），则
\[
\bigcap_{n=1}^{\infty}A_n\in\mathscr{F};
\]
\item 若 $A_1,\ldots,A_k\in\mathscr{F}$，则有限并与有限交都属于 $\mathscr{F}$；
\item 若 $A,B\in\mathscr{F}$，则 $A\setminus B\in\mathscr{F}$。
\end{enumerate}
\end{proposition}
\origin{原笔记第1.2节，命题“$\sigma$-域的基本封闭性”；源码第425行。}

\subsection{例题：\texorpdfstring{一些 $\sigma$-代数的例子}{一些 sigma-代数的例子}}\label{note:c1:e4}
\paragraph{题目}
设 $\Omega$ 为非空集合，$A\subseteq\Omega$。验证下列集合类是 $\sigma$-域，并验证前两项所述的最小性：
\begin{enumerate}
\item $\{\varnothing,\Omega\}$ 是 $\Omega$ 上所含事件最少的 $\sigma$-域；
\item 对任意 $A\subseteq\Omega$，集合类
\[
\{\varnothing,\Omega,A,A^c\}
\]
是包含 $A$ 的最小的 $\sigma$-域；
\item $\mathcal{P}(\Omega):=\{A\mid A\subseteq\Omega\}$ 是 $\Omega$ 的所有子集构成的 $\sigma$-域。
\end{enumerate}

\basis{原笔记第1.2节，例“一些 $\sigma$-代数的例子”；源码\texttt{chapters/chapter1.tex}第451行。}
\begin{proof}[验证]
\begin{enumerate}
\item 集合类 $\{\varnothing,\Omega\}$ 含 $\Omega$，取补仍在其中；其中任意可列个集合的并只能是 $\varnothing$ 或 $\Omega$，所以它是 $\sigma$-域。任意 $\sigma$-域都必须含 $\Omega$ 及其补集 $\varnothing$，故它所含事件最少。
\item 将 $\Omega$ 分为 $A$ 与 $A^c$ 两部分，所列集合恰为这两部分的所有可能并集，故取补和可列并仍为其中一个集合。它含 $\Omega$，所以是 $\sigma$-域。任意包含 $A$ 的 $\sigma$-域还必须包含 $A^c,\varnothing,\Omega$，故所列集合类是包含 $A$ 的最小 $\sigma$-域。若 $A$ 为空集或全集，重复项合并后即为第一项。
\item $\mathcal P(\Omega)$ 含 $\Omega$，其任意元素的补集和任意可列个元素的并仍是 $\Omega$ 的子集，故仍属于 $\mathcal P(\Omega)$。因此它是 $\sigma$-域。
\end{enumerate}
\end{proof}

\begin{proposition}{$\sigma$-域的交}{intersection-of-sigma-algebras}
设 $I$ 是一个指标集。若对任意 $\alpha\in I$，$\mathscr{F}_\alpha$ 都是 $\Omega$ 上的 $\sigma$-域，则
\[
\mathscr{F}=\bigcap_{\alpha\in I}\mathscr{F}_\alpha
\]
也是 $\Omega$ 上的 $\sigma$-域。
\end{proposition}
\origin{原笔记第1.2节，命题“$\sigma$-域的交”；源码第465行。}

\begin{mytheorem}{由集合类生成的 $\sigma$-域}{generated-sigma-algebra}
设 $\Omega$ 是非空集合，$\mathscr{E}$ 是由 $\Omega$ 的一些子集构成的非空集合类。则存在唯一的 $\sigma$-域 $\mathscr{F}$ 满足：
\begin{enumerate}
\item $\mathscr{E}\subseteq\mathscr{F}$；
\item 对任意包含 $\mathscr{E}$ 的 $\sigma$-域 $\mathcal{M}$，都有 $\mathscr{F}\subseteq\mathcal{M}$。
\end{enumerate}
此时称 $\mathscr{F}$ 为由 $\mathscr{E}$ \textbf{生成的 $\sigma$-域}，记为 $\sigma(\mathscr{E})$。
\end{mytheorem}
\origin{原笔记第1.2节，定理“由集合类生成的 $\sigma$-域”；源码第477行。}

\begin{mydefinition}{Borel $\sigma$-域}{borel-sigma-algebra}
记
\[
\mathscr{E}_{\mathbb{R}}=\{(a,b]\mid a,b\in\mathbb{R}\}.
\]
称 $\sigma(\mathscr{E}_{\mathbb{R}})$ 为 $\mathbb{R}$ 上的\textbf{Borel $\sigma$-域}，记为 $\mathscr{B}(\mathbb{R})$。

类似地，记
\[
\mathscr{E}_{\mathbb{R}^d}
=\left\{(a_1,b_1]\times\cdots\times(a_d,b_d]\mid a_i,b_i\in\mathbb{R},\ 1\leq i\leq d\right\}.
\]
称 $\sigma(\mathscr{E}_{\mathbb{R}^d})$ 为 $\mathbb{R}^d$ 上的 Borel $\sigma$-域，记为 $\mathscr{B}(\mathbb{R}^d)$。
\end{mydefinition}
\origin{原笔记第1.2节，定义“Borel $\sigma$-域”；源码第495行。}

\begin{proposition}{常用集合属于 $\mathscr{B}(\mathbb{R})$}{common-borel-sets}
对任意 $a<b$，开区间 $(a,b)$、闭区间 $[a,b]$、半开半闭区间 $[a,b)$ 以及单点集 $\{b\}$ 都属于 $\mathscr{B}(\mathbb{R})$。
\end{proposition}
\origin{原笔记第1.2节，命题“常用集合属于 $\mathscr{B}(\mathbb{R})$”；源码第510行。}

\begin{mydefinition}{概率与概率空间}{probability-space}
设 $\Omega$ 是非空集合，$\mathscr{F}$ 是 $\Omega$ 上给定的一个 $\sigma$-域。若映射
\[
P:\mathscr{F}\longrightarrow[0,+\infty)
\]
满足：
\begin{enumerate}
\item \textbf{非负性：}对任意 $A\in\mathscr{F}$，有 $P(A)\geq0$；
\item \textbf{规范性：}$P(\Omega)=1$；
\item \textbf{可列可加性：}若 $A_n\in\mathscr{F}$（$n\geq1$）两两互斥，则
\[
P\left(\bigcup_{n=1}^{\infty}A_n\right)
=\sum_{n=1}^{\infty}P(A_n),
\]
\end{enumerate}
则称 $P$ 为 $\mathscr{F}$ 上的\textbf{概率}，三元组 $(\Omega,\mathscr{F},P)$ 称为\textbf{概率空间}。
\end{mydefinition}
\origin{原笔记第1.2节，定义“概率与概率空间”；源码第536行。}

\begin{proposition}{空集的概率}{probability-empty-set}
\[
P(\varnothing)=0.
\]
\end{proposition}
\origin{原笔记第1.2节，命题“空集的概率”；源码第555行。}

\begin{proposition}{有限可加性}{finite-additivity}
若 $A_1,\ldots,A_n\in\mathscr{F}$ 两两互斥，则
\[
P\left(\bigcup_{i=1}^{n}A_i\right)
=\sum_{i=1}^{n}P(A_i).
\]
\end{proposition}
\origin{原笔记第1.2节，命题“有限可加性”；源码第576行。}

\begin{proposition}{单调性}{probability-monotonicity}
若 $A,B\in\mathscr{F}$ 且 $B\subseteq A$，则
\[
P(B)\leq P(A).
\]
\end{proposition}
\origin{原笔记第1.2节，命题“单调性”；源码第588行。}

\begin{mycorollary}{概率的上界}{probability-upper-bound}
对任意 $A\in\mathscr{F}$，有
\[
0\leq P(A)\leq1.
\]
\end{mycorollary}
\origin{原笔记第1.2节，推论“概率的上界”；源码第602行。}

\begin{proposition}{补事件的概率}{probability-complement}
对任意 $A\in\mathscr{F}$，有
\[
P(A^c)=1-P(A).
\]
\end{proposition}
\origin{原笔记第1.2节，命题“补事件的概率”；源码第613行。}

\begin{proposition}{一般加法公式}{probability-addition-formula}
对任意事件 $A,B\in\mathscr{F}$，有
\[
P(A\cup B)=P(A)+P(B)-P(A\cap B).
\]
\end{proposition}
\origin{原笔记第1.2节，命题“一般加法公式”；源码第627行。}

\begin{mytheorem}{容斥恒等式}{inclusion-exclusion}
对任意 $n$ 个事件 $A_1,A_2,\ldots,A_n$，有
\[
P\left(\bigcup_{i=1}^{n}A_i\right)
=\sum_{r=1}^{n}(-1)^{r+1}
\sum_{1\leq i_1<\cdots<i_r\leq n}
P\left(A_{i_1}\cap\cdots\cap A_{i_r}\right).
\]
\end{mytheorem}
\origin{原笔记第1.2节，定理“容斥恒等式”；源码第655行。}

\begin{mycorollary}{有限次次可加性}{finite-subadditivity}
对任意 $n$ 个事件 $A_1,\ldots,A_n$，有
\[
P\left(\bigcup_{i=1}^{n}A_i\right)
\leq\sum_{i=1}^{n}P(A_i).
\]
\end{mycorollary}
\origin{原笔记第1.2节，推论“有限次次可加性”；源码第677行。}

\subsection{例题：Kounias 不等式}\label{note:c1:e5}
\paragraph{题目}
证明：对任意 $n$ 个事件 $A_1,\ldots,A_n$，有
\[
P\left(\bigcup_{r=1}^{n}A_r\right)
\leq\min_{1\leq k\leq n}
\left\{
\sum_{r=1}^{n}P(A_r)-\sum_{\substack{1\leq r\leq n\\ r\ne k}}P(A_r\cap A_k)
\right\}.
\]

\basis{原笔记第1.2节，例“Kounias 不等式”；源码\texttt{chapters/chapter1.tex}第696行。}
\begin{proof}
固定 $k$。由有限次次可加性，
\[
\begin{aligned}
\sum_{r=1}^{n}P(A_r)-\sum_{r\ne k}P(A_r\cap A_k)
&=P(A_k)+\sum_{r\ne k}P(A_r\cap A_k^c)\\
&\geq P\left(A_k\cup\bigcup_{r\ne k}(A_r\cap A_k^c)\right)\\
&=P\left(\bigcup_{r=1}^{n}A_r\right).
\end{aligned}
\]
对 $k=1,\ldots,n$ 取最小值即得结论。
\end{proof}

\begin{mydefinition}{事件序列的极限}{limit-of-event-sequences}
若事件列 $\{A_n\}_{n\geq1}$ 满足
\[
A_1\subseteq A_2\subseteq\cdots\subseteq A_n\subseteq\cdots,
\]
则称其为\textbf{递增序列}（或不降序列），并定义
\[
\lim_{n\to\infty}A_n:=\bigcup_{n=1}^{\infty}A_n.
\]
若
\[
A_1\supseteq A_2\supseteq\cdots\supseteq A_n\supseteq\cdots,
\]
则称其为\textbf{递减序列}（或不增序列），并定义
\[
\lim_{n\to\infty}A_n:=\bigcap_{n=1}^{\infty}A_n.
\]
\end{mydefinition}
\origin{原笔记第1.2节，定义“事件序列的极限”；源码第720行。}

\begin{mytheorem}{概率的连续性}{continuity-of-probability}
\begin{enumerate}
\item 若 $\{A_n\}_{n\geq1}$ 为递增序列，则
\[
P\left(\lim_{n\to\infty}A_n\right)=\lim_{n\to\infty}P(A_n).
\]
称为概率的\textbf{下连续性}；
\item 若 $\{A_n\}_{n\geq1}$ 为递减序列，则
\[
P\left(\lim_{n\to\infty}A_n\right)=\lim_{n\to\infty}P(A_n).
\]
称为概率的\textbf{上连续性}。
\end{enumerate}
\end{mytheorem}
\origin{原笔记第1.2节，定理“概率的连续性”；源码第739行。}

\begin{proposition}{可列次可加性（布尔不等式）}{countable-subadditivity}
对任意可列个事件 $A_1,A_2,\ldots$，有
\[
P\left(\bigcup_{i=1}^{\infty}A_i\right)
\leq\sum_{i=1}^{\infty}P(A_i).
\]
\end{proposition}
\origin{原笔记第1.2节，命题“可列次可加性（布尔不等式）”；源码第779行。}

\subsection{例题：可列个零概率事件的并}\label{note:c1:e6}
\paragraph{题目}
证明：若事件列 $\{A_i\}_{i\geq1}$ 满足 $P(A_i)=0$（$i\geq1$），则
\[
P\left(\bigcup_{i=1}^{\infty}A_i\right)=0.
\]

\basis{原笔记第1.2节，例“可列个零概率事件的并”；源码\texttt{chapters/chapter1.tex}第803行。}
\begin{proof}
由布尔不等式，
\[
0\leq P\left(\bigcup_{i=1}^{\infty}A_i\right)
\leq\sum_{i=1}^{\infty}P(A_i)=0.
\]
\end{proof}

\begin{mydefinition}{事件列的上极限与下极限}{limsup-liminf-events}
设 $A_n\in\mathscr{F}$（$n\geq1$），不要求事件列 $\{A_n\}$ 单调。定义
\[
\limsup_{n\to\infty}A_n
:=\bigcap_{n=1}^{\infty}\bigcup_{k=n}^{\infty}A_k
\]
为事件列的\textbf{上极限}，定义
\[
\liminf_{n\to\infty}A_n
:=\bigcup_{n=1}^{\infty}\bigcap_{k=n}^{\infty}A_k
\]
为事件列的\textbf{下极限}。
\end{mydefinition}
\origin{原笔记第1.2节，定义“事件列的上极限与下极限”；源码第824行。}

\begin{mydefinition}{事件列极限的存在}{event-limit-exists}
	当一个事件列的上极限与下极限相等时，也即
	\[\liminf_{n\to\infty}A_n =\limsup_{n\to\infty}A_n\]
	时，我们定义这个事件列的极限为
	\[\lim\limits_{n\to\infty}A_n :=\liminf_{n\to\infty}A_n =\limsup_{n\to\infty}A_n\]
\end{mydefinition}
\origin{原笔记第1.2节，定义“事件列极限的存在”；源码第860行。}

\subsection{例题：伯努利概率空间}\label{note:c1:e7}
\paragraph{题目}
设 $\Omega$ 是非空集合，$\varnothing\subsetneq A\subsetneq\Omega$，并令
\[
\mathscr{F}=\{\varnothing,\Omega,A,A^c\}.
\]
给定 $0<p<1$，定义
\[
P(\varnothing)=0,\qquad P(\Omega)=1,
\qquad P(A)=p,\qquad P(A^c)=1-p.
\]
验证 $(\Omega,\mathscr F,P)$ 是概率空间，并说明此模型表示的试验。

\basis{原笔记第1.2节，例“伯努利概率空间”；源码\texttt{chapters/chapter1.tex}第868行。}
\begin{proof}[验证]
$\mathscr F$ 是由 $A,A^c$ 的所有并集构成的 $\sigma$-域。所给四个概率均非负，且 $P(\Omega)=1$。

任取 $\mathscr F$ 中两两互斥的可列事件，非空事件至多有两个。若只有一个非空事件，可列可加性直接成立；若有两个，则只能为 $A$ 与 $A^c$，且
\[P(A\cup A^c)=1=p+(1-p)=P(A)+P(A^c).\]
若全部为空集，两边均为零。因此 $P$ 满足可列可加性，$(\Omega,\mathscr F,P)$ 构成概率空间，称为\textbf{伯努利概率空间}（或伯努利概率模型），用于讨论单个事件 $A$ 是否发生。
\end{proof}

\subsection{例题：古典概率空间}\label{note:c1:e8}
\paragraph{题目}
设 $\Omega$ 是只含有限个样本点的非空集合，并取
\[
\mathscr{F}=2^\Omega=\{A:A\subseteq\Omega\}.
\]
若 $n(A)$ 表示集合 $A$ 中元素的个数，定义
\[
P(A)=\frac{n(A)}{n(\Omega)},\qquad A\in\mathscr{F},
\]
验证 $(\Omega,\mathscr F,P)$ 是概率空间，并求每个单点事件的概率。

\basis{原笔记第1.2节，例“古典概率空间”；源码\texttt{chapters/chapter1.tex}第881行。}
\begin{proof}[验证]
$\mathscr F=2^\Omega$ 是 $\Omega$ 上的 $\sigma$-域。由计数的非负性，$P(A)\geq0$；又有 $P(\Omega)=n(\Omega)/n(\Omega)=1$。

若 $A_1,A_2,\ldots$ 两两互斥，由 $\Omega$ 有限，至多有限个 $A_i$ 非空。各非空集合的元素不重复，故
\[P\left(\bigcup_{i=1}^\infty A_i\right)=\frac{n(\bigcup_i A_i)}{n(\Omega)}=\sum_{i=1}^\infty\frac{n(A_i)}{n(\Omega)}=\sum_{i=1}^\infty P(A_i).\]
因此它构成古典概率空间。每个 $\omega\in\Omega$ 都满足
\[P(\{\omega\})=\frac1{n(\Omega)}.\]
\end{proof}

\subsection{例题：最大号码问题}\label{note:c1:e9}
\paragraph{题目}
某城市有编号为 $1,2,\ldots,N$ 的 $N$ 部卡车。独立记录遇到的 $n$ 辆卡车的编号，允许重复，且每次记录到各编号的概率均为 $1/N$。求记录到的最大编号为 $k$（$1\leq k\leq N$）的概率。

\basis{原笔记第1.2节，例“最大号码问题”；源码\texttt{chapters/chapter1.tex}第895行。}
\begin{proof}[解答]
令 $B_k$ 表示最大编号恰为 $k$，$A_k$ 表示最大编号不超过 $k$。则
\[
B_1=A_1,\qquad B_k=A_k\setminus A_{k-1}\quad(2\leq k\leq N),
\]
且
\[
P(A_k)=\left(\frac{k}{N}\right)^n.
\]
因此
\[
P(B_1)=\frac{1}{N^n},
\]
而当 $2\leq k\leq N$ 时，
\[
P(B_k)=P(A_k)-P(A_{k-1})
=\frac{k^n-(k-1)^n}{N^n}.
\]
\end{proof}

\subsection{例题：配对问题}\label{note:c1:e10}
\paragraph{题目}
在有 $n$ 人参加的晚会上，每人带来一顶各不相同的帽子。离场时，每人从随机混放的帽子中无放回地抽取一顶，所有分配方式等可能。求：
\begin{enumerate}
\item 至少有一人抽到自己帽子的概率。
\item 当 $n\to\infty$ 时，至少有一人抽到自己帽子与无人抽到自己帽子的概率极限。
\end{enumerate}

\basis{原笔记第1.2节，例“配对问题”；源码\texttt{chapters/chapter1.tex}第920行。}
\begin{proof}[解答]
\begin{enumerate}
\item 给人和帽子分别编号 $1,2,\ldots,n$，令 $A_i$ 表示第 $i$ 人抽到自己的帽子。对任意
$1\leq i_1<i_2<\cdots<i_r\leq n$，有
\[
P(A_{i_1}\cap\cdots\cap A_{i_r})=\frac{(n-r)!}{n!}.
\]
同阶交事件共有 $\binom{n}{r}$ 个，故
\[
\sum_{1\leq i_1<\cdots<i_r\leq n}
P(A_{i_1}\cap\cdots\cap A_{i_r})
=\binom{n}{r}\frac{(n-r)!}{n!}=\frac{1}{r!}.
\]
由容斥恒等式，所求概率为
\[
P\left(\bigcup_{i=1}^{n}A_i\right)
=\sum_{r=1}^{n}\frac{(-1)^{r+1}}{r!}
=1-\sum_{r=0}^{n}\frac{(-1)^r}{r!}.
\]
\item 因此
\[
\lim_{n\to\infty}P\left(\bigcup_{i=1}^{n}A_i\right)
=1-e^{-1}\approx0.6321,
\]
而无人抽到自己帽子的概率趋于 $e^{-1}\approx0.3679$。
\end{enumerate}
\end{proof}

\subsection{例题：圆桌夫妻问题}\label{note:c1:e11}
\paragraph{题目}
十对夫妻随机地沿圆桌坐下，求没有妻子坐在丈夫身边的概率。

\basis{原笔记第1.2节，例“圆桌夫妻问题”；源码\texttt{chapters/chapter1.tex}第950行。}
\begin{proof}[解答]
令 $A_i$ 表示第 $i$ 对夫妻相邻（$i=1,2,\ldots,10$）。总坐法数为
$(20-1)!=19!$。指定 $r$ 对夫妻相邻时，将每对夫妻看作一个整体，并计入夫妻内部的两种次序，故
\[
P(A_{i_1}\cap\cdots\cap A_{i_r})
=\frac{2^r(20-r-1)!}{19!}
=\frac{2^r(19-r)!}{19!}.
\]
由容斥恒等式，至少有一对夫妻相邻的概率为
\[
P\left(\bigcup_{i=1}^{10}A_i\right)
=\sum_{r=1}^{10}(-1)^{r+1}\binom{10}{r}
\frac{2^r(19-r)!}{19!}
\approx0.6605.
\]
所以，没有夫妻相邻的概率约为
\[
1-0.6605=0.3395.
\]
\end{proof}

\subsection{例题：反复抽取扑克牌中的大王}\label{note:c1:e12}
\paragraph{题目}
从一副 $54$ 张的扑克牌中反复抽牌，每次等可能地抽取一张，记录后放回并混匀，各次抽取相互独立。求：
\begin{enumerate}
\item 一直抽不到大王的概率；
\item 无穷多次抽到大王的概率。
\end{enumerate}

\basis{原笔记第1.2节，例“反复抽取扑克牌中的大王”；源码\texttt{chapters/chapter1.tex}第975行。}
\begin{proof}[解答]
设 $E_n$ 表示第 $n$ 次抽到大王。

\begin{enumerate}
\item 令
\[
C_n=\bigcap_{k=1}^{n}E_k^c.
\]
则 $C_n$ 表示前 $n$ 次均未抽到大王，且 $\{C_n\}$ 为递减序列。令
\[
C=\bigcap_{n=1}^{\infty}C_n,
\]
则 $C$ 表示始终抽不到大王。由独立性和概率的上连续性，
\[
P(C)=\lim_{n\to\infty}P(C_n)
=\lim_{n\to\infty}\left(\frac{53}{54}\right)^n=0.
\]
因此，最终至少抽到一次大王的概率为 $1$。
\item 令 $B$ 表示无穷多次抽到大王，则
\[
B=\limsup_{n\to\infty}E_n
=\bigcap_{n=1}^{\infty}\bigcup_{k=n}^{\infty}E_k.
\]
再令
\[
B_n=\bigcup_{k=n}^{\infty}E_k.
\]
事件 $B_n$ 表示“从第 $n$ 次抽牌起至少抽到一次大王”。由第一个问题同样的计算，
$P(B_n^c)=0$，故 $P(B_n)=1$。由于 $\{B_n\}$ 为递减序列，
\[
P(B)=P\left(\bigcap_{n=1}^{\infty}B_n\right)
=\lim_{n\to\infty}P(B_n)=1.
\]
\end{enumerate}
\end{proof}

\begin{mydefinition}{测度、有限测度与 $\sigma$ 有限测度}{measure-space}
设 $\Omega$ 是非空集合，$\mathscr{F}$ 是 $\Omega$ 上给定的一个 $\sigma$-域。若映射
\[
\mu:\mathscr{F}\longrightarrow[0,+\infty]
\]
满足：
\begin{enumerate}
\item $\mu(\varnothing)=0$；
\item 对任意两两互斥的事件列 $A_1,A_2,\ldots\in\mathscr{F}$，有
\[
\mu\left(\bigcup_{n=1}^{\infty}A_n\right)
=\sum_{n=1}^{\infty}\mu(A_n),
\]
\end{enumerate}
则称 $\mu$ 为 $\mathscr{F}$ 上的\textbf{测度}，三元组 $(\Omega,\mathscr{F},\mu)$ 称为\textbf{测度空间}。

\end{mydefinition}
\origin{原笔记第1.2节，定义“测度、有限测度与 $\sigma$ 有限测度”；源码第1024行。}

若存在两两互斥的事件列 $A_1,A_2,\ldots\in\mathscr{F}$，使得
\[
\bigcup_{n=1}^{\infty}A_n=\Omega,
\qquad
\mu(A_n)<+\infty\quad(n\geq1),
\]
则称 $\mu$ 为\textbf{$\sigma$ 有限测度}。若 $\mu(\Omega)<+\infty$，则称 $\mu$ 为\textbf{有限测度}；若 $\mu(\Omega)=1$，则称 $\mu$ 为\textbf{概率测度}，也简称为\textbf{概率}。

\begin{proposition}{测度的有限可加性}{measure-finite-additivity}
若 $A_1,\ldots,A_n\in\mathscr{F}$ 两两互斥，则
\[
\mu\left(\bigcup_{i=1}^{n}A_i\right)=\sum_{i=1}^{n}\mu(A_i).
\]
\end{proposition}
\origin{原笔记第1.2节，命题“测度的有限可加性”；源码第1050行。}

\begin{proposition}{测度的单调性}{measure-monotonicity}
若 $A,B\in\mathscr{F}$ 且 $B\subseteq A$，则
\[
\mu(B)\leq\mu(A).
\]
\end{proposition}
\origin{原笔记第1.2节，命题“测度的单调性”；源码第1061行。}

\begin{proposition}{测度的可减性}{measure-subtractivity}
若 $A,B\in\mathscr{F}$，$A\subseteq B$，且 $\mu(A)<+\infty$，则
\[
\mu(B\setminus A)=\mu(B)-\mu(A).
\]
\end{proposition}
\origin{原笔记第1.2节，命题“测度的可减性”；源码第1075行。}

\begin{proposition}{测度的容斥恒等式}{measure-inclusion-exclusion}
若 $A_1,A_2,\ldots,A_n\in\mathscr{F}$，且对每个 $1\leq k\leq n$ 都有 $\mu(A_k)<+\infty$，则
\[
\mu\left(\bigcup_{k=1}^{n}A_k\right)
=\sum_{r=1}^{n}(-1)^{r+1}
\sum_{1\leq i_1<i_2<\cdots<i_r\leq n}
\mu(A_{i_1}\cap A_{i_2}\cap\cdots\cap A_{i_r}).
\]
\end{proposition}
\origin{原笔记第1.2节，命题“测度的容斥恒等式”；源码第1090行。}

\begin{mytheorem}{测度的可列次可加性}{measure-countable-subadditivity}
对任意事件列 $A_1,A_2,\ldots\in\mathscr{F}$，有
\[
\mu\left(\bigcup_{n=1}^{\infty}A_n\right)
\leq\sum_{n=1}^{\infty}\mu(A_n).
\]
\end{mytheorem}
\origin{原笔记第1.2节，定理“测度的可列次可加性”；源码第1104行。}

\begin{proposition}{测度的下连续性}{measure-lower-continuity}
若 $\{A_n\}_{n\geq1}\subseteq\mathscr{F}$ 为递增序列，则
\[
\mu\left(\lim_{n\to\infty}A_n\right)
=\lim_{n\to\infty}\mu(A_n).
\]
\end{proposition}
\origin{原笔记第1.2节，命题“测度的下连续性”；源码第1129行。}

\begin{proposition}{测度的上连续性}{measure-upper-continuity}
若 $\{A_n\}_{n\geq1}\subseteq\mathscr{F}$ 为递减序列，且 $\mu(A_1)<+\infty$，则
\[
\mu\left(\lim_{n\to\infty}A_n\right)
=\lim_{n\to\infty}\mu(A_n).
\]
\end{proposition}
\origin{原笔记第1.2节，命题“测度的上连续性”；源码第1141行。}

\subsection{例题：Lebesgue 测度}\label{note:c1:e13}
\paragraph{题目}
在 $\mathscr{B}(\mathbb{R}^d)$ 上存在唯一的测度 $\lambda$，使得对任意半开长方体
\[
(a_1,b_1]\times(a_2,b_2]\times\cdots\times(a_d,b_d]
\]
（其中 $a_i\leq b_i$，$1\leq i\leq d$），都有
\[
\lambda\left((a_1,b_1]\times\cdots\times(a_d,b_d]\right)
=(b_1-a_1)(b_2-a_2)\cdots(b_d-a_d).
\]
该测度称为 $\mathbb{R}^d$ 上的\textbf{Lebesgue 测度}，它给出了通常所说的长度、面积和体积。以该测度及其长方体公式为已知，验证它在 $d=1,2,3$ 时分别给出半开区间、矩形和长方体的通常长度、面积和体积。

\basis{原笔记第1.2节，例“Lebesgue 测度”；源码\texttt{chapters/chapter1.tex}第1167行。}
\begin{proof}[验证]
由题设公式，当 $d=1$ 时，
\[\lambda((a_1,b_1])=b_1-a_1,\]
正是区间的长度。当 $d=2$ 时，
\[\lambda((a_1,b_1]\times(a_2,b_2])=(b_1-a_1)(b_2-a_2),\]
正是矩形面积。当 $d=3$ 时，
\[\lambda((a_1,b_1]\times(a_2,b_2]\times(a_3,b_3])=(b_1-a_1)(b_2-a_2)(b_3-a_3),\]
正是长方体体积。若某一边长为零，公式同样给出零。
\end{proof}

\clearpage\section{条件概率}\label{k:conditional}

\subsection{例题：掷两枚骰子}\label{note:c1:e14}
\paragraph{题目}
独立地依次掷两枚均匀骰子。已知第一枚骰子的点数为 $4$，求两枚骰子点数之和为偶数的概率。

\basis{原笔记第1.3节，例“掷两枚骰子”；源码\texttt{chapters/chapter1.tex}第1186行。}
\begin{proof}[解答]
设 $B$ 表示“第一枚骰子的点数为 $4$”，$A$ 表示“两枚骰子点数之和为偶数”。在 $B$ 已经发生时，第二枚骰子共有 $6$ 种等可能结果；其中点数为 $2,4,6$ 时 $A$ 发生，故
\[
P(A\mid B)=\frac{3}{6}=\frac12.
\]
\end{proof}

\begin{mydefinition}{条件概率}{conditional-probability}
设 $(\Omega,\mathscr{F},P)$ 为概率空间，$A,B\in\mathscr{F}$，且 $P(B)>0$。称
\[
P(A\mid B)=\frac{P(A\cap B)}{P(B)}
\]
为在事件 $B$ 发生的条件下事件 $A$ 的\textbf{条件概率}。
\end{mydefinition}
\origin{原笔记第1.3节，定义“条件概率”；源码第1195行。}

\subsection{例题：寻找钥匙}\label{note:c1:e15}
\paragraph{题目}
一串钥匙可能落在左口袋或右口袋。设 $L$ 表示“钥匙在左口袋”，$R$ 表示“钥匙在右口袋”，且
\[
P(L)=0.4,\qquad P(R)=0.4,\qquad R\subseteq L^{\mathrm c}.
\]
检查左口袋后发现没有钥匙，即已知 $L^{\mathrm c}$ 发生。求此时钥匙在右口袋的条件概率 $P(R\mid L^{\mathrm c})$。

\basis{原笔记第1.3节，例“寻找钥匙”；源码\texttt{chapters/chapter1.tex}第1203行。}
\begin{proof}[解答]
由 $R\subseteq L^{\mathrm c}$，有 $R\cap L^{\mathrm c}=R$，因此
\[
P(R\mid L^{\mathrm c})
=\frac{P(R\cap L^{\mathrm c})}{P(L^{\mathrm c})}
=\frac{P(R)}{1-P(L)}
=\frac{0.4}{0.6}
=\frac23.
\]
\end{proof}

\begin{mytheorem}{条件概率的乘法公式}{conditional-multiplication-formula}
若 $A,B\in\mathscr{F}$ 且 $P(B)>0$，则
\[
P(A\cap B)=P(B)P(A\mid B).
\]

更一般地，若 $A_1,\ldots,A_n\in\mathscr{F}$ 且
\[
P(A_1\cap\cdots\cap A_{n-1})>0,
\]
则
\[
\begin{aligned}
P(A_1\cap\cdots\cap A_n)
={}&P(A_1)P(A_2\mid A_1)P(A_3\mid A_1\cap A_2)\cdots\\
&\cdot P(A_n\mid A_1\cap\cdots\cap A_{n-1}).
\end{aligned}
\]
\end{mytheorem}
\origin{原笔记第1.3节，定理“条件概率的乘法公式”；源码第1220行。}

\begin{mytheorem}{全概率公式}{total-probability-formula}
设 $A_1,A_2,\ldots$ 构成样本空间 $\Omega$ 的一个可列分割，即
\[
A_i\cap A_j=\varnothing\quad(i\ne j),
\qquad
\bigcup_{i=1}^{\infty}A_i=\Omega,
\]
并且 $P(A_i)>0$。则对任意 $B\in\mathscr{F}$，
\[
P(B)=\sum_{i=1}^{\infty}P(A_i)P(B\mid A_i).
\]
当分割只有有限个事件时，相应地使用有限求和。
\end{mytheorem}
\origin{原笔记第1.3节，定理“全概率公式”；源码第1254行。}

\subsection{例题：同卵双胞胎的比例}\label{note:c1:e16}
\paragraph{题目}
设 $A$ 表示“双胞胎为同卵”，$B$ 表示“双胞胎性别相同”。已知同卵双胞胎性别必然相同，异卵双胞胎两人的性别可视为独立且各以 $1/2$ 的概率为男或女，所以
\[
P(B\mid A)=1,
\qquad
P(B\mid A^{\mathrm c})=\frac12.
\]
总体中双胞胎性别相同的比例为 $P(B)=0.64$，求同卵双胞胎所占比例 $P(A)$。

\basis{原笔记第1.3节，例“同卵双胞胎的比例”；源码\texttt{chapters/chapter1.tex}第1283行。}
\begin{proof}[解答]
若总体中双胞胎性别相同的比例为 $P(B)=0.64$，记 $P(A)=x$。由全概率公式，
\[
0.64=x+\frac12(1-x),
\]
解得
\[
x=0.28.
\]
因此同卵双胞胎约占 $28\%$。
\end{proof}

\subsection{例题：Pólya 罐模型}\label{note:c1:e17}
\paragraph{题目}
罐中原有 $a$ 个红球与 $b$ 个蓝球，其中 $a,b$ 为正整数。每次随机取出一个球，放回后再加入 $c$ 个同色球，其中 $c$ 为非负整数。证明：任意第 $k$ 次取到红球的概率都等于
\[
\frac{a}{a+b}.
\]

\basis{原笔记第1.3节，例“Pólya 罐模型”；源码\texttt{chapters/chapter1.tex}第1301行。}
\begin{proof}
对所有允许的初始红、蓝球数同时作关于 $k$ 的归纳。$k=1$ 时，红球占全部球的比例为 $a/(a+b)$，结论成立。假设对任意初始红、蓝球数，结论对第 $k-1$ 次抽取成立。

以第一次抽取的颜色作分割：若第一次取到红球，则此后罐中有 $a+c$ 个红球与 $b$ 个蓝球；若第一次取到蓝球，则有 $a$ 个红球与 $b+c$ 个蓝球。

由归纳假设和全概率公式，
\[
\begin{aligned}
P(\text{第 $k$ 次为红球})
&=\frac{a}{a+b}\cdot\frac{a+c}{a+b+c}
+\frac{b}{a+b}\cdot\frac{a}{a+b+c}\\
&=\frac{a}{a+b}.
\end{aligned}
\]
故结论对所有 $k\geq1$ 成立。
\end{proof}

\subsection{例题：\texorpdfstring{配对问题的递推与恰有 $k$ 个配对}{配对问题的递推与恰有 k 个配对}}\label{note:c1:e18}
\paragraph{题目}
将 $n$ 封信等可能地一一放入写有相应地址的 $n$ 个信封中。令 $E_n$ 表示没有一封信放入正确信封，记
\[q_n=P(E_n),\qquad q_0=1,\quad q_1=0.\]
\begin{enumerate}
\item 证明当 $n\geq2$ 时
\[q_n=\frac1n q_{n-2}+\frac{n-1}{n}q_{n-1},\]
并据此求 $q_n$。
\item 求恰有 $k$ 封信放入正确信封的概率，其中 $0\leq k\leq n$。
\end{enumerate}

\basis{原笔记第1.3节，例“配对问题的递推与恰有 $k$ 个配对”；源码\texttt{chapters/chapter1.tex}第1321行。}
\begin{proof}[解答]
\begin{enumerate}
\item 记 $D_n=n!q_n$ 为 $n$ 个元素的错排数。第一封信只能放入其余 $n-1$ 个信封。固定其中一个信封 $j$，将它对应的信记为第 $j$ 封信。

若第 $j$ 封信放入第一封信的信封，则这两封信互换，余下 $n-2$ 封信须全部错排，有 $D_{n-2}$ 种。

若第 $j$ 封信不放入第一封信的信封，则删去第一封信及信封 $j$，并把第一封信的信封改记为信封 $j$。剩余放法与 $n-1$ 封信的错排一一对应：第 $j$ 封信不在改记后的信封 $j$，其他信也不在各自信封；反过来可恢复原放法。因此这一类有 $D_{n-1}$ 种。

两类互斥且穷尽所有错排，所以
\[D_n=(n-1)(D_{n-1}+D_{n-2}).\]
除以 $n!$ 得
\[q_n=\frac1n q_{n-2}+\frac{n-1}{n}q_{n-1}.\]
移项并用 $q_1-q_0=-1$，得到
\begin{align*}
q_n-q_{n-1}&=-\frac1n(q_{n-1}-q_{n-2})=\frac{(-1)^n}{n!},\\
q_n&=q_0+\sum_{r=1}^n(q_r-q_{r-1})=\sum_{r=0}^n\frac{(-1)^r}{r!}.
\end{align*}
\item 先选出放置正确的 $k$ 封信，有 $\binom nk$ 种选择，再将其余 $n-k$ 封信完全错排，有 $D_{n-k}=(n-k)!q_{n-k}$ 种。总放法有 $n!$ 种，因此
\[P(\text{恰有 $k$ 封信放置正确})=\binom nk\frac{(n-k)!q_{n-k}}{n!}=\frac1{k!}\sum_{r=0}^{n-k}\frac{(-1)^r}{r!}.\]
\end{enumerate}
\end{proof}

\begin{mytheorem}{贝叶斯公式}{bayes-formula}
设 $A_1,A_2,\ldots$ 构成 $\Omega$ 的一个可列分割，且 $P(A_i)>0$。若 $B\in\mathscr{F}$ 且 $P(B)>0$，则
\[
P(A_i\mid B)
=\frac{P(A_i)P(B\mid A_i)}
{\displaystyle\sum_{k=1}^{\infty}P(A_k)P(B\mid A_k)}.
\]
有限分割的情形只需将分母改为有限求和。
\end{mytheorem}
\origin{原笔记第1.3节，定理“贝叶斯公式”；源码第1362行。}

\subsection{例题：保险事故原因}\label{note:c1:e19}
\paragraph{题目}
设 $A$ 表示“事故较容易发生”，$B$ 表示“发生事故”。已知
\[
P(A)=0.3,
\qquad
P(B\mid A)=0.4,
\qquad
P(B\mid A^{\mathrm c})=0.2.
\]
求：
\begin{enumerate}
\item 发生事故的概率 $P(B)$。
\item 已发生事故时，属于事故较容易发生这一情形的条件概率 $P(A\mid B)$。
\end{enumerate}

\basis{原笔记第1.3节，例“保险事故原因”；源码\texttt{chapters/chapter1.tex}第1384行。}
\begin{proof}[解答]
\begin{enumerate}
\item 由全概率公式，
\[
P(B)=0.3\times0.4+0.7\times0.2=0.26.
\]
\item 再由贝叶斯公式，
\[
P(A\mid B)
=\frac{0.3\times0.4}{0.26}
=\frac6{13}
\approx0.4615.
\]
\end{enumerate}
\end{proof}

\subsection{例题：Monty Hall 问题}\label{note:c1:e20}
\paragraph{题目}
汽车以相同概率位于三扇门中的任意一扇门后；三扇门后分别藏有一辆汽车和两只山羊。参赛者先选择第 $1$ 扇门；主持人知道汽车的位置，随后一定在第 $2,3$ 扇门中打开一扇藏有山羊的门。若两扇门都可打开，主持人以相同概率选择其中一扇。已知主持人打开了第 $3$ 扇门，求不换门与换门的获胜概率，并判断是否应改选第 $2$ 扇门。

\basis{原笔记第1.3节，例“Monty Hall 问题”；源码\texttt{chapters/chapter1.tex}第1406行。}
\begin{proof}[解答]
令 $A_i$ 表示“汽车在第 $i$ 扇门后”，$B$ 表示“主持人打开第 $3$ 扇门”。则
\[
P(A_1)=P(A_2)=P(A_3)=\frac13,
\]
且
\[
P(B\mid A_1)=\frac12,
\qquad
P(B\mid A_2)=1,
\qquad
P(B\mid A_3)=0.
\]
因此
\[
P(B)=\frac13\cdot\frac12+\frac13\cdot1+\frac13\cdot0=\frac12.
\]
由贝叶斯公式，
\[
P(A_1\mid B)=\frac13,
\qquad
P(A_2\mid B)=\frac23.
\]
所以改选第 $2$ 扇门的获胜概率为 $2/3$，应当换门。
\end{proof}

\begin{mytheorem}{条件概率构成概率测度}{conditional-probability-measure}
固定 $B\in\mathscr{F}$ 且 $P(B)>0$，定义
\[
P_B(A)=P(A\mid B),\qquad A\in\mathscr{F}.
\]
则 $P_B$ 是可测空间 $(\Omega,\mathscr{F})$ 上的一个概率测度。
\end{mytheorem}
\origin{原笔记第1.3节，定理“条件概率构成概率测度”；源码第1436行。}

\begin{mydefinition}{条件概率空间中的再次条件化}{iterated-conditioning}
固定 $B\in\mathscr{F}$ 且 $P(B)>0$。若 $C\in\mathscr{F}$ 且 $P_B(C)>0$，定义
\[
P_B(A\mid C)=\frac{P_B(A\cap C)}{P_B(C)}.
\]
\end{mydefinition}
\origin{原笔记第1.3节，定义“条件概率空间中的再次条件化”；源码第1463行。}

\begin{proposition}{再次条件化等价于对交事件条件化}{nested-conditioning}
固定 $B\in\mathscr{F}$ 且 $P(B)>0$。对任意 $C\in\mathscr{F}$，有
\[
P_B(C)>0\quad\Longleftrightarrow\quad P(B\cap C)>0.
\]
在上述条件成立时，对任意 $A\in\mathscr{F}$，
\[
P_B(A\mid C)=P(A\mid B\cap C).
\]
\end{proposition}
\origin{原笔记第1.3节，命题“再次条件化等价于对交事件条件化”；源码第1470行。}

\subsection{例题：血液化验与重复检测}\label{note:c1:e21}
\paragraph{题目}
某项血液化验对患者有 $95\%$ 的概率检出阳性，对健康人有 $1\%$ 的概率出现伪阳性，患者占总人口的 $0.5\%$。记 $A$ 为患病，$B$ 为第一次阳性，$B_2$ 为第二次阳性。
\begin{enumerate}
\item 若一次化验呈阳性，求此人确实患病的条件概率 $P(A\mid B)$。
\item 再次接受同样化验，假设两次化验在给定患病状态后条件独立，即
\[P(B_2\mid A\cap B)=P(B_2\mid A)=0.95,\qquad P(B_2\mid A^{\mathrm c}\cap B)=P(B_2\mid A^{\mathrm c})=0.01.\]
求连续两次呈阳性后的患病概率 $P(A\mid B\cap B_2)$。
\end{enumerate}

\basis{原笔记第1.3节，例“血液化验与重复检测”；源码\texttt{chapters/chapter1.tex}第1515行。}
\begin{proof}[解答]
\begin{enumerate}
\item 设 $A$ 表示“此人患病”，$B$ 表示“第一次化验呈阳性”。则
\[
P(A)=0.005,
\qquad
P(B\mid A)=0.95,
\qquad
P(B\mid A^{\mathrm c})=0.01.
\]
由贝叶斯公式，
\[
\begin{aligned}
P(A\mid B)
&=\frac{P(A)P(B\mid A)}
{P(A)P(B\mid A)+P(A^{\mathrm c})P(B\mid A^{\mathrm c})}\\
&=\frac{0.005\times0.95}
{0.005\times0.95+(1-0.005)\times0.01}\\
&\approx0.323.
\end{aligned}
\]
因此，即使检验对患者的阳性检出率很高，一次阳性结果所对应的患病概率也只有约 $32.3\%$。原因在于患病率很低，健康人数量很大，伪阳性在全部阳性结果中占据了不可忽略的比例。
\item 由题设的条件独立性，连续两次呈阳性后，
\[
\begin{aligned}
P(A\mid B\cap B_2)
&=\frac{P(A)P(B\mid A)P(B_2\mid A)}
{P(A)P(B\mid A)P(B_2\mid A)
+P(A^{\mathrm c})P(B\mid A^{\mathrm c})P(B_2\mid A^{\mathrm c})}\\
&=\frac{0.005\times0.95^2}
{0.005\times0.95^2+0.995\times0.01^2}\\
&\approx0.978.
\end{aligned}
\]
连续两次阳性使后验患病概率提高到约 $97.8\%$。
\end{enumerate}
\end{proof}

\subsection{例题：飞机失事搜索}\label{note:c1:e22}
\paragraph{题目}
一架飞机失踪，并且等可能地落在第 $1,2,3$ 号区域。设 $B_i$ 表示“飞机落在第 $i$ 个区域”，$A_i$ 表示“搜索第 $i$ 个区域后没有找到飞机”，并设
\[
P(A_i\mid B_i)=\beta_i,
\qquad 0<\beta_i<1,
\qquad i=1,2,3.
\]
若飞机不在被搜索的区域，则在那里搜索必然找不到飞机。

\begin{enumerate}
\item 只搜索第 $1$ 号区域且未找到飞机，求 $P(B_i\mid A_1)$（$i=1,2,3$）。
\item 随后搜索第 $2$ 号区域仍未找到飞机，求 $P(B_i\mid A_1\cap A_2)$（$i=1,2,3$）。
\item 最后搜索第 $3$ 号区域仍未找到飞机，求 $P(B_i\mid A_1\cap A_2\cap A_3)$（$i=1,2,3$）。
\end{enumerate}

\basis{原笔记第1.3节，例“飞机失事搜索”；源码\texttt{chapters/chapter1.tex}第1560行。}
\begin{proof}[解答]
$B_1,B_2,B_3$ 构成完备事件组，且 $P(B_i)=1/3$。若给定 $B_i$，在其他区域搜索必然失败。因此计算多个搜索失败的交事件时，只留下在真实所在区域漏检的条件；不需要假设各次搜索相互独立。
\begin{enumerate}
\item 由 $P(A_1\mid B_1)=\beta_1$、$P(A_1\mid B_2)=P(A_1\mid B_3)=1$，全概率公式给出
\[P(A_1)=\frac{\beta_1+2}{3}.\]
用贝叶斯公式分别计算：
\begin{align*}
P(B_1\mid A_1)&=\frac{(1/3)\beta_1}{(\beta_1+2)/3}=\frac{\beta_1}{\beta_1+2},\\
P(B_2\mid A_1)&=P(B_3\mid A_1)=\frac{(1/3)\cdot1}{(\beta_1+2)/3}=\frac1{\beta_1+2}.
\end{align*}
\item 给定 $B_1$，事件 $A_2$ 必然发生，故 $P(A_1\cap A_2\mid B_1)=\beta_1$；给定 $B_2$，事件 $A_1$ 必然发生，故该条件概率为 $\beta_2$；给定 $B_3$，两次搜索均必然失败，故该条件概率为 $1$。于是
\[P(A_1\cap A_2)=\frac{\beta_1+\beta_2+1}{3}.\]
由贝叶斯公式，
\begin{align*}
P(B_1\mid A_1\cap A_2)&=\frac{(1/3)\beta_1}{(\beta_1+\beta_2+1)/3}=\frac{\beta_1}{\beta_1+\beta_2+1},\\
P(B_2\mid A_1\cap A_2)&=\frac{\beta_2}{\beta_1+\beta_2+1},\\
P(B_3\mid A_1\cap A_2)&=\frac1{\beta_1+\beta_2+1}.
\end{align*}
\item 给定 $B_i$，另外两个区域的搜索失败是必然事件，故
\[P(A_1\cap A_2\cap A_3\mid B_i)=P(A_i\mid B_i)=\beta_i.\]
因此
\[P(A_1\cap A_2\cap A_3)=\frac{\beta_1+\beta_2+\beta_3}{3},\]
进而
\[P(B_i\mid A_1\cap A_2\cap A_3)=\frac{(1/3)\beta_i}{(\beta_1+\beta_2+\beta_3)/3}=\frac{\beta_i}{\beta_1+\beta_2+\beta_3},\quad i=1,2,3.\]
每次搜索失败后，都可将当前后验概率作为下一次更新的先验概率。
\end{enumerate}
\end{proof}

\subsection{例题：赌徒破产模型}\label{note:c1:e23}
\paragraph{题目}
设 $a$ 为正整数，质点从整数点 $i\in\{0,1,\ldots,a\}$ 出发。只要质点位于 $1,2,\ldots,a-1$ 中，它就以概率 $p$ 向右移动一步，以概率 $q=1-p$ 向左移动一步，其中 $0<p<1$；一旦到达 $0$ 或 $a$，便停留在那里。

\begin{enumerate}
\item 分别求质点最终到达 $0$ 与到达 $a$ 的概率，并说明它是否以概率 $1$ 被端点吸收。
\item 将位置解释为甲的赌本：甲、乙初始赌本为 $i$ 元和 $a-i$ 元，每局胜者从败者处获得 $1$ 元，甲每局获胜概率为 $p$。固定 $i$，令 $a\to\infty$，求甲最终破产概率的极限。
\end{enumerate}

\basis{原笔记第1.3节，例“赌徒破产模型”；源码\texttt{chapters/chapter1.tex}第1594行。}
\begin{proof}[解答]
\begin{enumerate}
\item 设 $E_i$ 表示“质点从 $i$ 出发”，$F$ 表示“最终到达 $0$”，并记
\[
Q_i=P(F\mid E_i).
\]
按第一次移动的方向使用全概率公式，得到
\[
Q_i=pQ_{i+1}+qQ_{i-1},
\qquad 1\leq i\leq a-1,
\]
边界条件为
\[
Q_0=1,
\qquad
Q_a=0.
\]
递推式可改写为
\[
Q_{i+1}-Q_i=\frac{q}{p}(Q_i-Q_{i-1}).
\]

令 $d_j=Q_j-Q_{j-1}$。当 $p=q=1/2$ 时，$d_j=d_1$，故
\[Q_i=1+id_1,\qquad 0=Q_a=1+ad_1.\]
由此得到
\[
Q_i=\frac{a-i}{a}.
\]
当 $p\ne q$ 时，令 $\rho=q/p\ne1$，则 $d_j=\rho^{j-1}d_1$。对差分求和，
\begin{align*}
Q_i&=1+d_1\sum_{j=1}^i\rho^{j-1}=1+d_1\frac{1-\rho^i}{1-\rho},\\
0=Q_a&=1+d_1\frac{1-\rho^a}{1-\rho}.
\end{align*}
消去 $d_1$，得到
\[
Q_i
=\frac{\left(\dfrac qp\right)^i-\left(\dfrac qp\right)^a}
{1-\left(\dfrac qp\right)^a}.
\]
因此
\[
P(F\mid E_i)=
\begin{cases}
\dfrac{\left(\dfrac qp\right)^i-\left(\dfrac qp\right)^a}
{1-\left(\dfrac qp\right)^a},&p\ne q,\\[1.2em]
\dfrac{a-i}{a},&p=q=\dfrac12.
\end{cases}
\]

同理，若 $M$ 表示“最终到达 $a$”，则
\[
P(M\mid E_i)=
\begin{cases}
\dfrac{1-\left(\dfrac qp\right)^i}
{1-\left(\dfrac qp\right)^a},&p\ne q,\\[1.2em]
\dfrac{i}{a},&p=q=\dfrac12.
\end{cases}
\]
两式相加得到
\[
P(F\cup M\mid E_i)=P(F\mid E_i)+P(M\mid E_i)=1,
\]
即质点以概率 $1$ 被某一个端点吸收。
\item 这个模型也可解释为赌徒破产问题：甲、乙的初始赌本分别为 $i$ 元和 $a-i$ 元，每局胜者从败者处获得 $1$ 元，甲每局获胜的概率为 $p$。固定甲的初始赌本 $i$，令乙的初始赌本趋于无穷，即令 $a\to\infty$，则
\[
\lim_{a\to\infty}P(F\mid E_i)=
\begin{cases}
\left(\dfrac qp\right)^i,&p>\dfrac12,\\[0.8em]
1,&p\leq\dfrac12.
\end{cases}
\]
因此，当 $i>0$ 且 $p>1/2$ 时，甲有正概率永不破产；当 $p\leq1/2$ 时，面对拥有无限赌资的对手，甲以概率 $1$ 破产。
\end{enumerate}
\end{proof}

\clearpage\section{事件的独立性和试验的独立性}\label{k:independence}

\begin{mydefinition}{两个事件的独立性}{independence-two-events}
设 $(\Omega,\mathscr{F},P)$ 为概率空间，$A,B\in\mathscr{F}$。若
\[
P(A\cap B)=P(A)P(B),
\]
则称事件 $A$ 与 $B$\textbf{相互独立}。
\end{mydefinition}
\origin{原笔记第1.4节，定义“两个事件的独立性”；源码第1675行。}

\begin{proposition}{独立性对补事件保持}{independence-complements}
若事件 $A$ 与 $B$ 相互独立，则下面三对事件也分别相互独立：
\[
A\text{ 与 }B^{\mathrm c},
\qquad
A^{\mathrm c}\text{ 与 }B,
\qquad
A^{\mathrm c}\text{ 与 }B^{\mathrm c}.
\]
\end{proposition}
\origin{原笔记第1.4节，命题“独立性对补事件保持”；源码第1691行。}

\subsection{例题：独立与互斥的区别}\label{note:c1:e24}
\paragraph{题目}
向区间 $[0,1]$ 内均匀随机投入一个粒子，以下区间表示粒子落入其中的事件。分别判断每对事件是否互斥、是否相互独立：
\begin{enumerate}
\item \[
A_1=\left[0,\frac12\right],
\qquad
B_1=\left[\frac14,\frac34\right].
\]
\item \[
A_2=\left[\frac15,\frac13\right],
\qquad
B_2=\left[\frac14,1\right].
\]
\end{enumerate}

\basis{原笔记第1.4节，例“独立与互斥的区别”；源码\texttt{chapters/chapter1.tex}第1726行。}
\begin{proof}[解答]
\begin{enumerate}
\item 有
\[
A_1\cap B_1=\left[\frac14,\frac12\right],
\]
并且
\[
P(A_1)=\frac12,
\qquad
P(B_1)=\frac12,
\qquad
P(A_1\cap B_1)=\frac14=P(A_1)P(B_1).
\]
所以 $A_1$ 与 $B_1$ 不互斥，但相互独立。
\item 有
\[
A_2\cap B_2=\left[\frac14,\frac13\right],
\]
且
\[
P(A_2)=\frac{2}{15},
\qquad
P(B_2)=\frac34,
\qquad
P(A_2\cap B_2)=\frac1{12}\ne\frac1{10}=P(A_2)P(B_2).
\]
所以 $A_2$ 与 $B_2$ 既不互斥，也不相互独立。
\end{enumerate}
\end{proof}

\begin{mydefinition}{两两独立与相互独立}{pairwise-mutual-independence}
	设 $A_1,A_2,\ldots,A_n$ 为 $n\geq2$ 个事件。

	若对任意 $i\ne j$ 都有
	\[
	P(A_i\cap A_j)=P(A_i)P(A_j),
	\]
	则称 $A_1,A_2,\ldots,A_n$\textbf{两两独立}。

	若对任意 $k\in\{2,3,\ldots,n\}$ 以及任意不同的指标
	\[
	1\leq i_1<i_2<\cdots<i_k\leq n,
	\]
	都有
	\[
	P(A_{i_1}\cap A_{i_2}\cap\cdots\cap A_{i_k})
	=\prod_{r=1}^{k}P(A_{i_r}),
	\]
	则称 $A_1,A_2,\ldots,A_n$\textbf{相互独立}。
\end{mydefinition}
\origin{原笔记第1.4节，定义“两两独立与相互独立”；源码第1769行。}

\subsection{例题：两两独立不推出相互独立}\label{note:c1:e25}
\paragraph{题目}
独立地掷两枚均匀骰子。令 $E$ 表示“两枚骰子的点数之和为 $7$”，$F$ 表示“第一枚骰子的点数为 $4$”，$G$ 表示“第二枚骰子的点数为 $3$”。判断 $E,F,G$ 是否两两独立、是否相互独立。

\basis{原笔记第1.4节，例“两两独立不推出相互独立”；源码\texttt{chapters/chapter1.tex}第1791行。}
\begin{proof}[解答]
有
\[
P(E)=P(F)=P(G)=\frac16,
\]
并且
\[
P(E\cap F)=P(E\cap G)=P(F\cap G)=P(E\cap F\cap G)=\frac1{36}.
\]
因此
\[
P(E\cap F)=P(E)P(F),
\qquad
P(E\cap G)=P(E)P(G),
\qquad
P(F\cap G)=P(F)P(G),
\]
所以 $E,F,G$ 两两独立。但是
\[
P(E\cap F\cap G)=\frac1{36}
\ne\frac1{216}=P(E)P(F)P(G),
\]
故它们并不相互独立。
\end{proof}

\subsection{例题：三色球模型}\label{note:c1:e26}
\paragraph{题目}
袋中有 $4$ 个球：一个只涂白色，一个只涂蓝色，一个只涂红色，另一个同时涂有白、蓝、红三色。等可能地取出一个球。设 $A,B,C$ 分别表示“取出的球涂有白色、蓝色、红色”。判断 $A,B,C$ 是否两两独立、是否相互独立。

\basis{原笔记第1.4节，例“三色球模型”；源码\texttt{chapters/chapter1.tex}第1816行。}
\begin{proof}[解答]
有
\[
P(A)=P(B)=P(C)=\frac12,
\]
且
\[
P(A\cap B)=P(A\cap C)=P(B\cap C)=P(A\cap B\cap C)=\frac14.
\]
因此
\[
P(A\cap B)=P(A)P(B),
\quad
P(A\cap C)=P(A)P(C),
\quad
P(B\cap C)=P(B)P(C),
\]
所以 $A,B,C$ 两两独立；但
\[
P(A\cap B\cap C)=\frac14
\ne\frac18=P(A)P(B)P(C),
\]
故它们不相互独立。
\end{proof}

\begin{mydefinition}{可列个事件的相互独立}{countable-event-independence}
设 $A_1,A_2,\ldots$ 为可列个事件。若其中任意有限多个不同事件都相互独立，则称事件列 $\{A_n\}_{n\geq1}$\textbf{相互独立}。
\end{mydefinition}
\origin{原笔记第1.4节，定义“可列个事件的相互独立”；源码第1847行。}

\begin{mydefinition}{独立试验、伯努利试验与独立重复试验}{independent-trials}
设一个试验由一系列子试验组成。
\begin{enumerate}
    \item 若分别由不同子试验决定的任意有限个事件相互独立，则称这些子试验\textbf{相互独立}。
    \item 若一个试验由 $n$ 个相互独立且同分布的子试验组成，则称其为\textbf{$n$ 重独立试验}。若每个子试验只有“成功”和“失败”两种结果，且成功概率为 $p$，$0<p<1$，则称其为\textbf{$n$ 重伯努利试验}。
    \item 若一个试验由可列个相互独立且同分布的子试验组成，则称其为\textbf{独立重复试验}。
\end{enumerate}
\end{mydefinition}
\origin{原笔记第1.4节，定义“独立试验、伯努利试验与独立重复试验”；源码第1853行。}

\subsection{例题：伯努利试验中的常用计算}\label{note:c1:e27}
\paragraph{题目}
进行独立重复伯努利试验，每次成功的概率为 $p$，失败的概率为 $q=1-p$。
其中 $0<p<1$。求：
\begin{enumerate}
\item 前 $n$ 次中恰好成功 $k$ 次的概率（$0\leq k\leq n$）。
\item 第一次成功发生在第 $k$ 次的概率（$k\geq1$）。
\item 第 $r$ 次成功发生在第 $k$ 次的概率（$1\leq r\leq k$）。
\item 永远不成功及最终至少成功一次的概率。
\item 第 $n$ 次成功出现在第 $m$ 次失败之前的概率（$m,n\geq1$），分别按第 $n$ 次成功前的失败次数、第 $m$ 次失败前的成功次数、前 $m+n-1$ 次成功总数分类，写出三种等价形式。
\end{enumerate}

\basis{原笔记第1.4节，例“伯努利试验中的常用计算”；源码\texttt{chapters/chapter1.tex}第1863行。}
\begin{proof}[解答]
\begin{enumerate}
\item 指定某 $k$ 次成功、其余 $n-k$ 次失败，概率为 $p^kq^{n-k}$。成功位置有 $\binom nk$ 种，各种情况互斥，故概率为
\[\binom nk p^kq^{n-k}.\]
\item 前 $k-1$ 次必须失败，第 $k$ 次成功，由独立性得
\[q^{k-1}p.\]
\item 前 $k-1$ 次中须恰好有 $r-1$ 次成功，第 $k$ 次还须成功，故概率为
\[\binom{k-1}{r-1}p^{r-1}q^{k-r}p=\binom{k-1}{r-1}p^rq^{k-r}.\]
\item 令 $F_n$ 表示前 $n$ 次全部失败，则
\[F_n\downarrow\bigcap_{n=1}^\infty F_n=\{\text{永远不成功}\}.\]
由概率的上连续性，
\[P(\text{永远不成功})=\lim_{n\to\infty}P(F_n)=\lim_{n\to\infty}q^n=0.\]
所以最终至少成功一次的概率为 $1$。
\item 记所求事件为 $E$。若第 $n$ 次成功前恰有 $k$ 次失败，则 $0\leq k\leq m-1$；前 $n+k-1$ 次有 $n-1$ 次成功，下一次成功。因此
\[P(E)=\sum_{k=0}^{m-1}\binom{k+n-1}{n-1}p^nq^k.\]
若按第 $m$ 次失败前的成功次数 $k$ 分类，则 $E$ 等价于 $k\geq n$。第 $m$ 次失败发生在第 $k+m$ 次，故
\[P(E)=\sum_{k=n}^{\infty}\binom{k+m-1}{k}p^kq^m.\]
将成功与失败互换，上一项说明以后至少还会失败一次的概率为 $1$。每次失败后的剩余试验仍为同样的独立重复试验，重复此论证并作有限次归纳，可知第 $m$ 次失败以概率 $1$ 最终发生。因此上述分类遗漏的事件只有零概率。

也可考察前 $m+n-1$ 次试验：$E$ 发生当且仅当其中至少有 $n$ 次成功。若至少有 $n$ 次成功，则失败至多 $m-1$ 次，第 $n$ 次成功已发生而第 $m$ 次失败尚未发生；反之，若成功至多 $n-1$ 次，则已有至少 $m$ 次失败，第 $m$ 次失败先发生。因此
\[P(E)=\sum_{k=n}^{m+n-1}\binom{m+n-1}{k}p^kq^{m+n-1-k}.\]
\end{enumerate}
\end{proof}

\subsection{例题：分赌注问题}\label{note:c1:e28}
\paragraph{题目}
甲、乙约定先胜 $t$ 局者获得全部赌注。比赛在甲已胜 $r$ 局、乙已胜 $s$ 局时中止，其中 $r<t$、$s<t$。以后每局甲获胜的概率为 $p$，且各局相互独立。求甲、乙最终获胜的概率，并说明中止时应如何分配全部赌注。

\basis{原笔记第1.4节，例“分赌注问题”；源码\texttt{chapters/chapter1.tex}第1908行。}
\begin{proof}[解答]
有
\[
\begin{aligned}
P(\text{甲最终先胜 }t\text{ 局})
&=P(t-r\text{ 次成功发生在 }t-s\text{ 次失败之前})\\
&=\sum_{k=0}^{t-s-1}
\binom{k+t-r-1}{t-r-1}p^{t-r}(1-p)^k.
\end{aligned}
\]
合理的分配方式是按甲、乙最终先胜 $t$ 局的概率分配赌注。记上式结果为 $u$，则乙最终先胜 $t$ 局的概率为 $1-u$，甲、乙应分别取得赌注的 $u$ 与 $1-u$ 部分。
\end{proof}

\subsection{例题：\texorpdfstring{事件 $A$ 先于事件 $B$ 发生}{事件 A 先于事件 B 发生}}\label{note:c1:e29}
\paragraph{题目}
进行独立重复试验。每次试验中，互斥事件 $A,B$ 发生的概率分别为 $p_A,p_B$，并设 $p_A+p_B>0$。求 $A$ 先于 $B$ 发生的概率，分别用按首次出现时刻求和和按第一次试验分类两种方法计算。

\basis{原笔记第1.4节，例“事件 $A$ 先于事件 $B$ 发生”；源码\texttt{chapters/chapter1.tex}第1921行。}
\begin{proof}[解答]
$A$ 首次出现且早于 $B$ 发生在第 $n$ 次，当且仅当前 $n-1$ 次二者均未发生，第 $n$ 次发生 $A$。这些事件两两互斥，故

\[
\begin{aligned}
P(A\text{ 先于 }B\text{ 发生})
&=\sum_{n=1}^{\infty}(1-p_A-p_B)^{n-1}p_A\\
&=\frac{p_A}{p_A+p_B}.
\end{aligned}
\]
也可以令所求概率为 $x$。按第一次试验分类，
\[
x=p_A+(1-p_A-p_B)x,
\]
解得同一结果。
\end{proof}

\subsection{例题：由独立性得到欧拉乘积公式}\label{note:c1:e30}
\paragraph{题目}
设 $s>1$，在正整数集合 $\mathbb Z_{>0}$ 上定义
\[
P(A)=\frac{\displaystyle\sum_{k\in A}k^{-s}}
{\displaystyle\sum_{k=1}^{\infty}k^{-s}}.
\]
将所有素数依次记为 $p_1,p_2,\ldots$，令
\[
A_i=\{mp_i:m\in\mathbb Z_{>0}\}.
\]
\begin{enumerate}
\item 证明事件列 $A_1,A_2,\ldots$ 相互独立。
\item 利用概率的连续性证明欧拉乘积公式
\[\sum_{k=1}^{\infty}\frac1{k^s}=\prod_{i=1}^{\infty}\left(1-\frac1{p_i^s}\right)^{-1}.\]
\end{enumerate}

\basis{原笔记第1.4节，例“由独立性得到欧拉乘积公式”；源码\texttt{chapters/chapter1.tex}第1937行。}
\begin{proof}
\begin{enumerate}
\item 由定义，
\[P(A_i)=\frac{\sum_{m=1}^\infty(mp_i)^{-s}}{\sum_{k=1}^\infty k^{-s}}=p_i^{-s}.\]
任取不同的 $i_1,\ldots,i_r$，有
\[
A_{i_1}\cap\cdots\cap A_{i_r}
=\{mp_{i_1}\cdots p_{i_r}:m\in\mathbb Z_{>0}\},
\]
所以
\[
P(A_{i_1}\cap\cdots\cap A_{i_r})
=\frac1{p_{i_1}^s\cdots p_{i_r}^s}
=\prod_{j=1}^{r}P(A_{i_j}).
\]
因此 $A_1,A_2,\ldots$ 相互独立。
\item 记 $C_n=\bigcap_{i=1}^n A_i^{\mathrm c}$，则 $C_n$ 递减。对每个 $n$，由有限个独立事件的补事件仍相互独立，
\[P(C_n)=\prod_{i=1}^n(1-p_i^{-s}).\]
又因为每个大于 $1$ 的正整数至少有一个素因子，
\[
\{1\}=\bigcap_{i=1}^{\infty}A_i^{\mathrm c}.
\]
由概率的上连续性和独立性，
\[
P(\{1\})
=\prod_{i=1}^{\infty}\left(1-\frac1{p_i^s}\right).
\]
另一方面，
\[
P(\{1\})=\left(\sum_{k=1}^{\infty}\frac1{k^s}\right)^{-1}.
\]
故得到欧拉乘积公式
\[
\sum_{k=1}^{\infty}\frac1{k^s}
=\prod_{i=1}^{\infty}\left(1-\frac1{p_i^s}\right)^{-1}.
\]
\end{enumerate}
\end{proof}

\clearpage\section{关于试验独立性的补充}\label{k:product}

设子试验 $1$ 和子试验 $2$ 的概率空间分别为
\[
(\Omega_1,\mathscr F_1,P_1),
\qquad
(\Omega_2,\mathscr F_2,P_2).
\]
两个子试验组成的大试验以
\[
\Omega=\Omega_1\times\Omega_2
\]
为样本空间。子试验 $1$ 中的事件 $A\in\mathscr F_1$ 在大样本空间中表示为 $A\times\Omega_2$；子试验 $2$ 中的事件 $B\in\mathscr F_2$ 表示为 $\Omega_1\times B$。两者同时发生对应
\[
(A\times\Omega_2)\cap(\Omega_1\times B)=A\times B.
\]

\subsection{例题：掷骰子与独立取球}\label{note:c1:e31}
\paragraph{题目}
子试验 $1$ 为掷一枚均匀骰子，$\Omega_1=\{1,2,3,4,5,6\}$；子试验 $2$ 为从红、白、蓝三个球中等概率取一个，$\Omega_2=\{\text{红},\text{白},\text{蓝}\}$。两个子试验相互独立。
\begin{enumerate}
\item 求每个联合结果 $(i,c)\in\Omega_1\times\Omega_2$ 的概率。
\item 对任意 $A\subseteq\Omega_1$、$B\subseteq\Omega_2$，验证 $P(A\times B)=P_1(A)P_2(B)$。
\end{enumerate}

\basis{原笔记第1.5节，例“掷骰子与独立取球”；源码\texttt{chapters/chapter1.tex}第1999行。}
\begin{proof}[解答]
\begin{enumerate}
\item 每个联合结果 $(i,c)\in\Omega_1\times\Omega_2$ 的概率为
\[
P(\{(i,c)\})=\frac16\cdot\frac13=\frac1{18}.
\]
\item $A\times B$ 中共有 $n(A)n(B)$ 个联合结果，每个结果的概率为 $1/18$，故
\[P(A\times B)=\frac{n(A)n(B)}{18}=\frac{n(A)}6\frac{n(B)}3=P_1(A)P_2(B).\]
\end{enumerate}
\end{proof}

\begin{mydefinition}{两个独立子试验的乘积空间表述}{two-independent-experiments-product}
在 $\Omega_1\times\Omega_2$ 上取乘积 $\sigma$ 域
\[
\mathscr F_1\otimes\mathscr F_2
=\sigma\bigl(\{A\times B:A\in\mathscr F_1,\ B\in\mathscr F_2\}\bigr).
\]
若其上的概率 $P$ 对任意 $A\in\mathscr F_1$、$B\in\mathscr F_2$ 都满足
\[
P(A\times B)=P_1(A)P_2(B),
\]
则称这两个子试验相互独立，并将 $P$ 记为 $P_1\otimes P_2$。
\end{mydefinition}
\origin{原笔记第1.5节，定义“两个独立子试验的乘积空间表述”；源码第2010行。}

\begin{mydefinition}{有限个独立子试验}{finite-independent-experiments-product}
设第 $i$ 个子试验的概率空间为 $(\Omega_i,\mathscr F_i,P_i)$，$i=1,\ldots,n$。在
\[
\Omega=\prod_{i=1}^{n}\Omega_i,
\qquad
\mathscr F=\bigotimes_{i=1}^{n}\mathscr F_i
\]
上，若概率 $P$ 对任意 $A_i\in\mathscr F_i$ 都满足
\[
P(A_1\times\cdots\times A_n)
=\prod_{i=1}^{n}P_i(A_i),
\]
则称这 $n$ 个子试验相互独立，并记
\[
P=P_1\otimes\cdots\otimes P_n.
\]
若各个概率空间完全相同，则它们组成 $n$ 重独立试验。
\end{mydefinition}
\origin{原笔记第1.5节，定义“有限个独立子试验”；源码第2034行。}

\subsection{例题：掷骰子后按点数取球}\label{note:c1:e32}
\paragraph{题目}
先掷一枚均匀骰子。若点数为偶数，则从红、白两球中等概率取一个；若点数为奇数，则取蓝球。用
\[
p(\omega_1,\{\omega_2\})
=P(\text{第二步取出 }\omega_2\mid\text{第一步点数为 }\omega_1)
\]
表示转移概率。其中 $\Omega_1=\{1,2,3,4,5,6\}$，$\Omega_2=\{\text{红},\text{白},\text{蓝}\}$。
\begin{enumerate}
\item 列出各点数条件下取出各颜色球的转移概率。
\item 对任意 $A\subseteq\Omega_1$、$B\subseteq\Omega_2$，写出联合概率 $P(A\times B)$ 的表达式。
\item 求取出红球、骰子点数为 $2$、以及点数为 $2$ 且取出红球的概率，并判断两个子试验是否独立。
\end{enumerate}

\basis{原笔记第1.5节，例“掷骰子后按点数取球”；源码\texttt{chapters/chapter1.tex}第2055行。}
\begin{proof}[解答]
\begin{enumerate}
\item 若 $\omega_1=2,4,6$，则
\[
p(\omega_1,\{\text{红}\})=p(\omega_1,\{\text{白}\})=\frac12,
\qquad
p(\omega_1,\{\text{蓝}\})=0;
\]
若 $\omega_1=1,3,5$，则
\[
p(\omega_1,\{\text{红}\})=p(\omega_1,\{\text{白}\})=0,
\qquad
p(\omega_1,\{\text{蓝}\})=1.
\]
\item 由全概率公式，对事件 $A\subseteq\Omega_1$ 和颜色集合 $B\subseteq\Omega_2$，联合概率可写为
\[
P(A\times B)
=\sum_{\omega_1\in A}P_1(\{\omega_1\})
\sum_{\omega_2\in B}p(\omega_1,\{\omega_2\}).
\]
\item 有
\[
P(\text{取出红球})=\frac14,
\qquad
P(\text{骰子点数为 }2)=\frac16,
\]
而
\[
P(\text{点数为 }2\text{ 且取出红球})=\frac1{12}
\ne\frac16\cdot\frac14.
\]
所以两个子试验不独立。
\end{enumerate}
\end{proof}

\clearpage\section*{第一至第五周作业}\addcontentsline{toc}{section}{第一至第五周作业}
% !TEX root = 概率论期中复习.tex
\clearpage\section{第一周：事件运算与古典、几何概型}
\subsection{教材第9--10页第1题}\label{ex:w01-t1}
\paragraph{题目}
设$A,B,C$是三个事件。用它们给出下列事件的表示：
\begin{enumerate}[label=(\arabic*)]
\item 至少有一个事件发生；\item 至少有两个事件发生；\item 所有事件都不发生；\item 至多一个事件发生；\item 至多两个事件发生；\item 恰有两个事件发生。
\end{enumerate}
\basis{原笔记第1.1节“事件的并、交、差与补事件”；见\hyperref[k:events]{事件运算与概率模型}。}
\paragraph{思路}“至少”用并表示，“同时”用交表示；“至多”可以取相应事件的补集。
\begin{proof}[解答]
\begin{enumerate}[label=(\arabic*)]
\item 至少一个发生，即$A,B,C$中存在一个发生，故为$A\cup B\cup C$。
\item 至少两个发生时，三对$AB,AC,BC$中至少有一对同时发生，故为$AB\cup AC\cup BC$。三个都发生也包含在内。
\item 三个都不发生为$A^cB^cC^c=(A\cup B\cup C)^c$。
\item 至多一个发生等价于没有一对同时发生，故为
\[(AB\cup AC\cup BC)^c=A^cB^c\cup A^cC^c\cup B^cC^c.\]
\item 至多两个发生只需排除三个都发生，故为$(ABC)^c=A^c\cup B^c\cup C^c$。
\item 恰有两个发生，须区分不发生的是哪一个，故为
\[ABC^c\cup AB^cC\cup A^cBC=(AB\cup AC\cup BC)\setminus ABC.\]
三项两两互斥，已穷尽所有恰有两个发生的情形。
\end{enumerate}
\end{proof}
\paragraph{答案}依次为$A\cup B\cup C$，$AB\cup AC\cup BC$，$A^cB^cC^c$，$(AB\cup AC\cup BC)^c$，$(ABC)^c$，$ABC^c\cup AB^cC\cup A^cBC$。

\subsection{教材第10页第2题}\label{ex:w01-t2}
\paragraph{题目}
从$0,1,2,3,4,5$中任取4个数排成一行。求恰好构成四位偶数的概率。
\basis{原笔记第1.1节古典概型与计数乘法原理。}
\paragraph{思路}采用有序无放回抽取，按末位是否为0分类；首位不能为0。
\begin{proof}[解答]
等可能排列总数为$6\cdot5\cdot4\cdot3=360$。

末位为0时，首位有5种选法，中间两位依次有4、3种，共$5\cdot4\cdot3=60$种。

末位为2或4时，末位有2种选法；首位不能为0且不能等于末位，有4种选法；中间两位从余下4个数字中有序选取，共$4\cdot3$种。因此
\[P=\frac{60+2\cdot4\cdot4\cdot3}{360}=\frac{13}{30}.\]
\end{proof}
\paragraph{答案}$13/30$。

\subsection{教材第10页第4题}\label{ex:w01-t4}
\paragraph{题目}
信封里有8把钥匙，其中只有一把能打开房门，现取出钥匙试开。分有放回和无放回两种情形，求3次以内能打开房门的概率。
\basis{原笔记第1.1节古典概型；第1.4节独立重复试验。}
\paragraph{思路}有放回取“前三次均失败”的补事件；无放回看正确钥匙在随机排列中的位置。
\begin{proof}[解答]
有放回且每次重新等可能抽取时，各次失败的概率都是$7/8$且相互独立，故
\[P(\text{3次以内打开})=1-\left(\frac78\right)^3=\frac{169}{512}.\]
无放回时，8把钥匙的试用顺序是随机排列。正确钥匙在每个位置的概率均为$1/8$，前三个位置互斥，故所求概率为$3/8$。也可直接计算为
\[1-\frac78\frac67\frac56=\frac38.\]
\end{proof}
\paragraph{答案}有放回$169/512$；无放回$3/8$。

\subsection{教材第10页第5题}\label{ex:w01-t5}
\paragraph{题目}
袋子里有$m$个白球，$n$个黑球（$m,n\geq1$）。
\begin{enumerate}[label=(\arabic*)]
\item 从袋中随机取两个球，求两球同色的概率；
\item 从袋中随机取一个球，记录颜色后放回，从中再取一球，求两次取的球同色的概率；
\item 证明(2)的概率始终大于(1)的概率。
\end{enumerate}
\basis{原笔记第1.1节组合计数；第1.4节独立试验。}
\paragraph{思路}同色分为都白与都黑两种互斥情况，再通分比较。
\begin{proof}[解答]
记$M=m+n\geq2$。
\begin{enumerate}[label=(\arabic*)]
\item 不放回时，所有两球组合等可能，故
\[p_1=\frac{\binom m2+\binom n2}{\binom M2}=\frac{m(m-1)+n(n-1)}{M(M-1)}.\]
当$m=1$或$n=1$时，相应组合数取0。
\item 放回后各次独立，故
\[p_2=\left(\frac mM\right)^2+\left(\frac nM\right)^2=\frac{m^2+n^2}{M^2}.\]
\item 用$m^2+n^2=M^2-2mn$化简，得
\begin{align*}
p_2-p_1&=\frac{(m^2+n^2)(M-1)-M(m^2+n^2-M)}{M^2(M-1)}\\
&=\frac{M^2-m^2-n^2}{M^2(M-1)}=\frac{2mn}{M^2(M-1)}>0.
\end{align*}
分母为正，且$m,n\geq1$，所以严格大于。
\end{enumerate}
\end{proof}
\paragraph{答案}$p_1=[m(m-1)+n(n-1)]/[M(M-1)]$，$p_2=(m^2+n^2)/M^2$，且$p_2-p_1=2mn/[M^2(M-1)]>0$。

\subsection{教材第10页第6题}\label{ex:w01-t6}
\paragraph{题目}
在区间$(0,1)$中随机地取出三个数，以它们作为三条线段的长度。求这些线段能连接成一个三角形的概率。
\basis{原笔记第1.1节几何概率模型；按三个数相互独立且均匀抽取理解。}
\paragraph{思路}在单位立方体内去掉一条边不小于另外两边之和的区域。
\begin{proof}[解答]
记三个数为$x,y,z$。组成非退化三角形当且仅当
\[x<y+z,\quad y<x+z,\quad z<x+y.\]
除去零体积的边界，三个失败区域互不相交，而且由对称性体积相等。其中$x\geq y+z$的体积为
\begin{align*}
\int_0^1\int_0^x\int_0^{x-y} dz\,dy\,dx
&=\int_0^1\int_0^x(x-y)\,dy\,dx\\
&=\int_0^1\frac{x^2}{2}\,dx=\frac16.
\end{align*}
单位立方体体积为1，故所求概率为$1-3/6=1/2$。
\end{proof}
\paragraph{答案}$1/2$。

\subsection{教材第10页第7题}\label{ex:w01-t7}
\paragraph{题目}
在打好方格的白纸上投掷一枚直径为1的硬币。问方格的边长至少应为多少才能使硬币与方格的边不交的概率小于$1/100$。
\basis{原笔记第1.1节几何概率模型。}
\paragraph{思路}硬币中心到网格线的距离必须都大于半径$1/2$；先按原题求概率条件，再检查“至少”的方向。
\begin{proof}[解答]
设方格边长为$a>0$。假设硬币中心在一个基本方格内均匀分布。当$a\leq1$时，硬币无法完整落在方格内部，所求不交概率为0。当$a>1$时，中心必须位于边长$a-1$的内方形，故
\[p(a)=\frac{(a-1)^2}{a^2}=\left(1-\frac1a\right)^2.\]
于是对$a>1$，
\[p(a)<\frac1{100}\iff 1-\frac1a<\frac1{10}\iff a<\frac{10}{9}.\]
合并得$0<a<10/9$。这给出的是上界，而不是题面“至少”所暗示的下界；题面方向有误，不能给出一个正的最小边长。若原意是“不交概率大于$1/100$”，则应为$a>10/9$；若要求“至少$1/100$”，则为$a\geq10/9$。
\end{proof}
\paragraph{答案}按原题“小于”理解，$0<a<10/9$，不存在正的最小边长。

\subsection{作业第2题：可列德摩根律}\label{ex:w01-h2}
\paragraph{题目}
设$A_i$，$i\geq1$是一列事件。求证：
\[\left(\bigcup_{i=1}^{\infty}A_i\right)^c=\bigcap_{i=1}^{\infty}A_i^c,\qquad
\left(\bigcap_{i=1}^{\infty}A_i\right)^c=\bigcup_{i=1}^{\infty}A_i^c.\]
\basis{原笔记第1.1节事件运算与德摩根律；此处是作业要求证明的结论，故展开证明。}
\paragraph{思路}逐点把集合所属关系转为量词。
\begin{proof}
任取$\omega\in\Omega$，有
\begin{align*}
\omega\in\left(\bigcup_{i=1}^{\infty}A_i\right)^c
&\iff \text{不存在$i\geq1$使$\omega\in A_i$}\\
&\iff \text{对每个$i\geq1$，$\omega\in A_i^c$}\\
&\iff \omega\in\bigcap_{i=1}^{\infty}A_i^c.
\end{align*}
两集合含有相同的样本点，第一式得证。同理，
\begin{align*}
\omega\in\left(\bigcap_{i=1}^{\infty}A_i\right)^c
&\iff \text{并非对所有$i\geq1$都有$\omega\in A_i$}\\
&\iff \text{存在$i\geq1$使$\omega\in A_i^c$}\\
&\iff \omega\in\bigcup_{i=1}^{\infty}A_i^c.
\end{align*}
故第二式也成立。
\end{proof}
\paragraph{答案}两式均成立。

\subsection{作业第3题：恰有两个事件发生}\label{ex:w01-h3}
\paragraph{题目}
设$A,B,C$是三个事件。求证：
\[ABC^c\cup AB^cC\cup A^cBC=(AB\cup AC\cup BC)\setminus ABC.\]
\basis{原笔记第1.1节事件运算；与\hyperref[ex:w01-t1]{本周教材第1题(6)}相联系。}
\paragraph{思路}右侧是在“至少两个发生”中排除“三个都发生”。
\begin{proof}
利用分配律与德摩根律，
\begin{align*}
(AB\cup AC\cup BC)\setminus ABC
&=(AB\cap(ABC)^c)\cup(AC\cap(ABC)^c)\cup(BC\cap(ABC)^c)\\
&=ABC^c\cup AB^cC\cup A^cBC.
\end{align*}
第二步中，例如$AB\cap(ABC)^c=AB\cap(A^c\cup B^c\cup C^c)=ABC^c$，因为$AB$分别与$A^c,B^c$的交为空集。另两项同理，故等式成立。
\end{proof}
\paragraph{答案}两边均表示恰有两个事件发生。

\subsection{作业第4项：绳子连接问题（阅读）}\label{ex:w01-h4}
\paragraph{原作业要求}
阅读：自己去查找一些关于“$n$根绳子，$2n$个端头任意连接”的相关类型的古典概型的题目和解答。（自己找来看明白即可，不需要上交。）
\paragraph{整理说明}原作业这一项是开放阅读要求，没有指定待求事件。保留原项以便核对作业范围，不将任意选定的连接问题冒充原题。
% !TEX root = 概率论期中复习.tex
\clearpage\section{第二周：概率空间、概率界与事件极限}
\subsection{教材第19页第1题}\label{ex:w02-t1}
\paragraph{题目}
假设某个随机试验包含可数无穷多个不同的结果，且每个结果出现的概率均严格大于零。解释为何不会所有结果发生的概率都相同。
\basis{原笔记第1.2节概率的可列可加性。}
\paragraph{思路}若各个结果具有同一个正概率，有限多个互斥结果的总概率也将超过1。
\begin{proof}[解答]
把这些结果列为$\omega_1,\omega_2,\ldots$，单点事件两两互斥。若所有单点的概率同为$p>0$，则对任意正整数$N$，
\[P(\{\omega_1,\ldots,\omega_N\})=Np\leq1.\]
取$N>1/p$即矛盾。因此正概率不能全部相同。若试图令共同值为0，可列可加性又给出$P(\Omega)=\sum_{i=1}^\infty0=0$，也不可能。
\end{proof}
\paragraph{答案}不存在可数无穷等可能且每个结果有正概率的模型。

\subsection{教材第19页第2题(1)}\label{ex:w02-t2}
\paragraph{题目}
设$(\Omega,\mathscr F)$为可测空间。给定非空集$E\subseteq\Omega$，记$\mathscr E=\{E\cap A:A\in\mathscr F\}$。证明$\mathscr E$是$E$上的$\sigma$代数。（按本周要求不做第(2)小题。）
\basis{原笔记第1.2节“可测空间”中$\sigma$代数的定义。}
\paragraph{思路}以$E$为新的全集，分别验证包含全集、相对补封闭和可列并封闭。
\begin{proof}
因为$\Omega\in\mathscr F$，有$E=E\cap\Omega\in\mathscr E$。

若$B\in\mathscr E$，则存在$A\in\mathscr F$使$B=E\cap A$。在新全集$E$中，$B$的补集为
\[E\setminus B=E\setminus(E\cap A)=E\cap A^c\in\mathscr E,\]
这里用了$A^c\in\mathscr F$。

若$B_i\in\mathscr E$，分别写成$B_i=E\cap A_i$，其中$A_i\in\mathscr F$，则
\[\bigcup_{i=1}^{\infty}B_i=E\cap\left(\bigcup_{i=1}^{\infty}A_i\right)\in\mathscr E,\]
因为$\bigcup_iA_i\in\mathscr F$。三个条件均满足，故$\mathscr E$是$E$上的$\sigma$代数。
\end{proof}
\paragraph{答案}$\mathscr E$是$E$上的$\sigma$代数，补运算必须相对于$E$进行。

\subsection{教材第19页第4题}\label{ex:w02-t4}
\paragraph{题目}
设$(\Omega,\mathscr F,P)$为概率空间。证明对$A,B\in\mathscr F$有
\[P(A\cap B)\geq P(A)+P(B)-1.\]
\basis{原笔记第1.2节命题“一般加法公式”。}
\paragraph{思路}用并事件的概率不超过1。
\begin{proof}
由加法公式，
\[P(A\cap B)=P(A)+P(B)-P(A\cup B).\]
因为$A\cup B\subseteq\Omega$，单调性给出$P(A\cup B)\leq1$。代入即得
\[P(A\cap B)\geq P(A)+P(B)-1.\]
\end{proof}
\paragraph{答案}所给下界成立；结合非负性还可写成$P(AB)\geq\max\{0,P(A)+P(B)-1\}$。

\subsection{教材第19页第5题}\label{ex:w02-t5}
\paragraph{题目}
设事件$A,B$满足$P(A)=3/4$和$P(B)=1/3$。证明$1/12\leq P(A\cap B)\leq1/3$。
\basis{原笔记第1.2节单调性、加法公式；\hyperref[ex:w02-t4]{上一题}。}
\paragraph{思路}下界用$P(A)+P(B)-1$，上界用$A\cap B\subseteq B$。
\begin{proof}
由上一题，
\[P(A\cap B)\geq\frac34+\frac13-1=\frac1{12}.\]
另一方面，$A\cap B\subseteq B$，所以$P(A\cap B)\leq P(B)=1/3$，合并即得结论。
\end{proof}
\paragraph{答案}$P(A\cap B)\in[1/12,1/3]$。

\subsection{作业第2题：偏离独立乘积的界}\label{ex:w02-h2}
\paragraph{题目}
求证：对任意事件$A,B\in\mathscr F$，有
\[|P(AB)-P(A)P(B)|\leq\frac14.\]
提示：不妨设$P(A)\leq P(B)$。注意不等式$ab\leq((a+b)/2)^2$。
\basis{原笔记第1.2节概率的单调性、加法公式。}
\paragraph{思路}按差值的正负分别估计，并按提示设$a=P(A)\leq b=P(B)$。
\begin{proof}
记$q=P(AB)$。先估计正向差值。由$q\leq a$与$a\leq b$，有
\begin{align*}
q-ab&\leq a(1-b)\leq a(1-a)\\
&=\frac14-\left(a-\frac12\right)^2\leq\frac14.
\end{align*}
再估计反向差值。若$a+b\leq1$，由$q\geq0$及提示中的不等式，
\[ab-q\leq ab\leq\left(\frac{a+b}{2}\right)^2\leq\frac14.\]
若$a+b>1$，由$q\geq a+b-1$，得到
\begin{align*}
ab-q&\leq ab-a-b+1=(1-a)(1-b)\\
&\leq\left(\frac{2-a-b}{2}\right)^2<\frac14.
\end{align*}
两种情况穷尽全部可能，故$-1/4\leq q-ab\leq1/4$，结论成立。
\end{proof}
\paragraph{答案}$1/4$为可达到的界，例如$A=B$且$P(A)=1/2$时取等号。

\subsection{作业第3题：开区间的可列表示}\label{ex:w02-h3}
\paragraph{题目}
求证：对任意实数$a<b$，有
\[(a,b)=\bigcup_{k=1}^{\infty}(a,b-1/k].\]
\basis{原笔记第1.2节常用博雷尔集合与阿基米德性质。}
\paragraph{思路}先证每个右侧区间包含于左侧，再对左侧每个点选取足够大的$k$。
\begin{proof}
若$x\in(a,b-1/k]$，则$a<x\leq b-1/k<b$，所以右侧并集包含于$(a,b)$。这里当$b-1/k\leq a$时，相应区间为空集。

反过来，任取$x\in(a,b)$，有$b-x>0$。选择正整数$k$满足$k\geq1/(b-x)$，则$1/k\leq b-x$，从而
\[a<x\leq b-1/k.\]
因此$x$属于右侧某一个区间，故$(a,b)$包含于该并集。两边互相包含，等式成立。
\end{proof}
\paragraph{答案}该可列并表示成立。

\subsection{作业第4题：无限重复试验的事件极限}\label{ex:w02-h4}
\paragraph{题目}
不断重复地做一个试验。设$0<p<1$，$A_n$表示前$n$次试验都不成功，且$P(A_n)=(1-p)^n$，$n\geq1$。
\begin{enumerate}[label=(\arabic*)]
\item 判断$A_n$是单增序列还是单减序列，并说明理由。
\item 说明$\lim_{n\to\infty}A_n$的含义。
\item 令$A$表示不断重复试验中至少成功一次，求$P(A)$。
\end{enumerate}
\basis{原笔记第1.2节事件的单调极限、概率的上连续性；第1.2节例“反复抽取扑克牌中的大王”。}
\paragraph{思路}先用事件的含义判断包含关系，再用概率上连续性。
\begin{proof}[解答]
\begin{enumerate}[label=(\arabic*)]
\item 前$n+1$次都不成功必然意味着前$n$次都不成功，所以$A_{n+1}\subseteq A_n$；这是单减事件列，不能只根据概率大小反推集合包含关系。
\item 单减列的极限为
\[\lim_{n\to\infty}A_n=\bigcap_{n=1}^{\infty}A_n,\]
表示每一次试验都不成功，即永远不成功。
\item 至少成功一次是永远不成功的补事件，因此
\begin{align*}
P(A)&=1-P\left(\bigcap_{n=1}^{\infty}A_n\right)\\
&=1-\lim_{n\to\infty}P(A_n)=1-\lim_{n\to\infty}(1-p)^n=1.
\end{align*}
使用上连续性是因为$A_n$单减且$P(A_1)\leq1<\infty$。概率为1不等于该事件在每个样本点都发生。
\end{enumerate}
\end{proof}
\paragraph{答案}单减；极限为永远不成功；至少成功一次的概率为1。
% !TEX root = 概率论期中复习.tex
\clearpage\section{第三周：可列可加性、容斥与条件概率}
\subsection{教材第19页第3题}\label{ex:w03-t3}
\paragraph{题目}
设$(\Omega_1,\mathscr F_1,\mu_1),\ldots,(\Omega_n,\mathscr F_n,\mu_n)$是$n$个有限测度空间，其中$\Omega_1,\ldots,\Omega_n$两两不相交。令$\Omega=\bigcup_{k=1}^n\Omega_k$。
\begin{enumerate}[label=(\arabic*)]
\item 证明$\mathscr F=\{A\subseteq\Omega:A\cap\Omega_k\in\mathscr F_k,\ 1\leq k\leq n\}$是$\Omega$上的$\sigma$代数；
\item 对$A\in\mathscr F$令$\mu(A)=\sum_{k=1}^n\mu_k(A\cap\Omega_k)$，证明$\mu$是$(\Omega,\mathscr F)$上的有限测度。
\end{enumerate}
\basis{原笔记第1.2节$\sigma$代数与测度的定义。测度是取值于$[0,\infty]$、空集取0且可列可加的集函数；有限测度另要求全集测度有限。}
\paragraph{思路}所有运算都限制到各个互不相交的分块上验证。
\begin{proof}
\begin{enumerate}[label=(\arabic*)]
\item 对每个$k$，$\Omega\cap\Omega_k=\Omega_k\in\mathscr F_k$，故$\Omega\in\mathscr F$。

若$A\in\mathscr F$，则
\[(\Omega\setminus A)\cap\Omega_k=\Omega_k\setminus(A\cap\Omega_k)\in\mathscr F_k,\]
因右侧是$\mathscr F_k$中集合相对于$\Omega_k$的补集。因此$\Omega\setminus A\in\mathscr F$。

若$A_j\in\mathscr F$，则对每个$k$，
\[\left(\bigcup_{j=1}^{\infty}A_j\right)\cap\Omega_k
=\bigcup_{j=1}^{\infty}(A_j\cap\Omega_k)\in\mathscr F_k.\]
所以$\bigcup_jA_j\in\mathscr F$。$\sigma$代数的三个条件均成立。
\item 由各$\mu_k$的非负性，$\mu(A)\geq0$，且$\mu(\varnothing)=0$。若$A_j\in\mathscr F$两两不交，则固定$k$时，$A_j\cap\Omega_k$也两两不交，从而
\begin{align*}
\mu\left(\bigcup_{j=1}^{\infty}A_j\right)
&=\sum_{k=1}^n\mu_k\left(\bigcup_{j=1}^{\infty}(A_j\cap\Omega_k)\right)\\
&=\sum_{k=1}^n\sum_{j=1}^{\infty}\mu_k(A_j\cap\Omega_k)\\
&=\sum_{j=1}^{\infty}\sum_{k=1}^n\mu_k(A_j\cap\Omega_k)
=\sum_{j=1}^{\infty}\mu(A_j).
\end{align*}
交换求和次序可先在$j\leq M$的有限和中进行，再令$M\to\infty$；由于$k$仅有有限个且各项非负，极限与有限和可交换。

最后，
\[\mu(\Omega)=\sum_{k=1}^n\mu_k(\Omega_k)<\infty,\]
因为这是有限个有限数的和。故$\mu$为有限测度。
\end{enumerate}
\end{proof}
\paragraph{答案}两项均成立，且$\mu(\Omega)=\sum_{k=1}^n\mu_k(\Omega_k)$。

\subsection{教材第19页第6题}\label{ex:w03-t6}
\paragraph{题目}
对于随机事件序列$A_1,A_2,\ldots\in\mathscr F$，证明：若对一切$k\geq1$有$P(A_k)=1$，则
\[P\left(\bigcap_{k=1}^{\infty}A_k\right)=1.\]
\basis{原笔记第1.2节可列次可加性、例“可列个零概率事件的并”。}
\paragraph{思路}对交集取补，把满概率事件变为零概率事件。
\begin{proof}
对每个$k$，$P(A_k^c)=1-P(A_k)=0$。由可列次可加性，
\[0\leq P\left(\bigcup_{k=1}^{\infty}A_k^c\right)
\leq\sum_{k=1}^{\infty}P(A_k^c)=0.\]
故补集的概率为0。再由德摩根律，
\[P\left(\bigcap_{k=1}^{\infty}A_k\right)
=1-P\left(\bigcup_{k=1}^{\infty}A_k^c\right)=1.\]
\end{proof}
\paragraph{答案}可列个满概率事件的交仍为满概率事件。

\subsection{教材第24页第1题}\label{ex:w03-t241}
\paragraph{题目}
对于事件$A,B$，证明：若$P(A\mid B)>P(A)$，则$P(B\mid A)>P(B)$。
\basis{原笔记第1.3节条件概率的定义；见\hyperref[k:conditional]{条件概率}。}
\paragraph{思路}两个条件概率不等式都等价于$P(AB)>P(A)P(B)$，但要先确认分母为正。
\begin{proof}
题设$P(A\mid B)$有定义，故$P(B)>0$。由
\[P(A\mid B)=\frac{P(AB)}{P(B)}>P(A),\]
得到$P(AB)>P(A)P(B)\geq0$。因此$P(A)\geq P(AB)>0$，可除以$P(A)$，得到
\[P(B\mid A)=\frac{P(AB)}{P(A)}>P(B).\]
\end{proof}
\paragraph{答案}结论成立，且题设自动保证$P(A)>0$。

\subsection{教材第24页第2题}\label{ex:w03-t242}
\paragraph{题目}
假设有三个事件$A,B,C$。
\begin{enumerate}[label=(\arabic*)]
\item 已知$P(A\mid C)>P(B\mid C)$和$P(A\mid C^c)>P(B\mid C^c)$，是否一定有$P(A)>P(B)$？若是，加以证明；否则举出反例。
\item 已知$P(A\mid C)>P(A\mid C^c)$和$P(B\mid C)>P(B\mid C^c)$，是否一定有$P(A)>P(B)$？若是，加以证明；否则举出反例。
\end{enumerate}
\basis{原笔记第1.3节全概率公式。题中的条件概率要求$0<P(C)<1$。}
\paragraph{思路}第一问在同一个分割内比较两个事件；第二问只比较每个事件在不同条件下的概率，不能决定两个事件的大小。
\begin{proof}[解答]
\begin{enumerate}[label=(\arabic*)]
\item 由全概率公式，
\begin{align*}
P(A)-P(B)
&=P(C)\bigl[P(A\mid C)-P(B\mid C)\bigr]\\
&\quad+P(C^c)\bigl[P(A\mid C^c)-P(B\mid C^c)\bigr]>0.
\end{align*}
两个权重均为正，两个括号也均为正，故一定有$P(A)>P(B)$。
\item 取$\Omega=(0,1)$上的均匀概率，令
\[C=(0,1/2),\quad A=(0,1/5)\cup(1/2,3/5),\quad B=(0,2/5)\cup(1/2,4/5).\]
用区间长度计算得
\[P(A\mid C)=\frac25>\frac15=P(A\mid C^c),\]
\[P(B\mid C)=\frac45>\frac35=P(B\mid C^c).\]
然而$P(A)=3/10<7/10=P(B)$。所有条件均满足而结论不成立，因此第二问的答案是否定的。
\end{enumerate}
\end{proof}
\paragraph{答案}(1)一定成立；(2)不一定，以上反例中$P(A)=3/10<P(B)=7/10$。

\subsection{教材第25页第4题}\label{ex:w03-t254}
\paragraph{题目}
盒子里有15个乒乓球，其中9个是新球。第一场比赛时从中取出3个用完后放回盒中。第二场比赛时再从中取出3个，求这3个都是新球的概率。
\basis{原笔记第1.3节全概率公式与古典计数。}
\paragraph{思路}以第一场取出的新球数作分割；新球使用后变为旧球。也可固定第二场取出的三个球，要求它们原本为新球且第一场未用过。
\begin{proof}[解答]
设$K$是第一场取出的新球数，则$K=0,1,2,3$，且
\[P(K=k)=\frac{\binom9k\binom6{3-k}}{\binom{15}3}.\]
给定$K=k$后，第二场之前剩下$9-k$个新球，总球数仍为15。因此全概率公式给出
\[P=\sum_{k=0}^3\frac{\binom9k\binom6{3-k}}{\binom{15}3}
\frac{\binom{9-k}3}{\binom{15}3}.\]
为简化计算，考察第二场随机取出的球组。它们原本都是新球的概率为$\binom93/\binom{15}3$。给定这个三球组，第一场必须全部从其余12个球中取出，其概率为$\binom{12}3/\binom{15}3$。故
\[P=\frac{\binom93\binom{12}3}{\binom{15}3^2}
=\frac{84\cdot220}{455^2}=\frac{528}{5915}.\]
这也核对了上面的全概率求和。
\end{proof}
\paragraph{答案}$528/5915\approx0.08926$。

\subsection{作业第2题：一排夫妻均不相邻}\label{ex:w03-h2}
\paragraph{题目}
有$n$对夫妇随机地坐成一排，求没有任何妻子坐在她丈夫身边的概率。
\basis{原笔记第1.2节例“圆桌夫妻问题”的捆绑计数与容斥方法；本题为一排座位，不能使用圆排列总数。}
\paragraph{思路}令$A_i$为第$i$对夫妻相邻；指定$r$对相邻后把各对捆为一个整体。
\begin{proof}[解答]
总排法为$(2n)!$，并且各排法等可能。指定$i_1,\ldots,i_r$这$r$对夫妻相邻时，捆绑后共有$2n-r$个对象，各对内部有两种次序，所以
\[P(A_{i_1}\cap\cdots\cap A_{i_r})=\frac{2^r(2n-r)!}{(2n)!}.\]
这样的$r$对可以从$n$对中选出$\binom nr$种。由容斥公式，至少一对夫妻相邻的概率为
\[P\left(\bigcup_{i=1}^nA_i\right)
=\sum_{r=1}^n(-1)^{r+1}\binom nr\frac{2^r(2n-r)!}{(2n)!}.\]
取补，并把常数1记为$r=0$项，所求概率为
\[\boxed{\sum_{r=0}^n(-1)^r\binom nr\frac{2^r(2n-r)!}{(2n)!}}.\]
\end{proof}
\paragraph{答案}$\displaystyle\sum_{r=0}^n(-1)^r\binom nr2^r(2n-r)!/(2n)!$。

\subsection{作业第3题：已知补事件条件概率}\label{ex:w03-h3}
\paragraph{题目}
设$P(A)=0.7$，$P(B\mid\overline A)=0.1$。求$P(A\cup B)$。
\basis{原笔记第1.3节条件概率乘法公式。}
\paragraph{思路}将并事件拆为两个互斥事件$A$与$A^cB$。
\begin{proof}[解答]
因为$A\cup B=A\mathbin{\dot\cup}(A^cB)$，所以
\begin{align*}
P(A\cup B)&=P(A)+P(A^cB)\\
&=P(A)+P(A^c)P(B\mid A^c)\\
&=0.7+0.3\times0.1=0.73.
\end{align*}
\end{proof}
\paragraph{答案}$0.73$。

\subsection{作业第4题：玻璃器皿连续三次未摔坏}\label{ex:w03-h4}
\paragraph{题目}
某玻璃器皿第一次落下时摔坏的概率为0.3。若第一次落下没有摔坏，则第二次落下摔坏的条件概率为0.6；若前两次落下都没有摔坏，则第三次落下摔坏的条件概率为0.9。求该玻璃器皿三次落下都没有摔坏的概率。
\basis{原笔记第1.3节条件概率的乘法公式。}
\paragraph{思路}按存活到每一次的条件顺序相乘，无须假设三次独立。
\begin{proof}[解答]
令$S_i$表示第$i$次落下未摔坏。题设给出
\[P(S_1)=0.7,\quad P(S_2\mid S_1)=0.4,\quad P(S_3\mid S_1S_2)=0.1.\]
由乘法公式，
\[P(S_1S_2S_3)=P(S_1)P(S_2\mid S_1)P(S_3\mid S_1S_2)
=0.7\times0.4\times0.1=0.028.\]
\end{proof}
\paragraph{答案}$0.028$。

\subsection{作业第5题：两种天气的递推}\label{ex:w03-h5}
\paragraph{题目}
只考虑有雨、无雨两种天气。已知今天天气时，明天天气保持不变的条件概率为$p$，$0<p<1$，改变的条件概率为$1-p$。设第一天无雨，求第$n$天也无雨的概率，$n\geq1$。
\basis{原笔记第1.3节全概率公式；原笔记第1.5节“不独立子试验与转移概率”。}
\paragraph{思路}以第$n$天是否无雨作分割，写递推式后减去平衡值$1/2$。
\begin{proof}[解答]
记$q_n=P(\text{第$n$天无雨})$，则$q_1=1$。按题给的一步转移概率与全概率公式，
\begin{align*}
q_{n+1}&=pq_n+(1-p)(1-q_n)\\
&=(2p-1)q_n+1-p.
\end{align*}
因此
\[q_{n+1}-\frac12=(2p-1)\left(q_n-\frac12\right).\]
迭代$n-1$次得到
\[q_n-\frac12=(2p-1)^{n-1}\left(q_1-\frac12\right)=\frac12(2p-1)^{n-1}.\]
所以
\[q_n=\frac{1+(2p-1)^{n-1}}2.\]
$n=1$时将零次幂按1处理，结果为1；$p=1/2$时$n\geq2$的结果为$1/2$。
\end{proof}
\paragraph{答案}$\bigl[1+(2p-1)^{n-1}\bigr]/2$。
% !TEX root = 概率论期中复习.tex
\section{第四周：条件概率与独立性}

\subsection{教材第25页第3题}
\label{ex:w04-t25-3}
\paragraph{题目}
考试中的一道多选题共有5个选项，某同学知道正确答案的概率为$p$，他在不知道答案时猜中答案的概率为$1/5$。
\begin{enumerate}
\item 求该同学做对一道多选题的概率；
\item 设对某道题该同学选对了答案，求他确实知道正确答案的概率。
\end{enumerate}
\textbf{作业订正：}本题为单选题，题面中的“多选题”按“单选题”理解。
\basis{原笔记第1章第3节“全概率公式”“贝叶斯公式”；见\hyperref[k:conditional]{条件概率与贝叶斯公式}。}
\paragraph{思路}按是否知道正确答案分类；已答对后的概率再用贝叶斯公式。
\begin{proof}[解答]
记$K$为“知道正确答案”，$D$为“答对”。知道答案时必能答对，所以
\[
\P(K)=p,\qquad \P(D\mid K)=1,\qquad \P(D\mid K^c)=\frac15.
\]
\begin{enumerate}
\item 由全概率公式，
\[
\P(D)=\P(K)\P(D\mid K)+\P(K^c)\P(D\mid K^c)
=p+\frac{1-p}{5}=\frac{1+4p}{5}.
\]
\item 因$\P(D)>0$，由贝叶斯公式，
\[
\P(K\mid D)=\frac{\P(K)\P(D\mid K)}{\P(D)}=\frac{5p}{1+4p}.
\]
\end{enumerate}
当$p=0$或$p=1$时，直接判断也得到上述结果，不需要对零概率的$K$或$K^c$作条件化。
\end{proof}
\paragraph{答案}$\P(D)=(1+4p)/5$；$\P(K\mid D)=5p/(1+4p)$。

\subsection{教材第25页第5题}
\label{ex:w04-t25-5}
\paragraph{题目}
在信号传递中发送信号“0”或“1”占比分别为$1/3$和$2/3$。由于外界干扰等原因，接收到的信号可能出现错误。已知出错的概率为$0.05$。设接收到的信号为“1”，求发送的信号确实为“1”的概率。
\basis{原笔记第1章第3节“贝叶斯公式”；见\hyperref[k:conditional]{条件概率与贝叶斯公式}。}
\paragraph{思路}区分发送信号与接收信号，写出接收到“1”的两条途径。
\begin{proof}[解答]
记$S_i$为“发送信号$i$”，$R_i$为“接收到信号$i$”。按照本题通常采用的模型，“出错的概率为$0.05$”解释为发送0、1时各自的条件误码率均为$0.05$，即
\[
\P(R_1\mid S_0)=0.05,\qquad \P(R_0\mid S_1)=0.05.
\]
于是
\begin{align*}
\P(R_1)&=\frac13\cdot0.05+\frac23\cdot0.95=\frac{39}{60},\\
\P(S_1\mid R_1)&=\frac{\P(S_1)\P(R_1\mid S_1)}{\P(R_1)}
=\frac{(2/3)\cdot0.95}{(1/3)\cdot0.05+(2/3)\cdot0.95}=\frac{38}{39}.
\end{align*}
若题意仅给出总误码率$\P(R_1S_0\cup R_0S_1)=0.05$，而不对两种信号的条件误码率作规定，则答案不能唯一确定。事实上，设$e_0=\P(R_1\mid S_0)$、$e_1=\P(R_0\mid S_1)$，题设只给出$e_0+2e_1=0.15$，而待求概率为
\[
\frac{2(1-e_1)}{e_0+2(1-e_1)},
\]
其值随$e_0,e_1$改变。
\end{proof}
\paragraph{答案}两种信号的条件误码率均为$0.05$时，答案为$38/39$；只给总误码率则条件不足。

\subsection{教材第25页第6题}
\label{ex:w04-t25-6}
\paragraph{题目}
盒子里有标号为$0,1,\ldots,k$的$k+1$枚不均匀的硬币。已知投掷第$i$号硬币正面朝上的概率为$i/k$。从盒中随机取一枚硬币重复投掷。设前$n$次投掷的结果都是正面朝上，求第$n+1$次也是正面朝上的概率。
\textbf{作业说明：}取定一枚硬币后，重复独立地投掷这一枚硬币，不换硬币。
\basis{原笔记第1章第3节“全概率公式”“贝叶斯公式”和第4节“独立试验”。}
\paragraph{思路}投掷在给定硬币后独立，但在未给定硬币时不能直接认为独立。先分别求前$n$次、前$n+1$次全为正面的概率。
\begin{proof}[解答]
设$K_i$为“取到第$i$号硬币”，$H_n$为“前$n$次均为正面”，其中$k\geq1$、$n\geq1$。由取币等可能以及给定硬币后的独立性，
\[
\P(K_i)=\frac1{k+1},\qquad \P(H_n\mid K_i)=\left(\frac ik\right)^n.
\]
用全概率公式得到
\[
\P(H_n)=\frac1{k+1}\sum_{i=0}^k\left(\frac ik\right)^n>0.
\]
记第$n+1$次为正面的事件为$H$。由于$H\cap H_n=H_{n+1}$，
\begin{align*}
\P(H\mid H_n)
&=\frac{\P(H_{n+1})}{\P(H_n)}\\
&=\frac{\sum_{i=0}^k(i/k)^{n+1}}{\sum_{i=0}^k(i/k)^n}
=\frac{\sum_{i=1}^ki^{n+1}}{k\sum_{i=1}^ki^n}.
\end{align*}
\end{proof}
\paragraph{答案}$\displaystyle\frac{\sum_{i=1}^ki^{n+1}}{k\sum_{i=1}^ki^n}$。

\subsection{教材第32页第1题}
\label{ex:w04-t32-1}
\paragraph{题目}
证明：
\begin{enumerate}
\item 若事件$A$与自身独立，则$\P(A)=0$或$1$；
\item 若$\P(A)=0$或$1$，则$A$与任意事件独立。
\end{enumerate}
\basis{原笔记第1章第4节“两个事件的独立性”；见\hyperref[k:independence]{事件独立性}。}
\paragraph{思路}将独立性的定义写成概率等式，并分别处理概率为0与1的情形。
\begin{proof}
\begin{enumerate}
\item 因$A\cap A=A$，独立性给出
\[
\P(A)=\P(A)^2,\qquad \P(A)(1-\P(A))=0.
\]
故$\P(A)=0$或$1$。
\item 任取事件$B$。若$\P(A)=0$，由$A\cap B\subseteq A$可得
\[
0\leq\P(A\cap B)\leq\P(A)=0,
\]
所以$\P(A\cap B)=0=\P(A)\P(B)$。

若$\P(A)=1$，则$\P(A^c)=0$，从而$\P(A^c\cap B)=0$。将$B$分解成两个互斥事件，得到
\[
\P(B)=\P(A\cap B)+\P(A^c\cap B)=\P(A\cap B).
\]
故仍有$\P(A\cap B)=\P(A)\P(B)$。两种情形均满足独立性的定义。
\end{enumerate}
\end{proof}
\paragraph{答案}两个结论均成立；事件与自身独立当且仅当其概率为0或1。

\subsection{教材第32页第2题}
\label{ex:w04-t32-2}
\paragraph{题目}
对于某个家庭考察两个事件$A=\{\text{至多有一个女孩}\}$和$B=\{\text{男女孩都有}\}$。设男女孩的出生率都是$1/2$。证明：对3孩之家，事件$A$和$B$独立；而对4孩之家，事件$A$和$B$不独立。
\basis{原笔记第1章第1节“古典概率模型”和第4节“两个事件的独立性”。}
\paragraph{思路}按女孩人数计数，比较$\P(AB)$和$\P(A)\P(B)$。采用题目通常隐含的模型：不同孩子的性别独立。
\begin{proof}
按出生次序记录性别。$n$个孩子共有$2^n$个等可能结果。$A\cap B$等价于“恰有一个女孩”。

对3孩之家，
\[
\P(A)=\frac{\binom30+\binom31}{2^3}=\frac12,
\qquad \P(B)=1-\frac2{2^3}=\frac34,
\qquad \P(AB)=\frac{\binom31}{2^3}=\frac38.
\]
因此$\P(AB)=\P(A)\P(B)$，二者独立。

对4孩之家，
\[
\P(A)=\frac{\binom40+\binom41}{2^4}=\frac5{16},
\qquad \P(B)=1-\frac2{2^4}=\frac78,
\qquad \P(AB)=\frac{\binom41}{2^4}=\frac14.
\]
而$\P(A)\P(B)=35/128\ne1/4$，故二者不独立。
\end{proof}
\paragraph{答案}3孩时独立；4孩时不独立。

\subsection{教材第33页第6题}
\label{ex:w04-t33-6}
\paragraph{题目}
重复投掷硬币，设第一次出现正面的概率为$c$（$0\leq c\leq1$），之后每次与上次出现相同面的概率为$p$（$0<p<1$）。求第$n$次出现正面的概率$p_n$，并讨论当$n\to\infty$时$p_n$的极限。

\textbf{提示：}根据全概率公式有$p_n=(1-p)+(2p-1)p_{n-1}$，于是有
\[
\begin{aligned}
p_n&=(1-p)[1+(2p-1)+\cdots+(2p-1)^{n-2}]+c(2p-1)^{n-1}\\
&=\frac12-\frac12(2p-1)^{n-1}+c(2p-1)^{n-1}.
\end{aligned}
\]
\basis{原笔记第1章第3节“全概率公式”及第5节“不独立子试验与转移概率”。}
\paragraph{思路}沿教材提示，先按上一次的结果分类，再将递推式的常数平衡值$1/2$减去。
\begin{proof}[解答]
设$H_n$为“第$n$次出现正面”。题设转移概率给出
\[
\P(H_n\mid H_{n-1})=p,\qquad
\P(H_n\mid H_{n-1}^c)=1-p.
\]
因而$n\geq2$时，
\begin{align*}
p_n&=pp_{n-1}+(1-p)(1-p_{n-1})\\
&=(1-p)+(2p-1)p_{n-1},\\
p_n-\frac12&=(2p-1)\left(p_{n-1}-\frac12\right).
\end{align*}
结合$p_1=c$，反复使用最后一个等式，得到
\[
p_n=\frac12+\left(c-\frac12\right)(2p-1)^{n-1}.
\]
对$n=1$此式直接取值为$c$；$p=1/2$时从$n=2$开始$p_n=1/2$。由于$|2p-1|<1$，无论初始概率$c$如何，都有$p_n\to1/2$。
\end{proof}
\paragraph{答案}$p_n=\frac12+(c-\frac12)(2p-1)^{n-1}$，极限为$1/2$。

\subsection{作业第2题：讲义例1.101第(3)问}
\label{ex:w04-h2}
\paragraph{原作业要求}
9月21日讲义例1.101第(3)题，可直接用第(1)(2)小题结论。

本次提供材料中未找到可核对的题面，按要求略去。

\subsection{作业第3题：两个并事件的独立性}
\label{ex:w04-h3}
\paragraph{题目}
设事件$A_1,A_2,A_3,A_4$相互独立。求证：$A_1\cup A_2$和$A_3\cup A_4$相互独立。
\basis{原笔记第1章第4节“相互独立”与命题“独立性对补事件保持”；见\hyperref[k:independence]{事件独立性}。}
\paragraph{思路}将并事件取补；先利用容斥公式说明相互独立的事件可同时换成补事件。
\begin{proof}
记$p_i=\P(A_i)$。对任意有限指标集$J\subseteq\{1,2,3,4\}$，容斥公式与相互独立性给出
\begin{align*}
\P\left(\bigcap_{j\in J}A_j^c\right)
&=1-\P\left(\bigcup_{j\in J}A_j\right)\\
&=\sum_{S\subseteq J}(-1)^{|S|}\prod_{j\in S}p_j
=\prod_{j\in J}(1-p_j).
\end{align*}
这里空集对应的乘积为1。因此，令$C=A_1\cup A_2$、$D=A_3\cup A_4$，就有
\begin{align*}
\P(C^c\cap D^c)&=\prod_{i=1}^4(1-p_i),\\
\P(C^c)\P(D^c)&=[(1-p_1)(1-p_2)][(1-p_3)(1-p_4)]
=\prod_{i=1}^4(1-p_i).
\end{align*}
故$C^c,D^c$独立。引用独立性对补事件保持的结论，即得$C,D$独立。
\end{proof}
\paragraph{答案}$A_1\cup A_2$与$A_3\cup A_4$相互独立。

\subsection{作业第4题：失败无穷多次}
\label{ex:w04-h4}
\paragraph{题目}
做独立重复试验，每次试验成功的概率为$p$，$0<p<1$。设$A$表示试验失败无穷多次，求$\P(A)$。
\basis{原笔记第1章第2节概率的连续性、第4节“伯努利试验中的常用计算”；见\hyperref[k:events]{事件与概率}、\hyperref[k:independence]{独立试验}。}
\paragraph{思路}补事件“只失败有限次”表示从某一次开始一直成功；先证明每个这样的尾事件概率为0。
\begin{proof}[解答]
设$S_j$为“第$j$次成功”，并对$m\geq1$定义
\[
E_m=\bigcap_{j=m}^\infty S_j.
\]
事件$E_m$表示“从第$m$次开始全都成功”。当$N\to\infty$时，$\bigcap_{j=m}^{m+N-1}S_j\downarrow E_m$。由独立性和概率的上连续性，
\[
\P(E_m)=\lim_{N\to\infty}\P\left(\bigcap_{j=m}^{m+N-1}S_j\right)
=\lim_{N\to\infty}p^N=0.
\]
失败只有有限次，当且仅当存在某个$m$使第$m$次起不再失败，因此
\[
A^c=\bigcup_{m=1}^\infty E_m.
\]
由可列次可加性（布尔不等式），$\P(A^c)\leq\sum_{m=1}^\infty\P(E_m)=0$，故$\P(A)=1$。
\end{proof}
\paragraph{答案}$\P(A)=1$。
% !TEX root = 概率论期中复习.tex
\section{第五周：重复试验与随机变量}

\subsection{教材第33页第4题}
\label{ex:w05-t33-4}
\paragraph{题目}
独立地重复投掷两枚骰子，并计算点数之和。求“和数为5”出现在“和数为7”之前的概率。
\basis{原笔记第1章第4节例“事件$A$先于事件$B$发生”；见\hyperref[k:independence]{独立试验}。}
\paragraph{思路}在第一次出现5或7以前，其余点数和只使试验继续。
\begin{proof}[解答]
记$A$为“一次投掷的点数和为5”，$B$为“点数和为7”。两枚均匀骰子共有36个等可能的有序结果，和为5有
\[
(1,4),(2,3),(3,2),(4,1)
\]
四种，和为7有六种，故
\[
\P(A)=\frac19,\qquad \P(B)=\frac16,\qquad
\P((A\cup B)^c)=\frac{13}{18}.
\]
$A,B$互斥，且各次投掷独立。按首次出现$A\cup B$的次数分类，所求概率为
\[
\sum_{n=1}^\infty\left(\frac{13}{18}\right)^{n-1}\frac19
=\frac{1/9}{1-13/18}=\frac25.
\]
\end{proof}
\paragraph{答案}$2/5$。

\subsection{教材第44页第5题}
\label{ex:w05-t44-5}
\paragraph{题目}
设随机变量$\xi$满足$\P(\xi\geq a)\geq1-\alpha$和$\P(\xi\leq b)\geq1-\beta$，其中$a\leq b\in\R$，而$\alpha,\beta\in[0,1]$。证明
\[
\P(a\leq\xi\leq b)\geq1-\alpha-\beta.
\]
\basis{原笔记第1章第2节命题“一般加法公式”“可列次可加性”；见\hyperref[k:events]{事件与概率}。}
\paragraph{思路}区间事件的补由左右两端越界事件组成，分别用已知概率控制。
\begin{proof}
由题设，
\[
\P(\xi<a)=1-\P(\xi\geq a)\leq\alpha,
\qquad
\P(\xi>b)=1-\P(\xi\leq b)\leq\beta.
\]
因$a\leq b$，事件$\{\xi<a\}$和$\{\xi>b\}$互斥，且它们的并正是$\{a\leq\xi\leq b\}^c$。所以
\begin{align*}
\P(a\leq\xi\leq b)
&=1-\P(\xi<a)-\P(\xi>b)\\
&\geq1-\alpha-\beta.
\end{align*}
即使$\alpha+\beta>1$，右侧为负，以上不等式仍然成立。
\end{proof}
\paragraph{答案}$\P(a\leq\xi\leq b)\geq1-\alpha-\beta$。

\subsection{作业第2题：区间投点的重复试验}
\label{ex:w05-h2}
\paragraph{题目}
一个试验为向$[0,1]$区间内随机地投掷一个粒子，观察该粒子的位置。设事件$A$表示粒子位于$[0,1/3]$区间内，事件$B$表示粒子位于$[1/4,3/4]$区间内。现将该试验独立重复地做。设事件$C$表示前$k$次独立重复试验中$A$至少发生一次，事件$D$表示前$k$次独立重复试验中$B$至少发生一次。
\begin{enumerate}
\item 求概率$\P(C)$和$\P(D)$。
\item 求概率$\P(C\cup D)$。
\item 请问：事件$C$和事件$D$是否相互独立？请说明理由。若其相互不独立，保持$A$事件不变，请调整$B$事件中的区间，要求保持$0<\P(B)<1$，使得事件$C$与$D$相互独立。
\item 由(3)启发，得到“$A$与$B$的相互独立性”和“$C$与$D$的相互独立性”的关系，请说明理由。
\item 保持$A$与$B$如题干不变，不进行(3)(4)的调整，求将该试验独立重复地进行，事件$A$严格发生在事件$B$之前的概率。
\end{enumerate}
\basis{原笔记第1章第1节“几何概率模型”、第4节例“独立与互斥的区别”“事件$A$先于事件$B$发生”；见\hyperref[k:independence]{事件独立性与重复试验}。}
\paragraph{思路}“至少一次”取补转化为“每次都不发生”；独立性等价于两个补事件独立；严格先发生要求首次出现$A\cup B$时落在$A\setminus B$。
\begin{proof}[解答]
按均匀投点模型，概率为区间长度，故
\[
\P(A)=\frac13,\qquad \P(B)=\frac12,
\qquad \P(AB)=\frac13-\frac14=\frac1{12}.
\]
因此
\[
\P(A\cup B)=\frac13+\frac12-\frac1{12}=\frac34,
\qquad \P(A\setminus B)=\frac14.
\]
以下$k$为正整数。
\begin{enumerate}
\item 由各次试验独立，
\[
\P(C^c)=\left(1-\frac13\right)^k=\left(\frac23\right)^k,
\qquad
\P(D^c)=\left(1-\frac12\right)^k=\left(\frac12\right)^k.
\]
所以
\[
\P(C)=1-\left(\frac23\right)^k,
\qquad
\P(D)=1-\left(\frac12\right)^k.
\]
\item $C^c\cap D^c$要求每次试验中$A,B$均不发生，因此
\[
\P(C^c\cap D^c)=\P((A\cup B)^c)^k=\left(\frac14\right)^k,
\qquad \P(C\cup D)=1-\left(\frac14\right)^k.
\]
\item 由于
\[
\P(C^c)\P(D^c)=\left(\frac13\right)^k
\ne\left(\frac14\right)^k=\P(C^c\cap D^c),
\]
$C^c,D^c$不独立，从而$C,D$不独立。

保持$A$不变，改为$B'=[1/6,2/3]$。此时
\[
\P(B')=\frac12,\qquad
\P(AB')=\frac13-\frac16=\frac16=\P(A)\P(B').
\]
所以$A,B'$独立。相应地，若$D'$表示前$k$次中$B'$至少发生一次，则
\[
\P(C^c\cap (D')^c)
=\P(A^c\cap (B')^c)^k
=[\P(A^c)\P((B')^c)]^k
=\P(C^c)\P((D')^c).
\]
故$C,D'$独立，且调整后的$\P(B')$满足要求。
\item 一般地，记$a=\P(A)$、$b=\P(B)$、$r=\P(AB)$。则
\[
\P(C^c\cap D^c)=(1-a-b+r)^k,
\quad \P(C^c)\P(D^c)=[(1-a)(1-b)]^k.
\]
两式中底数均非负，正整数次幂在非负数上是单射。因此
\begin{align*}
C,D\text{独立}
&\Longleftrightarrow C^c,D^c\text{独立}\\
&\Longleftrightarrow1-a-b+r=(1-a)(1-b)\\
&\Longleftrightarrow r=ab\\
&\Longleftrightarrow A,B\text{独立}.
\end{align*}
\item 严格先发生意味着$A$首次发生的次数小于$B$首次发生的次数。同时落入$A\cap B$不算$A$严格先发生。令$E_n$为“前$n-1$次中$A,B$均不发生，第$n$次发生$A\setminus B$”。这些事件两两不交，其并即为待求事件，故
\begin{align*}
\P(A\text{严格先于}B)
&=\sum_{n=1}^\infty\P(E_n)\\
&=\sum_{n=1}^\infty\left(\frac14\right)^{n-1}\frac14
=\frac{1/4}{1-1/4}=\frac13.
\end{align*}
\end{enumerate}
\end{proof}
\paragraph{答案}(1) $1-(2/3)^k$、$1-(1/2)^k$；(2) $1-(1/4)^k$；(3) 不独立，可改取$B'=[1/6,2/3]$；(4) $C,D$独立当且仅当$A,B$独立；(5) $1/3$。

\subsection{作业第3题：示性函数的乘积与并}
\label{ex:w05-h3}
\paragraph{题目}
设概率空间$(\Omega,\mathcal F,\P)$，事件$A,B\in\mathcal F$。求证：
\begin{enumerate}
\item 对任意$\omega\in\Omega$，示性函数$I_{AB}(\omega)=I_A(\omega)\cdot I_B(\omega)$；
\item 对任意$\omega\in\Omega$，示性函数$I_{A\cup B}(\omega)=I_A(\omega)+I_B(\omega)-I_A(\omega)\cdot I_B(\omega)$。
\end{enumerate}
\basis{原笔记第2章第1节例“示性函数”；见\hyperref[k:cdf]{随机变量与分布函数}。}
\paragraph{思路}按$\omega$分别属于$A\cap B$、$A\setminus B$、$B\setminus A$、$(A\cup B)^c$分类。四类互斥且覆盖$\Omega$。
\begin{proof}
对事件$E$，其示性函数$I_E(\omega)$在$\omega\in E$时为1，在$\omega\notin E$时为0。
\begin{enumerate}
\item 若$\omega\in A\cap B$，则$I_A(\omega)=I_B(\omega)=1$，乘积为1，等于$I_{AB}(\omega)$。若$\omega\notin A\cap B$，则$I_A(\omega),I_B(\omega)$至少有一个为0，乘积为0，也等于$I_{AB}(\omega)$。
\item 若$\omega\in A\cap B$，右端为$1+1-1=1$；若$\omega\in A\setminus B$，右端为$1+0-0=1$；若$\omega\in B\setminus A$，右端为$0+1-0=1$；若$\omega\notin A\cup B$，右端为0。前面三种情形恰好对应$I_{A\cup B}(\omega)=1$，最后一种对应$I_{A\cup B}(\omega)=0$。所以等式对所有$\omega$成立。
\end{enumerate}
\end{proof}
\paragraph{答案}$I_{AB}=I_AI_B$，$I_{A\cup B}=I_A+I_B-I_AI_B$，均为逐点恒等式。

\subsection{作业第4题：随机变量与常数取最小值}
\label{ex:w05-h4}
\paragraph{题目}
设$(\Omega,\mathcal F,\P)$是一个概率空间，$X$是该概率空间上的一个随机变量，$a$是一个常数。对任意$\omega\in\Omega$，令$Y(\omega):=\min\{X(\omega),a\}=X(\omega)\wedge a$。求证：$Y$也是该概率空间上的一个随机变量。
\basis{原笔记第2章第1节命题“随机变量的半直线判别”；见\hyperref[k:cdf]{随机变量与分布函数}。}
\paragraph{思路}不需要先求$Y$的分布，只需检查每个半直线的原像是否为事件。
\begin{proof}
由于$X(\omega)$与$a$均为有限实数，$Y$是$\Omega$上的实值函数。任取$t\in\R$。若$t<a$，则
\[
\min\{X(\omega),a\}\leq t\quad\Longleftrightarrow\quad X(\omega)\leq t.
\]
其中，若$X(\omega)>t$，因为$a>t$，两者的最小值也大于$t$。若$t\geq a$，则对每个$\omega$均有$Y(\omega)\leq a\leq t$。因此
\[
\{Y\leq t\}=
\begin{cases}
\{X\leq t\},&t<a,\\
\Omega,&t\geq a.
\end{cases}
\]
$X$为随机变量，故$\{X\leq t\}\in\mathcal F$；又$\Omega\in\mathcal F$。于是所有$\{Y\leq t\}$都属于$\mathcal F$，由半直线判别，$Y$是随机变量。
\end{proof}
\paragraph{答案}$Y=X\wedge a$是原概率空间上的随机变量。

\clearpage\part{一维随机变量及其分布}
% !TEX root = 概率论期中复习.tex
% 自动抽取自 chapters/chapter2.tex；原例题及其解答保留。

\clearpage\section{随机变量及其分布函数}\label{k:cdf}
\label{sec:random-variables-and-distribution-functions}

\begin{mydefinition}{随机变量}{random-variable}
给定概率空间 $(\Omega,\mathscr F,P)$。若映射 $X:\Omega\to\mathbb R$ 满足：对任意 $B\in\mathscr B(\mathbb R)$，都有
\[
X^{-1}(B):=\{\omega\in\Omega:X(\omega)\in B\}\in\mathscr F,
\]
则称 $X$ 为该概率空间上的\textbf{随机变量}。随机变量通常用 $X,Y,Z$ 或 $\xi,\eta$ 等表示。
\end{mydefinition}
\origin{原笔记第2.1节，定义“随机变量”；源码第9行。}

\begin{proposition}{随机变量的半直线判别}{random-variable-halfline-criterion}
映射 $X:\Omega\to\mathbb R$ 是随机变量，当且仅当对任意 $t\in\mathbb R$ 都有
\[
\{\omega\in\Omega:X(\omega)\leq t\}\in\mathscr F.
\]
\end{proposition}
\origin{原笔记第2.1节，命题“随机变量的半直线判别”；源码第19行。}

\subsection{例题：示性函数}\label{note:c2:e1}
\basis{原笔记第2.1节，例“示性函数”；源码\texttt{chapters/chapter2.tex}第26行。}
设 $A\in\mathscr F$，定义
\[
I_A(\omega)=
\begin{cases}
1,&\omega\in A,\\
0,&\omega\notin A.
\end{cases}
\]
称 $I_A$ 为事件 $A$ 的\textbf{示性函数}。由于
\[
\{I_A\leq t\}=
\begin{cases}
\varnothing,&t<0,\\
A^{\mathrm c},&0\leq t<1,\\
\Omega,&t\geq1,
\end{cases}
\]
上述集合均属于 $\mathscr F$，所以 $I_A$ 是随机变量。并且
\[
P(I_A=1)=P(A),\qquad P(I_A=0)=P(A^{\mathrm c}).
\]

\begin{mydefinition}{分布函数}{distribution-function}
设 $X$ 是概率空间 $(\Omega,\mathscr F,P)$ 上的随机变量。函数
\[
F(t):=P(X\leq t),\qquad t\in\mathbb R,
\]
称为随机变量 $X$ 的\textbf{分布函数}。需要区分不同随机变量时，也记为 $F_X(t)$。
\end{mydefinition}
\origin{原笔记第2.1节，定义“分布函数”；源码第73行。}

\begin{mytheorem}{分布函数的基本性质}{distribution-function-properties}
设 $F$ 为随机变量 $X$ 的分布函数，则：
\begin{enumerate}
    \item $F$ 单调不减，即 $a\leq b$ 时有 $F(a)\leq F(b)$；
    \item $\displaystyle\lim_{t\to-\infty}F(t)=0$，$\displaystyle\lim_{t\to+\infty}F(t)=1$；
    \item $F$ 右连续，即对任意 $a\in\mathbb R$，
    \[
    F(a)=\lim_{t\to a^+}F(t);
    \]
    \item 对任意 $a\in\mathbb R$，记
    \[
    F(a^-):=\lim_{t\to a^-}F(t),
    \]
    则
    \[
    P(X=a)=F(a)-F(a^-).
    \]
\end{enumerate}
\end{mytheorem}
\origin{原笔记第2.1节，定理“分布函数的基本性质”；源码第85行。}

\subsection{例题：含跳跃的分布函数}\label{note:c2:e2}
\basis{原笔记第2.1节，例“含跳跃的分布函数”；源码\texttt{chapters/chapter2.tex}第169行。}
随机变量 $X$ 的分布函数为
\[
F(t)=
\begin{cases}
0,&t<0,\\[0.3em]
\dfrac t2,&0\leq t<1,\\[0.5em]
\dfrac23,&1\leq t<2,\\[0.5em]
\dfrac{11}{12},&2\leq t<3,\\[0.5em]
1,&t\geq3.
\end{cases}
\]
分别计算：
\[
P(X<3)=F(3^-)=\frac{11}{12},
\]
\[
P(X=1)=F(1)-F(1^-)=\frac23-\frac12=\frac16,
\]
\[
P\left(X>\frac12\right)=1-F\left(\frac12\right)
=1-\frac14=\frac34,
\]
\[
P(2<X\leq4)=F(4)-F(2)=1-\frac{11}{12}=\frac1{12},
\]
\[
P\left(\frac12\leq X<2\right)
=F(2^-)-F\left(\left(\frac12\right)^-\right)
=\frac23-\frac14=\frac5{12}.
\]

\subsection{例题：可列个分布函数的凸组合}\label{note:c2:e3}
\basis{原笔记第2.1节，例“可列个分布函数的凸组合”；源码\texttt{chapters/chapter2.tex}第204行。}
设 $\{F_m\}_{m\geq1}$ 是一列分布函数，$a_m\geq0$ 且
\[
\sum_{m=1}^{\infty}a_m=1.
\]
令
\[
F(t)=\sum_{m=1}^{\infty}a_mF_m(t).
\]
称 $F$ 为这些分布函数的\textbf{凸组合}。证明 $F$ 也是分布函数。

首先，
\[
0\leq\sum_{m=1}^{\infty}a_mF_m(t)
\leq\sum_{m=1}^{\infty}a_m=1,
\]
所以级数收敛，且 $F$ 定义良好。

（1）若 $a\leq b$，则每个 $F_m(a)\leq F_m(b)$，从而
\[
F(a)=\sum_{m=1}^{\infty}a_mF_m(a)
\leq\sum_{m=1}^{\infty}a_mF_m(b)=F(b).
\]

（2）固定 $a\in\mathbb R$ 和 $\varepsilon>0$。选择正整数 $M$，使
\[
\sum_{m=M+1}^{\infty}a_m<\frac{\varepsilon}{2}.
\]
由 $F_1,\ldots,F_M$ 右连续，存在 $\delta>0$，使 $a<t<a+\delta$ 时，对所有 $1\leq m\leq M$ 都有
\[
0\leq F_m(t)-F_m(a)<\frac{\varepsilon}{2}.
\]
因此
\[
\begin{aligned}
0\leq F(t)-F(a)
&=\sum_{m=1}^{M}a_m\bigl(F_m(t)-F_m(a)\bigr)\\
&\quad+\sum_{m=M+1}^{\infty}a_m\bigl(F_m(t)-F_m(a)\bigr)\\
&\leq\frac{\varepsilon}{2}\sum_{m=1}^{M}a_m
+\sum_{m=M+1}^{\infty}a_m\\
&<\varepsilon.
\end{aligned}
\]
故 $F$ 右连续。

（3）仍选取上述 $M$。由于有限个分布函数 $F_1,\ldots,F_M$ 在负无穷处均趋于 $0$，当 $t$ 足够小时，
\[
0\leq F(t)
\leq\frac{\varepsilon}{2}\sum_{m=1}^{M}a_m
+\sum_{m=M+1}^{\infty}a_m
<\varepsilon.
\]
故 $\lim_{t\to-\infty}F(t)=0$。同理，当 $t$ 足够大时，$1-F_m(t)<\varepsilon/2$ 对所有 $1\leq m\leq M$ 成立，于是
\[
0\leq1-F(t)
=\sum_{m=1}^{\infty}a_m\bigl(1-F_m(t)\bigr)
<\varepsilon.
\]
故 $\lim_{t\to+\infty}F(t)=1$。

综上，$F$ 满足分布函数的三个本质条件，因此 $F$ 是一个分布函数。

\clearpage\section{离散型随机变量}\label{k:discrete}

\begin{mydefinition}{离散型随机变量与分布列}{discrete-random-variable-pmf}
若随机变量 $X$ 的可能取值为有限个 $x_1,\ldots,x_s$ 或可列个 $x_1,x_2,\ldots$，则称 $X$ 为\textbf{离散型随机变量}。这些取值互不相同。令
\[
p(x_i):=P(X=x_i),
\]
称各取值及其对应概率构成的表为 $X$ 的\textbf{概率分布列}，简称\textbf{分布列}。函数 $p$ 也称为概率质量函数。

分布列可写成表格
\[
\begin{array}{c|ccccc}
X&x_1&x_2&\cdots&x_s&\cdots\\ \hline
P&p(x_1)&p(x_2)&\cdots&p(x_s)&\cdots
\end{array}
\]
或两行矩阵的形式
\[
X\sim
\begin{pmatrix}
x_1&x_2&\cdots&x_s&\cdots\\
p(x_1)&p(x_2)&\cdots&p(x_s)&\cdots
\end{pmatrix}.
\]
有限取值的情形去掉末尾的省略号。
\end{mydefinition}
\origin{原笔记第2.2节，定义“离散型随机变量与分布列”；源码第274行。}

\begin{proposition}{分布列的基本性质}{pmf-basic-properties}
离散型随机变量的分布列满足
\[
p(x_i)\geq0,\qquad \sum_i p(x_i)=1.
\]
其中求和遍历所有不同取值；有限情形为有限和，可列情形为无穷级数。
\end{proposition}
\origin{原笔记第2.2节，命题“分布列的基本性质”；源码第301行。}

\begin{proposition}{分布列与分布函数相互确定}{pmf-cdf-equivalence}
设 $X$ 为离散型随机变量，分布列为 $p(x_i)$，分布函数为 $F$。则
\[
F(t)=\sum_{i:x_i\leq t}p(x_i),\qquad t\in\mathbb R,
\]
并且
\[
p(x_i)=F(x_i)-F(x_i^-).
\]
\end{proposition}
\origin{原笔记第2.2节，命题“分布列与分布函数相互确定”；源码第326行。}

\subsection{例题：收集齐优惠券所需的张数}\label{note:c2:e4}
\basis{原笔记第2.2节，例“收集齐优惠券所需的张数”；源码\texttt{chapters/chapter2.tex}第341行。}
有 $r$ 种不同优惠券，每次收集一张，每一种被收集到的概率均为 $1/r$，各次收集相互独立。设 $T$ 为首次收集齐全部 $r$ 种优惠券时累计收集的张数，求 $P(T=n)$。

当 $r=1$ 时，$P(T=1)=1$。以下设 $r\geq2$；显然 $n<r$ 时，$P(T=n)=0$。

先求 $P(T>m)$。设 $A_j^{(m)}$ 表示前 $m$ 次收集中没有第 $j$ 种优惠券，则
\[
\{T>m\}=\bigcup_{j=1}^rA_j^{(m)}.
\]
对任意不同的 $j_1,\ldots,j_k$，由独立性，
\[
P\left(A_{j_1}^{(m)}\cap\cdots\cap A_{j_k}^{(m)}\right)
=\left(1-\frac kr\right)^m.
\]
当 $m\geq r-1\geq1$ 时，容斥公式给出
\[
\begin{aligned}
P(T>m)
&=\sum_{k=1}^r(-1)^{k+1}\binom rk
\left(1-\frac kr\right)^m\\
&=\sum_{k=1}^{r-1}(-1)^{k+1}\binom rk
\left(1-\frac kr\right)^m,
\end{aligned}
\]
其中 $k=r$ 对应的项为 $0$。于是对 $n\geq r$，
\[
\begin{aligned}
P(T=n)
&=P(T>n-1)-P(T>n)\\
&=\sum_{k=1}^{r-1}(-1)^{k+1}\binom rk
\left[
\left(1-\frac kr\right)^{n-1}
-\left(1-\frac kr\right)^n
\right].
\end{aligned}
\]

\begin{mydefinition}{二项分布}{binomial-distribution}
设 $n$ 为正整数，$0<p<1$。若随机变量 $X$ 取值于 $0,1,\ldots,n$，且
\[
P(X=k)=\binom nk p^k(1-p)^{n-k},
\qquad k=0,1,\ldots,n,
\]
则称 $X$ 服从参数为 $(n,p)$ 的\textbf{二项分布}，记为 $X\sim B(n,p)$。
\end{mydefinition}
\origin{原笔记第2.2节，定义“二项分布”；源码第381行。}

\begin{proposition}{二项分布的单调性与最可能取值}{binomial-modes}
设 $X\sim B(n,p)$，$0<p<1$，令 $a=(n+1)p$。随着 $k$ 从 $0$ 增加到 $n$，概率 $P(X=k)$ 先不减后不增，最大值的位置为：
\begin{enumerate}
    \item 若 $a$ 不是整数，则唯一的最可能取值为 $\lfloor a\rfloor$；
    \item 若 $a=m$ 是整数，则 $m-1$ 和 $m$ 是两个并列的最可能取值，且 $1\leq m\leq n$。
\end{enumerate}
其中 $\lfloor a\rfloor$ 表示不大于 $a$ 的最大整数。
\end{proposition}
\origin{原笔记第2.2节，命题“二项分布的单调性与最可能取值”；源码第407行。}

\subsection{例题：均匀硬币的中心项}\label{note:c2:e5}
\basis{原笔记第2.2节，例“均匀硬币的中心项”；源码\texttt{chapters/chapter2.tex}第436行。}
独立地抛一枚均匀硬币 $2n$ 次，设 $X_n$ 为正面出现的次数，则
\[
X_n\sim B\left(2n,\frac12\right).
\]
由于 $(2n+1)/2=n+1/2$ 不是整数，唯一的最可能正面次数为 $n$，且
\[
P(X_n=n)=\binom{2n}{n}\frac1{2^{2n}}
=\frac{(2n)!}{(n!)^2}\frac1{2^{2n}}.
\]
由斯特林（Stirling）公式
\[
n!\sim\sqrt{2\pi n}\left(\frac ne\right)^n,
\qquad n\to\infty,
\]
可得
\[
\begin{aligned}
P(X_n=n)
&\sim
\frac{\sqrt{4\pi n}(2n/e)^{2n}}
{2\pi n(n/e)^{2n}}\frac1{2^{2n}}\\
&=\frac1{\sqrt{\pi n}}\longrightarrow0.
\end{aligned}
\]
这里 $u_n\sim v_n$ 表示 $u_n/v_n\to1$。对所有 $0\leq k\leq2n$，
\[
0\leq P(X_n=k)\leq P(X_n=n).
\]
因此，对每个固定的非负整数 $k$，由夹逼定理，
\[
\lim_{n\to\infty}P(X_n=k)=0.
\]

\begin{mydefinition}{泊松分布}{poisson-distribution}
设 $\lambda>0$。若随机变量 $X$ 取值于非负整数，且
\[
P(X=k)=\frac{\lambda^k}{k!}e^{-\lambda},
\qquad k=0,1,2,\ldots,
\]
则称 $X$ 服从参数为 $\lambda$ 的\textbf{泊松分布}，记为 $X\sim\mathcal P(\lambda)$。
\end{mydefinition}
\origin{原笔记第2.2节，定义“泊松分布”；源码第473行。}

\subsection{例题：一页上至少有一处印刷错误}\label{note:c2:e6}
\basis{原笔记第2.2节，例“一页上至少有一处印刷错误”；源码\texttt{chapters/chapter2.tex}第492行。}
设一本书某一页上的印刷错误数 $X$ 服从参数为 $\lambda=1/2$ 的泊松分布。则
\[
P(X\geq1)=1-P(X=0)=1-e^{-1/2}\approx0.393469.
\]

\subsection{例题：泊松粒子数的独立捕捉}\label{note:c2:e7}
\basis{原笔记第2.2节，例“泊松粒子数的独立捕捉”；源码\texttt{chapters/chapter2.tex}第499行。}
设某仪器在一定时间内发射的粒子总数 $\xi\sim\mathcal P(\lambda)$。给定 $\xi=n$ 时，每个粒子以相同概率 $p$ 被捕捉记录，且各粒子的捕捉结果相互独立，其中 $0<p<1$。记捕捉到的粒子总数为 $\eta$，证明 $\eta\sim\mathcal P(\lambda p)$。

令 $q=1-p$。给定 $\xi=n$ 时，$\eta$ 的条件分布为 $B(n,p)$。对任意非负整数 $k$，由全概率公式，
\[
\begin{aligned}
P(\eta=k)
&=\sum_{n=k}^{\infty}P(\xi=n)P(\eta=k\mid\xi=n)\\
&=\sum_{n=k}^{\infty}\frac{\lambda^n}{n!}e^{-\lambda}
\binom nk p^kq^{n-k}\\
&=\frac{e^{-\lambda}(\lambda p)^k}{k!}
\sum_{n=k}^{\infty}\frac{(\lambda q)^{n-k}}{(n-k)!}\\
&=\frac{e^{-\lambda}(\lambda p)^k}{k!}
\sum_{j=0}^{\infty}\frac{(\lambda q)^j}{j!}\\
&=\frac{(\lambda p)^k}{k!}e^{-\lambda}e^{\lambda q}
=\frac{(\lambda p)^k}{k!}e^{-\lambda p}.
\end{aligned}
\]
其中换元 $j=n-k$，所以 $\eta\sim\mathcal P(\lambda p)$。

\begin{proposition}{泊松分布的单调性与最可能取值}{poisson-modes}
设 $X\sim\mathcal P(\lambda)$。随着 $k$ 从 $0$ 增大，$P(X=k)$ 先不减后不增，并趋于 $0$。
\begin{enumerate}
    \item 若 $\lambda$ 不是整数，则唯一的最可能取值为 $\lfloor\lambda\rfloor$；
    \item 若 $\lambda=m$ 为正整数，则 $m-1$ 和 $m$ 是两个并列的最可能取值。
\end{enumerate}
\end{proposition}
\origin{原笔记第2.2节，命题“泊松分布的单调性与最可能取值”；源码第520行。}

\begin{mytheorem}{泊松定理}{poisson-limit-theorem}
设 $X_n\sim B(n,p_n)$，$0<p_n<1$，且
\[
np_n\longrightarrow\lambda>0.
\]
则对每个固定的非负整数 $k$，
\[
\lim_{n\to\infty}P(X_n=k)=\frac{\lambda^k}{k!}e^{-\lambda}.
\]
\end{mytheorem}
\origin{原笔记第2.2节，定理“泊松定理”；源码第546行。}

\subsection{例题：死亡人数不超过五人的泊松近似}\label{note:c2:e8}
\basis{原笔记第2.2节，例“死亡人数不超过五人的泊松近似”；源码\texttt{chapters/chapter2.tex}第588行。}
某患病人群中，每人因该病死亡的概率为 $0.2\%$。在 $1000$ 人的结果相互独立且具有相同概率的模型下，设死亡人数为 $X$，则
\[
X\sim B(1000,0.002).
\]
精确概率为
\[
P(X\leq5)
=\sum_{k=0}^{5}\binom{1000}{k}0.002^k0.998^{1000-k}.
\]
令 $Z\sim\mathcal P(2)$，因为 $np=1000\times0.002=2$，所以
\[
P(X\leq5)\approx P(Z\leq5)
=e^{-2}\sum_{k=0}^{5}\frac{2^k}{k!}
\approx0.983436\approx0.9834.
\]
这是对原二项分布概率的近似，不能将 $X$ 的实际分布直接改写为泊松分布。

\begin{mydefinition}{几何分布}{geometric-distribution}
设 $0<p<1$。若随机变量 $X$ 取值于正整数，且
\[
P(X=k)=(1-p)^{k-1}p,\qquad k=1,2,\ldots,
\]
则称 $X$ 服从参数为 $p$ 的\textbf{几何分布}，记为 $X\sim G(p)$。
\end{mydefinition}
\origin{原笔记第2.2节，定义“几何分布”；源码第618行。}

\begin{mytheorem}{几何分布的无记忆性}{geometric-memorylessness}
设 $X$ 是取正整数值的随机变量，且对每个非负整数 $k$ 都有 $P(X>k)>0$。则 $X$ 服从某个参数 $0<p<1$ 的几何分布，当且仅当对任意正整数 $n$ 和非负整数 $k$，
\[
P(X=n+k\mid X>k)=P(X=n).
\]
\end{mytheorem}
\origin{原笔记第2.2节，定理“几何分布的无记忆性”；源码第642行。}

\begin{mydefinition}{帕斯卡分布}{pascal-distribution}
设 $n$ 为正整数，$0<p<1$。若随机变量 $X$ 取值于 $n,n+1,\ldots$，且
\[
P(X=k)=\binom{k-1}{n-1}p^n(1-p)^{k-n},
\qquad k=n,n+1,\ldots,
\]
则称 $X$ 服从参数为 $(n,p)$ 的\textbf{帕斯卡分布}，记为 $X\sim\mathrm{Pa}(n,p)$。
\end{mydefinition}
\origin{原笔记第2.2节，定义“帕斯卡分布”；源码第686行。}

\begin{mydefinition}{负二项分布}{negative-binomial-distribution}
若 $X\sim\mathrm{Pa}(n,p)$，令 $Y=X-n$，则 $Y$ 表示取得第 $n$ 次成功时已经发生的\textbf{失败次数}，取值于非负整数。其分布列为
\[
\begin{aligned}
P(Y=k)&=P(X=n+k)\\
&=\binom{n+k-1}{n-1}p^n(1-p)^k\\
&=\binom{n+k-1}{k}p^n(1-p)^k,\qquad k=0,1,\ldots.
\end{aligned}
\]
称 $Y$ 服从参数为 $(n,p)$ 的\textbf{负二项分布}，记为 $Y\sim\mathrm{NB}(n,p)$。
\end{mydefinition}
\origin{原笔记第2.2节，定义“负二项分布”；源码第715行。}

\subsection{例题：巴拿赫火柴盒问题}\label{note:c2:e9}
\basis{原笔记第2.2节，例“巴拿赫火柴盒问题”；源码\texttt{chapters/chapter2.tex}第739行。}
一个人左右口袋各放一盒火柴，每盒最初有 $N$ 根，$N$ 为正整数。每次需要火柴时，相互独立地以概率 $1/2$ 选择一盒；若盒内有火柴，就取出一根。求他首次选到一个空盒、发现该盒已用光时，另一盒恰好剩下 $k$ 根的概率，其中 $0\leq k\leq N$。

设 $A_k$ 为首次发现右盒为空时，左盒恰好剩下 $k$ 根的事件。由于左右对称，两种首次发现空盒的事件互斥，所求概率为 $2P(A_k)$。

$A_k$ 发生时，这是第 $N+1$ 次选择右盒，而此前共选择左盒 $N-k$ 次。将选择右盒视为成功，则
\[
T\sim\mathrm{Pa}\left(N+1,\frac12\right),
\]
其中 $T$ 表示第 $N+1$ 次选择右盒的时刻。$A_k$ 等价于 $T=2N-k+1$；此前选择左盒的次数不超过 $N$，因而还没有发现左盒为空。于是
\[
\begin{aligned}
P(A_k)
&=\binom{2N-k}{N}
\left(\frac12\right)^{N+1}
\left(\frac12\right)^{N-k}\\
&=\binom{2N-k}{N}\frac1{2^{2N-k+1}}.
\end{aligned}
\]
因此所求概率为
\[
\boxed{\,P(\text{另一盒剩下 }k\text{ 根})
=\binom{2N-k}{N}\frac1{2^{2N-k}},\quad 0\leq k\leq N.\,}
\]

\textbf{注：}\quad 停止时刻是“首次选到空盒并发现已经用光”，不是“刚取走某盒的最后一根火柴”；区分这两个时刻才能得到正确的分布。

\clearpage\section{连续型随机变量}\label{k:continuous}

\begin{mydefinition}{连续型随机变量与概率密度}{continuous-random-variable-density}
设 $X$ 为随机变量。若存在非负可测函数 $f:\mathbb R\to[0,\infty)$，使得对任意 $B\in\mathscr B(\mathbb R)$，
\[
P(X\in B)=\int_B f(x)\,dx,
\]
则称 $X$ 为\textbf{连续型随机变量}，称 $f$ 为 $X$ 的\textbf{概率密度函数}，简称\textbf{密度函数}或\textbf{密度}。
\end{mydefinition}
\origin{原笔记第2.3节，定义“连续型随机变量与概率密度”；源码第794行。}

\begin{proposition}{概率密度的基本性质}{density-basic-properties}
若 $f$ 是连续型随机变量 $X$ 的密度，则
\[
f(x)\geq0,\qquad
\int_{\mathbb R}f(x)\,dx
=\int_{-\infty}^{+\infty}f(x)\,dx=1.
\]
\end{proposition}
\origin{原笔记第2.3节，命题“概率密度的基本性质”；源码第804行。}

\subsection{例题：分段密度的参数、区间概率与分布函数}\label{note:c2:e10}
\basis{原笔记第2.3节，例“分段密度的参数、区间概率与分布函数”；源码\texttt{chapters/chapter2.tex}第828行。}
设 $X$ 的概率密度为
\[
f(x)=
\begin{cases}
\dfrac18,&0<x<2,\\[0.4em]
kx,&2\leq x<4,\\
0,&\text{其他}.
\end{cases}
\]
求 $k$、$P(1<X\leq3)$ 以及分布函数 $F$。

（1）由密度的规范性，
\[
1=\int_0^2\frac18\,dx+\int_2^4kx\,dx
=\frac14+6k,
\]
故 $k=1/8$，也满足非负性要求。

（2）区间跨越密度的分段点 $2$，因此
\[
P(1<X\leq3)
=\int_1^2\frac18\,dx+\int_2^3\frac{x}{8}\,dx
=\frac18+\frac5{16}=\frac7{16}.
\]

（3）由 $F(t)=\int_{-\infty}^t f(x)\,dx$ 分段计算。特别地，$2\leq t<4$ 时，
\[
F(t)=\int_0^2\frac18\,dx+\int_2^t\frac{x}{8}\,dx
=\frac14+\frac{t^2-4}{16}=\frac{t^2}{16}.
\]
综上，
\[
F(t)=
\begin{cases}
0,&t<0,\\[0.3em]
\dfrac t8,&0\leq t<2,\\[0.5em]
\dfrac{t^2}{16},&2\leq t<4,\\[0.5em]
1,&t\geq4.
\end{cases}
\]
在 $0,2,4$ 处左右极限与函数值均相同，故 $F$ 连续。

\textbf{注：}\quad “几乎处处”是指除去一个勒贝格测度为 $0$ 的集合后成立；这里的测度是区间长度概念的推广。改变密度在单个点或零测集上的非负取值，不改变其所描述的分布。求导等式须按上述条件和几乎处处的意义理解。

\begin{proposition}{连续型分布的基本性质}{continuous-distribution-properties}
设 $X$ 为连续型随机变量，密度与分布函数分别为 $f$、$F$。则：
\begin{enumerate}
    \item 对任意 $t\in\mathbb R$，
    \[
    F(t)=\int_{-\infty}^t f(x)\,dx;
    \]
    \item $F$ 连续，且几乎处处可微，并且
    \[
    F'(t)=f(t)\qquad\text{几乎处处成立};
    \]
    在 $f$ 连续的点 $t$，由积分上限函数求导可得 $F'(t)=f(t)$；
    \item 若 Borel 集 $B$ 的勒贝格测度 $m(B)=0$，则
    \[
    P(X\in B)=0.
    \]
    特别地，对任意实数 $c$，$P(X=c)=0$；
    \item 对给定的随机变量 $X$，分布函数唯一，密度则只在相差一个零测集的意义下唯一。即若 $f_1,f_2$ 都是其密度，则
    \[
    m\bigl(\{x:f_1(x)\ne f_2(x)\}\bigr)=0.
    \]
\end{enumerate}
\end{proposition}
\origin{原笔记第2.3节，命题“连续型分布的基本性质”；源码第874行。}

\subsection{例题：五个电子管中恰有两个需要更换}\label{note:c2:e11}
\basis{原笔记第2.3节，例“五个电子管中恰有两个需要更换”；源码\texttt{chapters/chapter2.tex}第934行。}
某类电子管的寿命 $X$（单位：小时）具有密度
\[
f(x)=
\begin{cases}
\dfrac{100}{x^2},&x>100,\\[0.4em]
0,&x\leq100.
\end{cases}
\]
现有 $5$ 个具有相同寿命分布且相互独立的电子管，求 $150$ 小时内恰有 $2$ 个原装电子管因寿命耗尽需要更换的概率。

单个电子管在 $150$ 小时内需要更换的概率为
\[
p=P(X\leq150)
=\int_{100}^{150}\frac{100}{x^2}\,dx
=\left[-\frac{100}{x}\right]_{100}^{150}
=\frac13.
\]
记需要更换的原装电子管个数为 $Y$。由寿命的独立性，
\[
Y\sim B\left(5,\frac13\right),
\]
所以
\[
P(Y=2)=\binom52\left(\frac13\right)^2
\left(\frac23\right)^3=\frac{80}{243}.
\]

\subsection{例题：由分布函数求密度}\label{note:c2:e12}
\basis{原笔记第2.3节，例“由分布函数求密度”；源码\texttt{chapters/chapter2.tex}第963行。}
设连续型随机变量 $X$ 的分布函数为
\[
F(t)=
\begin{cases}
0,&t<0,\\
t^2,&0\leq t<1,\\
1,&t\geq1.
\end{cases}
\]
在各个开区间上求导，可选择一个密度版本为
\[
f(x)=
\begin{cases}
2x,&0<x<1,\\
0,&\text{其他}.
\end{cases}
\]
因为 $F$ 在 $1$ 处不可微，$f(1)=0$ 是选定的密度值，不是该点的导数；零测集上的取值可以另作非负指定。并且
\[
\int_0^1 2x\,dx=1,
\]
将此密度从 $-\infty$ 积分到 $t$，恰好还原上面的 $F(t)$。

\textbf{注：}\quad 可简写为“$f(x)=2x,\ 0<x<1$，其余处为 $0$”；不能把区间外的零值条件省略成对所有实数都有 $f(x)=2x$。

\begin{mydefinition}{均匀分布}{uniform-distribution}
设 $a<b$。若随机变量 $X$ 的密度可以取为
\[
f(x)=
\begin{cases}
\dfrac1{b-a},&a<x<b,\\[0.4em]
0,&\text{其他},
\end{cases}
\]
则称 $X$ 服从区间 $(a,b)$ 上的\textbf{均匀分布}，记为 $X\sim U(a,b)$。
\end{mydefinition}
\origin{原笔记第2.3节，定义“均匀分布”；源码第1005行。}

\subsection{例题：等公交车不超过五分钟}\label{note:c2:e13}
\basis{原笔记第2.3节，例“等公交车不超过五分钟”；源码\texttt{chapters/chapter2.tex}第1031行。}
公交车从 $7$ 点起每隔 $15$ 分钟到站一次。乘客到站时间服从 $7$ 点至 $7$ 点 $30$ 分之间的均匀分布，求等车时间不超过 $5$ 分钟的概率。

设 $X$ 为乘客到站时距 $7$ 点的分钟数，则 $X\sim U(0,30)$。在本时间段内公交车于 $0,15,30$ 分钟到站，所求事件为
\[
A=\{X=0\}\cup\{10\leq X\leq15\}\cup\{25\leq X\leq30\}.
\]
上述事件互斥，且 $P(X=0)=0$，因此
\[
\begin{aligned}
P(A)
&=P(X=0)+P(10\leq X\leq15)+P(25\leq X\leq30)\\
&=0+\int_{10}^{15}\frac1{30}\,dx+\int_{25}^{30}\frac1{30}\,dx\\
&=\frac5{30}+\frac5{30}=\frac13.
\end{aligned}
\]

\begin{mydefinition}{指数分布}{exponential-distribution}
设 $\lambda>0$。若随机变量 $X$ 的密度可以取为
\[
f(x)=
\begin{cases}
\lambda e^{-\lambda x},&x\geq0,\\
0,&x<0,
\end{cases}
\]
则称 $X$ 服从参数为 $\lambda$ 的\textbf{指数分布}，记为 $X\sim\operatorname{Exp}(\lambda)$。这里 $\lambda$ 使用率参数的约定。
\end{mydefinition}
\origin{原笔记第2.3节，定义“指数分布”；源码第1051行。}

\begin{proposition}{指数分布的分布函数与尾概率}{exponential-cdf}
对连续型随机变量 $X$，$X\sim\operatorname{Exp}(\lambda)$ 当且仅当
\[
F(t)=
\begin{cases}
1-e^{-\lambda t},&t\geq0,\\
0,&t<0.
\end{cases}
\]
等价地，其尾概率为
\[
P(X>t)=
\begin{cases}
e^{-\lambda t},&t\geq0,\\
1,&t<0.
\end{cases}
\]
\end{proposition}
\origin{原笔记第2.3节，命题“指数分布的分布函数与尾概率”；源码第1073行。}

\begin{mytheorem}{指数分布的无记忆性}{exponential-memorylessness}
设 $X$ 为取值非负的连续型随机变量，且对每个 $s>0$，$P(X>s)>0$。则 $X$ 服从某个参数 $\lambda>0$ 的指数分布，当且仅当对任意 $s,t>0$，
\[
P(X>s+t\mid X>s)=P(X>t).
\]
这一性质也称为\textbf{无效性}。
\end{mytheorem}
\origin{原笔记第2.3节，定理“指数分布的无记忆性”；源码第1108行。}

\begin{mylemma}{连续乘法方程}{continuous-multiplicative-equation}
若连续函数 $G:(0,\infty)\to(0,\infty)$ 满足
\[
G(s+t)=G(s)G(t),\qquad s,t>0,
\]
则存在 $b\in\mathbb R$，使 $G(t)=e^{bt}$ 对所有 $t>0$ 成立。
\end{mylemma}
\origin{原笔记第2.3节，引理“连续乘法方程”；源码第1151行。}

\subsection{例题：邮局中最后办完事情的人}\label{note:c2:e14}
\basis{原笔记第2.3节，例“邮局中最后办完事情的人”；源码\texttt{chapters/chapter2.tex}第1200行。}
邮局有两位职员。$A$ 到达时，$B$、$C$ 正分别接受服务；其中一人办完后，$A$ 立即开始接受服务。假定各人的服务时间相互独立，均服从同一 $\operatorname{Exp}(\lambda)$ 分布，求 $A$ 在三人中最后办完的概率。

由无记忆性，$A$ 开始服务时，另一位尚未办完者的剩余服务时间 $R$ 仍服从 $\operatorname{Exp}(\lambda)$。$A$ 的服务时间 $T$ 也服从同一分布，且与 $R$ 独立。此时已有一人办完，因此
\[
\{A\text{ 最后办完}\}=\{T>R\}.
\]
独立同分布给出对称性，
\[
P(T>R)=P(R>T).
\]
两者独立且为连续型，$P(T=R)=0$，所以
\[
P(T>R)+P(R>T)=1,
\qquad
P(A\text{ 最后办完})=\frac12.
\]

\textbf{注：}\quad 该模型要求服务时间相互独立；仅有相同的指数边缘分布，并不足以推出上述比较结论。

\subsection{例题：第一位顾客的到达时刻}\label{note:c2:e15}
\basis{原笔记第2.3节，例“第一位顾客的到达时刻”；源码\texttt{chapters/chapter2.tex}第1221行。}
设 $N(t)$ 表示商场从开门时刻 $0$ 到时刻 $t$ 的顾客到达人数，且对每个 $t>0$，
\[
N(t)\sim\mathcal P(\lambda t),\qquad \lambda>0.
\]
令 $S$ 为第一位顾客的到达时刻。对 $t>0$，
\[
\{S>t\}=\{N(t)=0\},
\]
所以
\[
\begin{aligned}
P(S>t)
&=P(N(t)=0)\\
&=e^{-\lambda t}.
\end{aligned}
\]
结合到达时刻非负可知，$S\sim\operatorname{Exp}(\lambda)$。

\begin{mydefinition}{Gamma 函数与 Gamma 分布}{gamma-function-and-distribution}
对 $\alpha>0$，定义
\[
\Gamma(\alpha)=\int_0^\infty x^{\alpha-1}e^{-x}\,dx,
\]
称为\textbf{Gamma 函数}或 $\Gamma$ 函数。

设 $\alpha>0$、$\lambda>0$。若随机变量 $X$ 的密度可以取为
\[
f(x)=
\begin{cases}
\dfrac{\lambda^\alpha}{\Gamma(\alpha)}x^{\alpha-1}e^{-\lambda x},
&x>0,\\[0.5em]
0,&x\leq0,
\end{cases}
\]
则称 $X$ 服从参数为 $(\alpha,\lambda)$ 的\textbf{Gamma 分布}，记为 $X\sim\Gamma(\alpha,\lambda)$。其中 $\alpha$ 为形状参数，$\lambda$ 为率参数；正半轴上的密度也可写成
\[
f(x)=\frac{(\lambda x)^{\alpha-1}}{\Gamma(\alpha)}
\lambda e^{-\lambda x},\qquad x>0.
\]
\end{mydefinition}
\origin{原笔记第2.3节，定义“Gamma 函数与 Gamma 分布”；源码第1243行。}

\begin{proposition}{Gamma 函数的基本性质}{gamma-function-properties}
Gamma 函数满足：
\begin{enumerate}
    \item 对 $\alpha>0$，$\Gamma(\alpha+1)=\alpha\Gamma(\alpha)$；
    \item $\Gamma(1)=1$，$\Gamma(1/2)=\sqrt\pi$；
    \item 对正整数 $n$，$\Gamma(n)=(n-1)!$。
\end{enumerate}
\end{proposition}
\origin{原笔记第2.3节，命题“Gamma 函数的基本性质”；源码第1280行。}

\begin{proposition}{Gamma 密度的归一化}{gamma-density-normalization}
上述 Gamma 分布定义中的 $f$ 确实是概率密度函数。
\end{proposition}
\origin{原笔记第2.3节，命题“Gamma 密度的归一化”；源码第1363行。}

\subsection{例题：\texorpdfstring{第 $n$ 位顾客的到达时刻}{第 n 位顾客的到达时刻}}\label{note:c2:e16}
\basis{原笔记第2.3节，例“第 $n$ 位顾客的到达时刻”；源码\texttt{chapters/chapter2.tex}第1381行。}
仍设商场在 $[0,t]$ 内的到达人数满足
\[
N(t)\sim\mathcal P(\lambda t),\qquad t>0.
\]
令 $S_n$ 为第 $n$ 位顾客的到达时刻，$n\geq1$。证明 $S_n\sim\Gamma(n,\lambda)$。

先求分布函数，再求导。对 $t>0$，
\[
\{S_n\leq t\}=\{N(t)\geq n\}.
\]
所以
\[
\begin{aligned}
F_{S_n}(t)
&=P(N(t)\geq n)\\
&=1-P(N(t)\leq n-1)\\
&=1-e^{-\lambda t}\sum_{k=0}^{n-1}\frac{(\lambda t)^k}{k!}.
\end{aligned}
\]
$t\leq0$ 时取 $F_{S_n}(t)=0$，与右侧在 $0$ 处的极限一致。

对 $t>0$ 求导，将乘积的两部分分别列出：
\[
\begin{aligned}
f_{S_n}(t)
&=e^{-\lambda t}\sum_{k=0}^{n-1}
\frac{\lambda^{k+1}t^k}{k!}
-e^{-\lambda t}\sum_{k=1}^{n-1}
\frac{\lambda^kt^{k-1}}{(k-1)!}\\
&=e^{-\lambda t}\left[
\sum_{k=0}^{n-1}\frac{\lambda^{k+1}t^k}{k!}
-\sum_{j=0}^{n-2}\frac{\lambda^{j+1}t^j}{j!}
\right]\\
&=\frac{\lambda^nt^{n-1}}{(n-1)!}e^{-\lambda t}\\
&=\frac{(\lambda t)^{n-1}}{\Gamma(n)}
\lambda e^{-\lambda t}.
\end{aligned}
\]
前两式中两个求和的共有项逐项抵消；$n=1$ 时第二个求和为空和，按 $0$ 处理。取 $t\leq0$ 时密度为 $0$，便得到 Gamma 密度，故
\[
S_n\sim\Gamma(n,\lambda).
\]
特别地，$S_1\sim\Gamma(1,\lambda)=\operatorname{Exp}(\lambda)$，与第一位顾客的结果一致。

\begin{mydefinition}{贝塔函数与贝塔分布}{beta-function-and-distribution}
设 $a,b>0$，定义\textbf{贝塔函数}
\[
B(a,b)=\int_0^1x^{a-1}(1-x)^{b-1}\,dx.
\]
若随机变量 $X$ 的密度可以取为
\[
f(x)=
\begin{cases}
\dfrac{x^{a-1}(1-x)^{b-1}}{B(a,b)},&0<x<1,\\[0.5em]
0,&\text{其他},
\end{cases}
\]
则称 $X$ 服从参数为 $(a,b)$ 的\textbf{贝塔分布}。
\end{mydefinition}
\origin{原笔记第2.3节，定义“贝塔函数与贝塔分布”；源码第1429行。}

\begin{proposition}{贝塔函数的基本性质}{beta-function-properties}
对 $a,b>0$，
\[
B(a,b)=B(b,a),
\qquad
B(a,b)=\frac{\Gamma(a)\Gamma(b)}{\Gamma(a+b)}.
\]
\end{proposition}
\origin{原笔记第2.3节，命题“贝塔函数的基本性质”；源码第1452行。}

\begin{mydefinition}{正态分布}{normal-distribution}
设 $\mu\in\mathbb R$、$\sigma>0$。若随机变量 $X$ 的密度为
\[
f(x)=\frac1{\sqrt{2\pi}\sigma}
\exp\left(-\frac{(x-\mu)^2}{2\sigma^2}\right),
\qquad x\in\mathbb R,
\]
则称 $X$ 服从参数为 $(\mu,\sigma^2)$ 的\textbf{正态分布}或\textbf{高斯分布}，记为 $X\sim N(\mu,\sigma^2)$。特别地，$N(0,1)$ 称为\textbf{标准正态分布}。
\end{mydefinition}
\origin{原笔记第2.3节，定义“正态分布”；源码第1504行。}

\begin{proposition}{正态密度的归一化}{normal-density-normalization}
上述 $f$ 非负且在实数轴上的积分为 $1$，因而确实是概率密度。
\end{proposition}
\origin{原笔记第2.3节，命题“正态密度的归一化”；源码第1514行。}

\begin{proposition}{正态密度的图像性质}{normal-density-shape}
正态密度关于 $x=\mu$ 对称，在 $\mu$ 处唯一达到最大值 $1/(\sqrt{2\pi}\sigma)$；以 $x$ 轴为水平渐近线，拐点为 $x=\mu\pm\sigma$，图像呈钟形。

固定 $\sigma$ 改变 $\mu$ 时，图像水平平移；固定 $\mu$ 增大 $\sigma$ 时，图像展宽、峰值降低，减小 $\sigma$ 时相反。
\end{proposition}
\origin{原笔记第2.3节，命题“正态密度的图像性质”；源码第1535行。}

标准正态分布的密度和分布函数分别记为
\[
\varphi(x)=\frac1{\sqrt{2\pi}}e^{-x^2/2},
\qquad
\Phi(t)=\int_{-\infty}^t\varphi(x)\,dx.
\]

\begin{proposition}{标准正态分布的对称公式}{standard-normal-symmetry}
对任意 $x\in\mathbb R$，
\[
\Phi(x)+\Phi(-x)=1.
\]
\end{proposition}
\origin{原笔记第2.3节，命题“标准正态分布的对称公式”；源码第1577行。}

\subsection{例题：对两侧尾部作符号翻转}\label{note:c2:e17}
\basis{原笔记第2.3节，例“对两侧尾部作符号翻转”；源码\texttt{chapters/chapter2.tex}第1600行。}
设 $X\sim N(0,1)$，$a>0$，定义
\[
Y=
\begin{cases}
X,&|X|\leq a,\\
-X,&|X|>a.
\end{cases}
\]
证明 $Y\sim N(0,1)$。

先写出分布函数：
\[
\begin{aligned}
F_Y(y)
&=P(X\leq y,\ |X|\leq a)
+P(-X\leq y,\ |X|>a)\\
&=P(X\leq y,\ |X|\leq a)
+P(X\geq-y,\ |X|>a).
\end{aligned}
\]
分三种情况计算。

（1）若 $y>a$，三个相应区间互不相交，
\[
\begin{aligned}
F_Y(y)
&=P(|X|\leq a)+P(-y\leq X<-a)+P(X>a)\\
&=P(X\geq-y)=\Phi(y).
\end{aligned}
\]

（2）若 $-a\leq y\leq a$，利用两侧尾概率相等，
\[
\begin{aligned}
F_Y(y)
&=P(-a\leq X\leq y)+P(X>a)\\
&=P(-a\leq X\leq y)+P(X<-a)\\
&=P(X\leq y)=\Phi(y).
\end{aligned}
\]

（3）若 $y<-a$，中间区间不贡献概率，
\[
F_Y(y)=0+P(X\geq-y)=\Phi(y).
\]
综上 $F_Y=\Phi$，故 $Y\sim N(0,1)$。

\begin{mytheorem}{正态分布的线性性质}{normal-affine-transformation}
若 $X\sim N(\mu,\sigma^2)$，$a,b\in\mathbb R$ 且 $a\ne0$，则
\[
aX+b\sim N(a\mu+b,a^2\sigma^2).
\]
\end{mytheorem}
\origin{原笔记第2.3节，定理“正态分布的线性性质”；源码第1651行。}

\begin{mycorollary}{正态分布的标准化}{normal-standardization}
若 $X\sim N(\mu,\sigma^2)$，则
\[
\frac{X-\mu}{\sigma}\sim N(0,1).
\]
\end{mycorollary}
\origin{原笔记第2.3节，推论“正态分布的标准化”；源码第1700行。}

\subsection{例题：三倍尺度区间的概率}\label{note:c2:e18}
\basis{原笔记第2.3节，例“三倍尺度区间的概率”；源码\texttt{chapters/chapter2.tex}第1727行。}
设 $X\sim N(\mu,\sigma^2)$，则
\[
\begin{aligned}
P(|X-\mu|<3\sigma)
&=P\left(-3<\frac{X-\mu}{\sigma}<3\right)\\
&=\Phi(3)-\Phi(-3)\\
&=2\Phi(3)-1\approx0.9973002.
\end{aligned}
\]
若先取 $\Phi(3)\approx0.9987$，再计算得到 $0.9974$，这是粗查表近似；较精细的概率约为 $99.73\%$。

\textbf{注：}\quad 这称为 $3\sigma$ 原则，可用于理解质量控制中的异常提示。但区间外仍有约 $0.27\%$ 的概率，并非概率为 $0$；一次观测落在区间外，不能单凭这一点就断言随机变量不服从该正态分布。

\subsection{例题：由两端比例估计录取分数线}\label{note:c2:e19}
\basis{原笔记第2.3节，例“由两端比例估计录取分数线”；源码\texttt{chapters/chapter2.tex}第1742行。}
共有 $10000$ 人报考，成绩大致服从正态分布；$90$ 分以上有 $359$ 人，$60$ 分以下有 $1151$ 人，现从高到低录取 $2500$ 人，估计最低录取分数。

设成绩模型为 $X\sim N(\mu,\sigma^2)$，将人数比例近似为概率。首先，
\[
\begin{aligned}
P(X\geq90)&\approx\frac{359}{10000}=0.0359,\\
\Phi\left(\frac{90-\mu}{\sigma}\right)&\approx0.9641.
\end{aligned}
\]
查表得到
\[
\frac{90-\mu}{\sigma}\approx1.8.
\]
另一方面，
\[
\begin{aligned}
\Phi\left(\frac{60-\mu}{\sigma}\right)
&\approx0.1151,\\
\Phi\left(\frac{\mu-60}{\sigma}\right)
&\approx1-0.1151=0.8849.
\end{aligned}
\]
所以 $(\mu-60)/\sigma\approx1.2$。联立两式，
\[
\begin{cases}
90-\mu\approx1.8\sigma,\\
\mu-60\approx1.2\sigma,
\end{cases}
\qquad
\sigma\approx10,\quad \mu\approx72.
\]

设最低录取分为 $c$。取前 $25\%$，意味着
\[
\begin{aligned}
P(X\geq c)&\approx0.25,\\
\Phi\left(\frac{c-72}{10}\right)&\approx0.75.
\end{aligned}
\]
按粗查表值 $(c-72)/10\approx0.68$，得到 $c\approx78.8$ 分。若用较精细的 $0.75$ 分位数 $0.67449$，在同一 $\mu\approx72,\sigma\approx10$ 的近似下，
\[
c\approx72+10\times0.67449\approx78.74\text{ 分}.
\]
$78.8$ 分是粗查表结果，不是精确分位数。

\clearpage\section{随机变量函数的分布}\label{k:transform}

\begin{mydefinition}{博雷尔函数}{borel-function}
若函数 $g:\mathbb R\to\mathbb R$ 满足
\[
g^{-1}(B)\in\mathcal B(\mathbb R),
\qquad \forall B\in\mathcal B(\mathbb R),
\]
则称 $g$ 为\textbf{博雷尔函数}，也称 $(\mathbb R,\mathcal B(\mathbb R))$ 上的可测函数。
\end{mydefinition}
\origin{原笔记第2.4节，定义“博雷尔函数”；源码第1803行。}

\subsection{例题：离散型随机变量的多项式变换}\label{note:c2:e20}
\basis{原笔记第2.4节，例“离散型随机变量的多项式变换”；源码\texttt{chapters/chapter2.tex}第1818行。}
已知 $X$ 的分布列为
\[
\begin{array}{c|ccccc}
X&-2&-1&0&1&2\\ \hline
P&0.2&0.1&0.1&0.3&0.3
\end{array},
\]
求 $Y=X^2+X$ 的分布列。

逐项计算并保留原来的概率：
\[
\begin{array}{c|ccccc}
X&-2&-1&0&1&2\\ \hline
Y=X^2+X&2&0&0&2&6\\
P&0.2&0.1&0.1&0.3&0.3
\end{array}.
\]
$Y=0$ 对应 $X=-1,0$，$Y=2$ 对应 $X=-2,1$，$Y=6$ 对应 $X=2$。合并后得到
\[
\begin{array}{c|ccc}
Y&0&2&6\\ \hline
P&0.2&0.5&0.3
\end{array}.
\]
三个概率之和为 $1$。

\subsection{例题：一般连续型随机变量的线性变换}\label{note:c2:e21}
\basis{原笔记第2.4节，例“一般连续型随机变量的线性变换”；源码\texttt{chapters/chapter2.tex}第1848行。}
设 $X$ 的密度为 $f_X$，求 $Y=aX+b$（$a\ne0$）的密度。

先按 $a$ 的符号求分布函数。$a>0$ 时，
\[
\begin{aligned}
F_Y(t)
&=P\left(X\leq\frac{t-b}{a}\right)\\
&=F_X\left(\frac{t-b}{a}\right).
\end{aligned}
\]
$a<0$ 时，利用 $X$ 无单点质量，
\[
\begin{aligned}
F_Y(t)
&=P\left(X\geq\frac{t-b}{a}\right)\\
&=1-F_X\left(\frac{t-b}{a}\right).
\end{aligned}
\]
分别对上式求导，合并得到一个密度版本
\[
\boxed{\quad f_Y(t)=\frac1{|a|}f_X\left(\frac{t-b}{a}\right).\quad}
\]

\begin{mydefinition}{分布函数的广义逆}{generalized-inverse-cdf}
设 $G$ 单调不减、右连续，且
\[
\lim_{x\to-\infty}G(x)=0,\qquad
\lim_{x\to+\infty}G(x)=1.
\]
对 $p\in(0,1)$，定义其\textbf{广义逆}
\[
G^{-1}(p)=\inf\{x\in\mathbb R:G(x)\geq p\}.
\]
这里不要求 $G$ 严格单调，也不要求它有通常意义的反函数。
\end{mydefinition}
\origin{原笔记第2.4节，定义“分布函数的广义逆”；源码第1887行。}

\begin{mylemma}{广义逆与分布函数的不等式}{generalized-inverse-equivalence}
对任意 $p\in(0,1)$ 和 $t\in\mathbb R$，
\[
G^{-1}(p)\leq t\quad\Longleftrightarrow\quad p\leq G(t).
\]
\end{mylemma}
\origin{原笔记第2.4节，引理“广义逆与分布函数的不等式”；源码第1906行。}

\begin{mytheorem}{由均匀分布构造给定分布}{inverse-transform-construction}
设 $X\sim U(0,1)$，$G$ 及其广义逆满足前述定义。令
\[
Y=G^{-1}(X),
\]
则 $Y$ 的分布函数就是 $G$。
\end{mytheorem}
\origin{原笔记第2.4节，定理“由均匀分布构造给定分布”；源码第1935行。}

\begin{mytheorem}{分支反函数的密度变换公式}{inverse-branches-density}
设 $X$ 有密度 $f_X$，$g$ 为博雷尔函数，$Y=g(X)$。设存在有限个或可数个互不相交的开区间 $I_i$，满足
\[
P\left(X\in\bigcup_i I_i\right)=1.
\]
在每个 $I_i$ 上，$g$ 连续、严格单调，像 $J_i=g(I_i)$ 为开区间；其反函数 $h_i:J_i\to I_i$ 连续可微。这里 $h_i$ 称为第 $i$ 个\textbf{反函数分支}，即
\[
g(h_i(y))=y,\qquad y\in J_i.
\]
则 $Y$ 有一个密度版本
\[
f_Y(y)=\sum_{\{i:y\in J_i\}}f_X(h_i(y))\,|h_i'(y)|.
\]
没有分支贡献时，该和取 $0$。非负级数在几乎处处有限；在可能出现的零测异常点上，可另赋有限值而不改变分布。
\end{mytheorem}
\origin{原笔记第2.4节，定理“分支反函数的密度变换公式”；源码第1977行。}

\subsection{例题：用反函数公式求线性变换的密度}\label{note:c2:e22}
\basis{原笔记第2.4节，例“用反函数公式求线性变换的密度”；源码\texttt{chapters/chapter2.tex}第1993行。}
设 $Y=aX+b$，$a\ne0$。变换在整个实轴上只有一个反函数分支
\[
h(y)=\frac{y-b}{a},\qquad |h'(y)|=\frac1{|a|}.
\]
所以
\[
f_Y(y)=\frac1{|a|}f_X\left(\frac{y-b}{a}\right).
\]
这与分布函数法的结果一致。

\subsection{例题：平方变换与不同分支的取值范围}\label{note:c2:e23}
\basis{原笔记第2.4节，例“平方变换与不同分支的取值范围”；源码\texttt{chapters/chapter2.tex}第2005行。}
设随机变量 $X$ 的密度为
\[
f_X(x)=
\begin{cases}
\dfrac14|x|,&-2\leq x<0,\\
x,&0\leq x\leq1,\\
0,&\text{其他}.
\end{cases}
\]
求 $Y=X^2$ 的密度。

负半轴的概率为 $\int_{-2}^0|x|/4\,dx=1/2$，正半轴的概率为 $\int_0^1x\,dx=1/2$，所以给定密度归一化。由 $-2\leq X\leq1$ 以概率 $1$ 成立，$Y$ 可限制在 $[0,4]$。

对 $y>0$，两个反函数分支为 $h_1(y)=\sqrt y$ 和 $h_2(y)=-\sqrt y$，它们的导数绝对值都为 $1/(2\sqrt y)$，因此
\[
f_Y(y)=\frac{f_X(\sqrt y)+f_X(-\sqrt y)}{2\sqrt y}.
\]
当 $0<y<1$ 时，正、负两个原像均落在原密度的非零区间内，故
\[
f_Y(y)=\frac{\sqrt y+\sqrt y/4}{2\sqrt y}=\frac58.
\]
当 $1<y<4$ 时，正原像超过 $1$，只有负分支贡献，故
\[
f_Y(y)=\frac{0+\sqrt y/4}{2\sqrt y}=\frac18.
\]
当 $y<0$ 或 $y>4$ 时，概率密度可取 $0$。在 $y=0$ 处不能代入含 $1/\sqrt y$ 的公式；但
\[
P(Y=0)=P(X=0)=0,
\]
可另指定该点的密度值。取密度版本为
\[
f_Y(y)=
\begin{cases}
\dfrac58,&0\leq y\leq1,\\
\dfrac18,&1<y\leq4,\\
0,&\text{其他}.
\end{cases}
\]
检查全积分：
\[
\int_{\mathbb R}f_Y(y)\,dy=\frac58+\frac18(4-1)=1.
\]

\subsection{例题：均匀随机变量的余弦变换}\label{note:c2:e24}
\basis{原笔记第2.4节，例“均匀随机变量的余弦变换”；源码\texttt{chapters/chapter2.tex}第2050行。}
设 $X\sim U(0,6\pi)$，求 $Y=\cos X$ 的密度。

$Y\in[-1,1]$。对 $-1<y<1$，令 $\theta=\arccos y\in(0,\pi)$。区间 $(0,6\pi)$ 包含三个完整周期，每个周期上余弦先减后增，所以所有原像恰为六个互不相同的点
\[
h_{k,1}(y)=\arccos y+2k\pi,\qquad
h_{k,2}(y)=2\pi-\arccos y+2k\pi,
\quad k=0,1,2.
\]
六个反函数分支的导数绝对值均为 $1/\sqrt{1-y^2}$，原密度均为 $1/(6\pi)$，于是
\begin{align*}
f_Y(y)
&=\sum_{k=0}^2
\left[
\frac1{6\pi\sqrt{1-y^2}}+
\frac1{6\pi\sqrt{1-y^2}}
\right]\\
&=\frac1{\pi\sqrt{1-y^2}},\qquad -1<y<1.
\end{align*}
在 $y=\pm1$ 处，原像为有限个点，概率为零；分支导数公式在端点不适用，可将密度另取为 $0$。所以一个完整版本为
\[
f_Y(y)=
\begin{cases}
\dfrac1{\pi\sqrt{1-y^2}},&-1<y<1,\\
0,&\text{其他}.
\end{cases}
\]
虽然密度在两端附近无界，但积分收敛，并且
\[
\int_{-1}^1\frac{dy}{\pi\sqrt{1-y^2}}
=\frac1\pi\left[\arcsin y\right]_{-1}^1=1.
\]

\subsection{例题：倒数的小数部分与可数个反函数分支}\label{note:c2:e25}
\basis{原笔记第2.4节，例“倒数的小数部分与可数个反函数分支”；源码\texttt{chapters/chapter2.tex}第2084行。}
设 $X\sim U(0,1)$，令 $Y$ 为 $1/X$ 的小数部分。记 $\lfloor u\rfloor$ 为不超过 $u$ 的最大整数，即向下取整，因而
\[
Y=\frac1X-\left\lfloor\frac1X\right\rfloor\in[0,1).
\]
在零概率事件 $\{X\notin(0,1)\}$ 上，可另定义 $Y=0$。

对 $0<y<1$，若 $Y=y$，则唯一的整数部分 $k=\lfloor1/X\rfloor$ 满足 $k\geq1$，且
\[
\frac1X=k+y
\quad\Longleftrightarrow\quad X=h_k(y)=\frac1{k+y}.
\]
反过来，每个 $k\geq1$ 给出的 $h_k(y)\in(0,1)$ 都是有效原像。各分支对应互不相交的区间
\[
I_k=\left(\frac1{k+1},\frac1k\right).
\]
这些区间只遗漏可数个点 $1/k$，其总概率为零，满足分支公式的全概率条件。由于
\[
|h_k'(y)|=\frac1{(k+y)^2},
\]
得到一个密度版本
\[
f_Y(y)=
\begin{cases}
\displaystyle\sum_{k=1}^{\infty}\frac1{(k+y)^2},&0\leq y<1,\\
0,&\text{其他}.
\end{cases}
\]
这里在 $y=0$ 处的值是另选的密度值；该点的原像是可数集，$P(Y=0)=0$。对 $0\leq y<1$，各项不超过 $1/k^2$，所以级数收敛。

利用非负求和与积分交换，再作望远镜求和，核验
\begin{align*}
\int_0^1 f_Y(y)\,dy
&=\sum_{k=1}^{\infty}\int_0^1\frac{dy}{(k+y)^2}\\
&=\sum_{k=1}^{\infty}\left(\frac1k-\frac1{k+1}\right)
=1.
\end{align*}

\clearpage\section*{第六至第九周作业}\addcontentsline{toc}{section}{第六至第九周作业}
% !TEX root = 概率论期中复习.tex
\section{第六周：分布函数与端点概率}

\subsection{作业第1题：取出号码的中间值}
\label{ex:w06-h1}
不透明的口袋中有5个球，编号为1、2、3、4、5号。现从袋中不放回任取3个球。设$X$表示取出的3个球中编号居中的号码，求随机变量$X$的分布函数$F(t)$。
\basis{原笔记第1章第1节“古典概率模型”、第2章第2节命题“分布列与分布函数相互确定”；见\hyperref[k:discrete]{离散型分布}。}
\paragraph{思路}固定中间号码$j$后，要在其左右各选一个号码；先得分布列，再累加。
\begin{proof}[解答]
全部无序取球结果共有$\binom53=10$种，且等可能。最小号码1不可能居中，最大号码5也不可能居中，所以$X$只取2、3、4。

若中间号码为$j$，则必须取到$j$，在$\{1,\ldots,j-1\}$中选一个号码，在$\{j+1,\ldots,5\}$中选一个号码。这给出$(j-1)(5-j)$种不同结果，且每种中间号码为$j$的结果都恰好由此构成。因此
\[
\P(X=j)=\frac{(j-1)(5-j)}{10},\qquad j=2,3,4.
\]
分布列为
\[
\begin{array}{c|ccc}
j&2&3&4\\\hline
\P(X=j)&3/10&2/5&3/10
\end{array}.
\]
由$F(t)=\P(X\leq t)$逐段累加，得到
\[
F(t)=\begin{cases}
0,&t<2,\\
3/10,&2\leq t<3,\\
7/10,&3\leq t<4,\\
1,&t\geq4.
\end{cases}
\]
\end{proof}
\paragraph{答案}$X$在2、3、4处的概率分别为$3/10,2/5,3/10$，分布函数如上。

\subsection{作业第2题：随机选盒后取球}
\label{ex:w06-h2}
有3个外表一样的不透明的盒子。第一个盒子里装有1个红球和2个蓝球，第二个盒子里装有2个红球和2个蓝球，第3个盒子里装有3个红球和1个蓝球。现任取一个盒子，再从该盒子中任取两个球。设$X$表示所取到的红球个数。
\begin{enumerate}
\item 求概率$\P(X=0),\P(X=1),\P(X=2)$。
\item 求随机变量$X$的分布函数。
\end{enumerate}
\basis{原笔记第1章第3节“全概率公式”、第2章第2节“超几何分布”与“分布列与分布函数相互确定”。}
\paragraph{思路}按盒子分类，在每个盒内按无放回抽取计数，然后取三种条件分布的平均。
\begin{proof}[解答]
记$B_i$为“取到第$i$个盒子”，则$\P(B_i)=1/3$。题中的“任取两个球”按一次取出两个不同球理解。
\begin{enumerate}
\item 从第一个盒子中抽取两个球有$\binom32=3$种；第二、三个盒子均有$\binom42=6$种。分别按红球个数计数，得到
\[
\begin{array}{c|ccc}
&X=0&X=1&X=2\\\hline
B_1&\binom22/\binom32&\binom11\binom21/\binom32&0\\[0.3em]
B_2&\binom22/\binom42&\binom21\binom21/\binom42&\binom22/\binom42\\[0.3em]
B_3&0&\binom31\binom11/\binom42&\binom32/\binom42
\end{array}.
\]
表内各行表示给定相应盒子后的条件概率。由全概率公式，
\begin{align*}
\P(X=0)&=\frac13\left(\frac13+\frac16+0\right)=\frac16,\\
\P(X=1)&=\frac13\left(\frac23+\frac46+\frac36\right)=\frac{11}{18},\\
\P(X=2)&=\frac13\left(0+\frac16+\frac36\right)=\frac29.
\end{align*}
三者之和为$3/18+11/18+4/18=1$。
\item 将以上分布列累加，得到
\[
F(t)=\begin{cases}
0,&t<0,\\
1/6,&0\leq t<1,\\
7/9,&1\leq t<2,\\
1,&t\geq2.
\end{cases}
\]
\end{enumerate}
\end{proof}
\paragraph{答案}(1) $1/6,11/18,2/9$；(2) 分布函数如上。

\subsection{作业第3题：分段分布函数的跳跃与连续部分}
\label{ex:w06-h3}
设随机变量$X$的分布函数为
\[
F(t)=\begin{cases}
0,&t<0,\\
\dfrac1{10}(t+1),&0\leq t<1,\\
\dfrac15(t+1),&1\leq t<3,\\
\dfrac{11}{12},&3\leq t<4,\\
1,&t\geq4.
\end{cases}
\]
\begin{enumerate}
\item 求概率$\P(X=i)$，$i=0,1,3,4$，并判断概率$\P(\{X=0\}\cup\{X=1\}\cup\{X=3\}\cup\{X=4\})$是否等于1。
\item 求概率$\P(0<X<1),\P(1<X<3),\P(3<X<4)$。
\end{enumerate}
\basis{原笔记第2章第1节定理“分布函数的基本性质”与“用分布函数计算概率”；见\hyperref[k:cdf]{分布函数与端点概率}。}
\paragraph{思路}单点概率是跳跃量$F(a)-F(a^-)$；开区间概率为$F(b^-)-F(a)$，不能把端点跳跃误算进去。
\begin{proof}[解答]
\begin{enumerate}
\item 根据每段的端点值和左极限，
\begin{align*}
\P(X=0)&=F(0)-F(0^-)=\frac1{10}-0=\frac1{10},\\
\P(X=1)&=F(1)-F(1^-)=\frac25-\frac15=\frac15,\\
\P(X=3)&=F(3)-F(3^-)=\frac{11}{12}-\frac45=\frac7{60},\\
\P(X=4)&=F(4)-F(4^-)=1-\frac{11}{12}=\frac1{12}.
\end{align*}
这四个事件两两互斥，故其并的概率为
\[
\frac1{10}+\frac15+\frac7{60}+\frac1{12}=\frac12\ne1.
\]
因此，$X$并不以概率1只取这四个值，不能仅看分段的断点就把它当作四点离散分布。
\item 由开区间公式，
\begin{align*}
\P(0<X<1)&=F(1^-)-F(0)=\frac15-\frac1{10}=\frac1{10},\\
\P(1<X<3)&=F(3^-)-F(1)=\frac45-\frac25=\frac25,\\
\P(3<X<4)&=F(4^-)-F(3)=\frac{11}{12}-\frac{11}{12}=0.
\end{align*}
三个开区间的概率之和为$1/2$，与四个单点的总概率相加为1，与全部分布相符。
\end{enumerate}
\end{proof}
\paragraph{答案}(1) $1/10,1/5,7/60,1/12$，四点总概率为$1/2$；(2) $1/10,2/5,0$。

\subsection{作业第4题：由单点概率确定分布函数参数}
\label{ex:w06-h4}
设随机变量$X$的分布函数为
\[
F(t)=\begin{cases}
0,&t<-2,\\
\dfrac15,&-2\leq t<-1,\\
Bt+A,&-1\leq t<0,\\
1-\dfrac12e^{-t},&t\geq0,
\end{cases}
\]
其中$A$和$B$是未知常数，并且$\P(X=-1)=1/5$，$\P(X=0)=0$。
\begin{enumerate}
\item 求常数$A$和$B$。
\item 求概率$\P(-1\leq X\leq-1/2)$。
\end{enumerate}
\basis{原笔记第2章第1节定理“分布函数的基本性质”及区间概率公式；见\hyperref[k:cdf]{分布函数与端点概率}。}
\paragraph{思路}以0点连续和$-1$点跳跃量列方程，闭区间左端要取左极限。
\begin{proof}[解答]
\begin{enumerate}
\item 在0处，$F(0)=1/2$、$F(0^-)=A$，所以
\[
0=\P(X=0)=F(0)-F(0^-)=\frac12-A,
\]
得到$A=1/2$。在$-1$处，$F(-1)=A-B$、$F(-1^-)=1/5$，因此
\[
\frac15=\P(X=-1)=A-B-\frac15.
\]
由此$A-B=2/5$，得到$B=1/10$。

所得函数在$[-1,0)$上斜率为正，在$[0,\infty)$上也单调增加；它在0处连续，在$-2,-1$处向上跳跃，且两端极限为0、1，符合分布函数的性质。
\item 因区间包含左端点$-1$，
\begin{align*}
\P\left(-1\leq X\leq-\frac12\right)
&=F\left(-\frac12\right)-F(-1^-)\\
&=\left(-\frac1{20}+\frac12\right)-\frac15=\frac14.
\end{align*}
\end{enumerate}
\end{proof}
\paragraph{答案}$A=1/2$，$B=1/10$；所求概率为$1/4$。

\subsection{作业第5题（思考题，选做）：分布函数的一致收敛}
\label{ex:w06-h5}
设$F(t),F_n(t)$，$n\geq1$都是分布函数。已知$F(t)$在$\R$上连续，并且对任意$t\in\R$，有
\[
\lim_{n\to+\infty}F_n(t)=F(t).
\]
求证：当$n\to+\infty$时，$F_n(t)$在$\R$上一致收敛于$F(t)$。
\basis{原笔记第2章第1节“分布函数的基本性质”（单调性与两端极限）；连续函数在闭有界区间上的一致连续性用于有限分点。}
\paragraph{思路}先用两端极限控制左右尾部，再在中间闭区间上取有限分点，借助单调性把分点上的收敛控制推广到整个实轴。
\begin{proof}
一致收敛的目标是：对任意$\varepsilon>0$，找到与$t$无关的$N$，使$n\geq N$时所有$t\in\R$均满足$|F_n(t)-F(t)|<\varepsilon$。

固定$\varepsilon>0$，令$\delta=\varepsilon/3$。利用$F$在两端的极限，选取$M>0$使
\[
F(-M)<\delta,\qquad 1-F(M)<\delta.
\]
因为$F$在$[-M,M]$上连续，所以在该闭区间上一致连续。于是可以选取有限分点
\[
-M=t_0<t_1<\cdots<t_r=M,
\]
使每个相邻分点满足$F(t_{j+1})-F(t_j)<\delta$。具体地，先取使$|s-t|$充分小时$|F(s)-F(t)|<\delta$的统一长度，再将$[-M,M]$等分到每段短于该长度。

每个固定分点$t_j$上都有$F_n(t_j)\to F(t_j)$。分点只有有限个，因此存在统一的$N$使$n\geq N$时
\[
|F_n(t_j)-F(t_j)|<\delta,\qquad 0\leq j\leq r.
\]
以下固定$n\geq N$，分三类控制$t$。
\begin{enumerate}
\item 若$-M\leq t\leq M$，则$t$位于某个$[t_j,t_{j+1}]$中。$F_n,F$均单调不减，故
\begin{align*}
F_n(t)-F(t)
&\leq F_n(t_{j+1})-F(t_j)\\
&=[F_n(t_{j+1})-F(t_{j+1})]+[F(t_{j+1})-F(t_j)]<2\delta,\\
F(t)-F_n(t)
&\leq F(t_{j+1})-F_n(t_j)\\
&=[F(t_{j+1})-F(t_j)]+[F(t_j)-F_n(t_j)]<2\delta.
\end{align*}
所以$|F_n(t)-F(t)|<2\delta$。
\item 若$t<-M$，利用概率非负与单调性，
\[
0\leq F(t)\leq F(-M)<\delta,
\qquad 0\leq F_n(t)\leq F_n(-M)<F(-M)+\delta<2\delta.
\]
因此$|F_n(t)-F(t)|<2\delta$。
\item 若$t>M$，同理
\[
0\leq1-F(t)\leq1-F(M)<\delta,
\qquad0\leq1-F_n(t)\leq1-F_n(M)<2\delta.
\]
从而$|F_n(t)-F(t)|<2\delta$。
\end{enumerate}
三个范围覆盖$\R$。于是$n\geq N$时，
\[
\sup_{t\in\R}|F_n(t)-F(t)|\leq2\delta=\frac{2\varepsilon}{3}<\varepsilon,
\]
证明了在整个实轴上的一致收敛。
\end{proof}
\paragraph{答案}$\displaystyle\sup_{t\in\R}|F_n(t)-F(t)|\to0$。
% !TEX root = 概率论期中复习.tex
\section{第七周：离散分布与密度计算}

\subsection{教材第52页第1题}\label{ex:w07-b52-1}
设某市不同路口的红绿灯独立工作，每个路口的红灯、绿灯出现的时间相同。用 $X$ 表示一辆汽车首次遇到红灯时已经经过的路口数。求 $X$ 的概率分布。
\basis{原笔记第2章第2节“几何分布”；第1章第4节“独立试验”。}
\paragraph{思路} 将经过一个绿灯路口看成一次失败，第一次遇红灯看成首次成功；先明确计数是否包含当前红灯路口。
\begin{proof}[解答]
按“已经经过”指此前已通过的路口理解，$X=k$ 表示前 $k$ 个路口都是绿灯，而第 $k+1$ 个路口是红灯。每处红、绿灯概率各为 $1/2$，不同路口独立，所以
\[
\P(X=k)=\left(\frac12\right)^k\frac12=2^{-(k+1)},\qquad k=0,1,2,\ldots.
\]
这些概率之和为 $\sum_{k=0}^{\infty}2^{-(k+1)}=1$，故给出了完整分布。

若将首次遇到红灯的那个路口也计入，则相应变量为 $N=X+1$，其分布为 $\P(N=n)=2^{-n}$，$n=1,2,\ldots$。两种计数仅相差 $1$。
\end{proof}
\paragraph{答案} 已通过的路口数满足 $\P(X=k)=2^{-(k+1)}$，$k\geq0$；包含当前路口的计数服从参数 $1/2$ 的几何分布。

\subsection{教材第52页第2题}\label{ex:w07-b52-2}
用 $\xi$ 表示成功概率为 $p$（$0\leq p\leq1$）的伯努利实验中第一个游程（即连续成功或连续失败）的长度。求 $\xi$ 的分布。

作业说明：用 $1$ 表示成功、$0$ 表示失败。对 $11000111$，第一个游程为最左端的 $11$，长度为 $2$；对 $010000111$，第一个游程为最左端的 $0$，长度为 $1$；对 $00011100$，第一个游程为最左端的 $000$，长度为 $3$。
\basis{原笔记第1章第4节“独立试验”；第2章第2节“几何分布”。}
\paragraph{思路} 按第一个结果是成功还是失败，将长度为 $k$ 的事件拆为两个不交事件。
\begin{proof}[解答]
先设 $0<p<1$，记 $q=1-p$。$\xi=k$ 当且仅当出现下列两种情形之一：前 $k$ 次均成功、第 $k+1$ 次失败；前 $k$ 次均失败、第 $k+1$ 次成功。由独立性及有限可加性，
\[
\P(\xi=k)=p^kq+q^kp,\qquad k=1,2,\ldots.
\]
同时，$\xi>k$ 等价于前 $k+1$ 次结果全部相同，故
\[
\P(\xi>k)=p^{k+1}+q^{k+1}\longrightarrow0.
\]
因而 $\xi$ 以概率 $1$ 有限，所列概率构成完整分布。也可直接核对
\[
\sum_{k=1}^{\infty}(p^kq+q^kp)=q\frac{p}{1-p}+p\frac{q}{1-q}=1.
\]
若 $p=0$ 或 $p=1$，所有试验的结果几乎必然相同，第一个游程不会结束。此时应允许取值 $\xi=\infty$，且 $\P(\xi=\infty)=1$；它不是取有限实值的随机变量。
\end{proof}
\paragraph{答案} $0<p<1$ 时，$\P(\xi=k)=p^k(1-p)+(1-p)^kp$，$k\geq1$；$p=0,1$ 时游程无限长。

\subsection{教材第53页第3题}\label{ex:w07-b53-3}
连续射击某个目标，直至第一次击中目标或者射击次数达到 $5$ 次就停止。假定一次射击击中目标的概率为 $p$（$0\leq p\leq1$）且各次射击结果独立。用 $X$ 表示停止时的射击次数。求 $X$ 的概率分布。
\basis{原笔记第2章第2节“几何分布”；第1章第4节“独立试验”。}
\paragraph{思路} 前四次停止必须因为当次首次击中；第五次停止则只要求前四次未击中。
\begin{proof}[解答]
记 $q=1-p$。对 $k=1,2,3,4$，停止于第 $k$ 次等价于前 $k-1$ 次未中、第 $k$ 次击中，故
\[
\P(X=k)=q^{k-1}p.
\]
而 $X=5$ 只要求前四次未击中。无论第五次是否击中，都会停止，所以
\[
\P(X=5)=q^4p+q^5=q^4.
\]
当 $p=0$ 时必有 $X=5$；当 $p=1$ 时必有 $X=1$。上述分布包含这两种端点情形。
\end{proof}
\paragraph{答案} $\P(X=k)=(1-p)^{k-1}p$（$1\leq k\leq4$），$\P(X=5)=(1-p)^4$。

\subsection{教材第53页第4题}\label{ex:w07-b53-4}
某种昆虫在单位时间内产卵的个数 $\xi$ 服从泊松分布 $P(\lambda)$，其中每个卵孵化出幼虫的概率为 $p$，各个卵是否孵化相互独立，并且虫子产卵的个数和产出卵是否孵化相互独立。以 $\eta$ 表示单位时间内幼虫的个数。求 $\eta$ 的概率分布。
\basis{原笔记第2章第2节“二项分布”“泊松分布”，例“泊松粒子数的独立捕捉”。}
\paragraph{思路} 给定产卵数后，幼虫数是二项分布；再对产卵数用全概率公式求和。
\begin{proof}[解答]
给定 $\xi=n$，每个卵独立地以概率 $p$ 孵化，因此
\[
\P(\eta=k\mid\xi=n)=\binom nkp^k(1-p)^{n-k},\qquad 0\leq k\leq n.
\]
对 $k=0,1,2,\ldots$，全概率公式给出
\begin{align*}
\P(\eta=k)
&=\sum_{n=k}^{\infty}\binom nkp^k(1-p)^{n-k}e^{-\lambda}\frac{\lambda^n}{n!}\\
&=e^{-\lambda}\frac{(\lambda p)^k}{k!}\sum_{n=k}^{\infty}\frac{[\lambda(1-p)]^{n-k}}{(n-k)!}\\
&=e^{-\lambda}\frac{(\lambda p)^k}{k!}e^{\lambda(1-p)}
=e^{-\lambda p}\frac{(\lambda p)^k}{k!}.
\end{align*}
这正是参数为 $\lambda p$ 的泊松分布列。当 $p=0$ 时，$\eta=0$ 几乎必然；当 $p=1$ 时，$\eta=\xi$。
\end{proof}
\paragraph{答案} $\eta\sim P(\lambda p)$，即 $\P(\eta=k)=e^{-\lambda p}(\lambda p)^k/k!$。

\subsection{教材第53页第5题}\label{ex:w07-b53-5}
设随机变量 $X$ 服从泊松分布 $P(\lambda)$。记 $p=\P(X\text{ 为奇数})$ 和 $q=\P(X\text{ 为偶数})$。证明
\[
p=e^{-\lambda}\sinh\lambda=\frac12(1-e^{-2\lambda}),\qquad
q=e^{-\lambda}\cosh\lambda=\frac12(1+e^{-2\lambda}).
\]
提示：注意 $q+p=1$ 和 $q-p=e^{-2\lambda}$。
\basis{原笔记第2章第2节“泊松分布”；指数函数的幂级数。}
\paragraph{思路} 沿教材提示，对分布列作交错求和，再联立奇、偶概率之和为 $1$。
\begin{proof}
$X$ 只取非负整数，其中奇数与偶数划分了全部可能取值，因此 $p+q=1$。泊松分布列的绝对收敛性允许作下列求和：
\begin{align*}
q-p
&=\sum_{k=0}^{\infty}(-1)^k\P(X=k)\\
&=e^{-\lambda}\sum_{k=0}^{\infty}\frac{(-\lambda)^k}{k!}
=e^{-\lambda}e^{-\lambda}=e^{-2\lambda}.
\end{align*}
于是
\[
p=\frac{(p+q)-(q-p)}2=\frac{1-e^{-2\lambda}}2,\qquad
q=\frac{(p+q)+(q-p)}2=\frac{1+e^{-2\lambda}}2.
\]
再由 $\sinh t=(e^t-e^{-t})/2$、$\cosh t=(e^t+e^{-t})/2$，分别得到题中另两种表示。
\end{proof}
\paragraph{答案} 奇数概率为 $(1-e^{-2\lambda})/2$，偶数概率为 $(1+e^{-2\lambda})/2$。

\subsection{教材第58页第3题}\label{ex:w07-b58-3}
设绝对连续型随机变量 $X$ 的密度函数 $p$ 满足
\[
p(x)=\begin{cases}2ax,&0\leq x\leq1,\\2-x,&1<x\leq2,\\0,&\text{其他。}\end{cases}
\]
求常数 $a$ 和概率 $\P(X<3/2)$。
\basis{原笔记第2章第3节命题“概率密度的基本性质”，例“分段密度的参数、区间概率与分布函数”。}
\paragraph{思路} 先用总积分等于 $1$ 求参数，再在分段点 $1$ 处分开积分。
\begin{proof}[解答]
密度归一化要求
\[
1=\int_0^12ax\,dx+\int_1^2(2-x)\,dx=a+\frac12,
\]
故 $a=1/2$，此时各段密度均非负。于是
\begin{align*}
\P(X<3/2)
&=\int_0^1x\,dx+\int_1^{3/2}(2-x)\,dx\\
&=\frac12+\left[2x-\frac{x^2}2\right]_1^{3/2}
=\frac12+\frac38=\frac78.
\end{align*}
\end{proof}
\paragraph{答案} $a=1/2$，$\P(X<3/2)=7/8$。

\subsection{作业第2题：三角形密度的区间概率}\label{ex:w07-h2}
设连续型随机变量 $X$ 的密度函数为
\[
f(x)=\begin{cases}1-|x|,&-1\leq x\leq1,\\0,&\text{其他。}\end{cases}
\]
求概率 $\P(-\tfrac12<X\leq\tfrac23)$。
\basis{原笔记第2章第3节“利用密度与分布函数计算概率”。}
\paragraph{思路} 绝对值在 $0$ 两侧有不同表达，应在 $0$ 处分段积分。
\begin{proof}[解答]
由于 $X$ 连续型，端点取或不取不改变概率。在 $x\leq0$ 时 $1-|x|=1+x$，在 $x\geq0$ 时 $1-|x|=1-x$，所以
\begin{align*}
\P(-\tfrac12<X\leq\tfrac23)
&=\int_{-1/2}^0(1+x)\,dx+\int_0^{2/3}(1-x)\,dx\\
&=\left[x+\frac{x^2}2\right]_{-1/2}^0
+\left[x-\frac{x^2}2\right]_0^{2/3}\\
&=\frac38+\frac49=\frac{59}{72}.
\end{align*}
\end{proof}
\paragraph{答案} $59/72$。

\subsection{作业第3题：储油量}\label{ex:w07-h3}
一个加油站每周补给一次油。如果它每周的销售量（单位：千升）为随机变量 $X$，且 $X$ 的密度为
\[
f(x)=\begin{cases}0.05(1-x/100)^4,&0\leq x\leq100,\\0,&\text{其他。}\end{cases}
\]
设 $a$ 表示储油罐的最大储油量（单位：千升）。问 $a$ 至少取多大整数，才能把一周内缺油的概率控制到 $0.01$ 以下？
\basis{原笔记第2章第3节“利用密度与分布函数计算概率”。}
\paragraph{思路} 缺油事件是销售量超过储油量，先求尾概率，再按严格不等式取整数。
\begin{proof}[解答]
每周开始时储满油，缺油事件为 $\{X>a\}$。当 $0\leq a\leq100$ 时，令 $u=1-x/100$，得到
\begin{align*}
\P(X>a)
&=\int_a^{100}0.05(1-x/100)^4\,dx\\
&=5\int_0^{1-a/100}u^4\,du
=(1-a/100)^5.
\end{align*}
因此要求
\[
(1-a/100)^5<0.01
\quad\Longleftrightarrow\quad
a>100(1-0.01^{1/5})\approx60.1893.
\]
满足条件的最小整数为 $61$。其中 $a=60$ 时缺油概率为 $0.4^5=0.01024>0.01$，而 $a=61$ 时为 $0.39^5\approx0.00902<0.01$。
\end{proof}
\paragraph{答案} 至少 $61$ 千升。

\subsection{作业第4题：由归一化求密度参数}\label{ex:w07-h4}
设连续型随机变量 $X$ 的密度函数为
\[
f(x)=\begin{cases}Cxe^{-x/2},&x>0,\\0,&\text{其他。}\end{cases}
\]
求常数 $C$，求概率 $\P(-1<X\leq1)$。
\basis{原笔记第2章第3节命题“概率密度的基本性质”；分部积分。}
\paragraph{思路} 用分部积分算总积分，再在实际支撑集内计算概率。
\begin{proof}[解答]
由分部积分，
\[
\int_0^{\infty}xe^{-x/2}\,dx
=\left[-2xe^{-x/2}-4e^{-x/2}\right]_0^{\infty}=4.
\]
因此 $4C=1$，即 $C=1/4$。因为 $X>0$ 几乎必然，
\begin{align*}
\P(-1<X\leq1)
&=\frac14\int_0^1xe^{-x/2}\,dx\\
&=\frac14\left[-2xe^{-x/2}-4e^{-x/2}\right]_0^1
=1-\frac32e^{-1/2}.
\end{align*}
\end{proof}
\paragraph{答案} $C=1/4$，所求概率为 $1-\tfrac32e^{-1/2}$。
% !TEX root = 概率论期中复习.tex
\section{第八周：连续分布、对称性与无记忆性}

\subsection{教材第57页第2题}\label{ex:w08-b57-2}
设 $f,g$ 都是 $\R$ 上的概率分布密度函数，而 $p,q\in[0,1]$ 满足 $p+q=1$。证明 $pf+qg$ 是 $\R$ 上的概率分布密度函数。
\basis{原笔记第2章第3节命题“概率密度的基本性质”。}
\paragraph{思路} 分别核对非负性和归一化；这两项就是成为密度所需的条件。
\begin{proof}
记 $h=pf+qg$。由于 $p,q\geq0$ 且 $f,g$ 非负，有 $h(x)\geq0$。此外，$f,g$ 可积，故 $h$ 也可积，且
\begin{align*}
\int_{\R}h(x)\,dx
&=p\int_{\R}f(x)\,dx+q\int_{\R}g(x)\,dx\\
&=p+q=1.
\end{align*}
因此 $h$ 满足密度的非负性和归一化条件，是 $\R$ 上的概率分布密度函数。
\end{proof}
\paragraph{答案} $pf+qg$ 是密度；其分布函数为 $pF_f+qF_g$，其中 $F_f,F_g$ 分别是 $f,g$ 对应的分布函数。

\subsection{教材第58页第4题}\label{ex:w08-b58-4}
设 $\R$ 上的绝对连续型随机变量 $\xi$ 有密度函数
\[
f(x)=e^{-x-e^{-x}}.
\]
求 $\xi$ 的分布函数 $F$（称为极值分布）。
\basis{原笔记第2章第3节“密度与分布函数的关系”。}
\paragraph{思路} 将密度写成 $e^{-x}\exp(-e^{-x})$，用 $u=e^{-x}$ 换元。
\begin{proof}[解答]
对任意 $t\in\R$，
\begin{align*}
F(t)
&=\int_{-\infty}^t e^{-x}\exp(-e^{-x})\,dx\\
&=\int_{e^{-t}}^{\infty}e^{-u}\,du
=\exp(-e^{-t}).
\end{align*}
这里 $x\to-\infty$ 时 $u\to\infty$，且 $du=-e^{-x}\,dx$，因此换元后积分上下限及符号如上。所得函数在 $t\to-\infty$ 时趋于 $0$，在 $t\to\infty$ 时趋于 $1$。
\end{proof}
\paragraph{答案} $F(t)=\exp(-e^{-t})$，$t\in\R$。

\subsection{作业第2题：由三角形密度求分布函数}\label{ex:w08-h2}
设连续型随机变量 $X$ 的密度函数为
\[
f(x)=\begin{cases}x,&0\leq x\leq1,\\2-x,&1\leq x<2,\\0,&\text{其他。}\end{cases}
\]
求随机变量 $X$ 的分布函数 $F(t)$。
\basis{原笔记第2章第3节例“分段密度的参数、区间概率与分布函数”。}
\paragraph{思路} 按积分上限所在的密度分段分别计算；$t>1$ 时须包含第一段的全部积分。
\begin{proof}[解答]
$t<0$ 时，支撑集在 $t$ 的右侧，故 $F(t)=0$。$0\leq t\leq1$ 时，
\[
F(t)=\int_0^t x\,dx=\frac{t^2}2.
\]
$1<t<2$ 时，
\begin{align*}
F(t)&=\int_0^1x\,dx+\int_1^t(2-x)\,dx\\
&=\frac12+2t-\frac{t^2}2-\frac32
=1-\frac{(2-t)^2}2.
\end{align*}
$t\geq2$ 时已积尽全部概率，$F(t)=1$。两条中间表达在 $t=1$ 处均为 $1/2$，各分段相接连续。
\end{proof}
\paragraph{答案}
\[
F(t)=\begin{cases}0,&t<0,\\t^2/2,&0\leq t\leq1,\\1-(2-t)^2/2,&1<t<2,\\1,&t\geq2.\end{cases}
\]

\subsection{作业第3题：对称密度的分布函数}\label{ex:w08-h3}
设连续函数 $X$ 的密度函数 $f(x)$ 关于 $x=c$ 对称。求证：其分布函数 $F(t)$ 满足
\[
F(c-t)+F(c+t)=1,\qquad\forall t\in\R.
\]
注：不能只画图，要用关于密度的积分进行数学推导。
\basis{原笔记第2章第3节“密度与分布函数的关系”，命题“标准正态分布的对称公式”的积分思路。}
\paragraph{思路} 关于 $c$ 对称表示 $f(2c-x)=f(x)$，对左尾积分作反射换元即可得到右尾。
\begin{proof}
题面“连续函数 $X$”按具有所给密度的连续型随机变量理解。对称性给出 $f(2c-u)=f(u)$。对任意实数 $t$，在积分中令 $u=2c-x$，有
\begin{align*}
F(c-t)
&=\int_{-\infty}^{c-t}f(x)\,dx\\
&=\int_{c+t}^{\infty}f(2c-u)\,du\\
&=\int_{c+t}^{\infty}f(u)\,du
=1-F(c+t).
\end{align*}
最后一步使用密度的总积分为 $1$。将 $F(c+t)$ 移到等式左侧即得结论；这一换元对正、负及零的 $t$ 都成立。
\end{proof}
\paragraph{答案} $F(c-t)+F(c+t)=1$ 对全部 $t\in\R$ 成立；特别地 $F(c)=1/2$。

\subsection{作业第4题：先到列车的方向}\label{ex:w08-h4}
从早上7点开始，每隔15分钟都有一趟列车开往A地，而从7:05开始，每隔15分钟有一趟列车开往B地。
\begin{enumerate}
\item 如果某位乘客在7点到8点之间到达车站的时间服从均匀分布。他到达车站以后，无论下一趟进站的列车是开往A地还是开往B地的，他立刻上那一列车。那么他乘上开往A地的列车的概率是多大？
\item 如果该乘客在早上7:10到8:10之间到达车站的时间服从均匀分布，请回答(a)中相同的问题。
\end{enumerate}
\basis{原笔记第2章第3节“均匀分布”，例“等公交车不超过五分钟”。}
\paragraph{思路} 列出乘客会等到A车的到达时间段，以总长度除以均匀到达区间长度。
\begin{proof}[解答]
用从7点开始的分钟数表示时刻。A车在 $0,15,30,45,60,\ldots$ 到达，B车在 $5,20,35,50,65,\ldots$ 到达。乘客到达时刻恰为列车时刻的概率是 $0$，因此各段端点是否包含不影响答案。
\begin{enumerate}
\item 在 $[0,60]$ 内，乘客先等到A车的时间段为
\[
(5,15),\quad(20,30),\quad(35,45),\quad(50,60).
\]
总长度为 $4\times10=40$ 分钟。因此所求概率为 $40/60=2/3$。
\item 在 $[10,70]$ 内，相应时间段为
\[
(10,15),\quad(20,30),\quad(35,45),\quad(50,60),\quad(65,70).
\]
总长度为 $5+10+10+10+5=40$ 分钟，均匀到达区间长度仍为 $60$ 分钟，所以概率仍为 $2/3$。
\end{enumerate}
\end{proof}
\paragraph{答案} 两问均为 $2/3$。

\subsection{作业第5题：修理时间}\label{ex:w08-h5}
修理某机器所需的时间（单位：小时）服从参数为 $1/2$ 的指数分布。
\begin{enumerate}
\item 修理时间超过2小时的概率是多少？（结果保留3位小数）
\item 若已经持续修理了9小时的条件下，总共需要10小时才能修好的条件概率是多少？（结果保留3位小数）
\end{enumerate}
\basis{原笔记第2章第3节命题“指数分布的分布函数与尾概率”、定理“指数分布的无记忆性”。}
\paragraph{思路} 尾概率为 $e^{-t/2}$；第二问需区分“恰好10小时”和“10小时内修好”。
\begin{proof}[解答]
记总修理时间为 $T\sim\operatorname{Exp}(1/2)$。
\begin{enumerate}
\item 直接使用尾概率，
\[
\P(T>2)=e^{-2/2}=e^{-1}\approx0.368.
\]
\item 已持续修理9小时且尚未修好对应 $\{T>9\}$。若按题面字面“总共需要10小时”理解为恰好 $T=10$，连续型随机变量的单点概率为零，故
\[
\P(T=10\mid T>9)=\frac{\P(T=10)}{\P(T>9)}=0.000.
\]
若题意是“总修理时间不超过10小时”，则所求为
\begin{align*}
\P(T\leq10\mid T>9)
&=1-\P(T>10\mid T>9)\\
&=1-\frac{e^{-10/2}}{e^{-9/2}}
=1-e^{-1/2}\approx0.393.
\end{align*}
若题意是“至少还要修到第10小时”，相应的条件概率则为 $e^{-1/2}\approx0.607$。三者事件不同，不能混用。
\end{enumerate}
\end{proof}
\paragraph{答案} 第一问 $0.368$；第二问字面“恰好10小时”为 $0.000$，若指10小时内修好则为 $0.393$。

\subsection{作业第6题：旧车的剩余里程}\label{ex:w08-h6}
甲估计，一辆汽车在报废之前能行驶的里程数（单位：千里）服从参数为 $1/20$ 的指数分布。乙有一辆已经行驶了10公里的旧车。如果甲买下此车，那么该车至少还能行驶20千里的概率是多少？（结果保留3位小数）如果将以上的指数分布换为 $(0,40)$ 上的均匀分布，再求以上概率。
\basis{原笔记第2章第3节“指数分布的无记忆性”“均匀分布”。}
\paragraph{思路} 先统一单位，以“总里程已超过旧车里程”为条件计算剩余寿命的尾概率。
\begin{proof}[解答]
令 $L$ 表示总里程，以千里为单位，令 $d$ 表示旧车已经行驶的里程换算为相同单位后的数值。剩余里程至少为 $20$ 的条件概率是
\[
\P(L\geq d+20\mid L>d).
\]
对指数分布，只需 $d\geq0$，由无记忆性有
\[
\P(L\geq d+20\mid L>d)
=\frac{e^{-(d+20)/20}}{e^{-d/20}}
=e^{-1}\approx0.368.
\]
对 $L\sim U(0,40)$，当 $0\leq d<20$ 时，
\[
\P(L\geq d+20\mid L>d)
=\frac{40-d-20}{40-d}=\frac{20-d}{40-d}.
\]
题面总里程单位“千里”与已行驶“10公里”并不相同。按“里”为市里，$10$ 公里等于 $20$ 里，即 $d=0.02$ 千里，此时
\[
\frac{20-0.02}{40-0.02}=\frac{19.98}{39.98}\approx0.500.
\]
若题面“千里”实际用来指千公里，则 $d=0.01$，概率为 $19.99/39.99\approx0.500$。若原意是“已经行驶10千里”，则 $d=10$，均匀分布下的答案改为 $10/30\approx0.333$。指数分布的答案在这些解释下都为 $0.368$。
\end{proof}
\paragraph{答案} 指数分布下为 $0.368$；保留题面“10公里”并统一单位后，均匀分布下为 $0.500$；若原意为“10千里”，则为 $0.333$。

\subsection{作业第7题：正态分布的中心}\label{ex:w08-h7}
设 $X\sim N(\mu,4)$，并且 $\P(X\geq3)=0.5$。求 $\mu$ 是多少？
\basis{原笔记第2章第3节“标准正态分布与对称性”“正态分布的线性变换与标准化”。}
\paragraph{思路} 正态分布在其均值两侧各有一半概率，并利用其分布函数严格递增确定唯一的中心。
\begin{proof}[解答]
标准差为 $2$。令 $Z=(X-\mu)/2\sim N(0,1)$，以 $\Phi$ 表示标准正态分布函数，则
\[
\P(X\geq3)=1-\Phi\left(\frac{3-\mu}{2}\right)=\frac12.
\]
所以 $\Phi((3-\mu)/2)=1/2$。由于 $\Phi(0)=1/2$，且标准正态密度处处正，$\Phi$ 严格递增，得到 $(3-\mu)/2=0$，即 $\mu=3$。
\end{proof}
\paragraph{答案} $\mu=3$。
% !TEX root = 概率论期中复习.tex
\section{第九周：正态计算与随机变量函数的分布}

\subsection{教材第58页第4题（再次布置）}\label{ex:w09-b58-4}
设 $\R$ 上的绝对连续型随机变量 $\xi$ 有密度函数 $f(x)=e^{-x-e^{-x}}$。求 $\xi$ 的分布函数 $F$（称为极值分布）。
\basis{原笔记第2章第3节“密度与分布函数的关系”。}
\paragraph{思路} 本题在第八周已经布置，保留此次作业题号，并沿用同一换元。
\begin{proof}[解答]
令 $u=e^{-x}$，则 $du=-e^{-x}\,dx$。对 $t\in\R$，
\[
F(t)=\int_{-\infty}^{t}e^{-x}\exp(-e^{-x})\,dx
=\int_{e^{-t}}^{\infty}e^{-u}\,du
=\exp(-e^{-t}).
\]
完整换元说明见\hyperref[ex:w08-b58-4]{第八周教材第58页第4题}。
\end{proof}
\paragraph{答案} $F(t)=\exp(-e^{-t})$，$t\in\R$。

\subsection{教材第58页第5题}\label{ex:w09-b58-5}
设随机变量 $\xi$ 服从标准正态分布 $N(0,1)$。令 $X=\sigma\xi+a$，其中 $\sigma\ne0$ 和 $a$ 为常数。证明 $X$ 服从正态分布 $N(a,\sigma^2)$。
\basis{原笔记第2章第3节“正态分布的线性变换与标准化”；第4节例“一般连续型随机变量的线性变换”。}
\paragraph{思路} 沿原笔记的线性密度变换，绝对值 $|\sigma|$ 同时处理正、负的尺度。
\begin{proof}
记标准正态密度为 $\varphi(u)=(2\pi)^{-1/2}e^{-u^2/2}$。由线性变换的密度公式，
\begin{align*}
f_X(x)
&=\frac1{|\sigma|}\varphi\left(\frac{x-a}{\sigma}\right)\\
&=\frac1{\sqrt{2\pi}|\sigma|}
\exp\left(-\frac{(x-a)^2}{2\sigma^2}\right),\qquad x\in\R.
\end{align*}
这就是均值参数为 $a$、方差参数为 $\sigma^2$ 的正态密度，所以 $X\sim N(a,\sigma^2)$。这里即使 $\sigma<0$，标准差仍为正数 $|\sigma|$。
\end{proof}
\paragraph{答案} $X\sim N(a,\sigma^2)$，密度如上。

\subsection{教材第62页第1题}\label{ex:w09-b62-1}
设随机变量 $\xi$ 服从均匀分布 $U(0,1)$。求方程 $x^2+3\xi x+1=0$ 有实根的概率。
\basis{原笔记第2章第3节“均匀分布”。}
\paragraph{思路} 将有实根转化为判别式非负，再计算均匀随机变量落入相应区间的概率。
\begin{proof}[解答]
关于 $x$ 的方程有实根当且仅当
\[
\Delta=(3\xi)^2-4\geq0
\quad\Longleftrightarrow\quad |\xi|\geq\frac23.
\]
由于 $0<\xi<1$ 几乎必然，上述条件等价于 $\xi\geq2/3$。故
\[
\P(\Delta\geq0)=\P(\xi\geq2/3)=\int_{2/3}^{1}1\,du=\frac13.
\]
\end{proof}
\paragraph{答案} $1/3$。

\subsection{作业第2题：正态分布的区间概率}\label{ex:w09-h2}
设 $X\sim N(10,36)$，求以下概率（可以利用软件、计算器等，结果保留4位小数）：
\begin{enumerate}
\item $\P(X>5)$；
\item $\P(4\leq X\leq16)$；
\item $\P(X<16)$。
\end{enumerate}
\basis{原笔记第2章第3节“正态分布的线性变换与标准化”、命题“标准正态分布的对称公式”。}
\paragraph{思路} 方差为 $36$，标准差为 $6$；先标准化，再使用 $\Phi$。
\begin{proof}[解答]
令 $Z=(X-10)/6\sim N(0,1)$。
\begin{enumerate}
\item
\[
\P(X>5)=\P(Z>-5/6)
=1-\Phi(-5/6)=\Phi(5/6)\approx0.7977.
\]
\item
\[
\P(4\leq X\leq16)=\P(-1\leq Z\leq1)
=\Phi(1)-\Phi(-1)=2\Phi(1)-1\approx0.6827.
\]
\item
\[
\P(X<16)=\P(Z<1)=\Phi(1)\approx0.8413.
\]
\end{enumerate}
\end{proof}
\paragraph{答案} 依次为 $0.7977$、$0.6827$、$0.8413$。

\subsection{作业第3题：由分位点确定收入分布}\label{ex:w09-h3}
设某领域的物理学家的年收入近似服从正态分布。已知该群物理学家中年收入低于180000元的占25\%，而高于320000元的占25\%。
\begin{enumerate}
\item 近似求年收入少于200000元的概率。（结果保留4位小数）
\item 近似求年收入介入280000元和320000元之间的概率。（结果保留4位小数）
\end{enumerate}
\basis{原笔记第2章第3节例“由两端比例估计录取分数线”。}
\paragraph{思路} 两个已知分位点关于均值对称，先求均值、标准差，再代入目标区间。
\begin{proof}[解答]
记年收入为 $X\sim N(\mu,\sigma^2)$，$\sigma>0$，并记
\[
z_{0.75}=\Phi^{-1}(0.75)\approx0.67448975.
\]
已知条件等价于
\[
\frac{180000-\mu}{\sigma}=-z_{0.75},\qquad
\frac{320000-\mu}{\sigma}=z_{0.75}.
\]
两式相加得到 $\mu=250000$；相减得到
\[
\sigma=\frac{320000-180000}{2z_{0.75}}
=\frac{70000}{z_{0.75}}\approx103782.1553.
\]
\begin{enumerate}
\item
\[
\P(X<200000)=\Phi\left(\frac{-50000}{\sigma}\right)
=\Phi\left(-\frac57z_{0.75}\right)\approx0.3150.
\]
\item 将题面“介入”按“介于”理解，有
\begin{align*}
\P(280000<X<320000)
&=\Phi\left(\frac{70000}{\sigma}\right)
-\Phi\left(\frac{30000}{\sigma}\right)\\
&=0.75-\Phi\left(\frac37z_{0.75}\right)\approx0.1363.
\end{align*}
\end{enumerate}
\end{proof}
\paragraph{答案} $\mu=250000$、$\sigma\approx103782.1553$；两问概率分别为 $0.3150$、$0.1363$。

\subsection{作业第4题：正态分位点}\label{ex:w09-h4}
设 $X\sim N(12,4)$。若 $\P(X>c)=0.1$，求 $c$（结果保留2位小数）。
\basis{原笔记第2章第3节“正态分布的线性变换与标准化”。}
\paragraph{思路} 右尾为 $0.1$ 对应左尾分位点 $0.9$；标准差为 $2$。
\begin{proof}[解答]
由标准化，
\[
0.1=\P\left(\frac{X-12}{2}>\frac{c-12}{2}\right)
=1-\Phi\left(\frac{c-12}{2}\right).
\]
因此 $(c-12)/2=\Phi^{-1}(0.9)\approx1.28155157$，即
\[
c=12+2\Phi^{-1}(0.9)\approx14.56310313.
\]
\end{proof}
\paragraph{答案} $c\approx14.56$。

\subsection{作业第5题：线段分割的长短比}\label{ex:w09-h5}
从长为 $L$ 的线段上随机等可能地选取一点将线段截断。设 $X$ 表示从线段左端到该断点的距离。求短的那一截长度相对于长的那一截长度之比小于 $1/4$ 的概率。
\basis{原笔记第2章第3节“均匀分布”；第4节“先分布函数，后密度”中求取值集合的方法。}
\paragraph{思路} 在断点位于左半段、右半段时分别解比例不等式。
\begin{proof}[解答]
由随机等可能选点，$X\sim U(0,L)$。两截长度为 $X$ 与 $L-X$。

当 $0<X\leq L/2$ 时，短截为 $X$，故
\[
\frac{X}{L-X}<\frac14
\quad\Longleftrightarrow\quad 5X<L
\quad\Longleftrightarrow\quad X<L/5.
\]
当 $L/2<X<L$ 时，短截为 $L-X$，故
\[
\frac{L-X}{X}<\frac14
\quad\Longleftrightarrow\quad 5X>4L
\quad\Longleftrightarrow\quad X>4L/5.
\]
两个事件互斥，合并得到
\[
\P\left(\frac{\min(X,L-X)}{\max(X,L-X)}<\frac14\right)
=\frac{L/5+L/5}{L}=\frac25.
\]
\end{proof}
\paragraph{答案} $2/5$。

\subsection{作业第6题：由均匀分布构造指数分布}\label{ex:w09-h6}
设 $X\sim U(0,1)$，令 $Y=-\lambda^{-1}\ln X$，$\lambda>0$。求 $Y$ 的密度函数并判断 $Y$ 服从什么常用分布。
\basis{原笔记第2章第4节“求随机变量函数密度的一般方法”；第3节“指数分布”。}
\paragraph{思路} 变换严格递减，解 $Y\leq y$ 时必须反向改变关于 $X$ 的不等号。
\begin{proof}[解答]
因为 $0<X<1$ 几乎必然，$Y>0$ 几乎必然。对 $y<0$，$F_Y(y)=0$；对 $y\geq0$，
\begin{align*}
F_Y(y)
&=\P\left(-\frac1\lambda\ln X\leq y\right)\\
&=\P(X\geq e^{-\lambda y})
=1-e^{-\lambda y}.
\end{align*}
在各开区间求导，得到一个密度版本
\[
f_Y(y)=\begin{cases}\lambda e^{-\lambda y},&y>0,\\0,&y\leq0.\end{cases}
\]
这与参数为 $\lambda$ 的指数密度一致。
\end{proof}
\paragraph{答案} $Y\sim\operatorname{Exp}(\lambda)$，密度为 $\lambda e^{-\lambda y}$（$y>0$），其他处取 $0$。

\subsection{作业第7题：指数变量的尺度变换}\label{ex:w09-h7}
设 $X\sim\operatorname{Exp}(\lambda)$，常数 $c>0$。令 $Y=cX$，求 $Y$ 的密度函数，并判断 $Y$ 服从什么常用分布。
\basis{原笔记第2章第4节例“用反函数公式求线性变换的密度”；第3节指数分布的率参数约定。}
\paragraph{思路} 对线性变换使用密度公式，同时换算率参数。
\begin{proof}[解答]
反函数为 $x=y/c$，导数绝对值为 $1/c$。由于 $c>0$，$Y$ 只在正半轴有密度，于是
\[
f_Y(y)=\frac1c f_X(y/c)
=\begin{cases}\dfrac\lambda c e^{-\lambda y/c},&y>0,\\0,&y\leq0.\end{cases}
\]
令 $\lambda'=\lambda/c>0$，这正是率参数为 $\lambda'$ 的指数密度。
\end{proof}
\paragraph{答案} $Y\sim\operatorname{Exp}(\lambda/c)$。

\subsection{作业第8题：均匀变量的绝对值}\label{ex:w09-h8}
设 $X\sim U(-1/2,1/2)$。令 $Y=|X|$。求 $Y$ 的密度函数，并判断 $Y$ 服从什么常用分布。
\basis{原笔记第2章第4节“求随机变量函数密度的一般方法”；第3节“均匀分布”。}
\paragraph{思路} $|X|\leq y$ 等价于 $-y\leq X\leq y$，区间长度加倍。
\begin{proof}[解答]
对 $y<0$，$F_Y(y)=0$。当 $0\leq y<1/2$ 时，$[-y,y]$ 位于 $X$ 的均匀区间内，所以
\[
F_Y(y)=\P(-y\leq X\leq y)=\int_{-y}^y1\,dx=2y.
\]
对 $y\geq1/2$，$F_Y(y)=1$。因此
\[
F_Y(y)=\begin{cases}0,&y<0,\\2y,&0\leq y<1/2,\\1,&y\geq1/2,\end{cases}
\qquad
f_Y(y)=\begin{cases}2,&0<y<1/2,\\0,&\text{其他。}\end{cases}
\]
\end{proof}
\paragraph{答案} $Y\sim U(0,1/2)$。

\subsection{作业第9题：指数变量的余弦}\label{ex:w09-h9}
设 $X\sim\operatorname{Exp}(1)$，令 $Y=\cos X$。求 $Y$ 的密度函数。
\basis{原笔记第2章第4节定理“分支反函数的密度变换公式”，例“均匀随机变量的余弦变换”“倒数的小数部分与可数个反函数分支”。}
\paragraph{思路} 余弦在正半轴有无穷多个单调分支，必须将每个原像的密度贡献相加。
\begin{proof}[解答]
$Y$ 的值域包含于 $[-1,1]$。固定 $-1<y<1$，令 $\alpha=\arccos y\in(0,\pi)$。正半轴上满足 $\cos x=y$ 的全部原像为
\[
x=2k\pi+\alpha\quad(k=0,1,\ldots),\qquad
x=2k\pi-\alpha\quad(k=1,2,\ldots).
\]
这些原像分别位于余弦的各个单调区间，每一反函数分支的导数绝对值均为 $1/\sqrt{1-y^2}$。区间端点 $k\pi$ 构成可数集，$X$ 落在该集合的概率为零，因此这些开区间覆盖了全部概率。

由分支反函数公式，
\begin{align*}
f_Y(y)
&=\frac1{\sqrt{1-y^2}}
\left(\sum_{k=0}^{\infty}e^{-(2k\pi+\alpha)}
+\sum_{k=1}^{\infty}e^{-(2k\pi-\alpha)}\right)\\
&=\frac{e^{-\alpha}+e^{-2\pi+\alpha}}
{(1-e^{-2\pi})\sqrt{1-y^2}}.
\end{align*}
这是正项等比级数求和，因公比 $e^{-2\pi}<1$ 而收敛。$|y|>1$ 时密度为 $0$；在 $y=\pm1$ 处的密度值可另取 $0$，不影响分布。
\end{proof}
\paragraph{答案}
\[
f_Y(y)=\begin{cases}
\dfrac{e^{-\arccos y}+e^{-2\pi+\arccos y}}
{(1-e^{-2\pi})\sqrt{1-y^2}},&-1<y<1,\\[6pt]
0,&\text{其他。}
\end{cases}
\]

\subsection{作业第10题及思考题：概率积分变换}\label{ex:w09-h10}
设连续型随机变量 $X$ 的分布函数为 $F(t)$，且 $F(t)$ 严格单增。令 $Y=F(X)$（即 $Y$ 是随机变量 $X$ 和分布函数 $F$ 的复合）。求证：$Y\sim U(0,1)$。

思考题：如果去掉以上 $F(t)$ 严格单增的条件，结论对吗？
\basis{原笔记第2章第1节“分布函数的性质”；第4节“构造给定分布函数的随机变量”中广义逆的思想。}
\paragraph{思路} 严格单增时用通常反函数；去掉严格单增后，用同一水平集的最右端处理平坦区间。
\begin{proof}
先保留严格单增条件。因为 $X$ 连续型，$F$ 连续；结合两端极限为 $0,1$，对每个 $u\in(0,1)$ 恰有一个 $t_u\in\R$ 使 $F(t_u)=u$。严格单增性给出
\[
\{F(X)\leq u\}=\{X\leq t_u\},
\]
从而
\[
F_Y(u)=\P(F(X)\leq u)=F(t_u)=u,\qquad0<u<1.
\]
又因 $0\leq Y\leq1$，并由上式令 $u\downarrow0$ 知 $\P(Y=0)=0$，得到
\[
F_Y(u)=\begin{cases}0,&u<0,\\u,&0\leq u\leq1,\\1,&u>1.\end{cases}
\]
所以 $Y\sim U(0,1)$。

再去掉严格单增条件，保留 $X$ 连续型，故 $F$ 仍连续。固定 $u\in(0,1)$，令
\[
b_u=\sup\{x\in\R:F(x)\leq u\}.
\]
由 $F$ 的两端极限，集合非空且有上界，故 $b_u$ 为有限实数。由连续性和上确界定义，$F(b_u)\leq u$。若 $F(b_u)<u$，连续性保证某个 $b_u$ 右侧的点仍满足 $F(x)<u$，与上确界的定义矛盾。因此 $F(b_u)=u$。

由单调性，$x\leq b_u$ 时 $F(x)\leq u$；$x>b_u$ 时按上确界定义 $F(x)>u$。于是
\[
\P(F(X)\leq u)=\P(X\leq b_u)=F(b_u)=u.
\]
再如前处理 $u\leq0$ 与 $u\geq1$，得到同一个均匀分布函数。故严格单增条件可以去掉；连续性足以保证结论。
\end{proof}
\paragraph{答案} $F(X)\sim U(0,1)$；即使 $F$ 有平坦区间，结论仍成立。这里保留了题设中 $X$ 连续型的条件。

\clearpage\part{随机向量、条件分布与变量变换}
% !TEX root = 概率论期中复习.tex
% 自动抽取自 chapters/chapter3.tex；原例题及其解答保留。

\clearpage\section{联合分布函数和边缘分布函数}\label{k:joint}

\begin{mydefinition}{随机向量}{random-vector}
若随机变量 $X_1,X_2,\ldots,X_n$ 定义在同一个概率空间 $(\Omega,\mathcal F,P)$ 上，则称
\[
\boldsymbol X=(X_1,X_2,\ldots,X_n)
\]
为一个 \textbf{$n$ 维随机向量}，也称 $n$ 维随机变量。它把同一次试验结果 $\omega$ 映为
\[
\boldsymbol X(\omega)=\bigl(X_1(\omega),\ldots,X_n(\omega)\bigr)\in\mathbb R^n.
\]
\end{mydefinition}
\origin{原笔记第3.1节，定义“随机向量”；源码第9行。}

记 $\sigma(\mathcal C)$ 为包含集合族 $\mathcal C$ 的最小 $\sigma$-域。定义 $n$ 维博雷尔 $\sigma$-域
\[
\mathcal B(\mathbb R^n)
=\sigma\left(
\left\{\prod_{i=1}^n(a_i,b_i]:
a_i,b_i\in\mathbb R,\ a_i<b_i\right\}
\right).
\]
于是，随机向量的等价表述为：映射 $\boldsymbol X:\Omega\to\mathbb R^n$ 满足
\[
\boldsymbol X^{-1}(B)\in\mathcal F,
\qquad \forall B\in\mathcal B(\mathbb R^n).
\]

\begin{mydefinition}{联合分布函数}{joint-distribution-function}
设 $\boldsymbol X=(X_1,\ldots,X_n)$ 为 $n$ 维随机向量。对任意 $(t_1,\ldots,t_n)\in\mathbb R^n$，令
\[
F(t_1,\ldots,t_n)
=P(X_1\leq t_1,\ldots,X_n\leq t_n)
=P\left(\bigcap_{i=1}^n\{X_i\leq t_i\}\right).
\]
称 $F$ 为 $\boldsymbol X$ 的\textbf{联合分布函数}。逗号表示这些事件同时发生，即取交集。
\end{mydefinition}
\origin{原笔记第3.1节，定义“联合分布函数”；源码第56行。}

\begin{proposition}{联合分布函数的基本性质}{joint-cdf-properties}
设 $F$ 是 $n$ 维随机向量的联合分布函数，$n\geq2$，则：
\begin{enumerate}
\item 固定其他分量时，$F$ 关于每个分量单调不减。
\item 固定其他分量时，对每个 $i=1,\ldots,n$，
\[
\lim_{t_i\to-\infty}F(t_1,\ldots,t_i,\ldots,t_n)=0;
\]
当所有分量同时趋于 $+\infty$ 时，
\[
\lim_{\substack{t_1\to+\infty,\ldots,t_n\to+\infty}}
F(t_1,\ldots,t_n)=1.
\]
\item 固定其他分量时，$F$ 关于每个分量右连续。
\item 当 $n=2$ 时，对任意 $a<b$ 和 $c<d$，
\[
F(b,d)-F(a,d)-F(b,c)+F(a,c)\geq0.
\]
\end{enumerate}
\end{proposition}
\origin{原笔记第3.1节，命题“联合分布函数的基本性质”；源码第66行。}

\begin{mydefinition}{边缘分布函数}{marginal-distribution-function}
设 $\boldsymbol X=(X_1,\ldots,X_n)$ 为 $n$ 维随机向量。选取非空且不包含全部坐标的指标集
\[
I=\{i_1,\ldots,i_k\},\qquad
1\leq i_1<\cdots<i_k\leq n,\quad 1\leq k\leq n-1.
\]
子向量 $\boldsymbol X_I=(X_{i_1},\ldots,X_{i_k})$ 的联合分布函数
\[
F_I(t_{i_1},\ldots,t_{i_k})
=P(X_{i_1}\leq t_{i_1},\ldots,X_{i_k}\leq t_{i_k})
\]
称为原随机向量的一个\textbf{边缘分布函数}。
\end{mydefinition}
\origin{原笔记第3.1节，定义“边缘分布函数”；源码第175行。}

\begin{proposition}{由联合分布函数求边缘分布函数}{marginal-cdf-limits}
设 $(X,Y)$ 的联合分布函数为 $F(t,s)$，则
\[
F_X(t)=\lim_{s\to+\infty}F(t,s),\qquad
F_Y(s)=\lim_{t\to+\infty}F(t,s).
\]
一般地，设 $\boldsymbol X$ 的联合分布函数为 $F$，保留指标集 $I$ 对应的参数，将其他参数都趋于 $+\infty$，即得
\[
F_I((t_i)_{i\in I})
=\lim_{\substack{t_j\to+\infty\\j\notin I}}F(t_1,\ldots,t_n).
\]
其中 $(t_i)_{i\in I}$ 表示按指标从小到大排列的保留参数。
\end{proposition}
\origin{原笔记第3.1节，命题“由联合分布函数求边缘分布函数”；源码第195行。}

\begin{mydefinition}{离散型随机向量与联合分布列}{discrete-vector-joint-pmf}
若存在有限集或可数集 $S\subseteq\mathbb R^n$，使
\[
P(\boldsymbol X\in S)=1,
\]
则称 $\boldsymbol X=(X_1,\ldots,X_n)$ 为 \textbf{$n$ 维离散型随机向量}。也就是说，它以概率 $1$ 只取有限个或可数个向量值。

对 $\boldsymbol x=(x_1,\ldots,x_n)$，记
\[
p(\boldsymbol x)=p(x_1,\ldots,x_n)
=P(X_1=x_1,\ldots,X_n=x_n).
\]
这些单点概率构成 $\boldsymbol X$ 的\textbf{联合分布列}。在 $S$ 以外，$p(\boldsymbol x)=0$。
\end{mydefinition}
\origin{原笔记第3.1节，定义“离散型随机向量与联合分布列”；源码第244行。}

\begin{proposition}{联合分布列的基本性质}{joint-pmf-properties}
离散型随机向量的联合分布列满足
\[
p(x_1,\ldots,x_n)\geq0,\qquad
\sum_{x_1}\cdots\sum_{x_n}p(x_1,\ldots,x_n)=1.
\]
\end{proposition}
\origin{原笔记第3.1节，命题“联合分布列的基本性质”；源码第265行。}

\begin{proposition}{由联合分布列计算事件概率与分布函数}{joint-pmf-event-probability}
对任意 $B\in\mathcal B(\mathbb R^n)$，
\[
P(\boldsymbol X\in B)
=\sum_{\boldsymbol x\in S\cap B}p(\boldsymbol x).
\]
特别地，
\[
F(t_1,\ldots,t_n)
=\sum_{x_1\leq t_1}\cdots\sum_{x_n\leq t_n}
p(x_1,\ldots,x_n).
\]
\end{proposition}
\origin{原笔记第3.1节，命题“由联合分布列计算事件概率与分布函数”；源码第295行。}

\subsection{例题：指数随机变量经阈值变换得到的联合分布列}\label{note:c3:e1}
\paragraph{题目}
设 $Y\sim\operatorname{Exp}(1)$，令
\[
X_k=
\begin{cases}
0,&Y\leq k,\\
1,&Y>k,
\end{cases}
\qquad k=1,2.
\]
求 $(X_1,X_2)$ 的联合分布列。

\basis{原笔记第3.1节，例“指数随机变量经阈值变换得到的联合分布列”；源码\texttt{chapters/chapter3.tex}第324行。}
\begin{proof}[解答]
两个变量都由同一个 $Y$ 决定，四种组合必须先还原为 $Y$ 的事件，不能把各自的概率直接相乘。逐项计算：
\begin{align*}
P(X_1=0,X_2=0)
&=P(Y\leq1,Y\leq2)\\
&=P(Y\leq1)=\int_0^1e^{-y}\,dy
=1-e^{-1},\\
P(X_1=0,X_2=1)
&=P(Y\leq1,Y>2)=P(\varnothing)=0,\\
P(X_1=1,X_2=0)
&=P(Y>1,Y\leq2)\\
&=\int_1^2e^{-y}\,dy=e^{-1}-e^{-2},\\
P(X_1=1,X_2=1)
&=P(Y>1,Y>2)=P(Y>2)=e^{-2}.
\end{align*}
按行表示 $X_1$、按列表示 $X_2$，得到
\[
\begin{array}{c|cc}
X_1\backslash X_2&0&1\\ \hline
0&1-e^{-1}&0\\
1&e^{-1}-e^{-2}&e^{-2}
\end{array}.
\]
三段 $(0,1]$、$(1,2]$、$(2,+\infty)$ 覆盖 $Y$ 的全概率取值，故这些概率非负且总和为 $1$。
\end{proof}

\begin{mydefinition}{边缘分布列}{marginal-probability-mass}
设 $\boldsymbol X$ 为 $n$ 维离散型随机向量，$I=\{i_1,\ldots,i_k\}$ 为前述非空真子指标集。子向量
\[
\boldsymbol X_I=(X_{i_1},\ldots,X_{i_k})
\]
的联合分布列，称为原向量的一个\textbf{边缘分布列}。
\end{mydefinition}
\origin{原笔记第3.1节，定义“边缘分布列”；源码第361行。}

\begin{proposition}{由联合分布列求边缘分布列}{marginal-pmf-sums}
设二维离散型随机向量 $(X,Y)$ 的不同可能取值分别列为 $x_i$ 和 $y_j$，联合分布列为 $p(x_i,y_j)$，则
\[
p_X(x_i)=\sum_jp(x_i,y_j),\qquad
p_Y(y_j)=\sum_ip(x_i,y_j).
\]
求和遍历全部有限个或可数个相应取值。

一般地，固定保留坐标 $x_i$（$i\in I$），有
\[
p_I((x_i)_{i\in I})
=\sum_{\substack{x_j\in A_j\\j\notin I}}
p(x_1,\ldots,x_n).
\]
右侧表示对所有被舍去坐标的取值组合求和，保留坐标不参与求和。
\end{proposition}
\origin{原笔记第3.1节，命题“由联合分布列求边缘分布列”；源码第369行。}

\subsection{例题：由联合分布表求边缘分布列}\label{note:c3:e2}
\paragraph{题目}
已知 $(X,Y)$ 的联合分布列为
\[
\begin{array}{c|cc}
X\backslash Y&0&1\\ \hline
0&\dfrac13&\dfrac16\\[2pt]
1&\dfrac14&\dfrac14
\end{array}.
\]
求 $X,Y$ 的边缘分布列。

\basis{原笔记第3.1节，例“由联合分布表求边缘分布列”；源码\texttt{chapters/chapter3.tex}第415行。}
\begin{proof}[解答]
按行求和，
\begin{align*}
P(X=0)&=\frac13+\frac16=\frac12,\\
P(X=1)&=\frac14+\frac14=\frac12.
\end{align*}
按列求和，
\begin{align*}
P(Y=0)&=\frac13+\frac14=\frac7{12},\\
P(Y=1)&=\frac16+\frac14=\frac5{12}.
\end{align*}
将边缘概率写在联合分布表旁边：
\[
\begin{array}{c|cc|c}
X\backslash Y&0&1&P(X=x)\\ \hline
0&\dfrac13&\dfrac16&\dfrac12\\[2pt]
1&\dfrac14&\dfrac14&\dfrac12\\[2pt] \hline
P(Y=y)&\dfrac7{12}&\dfrac5{12}&1
\end{array}.
\]
两组边缘概率分别加和为 $1$，与联合分布总概率一致。
\end{proof}

记 $m_n(B)$ 为集合 $B$ 的 $n$ 维体积（严格地说是勒贝格测度）；$m_n(B)=0$ 的集合称为 $n$ 维零测集。“几乎处处”表示允许在一个零测集上例外。本小节的密度统一采用非负博雷尔可测版本，允许取 $+\infty$；由全积分为 $1$，密度几乎处处有限。

\begin{mydefinition}{连续型随机向量与联合密度}{continuous-vector-joint-density}
若随机向量 $\boldsymbol X=(X_1,\ldots,X_n)$ 存在非负博雷尔可测函数 $f:\mathbb R^n\to[0,+\infty]$，使对任意 $B\in\mathcal B(\mathbb R^n)$，
\[
P(\boldsymbol X\in B)=\int_B f(x_1,\ldots,x_n)\,dx_1\cdots dx_n,
\]
则称 $\boldsymbol X$ 为 \textbf{$n$ 维连续型随机向量}，$f$ 称为其\textbf{联合密度}。这里“连续型”指存在关于 $n$ 维体积的密度，不仅指分布函数连续。
\end{mydefinition}
\origin{原笔记第3.1节，定义“连续型随机向量与联合密度”；源码第448行。}

\begin{proposition}{联合密度的基本性质}{joint-density-properties}
联合密度非负，且
\[
\int_{\mathbb R^n}f(\boldsymbol x)\,d\boldsymbol x=1,
\qquad d\boldsymbol x=dx_1\cdots dx_n.
\]
\end{proposition}
\origin{原笔记第3.1节，命题“联合密度的基本性质”；源码第462行。}

\begin{proposition}{联合密度与联合分布函数}{joint-density-cdf}
设 $\boldsymbol X$ 有联合密度 $f$、联合分布函数 $F$，则：
\begin{enumerate}
\item 对任意 $\boldsymbol t=(t_1,\ldots,t_n)$，
\[
F(\boldsymbol t)=\int_{\{x_1\leq t_1,\ldots,x_n\leq t_n\}}
 f(\boldsymbol x)\,d\boldsymbol x.
\]
\item $F$ 是关于所有参数的连续函数。
\item 若博雷尔集 $B$ 满足 $m_n(B)=0$，则 $P(\boldsymbol X\in B)=0$。
\item 同一分布的两个联合密度在几乎处处相等；在零测集上改变密度的取值，只要仍保持非负和博雷尔可测，就不改变其分布。
\end{enumerate}
\end{proposition}
\origin{原笔记第3.1节，命题“联合密度与联合分布函数”；源码第475行。}

\subsection{例题：在抛物线围成的区域上归一化并计算概率}\label{note:c3:e3}
\paragraph{题目}
设联合密度在 $x^2\leq y\leq1$ 内为 $cx^2y$，区域外为 $0$。
\begin{enumerate}
\item 求归一化常数 $c$。

\item 求 $P(X>Y)$。
\end{enumerate}

\basis{原笔记第3.1节，例“在抛物线围成的区域上归一化并计算概率”；源码\texttt{chapters/chapter3.tex}第537行。}
\begin{proof}[解答]
\begin{enumerate}
\item 区域要求 $-1\leq x\leq1$。先固定 $x$，再沿 $y$ 积分：
\begin{align*}
1&=c\int_{-1}^1x^2\left(\int_{x^2}^1y\,dy\right)dx\\
&=\frac c2\int_{-1}^1(x^2-x^6)\,dx\\
&=\frac c2\left(\frac23-\frac27\right)=\frac{4c}{21}.
\end{align*}
所以 $c=21/4$。

\item 事件 $X>Y$ 还要求 $y<x$。结合 $y\geq x^2\geq0$，得到 $0<x<1$，且 $x^2\leq y<x$。边界曲线为零测集，不影响积分，故
\begin{align*}
P(X>Y)
&=\frac{21}4\int_0^1x^2\left(\int_{x^2}^{x}y\,dy\right)dx\\
&=\frac{21}8\int_0^1(x^4-x^6)\,dx
=\frac{21}8\left(\frac15-\frac17\right)=\frac3{20}.
\end{align*}
\end{enumerate}
\end{proof}

\begin{mydefinition}{平面区域上的均匀分布}{uniform-planar-region}
设 $G\subseteq\mathbb R^2$ 为有界博雷尔可测区域，面积 $m_2(G)>0$。若
\[
f(x,y)=
\begin{cases}
1/m_2(G),&(x,y)\in G,\\
0,&(x,y)\notin G,
\end{cases}
\]
则称 $(X,Y)$ 服从 $G$ 上的\textbf{均匀分布}，记为 $(X,Y)\sim U(G)$。对任意博雷尔集 $B$，
\[
P((X,Y)\in B)=\frac{m_2(B\cap G)}{m_2(G)}.
\]
\end{mydefinition}
\origin{原笔记第3.1节，定义“平面区域上的均匀分布”；源码第557行。}

\subsection{例题：圆盘上的均匀分布}\label{note:c3:e4}
\paragraph{题目}
设 $G=\{(x,y):x^2+y^2\leq R^2\}$，$R>0$，且 $(X,Y)\sim U(G)$。
\begin{enumerate}
\item 求 $P(|X|\leq R/2)$。

\item 令 $D=\sqrt{X^2+Y^2}$，求 $D$ 的分布函数。
\end{enumerate}

\basis{原笔记第3.1节，例“圆盘上的均匀分布”；源码\texttt{chapters/chapter3.tex}第572行。}
\begin{proof}[解答]
\begin{enumerate}
\item 联合密度在圆盘内为 $1/(\pi R^2)$，外部为 $0$。利用关于两轴的对称性，
\[
P(|X|\leq R/2)=\frac4{\pi R^2}\int_0^{R/2}\sqrt{R^2-x^2}\,dx.
\]
取 $x=R\sin\theta$，则 $dx=R\cos\theta\,d\theta$，上下限为 $0,\pi/6$，于是
\begin{align*}
P(|X|\leq R/2)
&=\frac4\pi\int_0^{\pi/6}\cos^2\theta\,d\theta\\
&=\frac4\pi\left[\frac\theta2+\frac{\sin2\theta}4\right]_0^{\pi/6}\\
&=\frac13+\frac{\sqrt3}{2\pi}.
\end{align*}

\item $t<0$ 时事件为空，$t\geq R$ 时覆盖整个圆盘。$0\leq t<R$ 时，事件对应半径为 $t$ 的同心圆盘，其面积与原圆盘面积之比为 $t^2/R^2$。所以
\[
F_D(t)=
\begin{cases}
0,&t<0,\\
t^2/R^2,&0\leq t<R,\\
1,&t\geq R.
\end{cases}
\]
它在 $0,R$ 两端连续。
\end{enumerate}
\end{proof}

\begin{mytheorem}{由联合密度求边缘密度}{marginal-density-integrals}
若 $(X,Y)$ 有联合密度 $f$，则两个分量都是一维连续型随机变量，其\textbf{边缘密度}可取为
\[
f_X(x)=\int_{\mathbb R}f(x,y)\,dy,\qquad
f_Y(y)=\int_{\mathbb R}f(x,y)\,dx.
\]
这里采用上述非负博雷尔可测的 $f$，故截面积分也是非负博雷尔可测函数，并且几乎处处有限；在取无穷大的博雷尔零测异常集上，可统一改取 $0$，或另取保持博雷尔可测的非负有限值，不影响分布。
\end{mytheorem}
\origin{原笔记第3.1节，定理“由联合密度求边缘密度”；源码第599行。}

\subsection{例题：固定坐标求积分上下限}\label{note:c3:e5}
\paragraph{题目}
设联合密度在 $x^2\leq y\leq x$ 内为 $6$，其余为 $0$。
\begin{enumerate}
\item 求边缘密度 $f_X,f_Y$。

\item 求 $P(-1<X\leq1/2)$。
\end{enumerate}

\basis{原笔记第3.1节，例“固定坐标求积分上下限”；源码\texttt{chapters/chapter3.tex}第641行。}
\begin{proof}[解答]
\begin{enumerate}
\item 由 $x^2\leq x$ 可得 $0\leq x\leq1$。固定 $x\in(0,1)$ 时，$y$ 从 $x^2$ 到 $x$，所以
\[
f_X(x)=
\begin{cases}
\displaystyle\int_{x^2}^x6\,dy=6(x-x^2),&0<x<1,\\
0,&\text{其他}.
\end{cases}
\]


固定 $y$ 时，$y\leq x$ 给出 $x\geq y$，$x^2\leq y$ 结合 $x\geq0$ 给出 $x\leq\sqrt y$。只有 $0<y<1$ 时存在非退化截面，所以
\[
f_Y(y)=
\begin{cases}
\displaystyle\int_y^{\sqrt y}6\,dx=6(\sqrt y-y),&0<y<1,\\
0,&\text{其他}.
\end{cases}
\]
端点取 $0$ 不影响积分。核验两个边缘的归一化：
\[
\int_0^1 6(x-x^2)\,dx=1,\qquad
\int_0^1 6(\sqrt y-y)\,dy=6\left(\frac23-\frac12\right)=1.
\]

\item \begin{align*}
P(-1<X\leq1/2)
&=\int_0^{1/2}6(x-x^2)\,dx\\
&=[3x^2-2x^3]_0^{1/2}=\frac12.
\end{align*}
\end{enumerate}
\end{proof}

\subsection{例题：各分量连续型不保证向量连续型}\label{note:c3:e6}
\paragraph{题目}
设 $X$ 是一维连续型随机变量，令 $Y=X$。
\begin{enumerate}
\item 证明 $Y$ 也是一维连续型随机变量，但 $(X,Y)$ 不存在二维联合密度。

\item 求 $(X,Y)$ 的联合分布函数，并说明联合分布函数连续是否足以保证存在联合密度。
\end{enumerate}

\basis{原笔记第3.1节，例“各分量连续型不保证向量连续型”；源码\texttt{chapters/chapter3.tex}第674行。}
\begin{proof}
\begin{enumerate}
\item 由于 $Y=X$，$Y$ 与 $X$ 同分布，所以也是一维连续型随机变量。

事实上，$(X,Y)$ 必落在直线 $L=\{(x,y):x=y\}$ 上，所以 $P((X,Y)\in L)=1$。这条直线的二维面积为 $0$。补充一个面积说明：在 $[-M,M]^2$ 内，将对角线分成 $N$ 段，每段可放入边长 $2M/N$ 的正方形，总面积不超过 $4M^2/N$；令 $N\to\infty$，得这段对角线零测。整条直线是 $M=1,2,\ldots$ 所对应线段的可数并，故仍零测。

若存在联合密度 $f$，则应有
\[
1=P((X,Y)\in L)=\int_Lf(x,y)\,dx\,dy=0,
\]
矛盾。因此各分量分别具有密度，不足以保证存在联合密度。

\item 联合分布函数为：
\begin{align*}
F(t,s)&=P(X\leq t,Y\leq s)\\
&=P(X\leq t,X\leq s)=F_X(\min\{t,s\}).
\end{align*}
由于 $\min\{t,s\}=(t+s-|t-s|)/2$ 连续，且 $F_X$ 连续，$F$ 也是连续函数。这说明即使联合分布函数连续，向量仍可能没有联合密度。
\end{enumerate}
\end{proof}

\clearpage\section{随机变量的独立性}\label{k:rvindependence}

\begin{mydefinition}{随机变量的相互独立}{mutually-independent-random-variables}
设 $X_1,\ldots,X_n$ 定义在同一概率空间上。若对任意 $A_1,\ldots,A_n\in\mathcal B(\mathbb R)$，
\[
P(X_1\in A_1,\ldots,X_n\in A_n)=\prod_{i=1}^nP(X_i\in A_i),
\]
则称这些随机变量\textbf{相互独立}。对无穷列 $X_1,X_2,\ldots$，若任取其中有限个都相互独立，则称该列相互独立。
\end{mydefinition}
\origin{原笔记第3.2节，定义“随机变量的相互独立”；源码第699行。}

\begin{proposition}{独立随机变量的函数}{functions-of-independent-variables}
设 $X_1,\ldots,X_n$ 相互独立，$g_i:\mathbb R\to\mathbb R$ 为博雷尔函数，$\varphi:\mathbb R^k\to\mathbb R$ 为博雷尔函数，$1\leq k<n$。则：
\begin{enumerate}
\item $g_1(X_1),\ldots,g_n(X_n)$ 相互独立；
\item $\varphi(X_1,\ldots,X_k),X_{k+1},\ldots,X_n$ 相互独立。
\end{enumerate}
多元博雷尔函数指任意一维博雷尔集的原像均为相应维数的博雷尔集。
\end{proposition}
\origin{原笔记第3.2节，命题“独立随机变量的函数”；源码第711行。}

\begin{mytheorem}{独立性的分布判据}{independence-distribution-criteria}
设 $F$ 为 $(X_1,\ldots,X_n)$ 的联合分布函数。
\begin{enumerate}
\item $X_1,\ldots,X_n$ 相互独立，当且仅当对所有实数 $t_i$，
\[
F(t_1,\ldots,t_n)=\prod_{i=1}^nF_{X_i}(t_i).
\]
\item 若向量为离散型，则独立当且仅当其联合分布列在所有取值组合上满足
\[
p(x_1,\ldots,x_n)=\prod_{i=1}^np_{X_i}(x_i).
\]
\item 若向量具有联合密度 $f$，则独立当且仅当
\[
f(x_1,\ldots,x_n)=\prod_{i=1}^nf_{X_i}(x_i)
\quad\text{在 }\mathbb R^n\text{ 上几乎处处成立}.
\]
密度均采用前述非负博雷尔可测版本，可在零测异常集上改取 $0$ 为有限版本。
\end{enumerate}
\end{mytheorem}
\origin{原笔记第3.2节，定理“独立性的分布判据”；源码第754行。}

\begin{mycorollary}{独立的连续型分量具有联合密度}{independent-continuous-joint-density}
若 $X,Y$ 都是一维连续型随机变量且相互独立，则 $(X,Y)$ 是二维连续型随机向量，其联合密度可取为 $f_X(x)f_Y(y)$。进一步，对相互独立的连续型变量 $X_1,\ldots,X_n$，联合密度可取为 $\prod_i f_{X_i}(x_i)$。
\end{mycorollary}
\origin{原笔记第3.2节，推论“独立的连续型分量具有联合密度”；源码第835行。}

\subsection{例题：泊松总人数的随机分流}\label{note:c3:e7}
\paragraph{题目}
设一天进入邮局的总人数 $Z\sim P(\lambda)$，$\lambda>0$。给定 $Z=m$ 时，将这 $m$ 人独立地以概率 $p$ 标记为男性，以概率 $1-p$ 标记为女性，$0<p<1$。这一条件模型由“各人的标记相互独立且独立于总人数”保证。记男性人数为 $X$、女性人数为 $Y$，则 $Z=X+Y$，且
\[
P(X=k\mid Z=m)=\binom mkp^k(1-p)^{m-k},\qquad 0\leq k\leq m.
\]
\begin{enumerate}
\item 求男性人数 $X$ 与女性人数 $Y$ 的分布。

\item 求 $(X,Y)$ 的联合分布列，判断 $X,Y$ 是否独立，并说明 $p=0$ 或 $1$ 时的退化情形。
\end{enumerate}

\basis{原笔记第3.2节，例“泊松总人数的随机分流”；源码\texttt{chapters/chapter3.tex}第856行。}
\begin{proof}[解答]
\begin{enumerate}
\item 对 $k=0,1,\ldots$，由全概率公式，
\begin{align*}
P(X=k)
&=\sum_{m=k}^{\infty}\binom mkp^k(1-p)^{m-k}
 \frac{\lambda^m}{m!}e^{-\lambda}\\
&=\frac{e^{-\lambda}(p\lambda)^k}{k!}
 \sum_{m=k}^{\infty}\frac{[(1-p)\lambda]^{m-k}}{(m-k)!}\\
&=\frac{e^{-\lambda}(p\lambda)^k}{k!}e^{(1-p)\lambda}
=\frac{(p\lambda)^k}{k!}e^{-p\lambda}.
\end{align*}
所以 $X\sim P(p\lambda)$。将“男性”与“女性”的标记交换，条件成功概率变为 $1-p$，同一求和给出
\[
P(Y=j)=\frac{[(1-p)\lambda]^j}{j!}e^{-(1-p)\lambda},
\qquad Y\sim P((1-p)\lambda).
\]

\item 还须计算联合分布，而不能仅凭两个边缘都是泊松分布推出独立。对 $k,j\geq0$，
\begin{align*}
P(X=k,Y=j)
&=P(X=k,Z=k+j)\\
&=\binom{k+j}{k}p^k(1-p)^j
  \frac{\lambda^{k+j}}{(k+j)!}e^{-\lambda}\\
&=\frac{(p\lambda)^k}{k!}e^{-p\lambda}
  \frac{[(1-p)\lambda]^j}{j!}e^{-(1-p)\lambda}\\
&=P(X=k)P(Y=j).
\end{align*}
故由分布列判据，$X,Y$ 相互独立。$p=0$ 或 $1$ 时，一类人数恒为零，另一类等于 $Z$，结论相应退化为常量与 $Z$ 的独立性。
\end{enumerate}
\end{proof}

\begin{proposition}{非负函数分解判据}{independence-factorization}
\begin{enumerate}
\item 对二维离散型向量，$X,Y$ 独立当且仅当联合分布列可写为
\[
p(x,y)=h(x)g(y)
\]
在全部取值组合上成立，其中 $h,g$ 为非负有限值函数。
\item 对具有联合密度 $f$ 的二维向量，$X,Y$ 独立当且仅当存在非负有限值博雷尔函数 $h,g$，使
\[
f(x,y)=h(x)g(y)\qquad\text{几乎处处}.
\]
$h,g$ 本身不必各自归一化。区域外的零值也须包含在分解中。
\end{enumerate}
\end{proposition}
\origin{原笔记第3.2节，命题“非负函数分解判据”；源码第889行。}

\subsection{例题：乘积区域与三角形区域中的独立性}\label{note:c3:e8}
\paragraph{题目}
记 $I_A$ 为集合 $A$ 的示性函数：在 $A$ 内取 $1$，外取 $0$。
在下列两种情形中，分别求边缘密度并判断 $X,Y$ 是否独立，说明理由。
\begin{enumerate}
\item 设联合密度为
\[
f(x,y)=6xy^2 I_{(0,1)\times(0,1)}(x,y).
\]

\item 设联合密度为
\[
f(x,y)=24xyI_T(x,y),\qquad
T=\{(x,y):x>0,\ y>0,\ x+y<1\}.
\]
\end{enumerate}

\basis{原笔记第3.2节，例“乘积区域与三角形区域中的独立性”；源码\texttt{chapters/chapter3.tex}第930行。}
\begin{proof}[解答]
\begin{enumerate}
\item 由于 $I_{(0,1)\times(0,1)}(x,y)=I_{(0,1)}(x)I_{(0,1)}(y)$，
\[
f(x,y)=[6xI_{(0,1)}(x)]\,[y^2I_{(0,1)}(y)].
\]
两个非负因子的积分分别为 $3$ 和 $1/3$，乘积为 $1$，故满足分解判据。更直接地，
\[
f_X(x)=2xI_{(0,1)}(x),\qquad
f_Y(y)=3y^2I_{(0,1)}(y),
\]
二者乘积正是联合密度，所以 $X,Y$ 独立。

\item 虽然 $24xy$ 是乘积，但三角形限制也必须计入。取
\[
A=\{X>1/2\},\qquad B=\{Y>1/2\}.
\]
两个事件的交集在 $T$ 上为空，故 $P(A\cap B)=0$。边缘密度为
\[
f_X(x)=\int_0^{1-x}24xy\,dy=12x(1-x)^2,\qquad 0<x<1.
\]
区域与密度交换 $x,y$ 后不变，故 $f_Y(y)=12y(1-y)^2$（$0<y<1$），区域外均为 $0$。因此
\[
P(A)=\int_{1/2}^1 12x(1-x)^2\,dx=\frac5{16},
\qquad P(B)=\frac5{16}.
\]
于是 $P(A\cap B)=0\ne25/256=P(A)P(B)$，所以 $X,Y$ 不独立。这也具体说明了区域限制不能忽略。
\end{enumerate}
\end{proof}

\begin{mydefinition}{独立同分布}{iid-random-variables}
若 $X_1,\ldots,X_n$ 相互独立，并且
\[
F_{X_1}(t)=\cdots=F_{X_n}(t),\qquad \forall t\in\mathbb R,
\]
则称它们\textbf{独立同分布}，简写为 i.i.d.。
\end{mydefinition}
\origin{原笔记第3.2节，定义“独立同分布”；源码第969行。}

\subsection{例题：独立同分布的指数变量的联合密度}\label{note:c3:e9}
\paragraph{题目}
设 $X,Y$ 独立同分布于 $\operatorname{Exp}(\lambda)$，$\lambda>0$。求 $(X,Y)$ 的联合密度并核验其归一化。

\basis{原笔记第3.2节，例“独立同分布的指数变量的联合密度”；源码\texttt{chapters/chapter3.tex}第979行。}
\begin{proof}[解答]
\[
f_X(x)=\lambda e^{-\lambda x}I_{(0,\infty)}(x),\qquad
f_Y(y)=\lambda e^{-\lambda y}I_{(0,\infty)}(y).
\]
由独立连续型分量的联合密度公式，
\[
f(x,y)=
\begin{cases}
\lambda^2e^{-\lambda(x+y)},&x>0,\ y>0,\\
0,&\text{其他}.
\end{cases}
\]
其全积分等于两个边缘密度积分之积 $1\cdot1=1$。
\end{proof}

\begin{mytheorem}{独立随机变量求和的卷积公式}{sum-convolution}
设 $X,Y$ 相互独立。
\begin{enumerate}
\item 若二者为离散型，不同可能取值分别为 $a_i,b_j$，则
\[
p_{X+Y}(z)=\sum_i p_X(a_i)p_Y(z-a_i)
=\sum_jp_X(z-b_j)p_Y(b_j).
\]
不属于相应取值集合的单点概率取 $0$，有限取值时求和亦为有限和。
\item 若二者为连续型，密度分别为 $f_X,f_Y$，则 $X+Y$ 也是连续型，一个密度版本为
\[
f_{X+Y}(z)=\int_{\mathbb R}f_X(x)f_Y(z-x)\,dx
=\int_{\mathbb R}f_X(z-y)f_Y(y)\,dy.
\]
积分几乎处处有限，异常点可另取 $0$。这些求和或积分公式称为\textbf{卷积公式}。
\end{enumerate}
\end{mytheorem}
\origin{原笔记第3.2节，定理“独立随机变量求和的卷积公式”；源码第998行。}

\begin{mytheorem}{二项分布的可加性}{binomial-additivity}
设 $X_1,\ldots,X_n$ 相互独立，且 $X_i\sim B(m_i,p)$，各变量具有相同的成功概率 $p$，则
\[
\sum_{i=1}^nX_i\sim B\left(\sum_{i=1}^nm_i,p\right).
\]
\end{mytheorem}
\origin{原笔记第3.2节，定理“二项分布的可加性”；源码第1051行。}

\begin{mytheorem}{泊松分布的可加性}{poisson-additivity}
设 $X_1,\ldots,X_n$ 相互独立，$X_i\sim P(\lambda_i)$，$\lambda_i>0$，则
\[
\sum_{i=1}^nX_i\sim P\left(\sum_{i=1}^n\lambda_i\right).
\]
\end{mytheorem}
\origin{原笔记第3.2节，定理“泊松分布的可加性”；源码第1084行。}

\begin{mytheorem}{Gamma 分布的可加性}{gamma-additivity}
设 $X_1,\ldots,X_n$ 相互独立，且 $X_i\sim\Gamma(\alpha_i,\lambda)$，其中 $\alpha_i>0$，公共率参数 $\lambda>0$，则
\[
\sum_{i=1}^nX_i\sim\Gamma\left(\sum_{i=1}^n\alpha_i,\lambda\right).
\]
\end{mytheorem}
\origin{原笔记第3.2节，定理“Gamma 分布的可加性”；源码第1105行。}

\begin{mydefinition}{卡方分布}{chi-square-distribution}
设 $Z_1,\ldots,Z_n$ 独立同分布于 $N(0,1)$，$n$ 为正整数。称
\[
Y=\sum_{i=1}^nZ_i^2
\]
服从自由度为 $n$ 的\textbf{卡方分布}，记为 $Y\sim\chi^2(n)$。
\end{mydefinition}
\origin{原笔记第3.2节，定义“卡方分布”；源码第1142行。}

\begin{proposition}{卡方分布与 Gamma 分布}{chi-square-gamma}
自由度为 $n$ 的卡方分布就是 $\Gamma(n/2,1/2)$，其中第二参数为率参数。
\end{proposition}
\origin{原笔记第3.2节，命题“卡方分布与 Gamma 分布”；源码第1150行。}

\begin{mytheorem}{独立正态随机变量的线性组合}{normal-linear-combination}
设 $X_i\sim N(\mu_i,\sigma_i^2)$ 相互独立，$\sigma_i>0$，$a_i,c\in\mathbb R$。记
\[
M=\sum_{i=1}^na_i\mu_i+c,\qquad V=\sum_{i=1}^na_i^2\sigma_i^2.
\]
若 $V>0$，则
\[
\sum_{i=1}^na_iX_i+c\sim N(M,V).
\]
若 $V=0$，则所有 $a_i=0$，该线性组合恒等于 $c$。
\end{mytheorem}
\origin{原笔记第3.2节，定理“独立正态随机变量的线性组合”；源码第1169行。}

\subsection{例题：两个独立均匀变量之和}\label{note:c3:e10}
\paragraph{题目}
设 $X,Y$ 独立同分布于 $U(0,1)$。求 $X+Y$ 的密度。

\basis{原笔记第3.2节，例“两个独立均匀变量之和”；源码\texttt{chapters/chapter3.tex}第1209行。}
\begin{proof}[解答]
卷积的被积函数在
\[
0<x<1,\qquad 0<z-x<1
\]
时为 $1$，其余为 $0$。两个不等式同时成立等价于
\[
\max(0,z-1)<x<\min(1,z).
\]
$0<z\leq1$ 时截面为 $(0,z)$，$1<z<2$ 时截面为 $(z-1,1)$，其他情形为空或退化。因此
\[
f_{X+Y}(z)=
\begin{cases}
\displaystyle\int_0^z1\,dx=z,&0<z\leq1,\\
\displaystyle\int_{z-1}^11\,dx=2-z,&1<z<2,\\
0,&\text{其他}.
\end{cases}
\]
这个分布称为三角分布，密度全积分为 $1/2+1/2=1$。
\end{proof}

\clearpage\section{条件分布}\label{k:conditionaldist}

\begin{mydefinition}{给定事件的条件分布函数}{conditional-cdf-event}
对事件 $B$，若 $P(B)>0$，定义
\[
F_X(t\mid B)=P(X\leq t\mid B),\qquad t\in\mathbb R,
\]
称为 $B$ 发生条件下 $X$ 的\textbf{条件分布函数}。
\end{mydefinition}
\origin{原笔记第3.3节，定义“给定事件的条件分布函数”；源码第1239行。}

\begin{mydefinition}{离散型条件分布列与条件分布函数}{discrete-conditional-distribution}
设 $(X,Y)$ 的不同可能坐标值分别列为 $x_i,y_j$。
\begin{enumerate}
\item 固定满足 $P(Y=y_j)>0$ 的 $y_j$，定义
\[
p_{X\mid Y}(x_i\mid y_j)=P(X=x_i\mid Y=y_j)
=\frac{P(X=x_i,Y=y_j)}{P(Y=y_j)}.
\]
这些概率构成 $Y=y_j$ 条件下 $X$ 的条件分布列，其条件分布函数为
\[
F_{X\mid Y}(t\mid y_j)=\sum_{x_i\leq t}p_{X\mid Y}(x_i\mid y_j).
\]
\item 固定满足 $P(X=x_i)>0$ 的 $x_i$，定义
\[
p_{Y\mid X}(y_j\mid x_i)=P(Y=y_j\mid X=x_i)
=\frac{P(X=x_i,Y=y_j)}{P(X=x_i)}.
\]
相应条件分布函数为
\[
F_{Y\mid X}(t\mid x_i)=\sum_{y_j\leq t}p_{Y\mid X}(y_j\mid x_i).
\]
\end{enumerate}
求和遍历有限个或可数个取值，分母为零时不使用这些比值公式。
\end{mydefinition}
\origin{原笔记第3.3节，定义“离散型条件分布列与条件分布函数”；源码第1249行。}

\begin{proposition}{条件分布列的归一化}{conditional-pmf-normalization}
上述两种条件分布列，分别是概率空间 $(\Omega,\mathcal F,P(\cdot\mid Y=y_j))$ 下 $X$ 的分布列，以及 $(\Omega,\mathcal F,P(\cdot\mid X=x_i))$ 下 $Y$ 的分布列。
\end{proposition}
\origin{原笔记第3.3节，命题“条件分布列的归一化”；源码第1274行。}

\subsection{例题：给定独立泊松变量的和}\label{note:c3:e11}
\paragraph{题目}
设 $X\sim P(\lambda_1)$、$Y\sim P(\lambda_2)$ 相互独立，$\lambda_1,\lambda_2>0$，求给定 $X+Y=n$ 时 $X$ 的分布，$n$ 为非负整数。

\basis{原笔记第3.3节，例“给定独立泊松变量的和”；源码\texttt{chapters/chapter3.tex}第1298行。}
\begin{proof}[解答]
泊松可加性给出 $X+Y\sim P(\lambda_1+\lambda_2)$，故该条件事件具有正概率。对 $0\leq k\leq n$，
\begin{align*}
P(X=k\mid X+Y=n)
&=\frac{P(X=k,Y=n-k)}{P(X+Y=n)}\\
&=\frac{\dfrac{\lambda_1^k}{k!}e^{-\lambda_1}
 \dfrac{\lambda_2^{n-k}}{(n-k)!}e^{-\lambda_2}}
 {\dfrac{(\lambda_1+\lambda_2)^n}{n!}e^{-(\lambda_1+\lambda_2)}}\\
&=\binom nk\left(\frac{\lambda_1}{\lambda_1+\lambda_2}\right)^k
 \left(\frac{\lambda_2}{\lambda_1+\lambda_2}\right)^{n-k}.
\end{align*}
其他 $k$ 的条件概率为 $0$。所以条件分布为
\[
X\mid(X+Y=n)\sim B\left(n,\frac{\lambda_1}{\lambda_1+\lambda_2}\right).
\]
当 $n=0$ 时，条件下 $X=0$ 的概率为 $1$。
\end{proof}

\begin{mydefinition}{连续型条件概率的邻域极限}{continuous-conditional-neighborhood}
固定 $y$ 及博雷尔集 $A$。若充分小的 $\varepsilon>0$ 均满足
\[
P(y-\varepsilon<Y\leq y+\varepsilon)>0,
\]
且下列极限存在，则用邻域极限记
\[
P(X\in A\mid Y=y)
=\lim_{\varepsilon\downarrow0}
\frac{P(X\in A,\ y-\varepsilon<Y\leq y+\varepsilon)}
 {P(y-\varepsilon<Y\leq y+\varepsilon)}.
\]
即先计算 $P(X\in A\mid y-\varepsilon<Y\leq y+\varepsilon)$，再令 $\varepsilon\downarrow0$。
\end{mydefinition}
\origin{原笔记第3.3节，定义“连续型条件概率的邻域极限”；源码第1326行。}

\begin{mytheorem}{由联合密度求条件密度}{conditional-density-ratio}
设 $(X,Y)$ 的联合密度为 $f$，边缘密度选为截面积分
\[
f_Y(y)=\int_{\mathbb R}f(x,y)\,dx,\qquad
f_X(x)=\int_{\mathbb R}f(x,y)\,dy.
\]
\begin{enumerate}
\item 对 $0<f_Y(y)<\infty$ 的 $y$，定义
\[
f_{X\mid Y}(x\mid y)=\frac{f(x,y)}{f_Y(y)},\qquad
P(X\in A\mid Y=y)=\int_Af_{X\mid Y}(x\mid y)\,dx.
\]
它给出 $X$ 在 $Y=y$ 条件下的一个条件密度版本。
\item 对 $0<f_X(x)<\infty$ 的 $x$，相应定义为
\[
f_{Y\mid X}(y\mid x)=\frac{f(x,y)}{f_X(x)},\qquad
P(Y\in A\mid X=x)=\int_Af_{Y\mid X}(y\mid x)\,dy.
\]
\end{enumerate}
进一步，固定 $A$，若 $g_A(v)=\int_Af(u,v)\,du$ 及 $f_Y(v)$ 在 $v=y$ 连续，且 $f_Y(y)>0$，则第（1）项与该集合 $A$ 的邻域极限一致。第（2）项有交换坐标后的对应结论。
\end{mytheorem}
\origin{原笔记第3.3节，定理“由联合密度求条件密度”；源码第1350行。}

\begin{mycorollary}{条件分布函数的积分公式}{conditional-cdf-density}
在上述有效条件值处，
\begin{enumerate}
\item $F_{X\mid Y}(t\mid y)=\displaystyle\int_{-\infty}^t
f_{X\mid Y}(x\mid y)\,dx$；
\item $F_{Y\mid X}(t\mid x)=\displaystyle\int_{-\infty}^t
f_{Y\mid X}(y\mid x)\,dy$。
\end{enumerate}
\end{mycorollary}
\origin{原笔记第3.3节，推论“条件分布函数的积分公式”；源码第1422行。}

\subsection{例题：固定条件值与区间条件事件}\label{note:c3:e12}
\paragraph{题目}
设
\[
f(x,y)=
\begin{cases}
xe^{-y},&0<x<y,\\
0,&\text{其他}.
\end{cases}
\]
\begin{enumerate}
\item 求 $f_{X\mid Y}(x\mid y)$（$y>0$）。

\item 求 $f_{X\mid Y}(x\mid1/2)$。

\item 求 $P(-1<X\leq1/4\mid Y=1/2)$。

\item 求 $P(-1<X\leq1/4\mid Y\leq1/2)$。
\end{enumerate}

\basis{原笔记第3.3节，例“固定条件值与区间条件事件”；源码\texttt{chapters/chapter3.tex}第1445行。}
\begin{proof}[解答]
\begin{enumerate}
\item 固定 $y>0$，有
\[
f_Y(y)=\int_0^yxe^{-y}\,dx=\frac{y^2}2e^{-y},
\]
$y\leq0$ 时取 $0$。其积分为 $\Gamma(3)/2=1$，故给定联合密度也已归一化。因此
\[
f_{X\mid Y}(x\mid y)=
\begin{cases}
2x/y^2,&0<x<y,\\
0,&\text{其他},
\end{cases}\qquad y>0.
\]

\item 代入 $y=1/2$ 得
\[
f_{X\mid Y}(x\mid1/2)=
\begin{cases}
8x,&0<x<1/2,\\
0,&\text{其他}.
\end{cases}
\]

\item \[
P(-1<X\leq1/4\mid Y=1/2)=\int_0^{1/4}8x\,dx=\frac14.
\]

\item 对于条件事件 $Y\leq1/2$，应对整个条件区域积分。固定 $y$ 时，$x$ 的上限为 $\min(y,1/4)$，需在 $y=1/4$ 处分段：
\begin{align*}
&P(-1<X\leq1/4\mid Y\leq1/2)\\
&=\frac{\displaystyle
 \int_0^{1/4}\int_0^yxe^{-y}\,dx\,dy+
 \int_{1/4}^{1/2}\int_0^{1/4}xe^{-y}\,dx\,dy}
 {\displaystyle\int_0^{1/2}\frac{y^2}2e^{-y}\,dy}\\
&=\frac{\displaystyle
 \int_0^{1/4}\frac{y^2}2e^{-y}\,dy+
 \frac1{32}\int_{1/4}^{1/2}e^{-y}\,dy}
 {\displaystyle\int_0^{1/2}\frac{y^2}2e^{-y}\,dy}\\
&=\frac{1-\frac54e^{-1/4}-\frac1{32}e^{-1/2}}
 {1-\frac{13}8e^{-1/2}}
\approx0.524403.
\end{align*}
它与给定单点条件值 $Y=1/2$ 时的结果不同。
\end{enumerate}
\end{proof}

\subsection{例题：\texorpdfstring{$t$ 分布的密度}{t 分布的密度}}\label{note:c3:e13}
\paragraph{题目}
设 $Z\sim N(0,1)$、$Y\sim\chi^2(n)$ 相互独立，$n$ 为正整数。定义
\[
T=\frac{Z}{\sqrt{Y/n}}.
\]
称 $T$ 服从自由度为 $n$ 的 \textbf{$t$ 分布}。由于 $Y>0$ 以概率 $1$ 成立，该比值在全概率事件上有定义；其他样本点上可另定义为 $0$。
求 $T$ 的密度。

\basis{原笔记第3.3节，例“$t$ 分布的密度”；源码\texttt{chapters/chapter3.tex}第1500行。}
\begin{proof}[解答]
固定 $y>0$，由独立性，$Z$ 在条件 $Y=y$ 下仍服从 $N(0,1)$，故可取
\[
T\mid(Y=y)\sim N(0,n/y),\qquad
f_{T\mid Y}(t\mid y)=\sqrt{\frac{y}{2\pi n}}e^{-yt^2/(2n)}.
\]
又因 $\chi^2(n)=\Gamma(n/2,1/2)$，
\[
f_Y(y)=\frac{y^{n/2-1}e^{-y/2}}{2^{n/2}\Gamma(n/2)},\qquad y>0.
\]
条件密度与边缘密度相乘，得到
\[
f_{T,Y}(t,y)=
\frac{y^{(n-1)/2}e^{-\frac y2(1+t^2/n)}}
 {2^{(n+1)/2}\sqrt{n\pi}\Gamma(n/2)},
\qquad t\in\mathbb R,\ y>0,
\]
其他地方取 $0$。这一乘法也可直接从变换 $z=t\sqrt{y/n}$ 核验：其雅可比绝对值为 $\sqrt{y/n}$，原联合密度为 $f_Z(z)f_Y(y)$。

令 $a=(n+1)/2$、$b=(1+t^2/n)/2>0$，再作 $u=by$ 的换元，得
\[
\int_0^\infty y^{a-1}e^{-by}\,dy=b^{-a}\Gamma(a).
\]
所以
\begin{align*}
f_T(t)
&=\int_0^\infty f_{T,Y}(t,y)\,dy\\
&=\frac{\Gamma((n+1)/2)}{\sqrt{n\pi}\Gamma(n/2)}
 \left(1+\frac{t^2}n\right)^{-(n+1)/2},
\qquad t\in\mathbb R.
\end{align*}
\end{proof}

\begin{proposition}{离散与连续变量的条件分布乘法公式}{mixed-conditional-bayes}
设 $X$ 连续型，$Y$ 离散型。对 $p_j=P(Y=y_j)>0$ 的取值，设 $q_j(x)=f_{X\mid Y}(x\mid y_j)$ 是条件密度。则 $X$ 的边缘密度可取为
\[
f_X(x)=\sum_jp_jq_j(x).
\]
在 $0<f_X(x)<\infty$ 处，可取
\[
P(Y=y_j\mid X=x)=\frac{p_jq_j(x)}{f_X(x)}.
\]
从而
\[
P(Y=y_j\mid X=x)f_X(x)
=f_{X\mid Y}(x\mid y_j)P(Y=y_j)
\quad\text{几乎处处}.
\]
求和仅遍历正概率取值；$p_j=0$ 的项可按零项处理。
\end{proposition}
\origin{原笔记第3.3节，命题“离散与连续变量的条件分布乘法公式”；源码第1549行。}

\subsection{例题：均匀先验与二项抽样}\label{note:c3:e14}
\paragraph{题目}
设 $X\sim U(0,1)$，给定 $X=x$ 时 $Y\sim B(M,x)$，其中 $M$ 为正整数。求给定 $Y=n$ 时 $X$ 的分布，$0\leq n\leq M$。

\basis{原笔记第3.3节，例“均匀先验与二项抽样”；源码\texttt{chapters/chapter3.tex}第1588行。}
\begin{proof}[解答]
由乘法公式，对 $0<x<1$，
\[
f_{X\mid Y}(x\mid n)
=\frac{f_X(x)P(Y=n\mid X=x)}{P(Y=n)}
=\frac{\binom Mn x^n(1-x)^{M-n}}{P(Y=n)}.
\]
分母由积分归一化：
\begin{align*}
P(Y=n)
&=\binom Mn\int_0^1x^n(1-x)^{M-n}\,dx\\
&=\binom Mn B(n+1,M-n+1)\\
&=\frac{M!}{n!(M-n)!}\frac{n!(M-n)!}{(M+1)!}
=\frac1{M+1}>0.
\end{align*}
所以
\[
f_{X\mid Y}(x\mid n)=
\begin{cases}
\displaystyle\frac{x^n(1-x)^{M-n}}{B(n+1,M-n+1)},&0<x<1,\\
0,&\text{其他}.
\end{cases}
\]
即 $X\mid(Y=n)\sim\operatorname{Beta}(n+1,M-n+1)$。

这里 $X$ 是未知成功概率，$Y$ 是给定 $X$ 后 $M$ 次条件独立伯努利试验的成功次数。先验为 $U(0,1)=\operatorname{Beta}(1,1)$，观察到 $n$ 次成功后，后验参数变为 $1+n$、$1+M-n$。独立性是给定同一个成功概率 $X=x$ 后的条件独立性。
\end{proof}

\clearpage\section{次序统计量}\label{k:order}
本节公式中的样本$X_1,\ldots,X_n$独立同分布，公共分布函数为$F$；涉及密度时假设公共分布连续型且密度为$f$。

\begin{mydefinition}{次序统计量}{order-statistics}
对每个样本点 $\omega$，将实数 $X_1(\omega),\ldots,X_n(\omega)$ 按从小到大排列，保留重复值，得到
\[
X_{(1)}(\omega)\leq\cdots\leq X_{(n)}(\omega).
\]
对应的随机变量 $X_{(1)},\ldots,X_{(n)}$ 称为 $X_1,\ldots,X_n$ 的\textbf{次序统计量}。
\end{mydefinition}
\origin{原笔记第3.4节，定义“次序统计量”；源码第1621行。}

\subsection{例题：三个变量的排序}\label{note:c3:e15}
\paragraph{题目}
在样本点 $\omega$ 和 $\widetilde\omega$ 上，分别有
\[
(X_1(\omega),X_2(\omega),X_3(\omega))=(2,1,3),
\qquad
(X_1(\widetilde\omega),X_2(\widetilde\omega),X_3(\widetilde\omega))=(0,2,1).
\]
求这两个样本点上的次序统计量取值，并指出排序后的三个值分别来自哪个原变量。

\basis{原笔记第3.4节，例“三个变量的排序”；源码\texttt{chapters/chapter3.tex}第1635行。}
\begin{proof}[解答]
若在样本点 $\omega$ 上，
\[
(X_1(\omega),X_2(\omega),X_3(\omega))=(2,1,3),
\]
则
\[
(X_{(1)}(\omega),X_{(2)}(\omega),X_{(3)}(\omega))=(1,2,3),
\]
分别取自 $X_2,X_1,X_3$。若在另一样本点 $\widetilde\omega$ 上，原取值为 $(0,2,1)$，则排序取值为 $(0,1,2)$，分别取自 $X_1,X_3,X_2$。因此对应原变量的排列会随样本点改变。
\end{proof}

\begin{proposition}{最小值与最大值的分布}{minimum-maximum-cdf}
对任意 $t\in\mathbb R$，
\[
F_{X_{(1)}}(t)=1-(1-F(t))^n,\qquad
F_{X_{(n)}}(t)=F(t)^n.
\]
\end{proposition}
\origin{原笔记第3.4节，命题“最小值与最大值的分布”；源码第1649行。}

\begin{proposition}{连续独立样本无相同取值}{no-ties-continuous-sample}
对 $i\ne j$，$P(X_i=X_j)=0$。因此以概率 $1$ 有
\[
X_{(1)}<X_{(2)}<\cdots<X_{(n)}.
\]
\end{proposition}
\origin{原笔记第3.4节，命题“连续独立样本无相同取值”；源码第1673行。}

\begin{mytheorem}{次序统计量的联合密度}{joint-density-order-statistics}
次序统计量 $(X_{(1)},\ldots,X_{(n)})$ 的联合密度可取为
\[
f_{\mathrm{ord}}(t_1,\ldots,t_n)=
\begin{cases}
\displaystyle n!\prod_{i=1}^nf(t_i),&t_1<\cdots<t_n,\\
0,&\text{其他}.
\end{cases}
\]
\end{mytheorem}
\origin{原笔记第3.4节，定理“次序统计量的联合密度”；源码第1687行。}

\subsection{例题：路上随机位置的最小间距}\label{note:c3:e16}
\paragraph{题目}
三个人的位置 $X_1,X_2,X_3$ 相互独立且都服从 $U(0,1)$，距离以英里为单位。
\begin{enumerate}
\item 设 $0\leq d\leq1/2$，求任意两人之间的距离都大于 $d$ 的概率。

\item 推广到 $n\geq2$ 个人的位置独立同分布于 $U(0,1)$ 的情形，求 $d\geq0$ 时任意两人之间的距离都大于 $d$ 的概率。
\end{enumerate}

\basis{原笔记第3.4节，例“路上随机位置的最小间距”；源码\texttt{chapters/chapter3.tex}第1718行。}
\begin{proof}[解答]
\begin{enumerate}
\item 将位置排序后，任意两点的距离都大于 $d$，当且仅当两个相邻间距都大于 $d$。排序向量的密度在 $0<t_1<t_2<t_3<1$ 内为 $3!=6$，所以
\begin{align*}
P&=P(X_{(2)}>X_{(1)}+d,\ X_{(3)}>X_{(2)}+d)\\
&=6\int_0^{1-2d}\int_{t_1+d}^{1-d}\int_{t_2+d}^1
 dt_3\,dt_2\,dt_1\\
&=6\int_0^{1-2d}\int_{t_1+d}^{1-d}(1-t_2-d)\,dt_2\,dt_1\\
&=3\int_0^{1-2d}(1-t_1-2d)^2\,dt_1
=(1-2d)^3.
\end{align*}
$d=1/2$ 时概率为 $0$，$d=0$ 时由于无相同取值，概率为 $1$。

\item 若有 $n\geq2$ 个人，位置仍独立同分布于 $U(0,1)$，对 $0\leq d<1/(n-1)$，令
\[
s_i=t_i-(i-1)d,\qquad L=1-(n-1)d.
\]
相邻间距大于 $d$ 的区域恰好变为
\[
0<s_1<s_2<\cdots<s_n<L.
\]
平移的雅可比为 $1$；这个有序区域的体积是立方体 $(0,L)^n$ 体积的 $1/n!$，因为除了坐标相等的零测边界外，立方体分成 $n!$ 个等体积的排列区域。乘上排序密度 $n!$，所求概率为
\[
[1-(n-1)d]^n.
\]
当 $d\geq1/(n-1)$ 时，$n-1$ 个相邻间距之和无法严格小于 $1$，所求概率为 $0$。
\end{enumerate}
\end{proof}

\begin{mytheorem}{单个次序统计量的密度}{single-order-density}
设 $X_1,\ldots,X_n$ 独立同分布于连续型变量，共同分布函数为 $F$、密度为 $f$。第 $k$ 个次序统计量的密度可取为
\[
f_{X_{(k)}}(t)=
\frac{n!}{(k-1)!(n-k)!}
F(t)^{k-1}(1-F(t))^{n-k}f(t),
\qquad 1\leq k\leq n.
\]
指数为 $0$ 的因子按 $1$ 处理，包括 $F(t)=0$ 或 $1$ 的端点情形。
\end{mytheorem}
\origin{原笔记第3.4节，定理“单个次序统计量的密度”；源码第1747行。}

\begin{mytheorem}{两个次序统计量的联合密度}{pair-order-density}
对 $1\leq i<j\leq n$，记
\[
C_{ij}=\frac{n!}{(i-1)!(j-i-1)!(n-j)!}.
\]
则 $(X_{(i)},X_{(j)})$ 的联合密度在 $x<y$ 时可取为
\[
\begin{aligned}
f_{i,j}(x,y)
={}&C_{ij}\,F(x)^{i-1}[F(y)-F(x)]^{j-i-1}\\
&{}\cdot[1-F(y)]^{n-j}f(x)f(y),
\end{aligned}
\]
在 $x\geq y$ 时取 $0$。其中零次幂仍按 $1$ 处理。
\end{mytheorem}
\origin{原笔记第3.4节，定理“两个次序统计量的联合密度”；源码第1779行。}

\subsection{例题：样本极差的分布}\label{note:c3:e17}
\paragraph{题目}
设 $X_1,\ldots,X_n$ 独立同分布，公共分布为连续型，分布函数为 $F$、密度为 $f$，$n\geq2$。称
\[R=X_{(n)}-X_{(1)}\]
为样本的极差。
\begin{enumerate}
\item 求 $R$ 的分布函数。

\item 若原样本服从 $U(0,1)$，求 $R$ 的分布函数与密度，指出其分布，并说明 $n=1$ 时的情形。
\end{enumerate}

\basis{原笔记第3.4节，例“样本极差的分布”；源码\texttt{chapters/chapter3.tex}第1817行。}
\begin{proof}[解答]
\begin{enumerate}
\item 由两个次序统计量的密度公式，
\[
f_{1,n}(x,y)=n(n-1)[F(y)-F(x)]^{n-2}f(x)f(y),
\qquad x<y.
\]
$r<0$ 时 $F_R(r)=0$。对 $r\geq0$，积分区域为 $x<y\leq x+r$，所以
\begin{align*}
F_R(r)
&=n(n-1)\int_{\mathbb R}f(x)
 \int_x^{x+r}[F(y)-F(x)]^{n-2}f(y)\,dy\,dx\\
&=n\int_{\mathbb R}[F(x+r)-F(x)]^{n-1}f(x)\,dx.
\end{align*}
内层积分利用 $F$ 的绝对连续性及 $F'=f$ 几乎处处；原函数为
$[F(y)-F(x)]^{n-1}/(n-1)$，在下端点为 $0$。

\item 若原样本服从 $U(0,1)$，则 $r\geq1$ 时 $F_R(r)=1$。对 $0<r<1$，在 $x=1-r$ 处分段：
\begin{align*}
F_R(r)
&=n\int_0^1[F(x+r)-F(x)]^{n-1}\,dx\\
&=n\left[\int_0^{1-r}r^{n-1}\,dx+
 \int_{1-r}^1(1-x)^{n-1}\,dx\right]\\
&=nr^{n-1}(1-r)+r^n
=nr^{n-1}-(n-1)r^n.
\end{align*}
因此
\[
F_R(r)=
\begin{cases}
0,&r\leq0,\\
nr^{n-1}-(n-1)r^n,&0<r<1,\\
1,&r\geq1,
\end{cases}
\]
且一个密度版本为
\[
f_R(r)=
\begin{cases}
n(n-1)r^{n-2}(1-r),&0<r<1,\\
0,&\text{其他}.
\end{cases}
\]
因为 $B(n-1,2)=1/[n(n-1)]$，故 $R\sim\operatorname{Beta}(n-1,2)$。若样本只有一个变量，则极差恒为 $0$，不使用上述 $n\geq2$ 的密度公式。
\end{enumerate}
\end{proof}

\clearpage\section{随机变量函数的联合分布}\label{k:jacobian}

\begin{mytheorem}{二维一一变换的密度公式}{bivariate-density-transformation}
设 $(X,Y)$ 具有联合密度 $f_{X,Y}$，且以概率 $1$ 落在开区域 $D$。设
\[
g:D\to E,\qquad (u,v)=(g_1(x,y),g_2(x,y))
\]
为到开区域 $E$ 的一一映射，$g_1,g_2$ 有连续偏导数，且在 $D$ 上
\[
J_g(x,y)=\det
\begin{pmatrix}
\partial g_1/\partial x&\partial g_1/\partial y\\
\partial g_2/\partial x&\partial g_2/\partial y
\end{pmatrix}\ne0.
\]
记逆变换为 $h=(h_1,h_2)$。则 $(U,V)=g(X,Y)$ 具有联合密度
\[
f_{U,V}(u,v)=
\begin{cases}
f_{X,Y}(h_1(u,v),h_2(u,v))\,|J_h(u,v)|,&(u,v)\in E,\\
0,&(u,v)\notin E,
\end{cases}
\]
其中
\[
J_h(u,v)=\frac{\partial(h_1,h_2)}{\partial(u,v)},\qquad
|J_h(u,v)|=\frac1{|J_g(h_1(u,v),h_2(u,v))|}.
\]
\end{mytheorem}
\origin{原笔记第3.5节，定理“二维一一变换的密度公式”；源码第1869行。}

\subsection{例题：引入辅助变量求和的密度}\label{note:c3:e18}
\paragraph{题目}
设 $(X,Y)$ 的联合密度为 $f(x,y)$，令 $U=X+Y$。
\begin{enumerate}
\item 通过引入辅助变量求 $U$ 的密度。

\item 若 $X,Y$ 独立，将结果写为边缘密度的卷积公式。
\end{enumerate}

\basis{原笔记第3.5节，例“引入辅助变量求和的密度”；源码\texttt{chapters/chapter3.tex}第1915行。}
\begin{proof}[解答]
\begin{enumerate}
\item 增加辅助变量 $V=X$，则
\[
U=X+Y,\quad V=X
\quad\Longleftrightarrow\quad
X=V,\quad Y=U-V.
\]
逆雅可比为
\[
\frac{\partial(x,y)}{\partial(u,v)}
=\det\begin{pmatrix}0&1\\1&-1\end{pmatrix}=-1.
\]
故
\[
f_{U,V}(u,v)=f(v,u-v).
\]
对辅助变量积分，得到
\[
f_{X+Y}(u)=\int_{\mathbb R}f(v,u-v)\,dv.
\]

\item 若 $X,Y$ 独立，则可取 $f(x,y)=f_X(x)f_Y(y)$，于是
\[
f_{X+Y}(u)=\int_{\mathbb R}f_X(v)f_Y(u-v)\,dv
=\int_{\mathbb R}f_X(u-v)f_Y(v)\,dv,
\]
与卷积公式一致。
\end{enumerate}
\end{proof}

\subsection{例题：独立 Gamma 变量的和与比例}\label{note:c3:e19}
\paragraph{题目}
设 $X\sim\Gamma(\alpha,\lambda)$、$Y\sim\Gamma(\beta,\lambda)$ 相互独立，$\alpha,\beta,\lambda>0$，$\lambda$ 为公共率参数。令
\[
U=X+Y,\qquad V=\frac{X}{X+Y}.
\]
因 $X,Y>0$ 以概率 $1$ 成立，该变换在全概率区域上有定义。
\begin{enumerate}
\item 求 $(U,V)$ 的联合密度。

\item 求 $U,V$ 的边缘密度并指出其分布。

\item 判断 $U,V$ 是否独立。
\end{enumerate}

\basis{原笔记第3.5节，例“独立 Gamma 变量的和与比例”；源码\texttt{chapters/chapter3.tex}第1943行。}
\begin{proof}[解答]
\begin{enumerate}
\item 逆变换为
\[
x=uv,\qquad y=u(1-v),
\]
故
\[
\frac{\partial(x,y)}{\partial(u,v)}
=\det\begin{pmatrix}v&u\\1-v&-u\end{pmatrix}=-u.
\]
原区域 $x>0,y>0$ 等价于 $u>0,0<v<1$；该区域中 $u>0$，所以逆雅可比绝对值为 $u$。

由独立性，原联合密度在正象限为
\[
f_{X,Y}(x,y)=
\frac{\lambda^{\alpha+\beta}}{\Gamma(\alpha)\Gamma(\beta)}
x^{\alpha-1}y^{\beta-1}e^{-\lambda(x+y)}.
\]
代入逆变换并乘上 $u$：
\begin{align*}
f_{U,V}(u,v)
&=\frac{\lambda^{\alpha+\beta}
 (uv)^{\alpha-1}[u(1-v)]^{\beta-1}}
 {\Gamma(\alpha)\Gamma(\beta)}e^{-\lambda u}\,u\\
&=\frac{\lambda^{\alpha+\beta}
 u^{\alpha+\beta-1}e^{-\lambda u}}
 {\Gamma(\alpha+\beta)}
 \frac{v^{\alpha-1}(1-v)^{\beta-1}}{B(\alpha,\beta)},
\qquad u>0,\ 0<v<1.
\end{align*}
区域外取 $0$。最后一步使用
$B(\alpha,\beta)=\Gamma(\alpha)\Gamma(\beta)/\Gamma(\alpha+\beta)$。

\item 右侧已经分解成一个 Gamma 密度和一个 Beta 密度的乘积，两因子各自积分为 $1$。因此
\[
f_U(u)=
\begin{cases}
\displaystyle\frac{\lambda^{\alpha+\beta}}{\Gamma(\alpha+\beta)}
u^{\alpha+\beta-1}e^{-\lambda u},&u>0,\\
0,&\text{其他},
\end{cases}
\]
\[
f_V(v)=
\begin{cases}
\displaystyle\frac{v^{\alpha-1}(1-v)^{\beta-1}}{B(\alpha,\beta)},&0<v<1,\\
0,&\text{其他}.
\end{cases}
\]
所以 $U\sim\Gamma(\alpha+\beta,\lambda)$、$V\sim\operatorname{Beta}(\alpha,\beta)$。

\item 由第（1）问所得联合密度等于第（2）问两个边缘密度之积，故 $U,V$ 相互独立。
\end{enumerate}
\end{proof}

\subsection{例题：标准正态变量的极坐标变换}\label{note:c3:e20}
\paragraph{题目}
设 $X,Y$ 相互独立，均服从 $N(0,1)$。令
\[
R=\sqrt{X^2+Y^2},
\]
并以 $\Theta\in[0,2\pi)$ 表示点 $(X,Y)$ 的极角；在原点约定 $\Theta=0$。
\begin{enumerate}
\item 求 $R,\Theta$ 的联合密度及两个边缘密度，判断独立性，并求 $R^2$ 的分布。

\item 反过来，取相互独立的 $U_1,U_2\sim U(0,1)$，令
\[X=\sqrt{-2\ln U_1}\cos(2\pi U_2),\qquad Y=\sqrt{-2\ln U_1}\sin(2\pi U_2).\]
证明生成的 $X,Y$ 相互独立且均服从 $N(0,1)$。
\end{enumerate}

\basis{原笔记第3.5节，例“标准正态变量的极坐标变换”；源码\texttt{chapters/chapter3.tex}第2013行。}
\begin{proof}[解答]
\begin{enumerate}
\item 原联合密度为
\[
f_{X,Y}(x,y)=\frac{1}{2\pi}e^{-(x^2+y^2)/2}.
\]
去掉原点和非负横轴后，剩余区域的概率为 $1$。极坐标变换将其一一映到 $r>0,\ 0<\theta<2\pi$，逆变换为
\[
x=r\cos\theta,\qquad y=r\sin\theta.
\]
其雅可比为
\[
\left|\frac{\partial(x,y)}{\partial(r,\theta)}\right|
=\left|\det
\begin{pmatrix}
\cos\theta&-r\sin\theta\\
\sin\theta&r\cos\theta
\end{pmatrix}\right|=r.
\]
所以可取联合密度
\[
f_{R,\Theta}(r,\theta)=
\begin{cases}
\displaystyle\frac{r}{2\pi}e^{-r^2/2},&r>0,\ 0<\theta<2\pi,\\
0,&\text{其他}.
\end{cases}
\]
两个边缘密度分别为
\[
f_R(r)=
\begin{cases}re^{-r^2/2},&r>0,\\0,&\text{其他},\end{cases}
\qquad
f_\Theta(\theta)=
\begin{cases}\displaystyle\frac{1}{2\pi},&0<\theta<2\pi,\\0,&\text{其他}.\end{cases}
\]
联合密度等于边缘密度的乘积，故 $R,\Theta$ 相互独立，且 $\Theta\sim U(0,2\pi)$。$R$ 的分布称为\textbf{瑞利分布}。对于 $t\geq0$，
\[
P(R^2>t)=\int_{\sqrt t}^{\infty}re^{-r^2/2}\,dr=e^{-t/2},
\]
故 $R^2\sim\operatorname{Exp}(1/2)=\Gamma(1,1/2)=\chi^2(2)$，其中指数分布和 Gamma 分布均采用率参数。

\item 反过来，取相互独立的 $U_1,U_2\sim U(0,1)$，令
\[
R=\sqrt{-2\ln U_1},\qquad \Theta=2\pi U_2.
\]
由于这两个变量分别由独立的 $U_1,U_2$ 决定，它们仍相互独立。对于 $t\geq0$，
\[
P(-2\ln U_1>t)=P(U_1<e^{-t/2})=e^{-t/2},
\]
因而 $R,\Theta$ 具有上面求出的分布。

令
\[
X=\sqrt{-2\ln U_1}\cos(2\pi U_2),\qquad
Y=\sqrt{-2\ln U_1}\sin(2\pi U_2).
\]
在去掉非负横轴的平面上，正向极坐标变换的雅可比绝对值为 $1/r$，于是
\begin{align*}
f_{X,Y}(x,y)
&=\frac{r}{2\pi}e^{-r^2/2}\frac1r\\
&=\frac{1}{\sqrt{2\pi}}e^{-x^2/2}
  \frac{1}{\sqrt{2\pi}}e^{-y^2/2}.
\end{align*}
被去掉的边界对应 $\Theta=0$，其概率为 $0$，不影响分布。密度分解说明生成的 $X,Y$ 相互独立，且均服从 $N(0,1)$。
\end{enumerate}
\end{proof}

\textbf{注：}\quad 设随机向量 $\boldsymbol X=(X_1,\ldots,X_n)$ 的密度为 $f_{\boldsymbol X}$，开集 $D\subseteq\mathbb R^n$ 满足 $P(\boldsymbol X\in D)=1$。若 $g:D\to E$ 是开集之间的一一映射，连续可微且处处满足 $\det Dg\ne0$，则逆映射 $h=g^{-1}$ 连续可微。$\boldsymbol Y=g(\boldsymbol X)$ 的密度可取为
\[
f_{\boldsymbol Y}(\boldsymbol y)=
\begin{cases}
f_{\boldsymbol X}(h(\boldsymbol y))|\det Dh(\boldsymbol y)|,
&\boldsymbol y\in E,\\
0,&\boldsymbol y\notin E.
\end{cases}
\]
在 $D$ 外可将变换固定定义为某个向量；由于 $P(\boldsymbol X\notin D)=0$，这不影响分布。

\subsection{例题：独立指数变量的部分和}\label{note:c3:e21}
\paragraph{题目}
设 $X_1,\ldots,X_n$ 相互独立，均服从 $\operatorname{Exp}(\lambda)$，$\lambda>0$。令
\[
Y_i=\sum_{j=1}^{i}X_j,\qquad i=1,\ldots,n.
\]
\begin{enumerate}
\item 求 $(Y_1,\ldots,Y_n)$ 的联合密度。

\item 求 $Y_n$ 的边缘分布。

\item 设 $n\geq2$，求给定 $Y_n=z$ 时前 $n-1$ 个部分和的条件联合密度。
\end{enumerate}

\basis{原笔记第3.5节，例“独立指数变量的部分和”；源码\texttt{chapters/chapter3.tex}第2112行。}
\begin{proof}[解答]
\begin{enumerate}
\item 逆变换为 $x_1=y_1$、$x_i=y_i-y_{i-1}$（$i\geq2$），其导数矩阵是对角元为 $1$、紧邻对角线下方为 $-1$ 的下三角矩阵，行列式为 $1$。原区域 $x_i>0$ 等价于 $0<y_1<\cdots<y_n$，且 $\sum_i x_i=y_n$，故
\[
f_{Y_1,\ldots,Y_n}(y_1,\ldots,y_n)=
\begin{cases}
\lambda^ne^{-\lambda y_n},&0<y_1<\cdots<y_n,\\
0,&\text{其他}.
\end{cases}
\]

\item 当 $n\geq2$、$z>0$ 时，对前 $n-1$ 个坐标积分，得
\begin{align*}
f_{Y_n}(z)
&=\lambda^ne^{-\lambda z}
  \int_{0<y_1<\cdots<y_{n-1}<z}dy_1\cdots dy_{n-1}\\
&=\frac{\lambda^n}{(n-1)!}z^{n-1}e^{-\lambda z}.
\end{align*}
积分区域的体积为 $z^{n-1}/(n-1)!$：除去坐标相等的零体积集合，立方体 $(0,z)^{n-1}$ 按坐标的 $(n-1)!$ 种排列分成体积相同的区域。$z\leq0$ 时密度取 $0$，所以 $Y_n\sim\Gamma(n,\lambda)$；$n=1$ 时结论就是指数分布本身。

\item 对每个 $z>0$，$f_{Y_n}(z)>0$，因此
\[
f_{Y_1,\ldots,Y_{n-1}\mid Y_n}(y_1,\ldots,y_{n-1}\mid z)=
\begin{cases}
\displaystyle\frac{(n-1)!}{z^{n-1}},
&0<y_1<\cdots<y_{n-1}<z,\\
0,&\text{其他}.
\end{cases}
\]
这正是 $n-1$ 个相互独立的 $U(0,z)$ 变量的次序统计量联合密度。排序前的变量独立；当 $n-1\geq2$ 时，排序后的变量受到大小次序的约束，并不相互独立。
\end{enumerate}
\end{proof}

\clearpage\section*{第十至第十二周作业}\addcontentsline{toc}{section}{第十至第十二周作业}
% !TEX root = 概率论期中复习.tex
\section{第十周：联合分布与边缘分布}

\subsection{教材第72页第1题}\label{ex:w10-b72-1}
从 $1,2,\ldots,n$ 中随机选取一个数记为 $\eta$，再从 $1,2,\ldots,\eta$ 中随机选取一数记为 $\xi$。求 $(\xi,\eta)$ 的联合分布和边缘分布。
\basis{原笔记第3章第1节“离散型随机向量的联合分布列与边缘分布列”，第3节“离散型条件分布列与条件分布函数”。}
\paragraph{思路}先按选取顺序相乘得到联合概率，再沿另一坐标求和。
\begin{proof}[解答]
对 $1\leq i\leq j\leq n$，第一步选到 $j$ 的概率是 $1/n$，在此条件下第二步选到 $i$ 的概率是 $1/j$。因此
\[
P(\xi=i,\eta=j)=P(\eta=j)P(\xi=i\mid\eta=j)=\frac1{nj}.
\]
不满足上述整数范围时，联合概率为 $0$。固定 $i$ 时，只有 $j=i,\ldots,n$ 能使 $\xi=i$，故
\[
P(\xi=i)=\frac1n\sum_{j=i}^n\frac1j,\qquad i=1,\ldots,n.
\]
固定 $j$ 时，有 $j$ 个可能的 $i$，所以
\[
P(\eta=j)=\sum_{i=1}^j\frac1{nj}=\frac1n,\qquad j=1,\ldots,n.
\]
\end{proof}
\paragraph{答案}联合分布为 $p(i,j)=1/(nj)$（$1\leq i\leq j\leq n$）；两个边缘分布如上。

\subsection{教材第72页第2题}\label{ex:w10-b72-2}
设 $(\xi,\eta)$ 为绝对连续型随机向量。证明 $P(\xi=\eta)=0$。
\basis{原笔记第3章第1节“联合密度的基本性质”及例“各分量连续型不保证向量连续型”。}
\paragraph{思路}把等值事件写成平面上对角线的概率，对联合密度积分。
\begin{proof}
记联合密度为 $f$。对每个固定 $x$，单点集合 $\{x\}$ 的一维长度为 $0$，故
\[
\int_{\mathbb R}I_{\{y=x\}}f(x,y)\,dy=0.
\]
被积函数非负，可以逐次积分，得到
\[
P(\xi=\eta)=\iint_{\{x=y\}}f(x,y)\,dx\,dy
=\int_{\mathbb R}\left(\int_{\mathbb R}I_{\{y=x\}}f(x,y)\,dy\right)dx=0.
\]
这里使用的是联合分布具有密度；仅有两个边缘分布连续不足以推出结论，例如取 $\eta=\xi$。
\end{proof}
\paragraph{答案}$P(\xi=\eta)=0$。

\subsection{教材第72页第4题}\label{ex:w10-b72-4}
随机向量 $(\xi,\eta)$ 的联合分布密度为
\[
p(x,y)=ae^{-(2x+y)}I_{\{x>0\}}I_{\{y>0\}}.
\]
求：（1）常数 $a$；（2）$\xi$ 的边缘密度；（3）$P(\xi\leq2,\eta\leq1)$；（4）$P(\xi+\eta\leq2)$。
\basis{原笔记第3章第1节“联合密度的基本性质”“固定坐标求积分上下限”。}
\paragraph{思路}归一化确定常数；矩形区域与三角形区域分别积分。
\begin{proof}[解答]
\begin{enumerate}
\item 由总积分为 $1$，
\[
1=a\int_0^\infty e^{-2x}dx\int_0^\infty e^{-y}dy=\frac a2,
\qquad a=2.
\]
\item 对 $x>0$，
\[
f_\xi(x)=\int_0^\infty2e^{-2x-y}dy=2e^{-2x}.
\]
对 $x\leq0$ 取 $f_\xi(x)=0$。同理 $f_\eta(y)=e^{-y}$（$y>0$），区域外为 $0$。
\item 所求区域是 $0<x\leq2,0<y\leq1$，故
\[
P(\xi\leq2,\eta\leq1)=\int_0^2\int_0^1 2e^{-2x-y}\,dy\,dx
=(1-e^{-4})(1-e^{-1}).
\]
\item 在正象限内，$x+y\leq2$ 等价于 $0<x<2,0<y\leq2-x$。因此
\begin{align*}
P(\xi+\eta\leq2)
&=\int_0^2 2e^{-2x}(1-e^{-(2-x)})\,dx\\
&=(1-e^{-4})-2e^{-2}(1-e^{-2})
=(1-e^{-2})^2.
\end{align*}
\end{enumerate}
\end{proof}
\paragraph{答案}$a=2$；$f_\xi(x)=2e^{-2x}I_{\{x>0\}}$；后两问为 $(1-e^{-4})(1-e^{-1})$、$(1-e^{-2})^2$。

\subsection{教材第72页第5题}\label{ex:w10-b72-5}
设随机向量 $(\xi,\eta)$ 的联合分布密度为
\[
p(x,y)=3(1-x)I_{\{0<x<y<1\}}.
\]
求其边缘分布密度。
\basis{原笔记第3章第1节“固定坐标求积分上下限”。}
\paragraph{思路}固定 $x$ 时 $y\in(x,1)$；固定 $y$ 时 $x\in(0,y)$。
\begin{proof}[解答]
对 $0<x<1$，
\[
f_\xi(x)=\int_x^1 3(1-x)\,dy=3(1-x)^2.
\]
对 $0<y<1$，
\[
f_\eta(y)=\int_0^y3(1-x)\,dx=3y-\frac32y^2.
\]
相应区间之外的密度均为 $0$。两个表达式都非负，而且
\[
\int_0^1 3(1-x)^2dx=1,\qquad
\int_0^1\left(3y-\frac32y^2\right)dy=1.
\]
\end{proof}
\paragraph{答案}$f_\xi(x)=3(1-x)^2I_{\{0<x<1\}}$，$f_\eta(y)=(3y-\frac32y^2)I_{\{0<y<1\}}$。

\subsection{教材第72页第6题}\label{ex:w10-b72-6}
函数
\[
F(x,y)=(1-e^{-xy})I_{\{x>0,y>0\}}
\]
是否为二维概率分布函数？
\basis{原笔记第3章第1节“联合分布函数的基本性质”：每个矩形的分布函数增量必须非负。}
\paragraph{思路}检验矩形 $(1,2]\times(1,2]$ 对应的增量。
\begin{proof}[解答]
若 $F$ 是联合分布函数，则矩形概率应为
\begin{align*}
P(1<X\leq2,1<Y\leq2)
&=F(2,2)-F(1,2)-F(2,1)+F(1,1)\\
&=2e^{-2}-e^{-1}-e^{-4}.
\end{align*}
由于 $e>2$，有 $2e^{-2}<e^{-1}$，上述增量严格小于 $0$，与概率非负矛盾。
\end{proof}
\paragraph{答案}不是二维概率分布函数。

\subsection{教材第81页第4题}\label{ex:w10-b81-4}
设随机向量 $(\xi,\eta)$ 的联合分布密度为
\[
p(x,y)=4xyI_{\{0<x,y<1\}}.
\]
求 $\xi$ 和 $\eta$ 的边缘分布密度，并证明 $\xi$ 和 $\eta$ 相互独立。
\basis{原笔记第3章第1节边缘密度公式，第2节连续型随机变量的独立性判据。}
\paragraph{思路}求边缘密度后，检查联合密度是否等于二者乘积。
\begin{proof}[解答]
对 $0<x<1$，$f_\xi(x)=\int_0^14xy\,dy=2x$；对 $0<y<1$，$f_\eta(y)=\int_0^14xy\,dx=2y$。区间外均取 $0$。

当 $0<x,y<1$ 时，$p(x,y)=(2x)(2y)=f_\xi(x)f_\eta(y)$；当至少一个坐标不在 $(0,1)$ 内时，两边都为 $0$。因而联合密度分解为两个边缘密度的乘积，$\xi,\eta$ 相互独立。
\end{proof}
\paragraph{答案}$f_\xi(x)=2xI_{\{0<x<1\}}$、$f_\eta(y)=2yI_{\{0<y<1\}}$，且二者独立。

\subsection{作业第2题：由事件独立性确定联合分布}\label{ex:w10-h2}
已知 $(X,Y)$ 的联合分布列如下，其中 $a,b$ 是常数，事件 $\{X=1\}$ 和 $\{X+Y=1\}$ 相互独立。
\[
\begin{array}{c|cc}
X\backslash Y&0&1\\\hline
0&0.4&a\\
1&b&0.1
\end{array}
\]
（1）求 $a,b$；（2）求 $(X,Y)$ 的边缘分布列；（3）判断 $X$ 和 $Y$ 是否相互独立。
\basis{原笔记第3章第1节例“由联合分布表求边缘分布列”，第2节离散型独立性判据。}
\paragraph{思路}先将给定的两个事件写成表格中格子的并，再使用独立性方程。
\begin{proof}[解答]
\begin{enumerate}
\item 归一化给出 $a+b=0.5$。令 $A=\{X=1\}$、$B=\{X+Y=1\}$，则
\[
P(A)=b+0.1,\quad P(B)=a+b=0.5,\quad P(AB)=P(X=1,Y=0)=b.
\]
由独立性，$b=0.5(b+0.1)$，解得 $b=0.1,a=0.4$。
\item 沿行、列求和，有
\[
P(X=0)=0.8,\quad P(X=1)=0.2;\qquad
P(Y=0)=P(Y=1)=0.5.
\]
\item 四个联合概率分别为
\[
0.4=0.8\cdot0.5,\quad0.4=0.8\cdot0.5,\quad
0.1=0.2\cdot0.5,\quad0.1=0.2\cdot0.5.
\]
每一个可能取值对的联合概率均等于边缘概率的乘积，故 $X,Y$ 独立。
\end{enumerate}
\end{proof}
\paragraph{答案}$a=0.4,b=0.1$；$X$ 的分布为 $(0.8,0.2)$，$Y$ 的分布为 $(0.5,0.5)$；相互独立。

\subsection{作业第3题：三项分布的边缘分布}\label{ex:w10-h3}
设正整数 $n\geq1$，$0<p_1,p_2,p_3<1$，$p_1+p_2+p_3=1$。二维离散型随机变量 $(X_1,X_2)$ 服从三项分布，即对非负整数 $k_1,k_2$，$k_1+k_2\leq n$，
\[
P(X_1=k_1,X_2=k_2)=\frac{n!}{k_1!k_2!(n-k_1-k_2)!}
p_1^{k_1}p_2^{k_2}p_3^{n-k_1-k_2}.
\]
求 $(X_1,X_2)$ 的边缘分布列，并写出 $X_1,X_2$ 分别满足什么常用分布。
\basis{原笔记第3章第1节“由联合分布列求边缘分布列”，第2章二项分布。}
\paragraph{思路}固定一种结果的出现次数，将其余两种结果合并，用二项式公式求和。
\begin{proof}[解答]
固定 $0\leq k\leq n$，
\begin{align*}
P(X_1=k)
&=\sum_{j=0}^{n-k}\frac{n!}{k!j!(n-k-j)!}p_1^kp_2^jp_3^{n-k-j}\\
&=\binom nkp_1^k\sum_{j=0}^{n-k}\binom{n-k}{j}p_2^jp_3^{n-k-j}\\
&=\binom nkp_1^k(p_2+p_3)^{n-k}
=\binom nkp_1^k(1-p_1)^{n-k}.
\end{align*}
同样，固定 $X_2=k$ 时对 $X_1=0,\ldots,n-k$ 求和，得到
\[
P(X_2=k)=\binom nkp_2^k(p_1+p_3)^{n-k}
=\binom nkp_2^k(1-p_2)^{n-k}.
\]
因此 $X_1\sim B(n,p_1)$，$X_2\sim B(n,p_2)$。
\end{proof}
\paragraph{答案}$X_1\sim B(n,p_1)$，$X_2\sim B(n,p_2)$，分布列如上。

\subsection{作业第4题：双曲线之间的联合密度}\label{ex:w10-h4}
设 $(X,Y)$ 有联合密度函数
\[
f(x,y)=Ax^{-3}y^{-2},\qquad 1<x,\quad x^{-1}<y<x,
\]
区域外为 $0$。（1）求 $A$；（2）求边缘密度 $f_X(x),f_Y(y)$；（3）求 $P(1<X<2)$、$P(\frac12<Y<2)$。
\basis{原笔记第3章第1节“联合密度的基本性质”及例“固定坐标求积分上下限”。}
\paragraph{思路}先对 $y$ 积分；反过来固定 $y$ 时，须同时满足 $x>1,x>y,x>1/y$，在 $y=1$ 处分段。
\begin{proof}[解答]
\begin{enumerate}
\item 对固定 $x>1$，
\[
\int_{1/x}^xy^{-2}dy=x-\frac1x.
\]
于是
\[
1=A\int_1^\infty(x^{-2}-x^{-4})\,dx=A\left(1-\frac13\right),
\qquad A=\frac32.
\]
\item 由上一计算，
\[
f_X(x)=\begin{cases}\frac32(x^{-2}-x^{-4}),&x>1,\\0,&\text{其他}.\end{cases}
\]
固定 $y>0$，若 $0<y<1$，则 $x>1/y$；若 $y\geq1$，则 $x>y$。设 $L=\max(y,1/y)$，有
\[
f_Y(y)=\frac32y^{-2}\int_L^\infty x^{-3}dx=\frac34y^{-2}L^{-2}.
\]
因此可取
\[
f_Y(y)=\begin{cases}\frac34,&0<y<1,\\\frac3{4y^4},&y\geq1,\\0,&y\leq0.\end{cases}
\]
\item 使用边缘密度，
\begin{align*}
P(1<X<2)&=\frac32\int_1^2(x^{-2}-x^{-4})dx
=\left[-\frac3{2x}+\frac1{2x^3}\right]_1^2=\frac5{16},\\
P\left(\frac12<Y<2\right)
&=\int_{1/2}^1\frac34dy+\int_1^2\frac3{4y^4}dy
=\frac38+\frac14\left(1-\frac18\right)=\frac{19}{32}.
\end{align*}
\end{enumerate}
\end{proof}
\paragraph{答案}$A=3/2$；边缘密度如上；两概率为 $5/16$、$19/32$。
% !TEX root = 概率论期中复习.tex
\section{第十一周：独立性、卷积与条件分布}

\subsection{教材第80页第1题}\label{ex:w11-b80-1}
设矩形 $(0,\pi/2)^2$ 上的随机向量 $(\xi,\eta)$ 有联合分布密度
\[
p(x,y)=\frac12\sin(x+y).
\]
对于 $0<y<\pi/2$ 求给定 $\eta=y$ 时 $\xi$ 的条件密度函数。
\basis{原笔记第3章第3节“由联合密度求条件密度”。}
\paragraph{思路}先对 $x$ 积分求条件变量的边缘密度，再取比值。
\begin{proof}[解答]
对 $0<y<\pi/2$，
\begin{align*}
f_\eta(y)&=\int_0^{\pi/2}\frac12\sin(x+y)dx\\
&=\frac12\bigl(\cos y-\cos(\pi/2+y)\bigr)
=\frac12(\cos y+\sin y)>0.
\end{align*}
所以相应条件密度可取为
\[
f_{\xi\mid\eta}(x\mid y)=
\begin{cases}\dfrac{\sin(x+y)}{\sin y+\cos y},&0<x<\pi/2,\\0,&\text{其他}.\end{cases}
\]
分子在此区间非负，其对 $x$ 的积分正好等于分母，所以该条件密度已归一化。
\end{proof}
\paragraph{答案}$f_{\xi\mid\eta}(x\mid y)=\sin(x+y)/(\sin y+\cos y)$，$0<x<\pi/2$；区间外为 $0$。

\subsection{教材第80页第2题}\label{ex:w11-b80-2}
设 $\eta$ 服从几何分布 $G(\lambda)$。令 $\xi$ 是成功概率为 $p$（$0<p<1$）的 $\eta$ 重伯努利试验的成功次数。求离散型随机向量 $(\xi,\eta)$ 的联合分布列。

作业对题意的说明：当 $\eta=n$ 时，$\xi$ 的条件分布为 $B(n,p)$。几何分布采用取值 $1,2,\ldots$ 的约定，$0<\lambda<1$。
\basis{原笔记第2章几何分布、二项分布，第3章第3节条件分布列的乘法公式。}
\paragraph{思路}把条件二项分布与几何分布相乘。
\begin{proof}[解答]
对整数 $n\geq1$ 和 $k=0,\ldots,n$，
\[
P(\eta=n)=\lambda(1-\lambda)^{n-1},\qquad
P(\xi=k\mid\eta=n)=\binom nkp^k(1-p)^{n-k}.
\]
由乘法公式，
\[
P(\xi=k,\eta=n)=\binom nkp^k(1-p)^{n-k}\lambda(1-\lambda)^{n-1}.
\]
其他取值对的联合概率为 $0$。对固定 $n$ 先对 $k$ 求和得到 $P(\eta=n)$，再对 $n$ 求和得到 $1$，故该联合分布归一化。
\end{proof}
\paragraph{答案}联合分布列如上，支持为 $n\geq1,0\leq k\leq n$ 的整数对。

\subsection{教材第81页第3题}\label{ex:w11-b81-3}
设 $\eta$ 服从指数分布 $\operatorname{Exp}(\lambda)$，而在给定 $\eta$ 的条件下 $\xi$ 服从正态分布 $N(0,\eta^2)$。求绝对连续型随机向量 $(\xi,\eta)$ 的联合分布密度。
\basis{原笔记第3章第3节“由联合密度求条件密度”，第2章指数分布、正态分布密度。}
\paragraph{思路}联合密度等于条件密度乘以条件变量的边缘密度。
\begin{proof}[解答]
对 $y>0$，给定 $\eta=y$ 时，$\xi$ 的标准差为 $y$，因此
\[
f_{\xi\mid\eta}(x\mid y)=\frac1{\sqrt{2\pi}y}
\exp\left(-\frac{x^2}{2y^2}\right),\qquad x\in\R.
\]
乘以 $f_\eta(y)=\lambda e^{-\lambda y}$，可得
\[
f_{\xi,\eta}(x,y)=
\begin{cases}
\dfrac{\lambda}{\sqrt{2\pi}y}
\exp\left(-\lambda y-\dfrac{x^2}{2y^2}\right),&x\in\R,\ y>0,\\
0,&y\leq0.
\end{cases}
\]
对 $x$ 积分先得到 $f_\eta(y)$，再对 $y$ 积分得到 $1$，所以表达式确为联合密度。
\end{proof}
\paragraph{答案}联合密度如上。

\subsection{教材第81页第5题}\label{ex:w11-b81-5}
设随机向量 $(\xi,\eta)$ 的联合分布密度为
\[
p(x,y)=8xyI_{\{0<x<y<1\}}.
\]
求 $\xi$ 和 $\eta$ 的边缘分布密度，并说明 $\xi$ 和 $\eta$ 是否相互独立。
\basis{原笔记第3章第1节边缘密度公式，第2节例“乘积区域与三角形区域中的独立性”。}
\paragraph{思路}先求两个边缘密度，再利用 $\xi<\eta$ 必然成立这一限制排除独立。
\begin{proof}[解答]
对 $0<x<1$ 与 $0<y<1$，分别有
\begin{align*}
f_\xi(x)&=\int_x^18xy\,dy=4x(1-x^2),\\
f_\eta(y)&=\int_0^y8xy\,dx=4y^3.
\end{align*}
区间外密度为 $0$。

事件 $A=\{\xi>3/4\}$ 与 $B=\{\eta<1/4\}$ 的概率都严格大于 $0$：它们分别是上述正密度在非空区间上的积分。但 $AB$ 要求同时有 $\xi>\eta$，与联合密度的支持 $\xi<\eta$ 矛盾，故 $P(AB)=0$。于是 $P(AB)\ne P(A)P(B)$，二者不独立。
\end{proof}
\paragraph{答案}$f_\xi(x)=4x(1-x^2)I_{\{0<x<1\}}$，$f_\eta(y)=4y^3I_{\{0<y<1\}}$；不独立。

\subsection{作业第2题：中间次序与两个变量之和}\label{ex:w11-h2}
设 $X_1,X_2,X_3$ 独立同分布，共同密度函数为 $f(x)=2x$，$0<x<1$，其他处为 $0$。
（1）求 $X_2$ 的取值在 $X_1$ 和 $X_3$ 之间的概率；（2）求 $X_1+X_2$ 的密度函数。
\basis{原笔记第3章第2节“独立随机变量求和的卷积公式”，第4节“连续独立样本无相同取值”。}
\paragraph{思路}第一问用独立同分布的排列对称性；第二问按两个自变量均落在 $(0,1)$ 来确定卷积上下限。
\begin{proof}[解答]
\begin{enumerate}
\item 连续独立样本相同取值的概率为 $0$，所以三个值以概率 $1$ 恰有一个中间值。三个变量独立同分布，交换下标不改变联合分布，因而 $X_1,X_2,X_3$ 为中间值的概率相同。这三个互斥事件的概率和为 $1$，故所求概率为 $1/3$。
\item 令 $S=X_1+X_2$。卷积中须同时有 $0<x<1$ 与 $0<s-x<1$，即
\[
\max(0,s-1)<x<\min(1,s).
\]
当 $0<s\leq1$ 时，
\[
f_S(s)=\int_0^s4x(s-x)dx
=4\left[\frac{s x^2}{2}-\frac{x^3}{3}\right]_0^s
=\frac23s^3.
\]
当 $1<s<2$ 时，
\begin{align*}
f_S(s)&=\int_{s-1}^14x(s-x)dx\\
&=4\left[\frac{s x^2}{2}-\frac{x^3}{3}\right]_{s-1}^1
=-\frac23s^3+4s-\frac83.
\end{align*}
当 $s\notin(0,2)$ 时，积分区域为空，密度取 $0$。
\end{enumerate}
\end{proof}
\paragraph{答案}（1）$1/3$；（2）
\[
f_S(s)=\begin{cases}\frac23s^3,&0<s\leq1,\\-\frac23s^3+4s-\frac83,&1<s<2,\\0,&\text{其他}.\end{cases}
\]

\subsection{作业第3题：首次非零结果与等待次数}\label{ex:w11-h3}
考虑独立重复试验，每次可能结果为 $0,1,\ldots,k$，第 $i$ 个结果出现的概率为 $p_i$，其中 $0<p_i<1$，$\sum_{i=0}^kp_i=1$。
令 $N$ 为首次出现非 $0$ 结果时所需的试验次数，$X$ 为这次出现的具体结果，$X\in\{1,\ldots,k\}$。
（1）求 $N$ 的分布列；（2）求 $X$ 的分布列；（3）证明 $X$ 与 $N$ 独立。
\basis{原笔记第1章独立重复试验，第3章第1节边缘分布列，第2节独立性判据。}
\paragraph{思路}事件 $\{N=n,X=j\}$ 表示前 $n-1$ 次均为 $0$，第 $n$ 次为 $j$。
\begin{proof}[解答]
\begin{enumerate}
\item 对 $n\geq1$，由各次试验独立，
\[
P(N=n)=p_0^{n-1}(1-p_0).
\]
一直为 $0$ 的概率是 $\lim_{n\to\infty}p_0^n=0$，故 $N$ 以概率 $1$ 为有限正整数；$N\sim G(1-p_0)$。
\item 对 $n\geq1$、$j=1,\ldots,k$，
\[
P(N=n,X=j)=p_0^{n-1}p_j.
\]
事件 $\{N=n,X=j\}$ 对不同 $n$ 互斥，故
\[
P(X=j)=\sum_{n=1}^\infty p_0^{n-1}p_j=\frac{p_j}{1-p_0}.
\]
\item 对所有上述 $n,j$，
\[
P(N=n)P(X=j)=p_0^{n-1}(1-p_0)\frac{p_j}{1-p_0}
=P(N=n,X=j).
\]
联合概率处处分解为边缘概率乘积，故 $N,X$ 独立。
\end{enumerate}
\end{proof}
\paragraph{答案}$N\sim G(1-p_0)$，$P(X=j)=p_j/(1-p_0)$，且 $N,X$ 独立。

\subsection{作业第4题：保龄球得分}\label{ex:w11-h4}
甲的保龄球得分近似服从 $N(170,20^2)$，乙的得分近似服从 $N(160,15^2)$。甲、乙每人玩一局，得分相互独立。求以下事件的近似概率，保留4位小数：（1）甲得分高一些；（2）两人总得分超过 $350$。提示：正态分布的线性性。
\basis{原笔记第3章第2节独立正态随机变量的线性组合，第2章标准正态分布函数 $\Phi$。}
\paragraph{思路}分别把两人得分之差、之和化成一个正态变量，再标准化。
\begin{proof}[解答]
记甲、乙得分为 $A,B$。由独立正态变量线性组合的结论，
\[
A-B\sim N(10,20^2+15^2)=N(10,25^2),\qquad
A+B\sim N(330,25^2).
\]
于是
\begin{align*}
P(A>B)&=P\left(\frac{A-B-10}{25}>-0.4\right)
=\Phi(0.4)\approx0.6554,\\
P(A+B>350)&=P\left(\frac{A+B-330}{25}>0.8\right)
=1-\Phi(0.8)\approx0.2119.
\end{align*}
\end{proof}
\paragraph{答案}（1）$0.6554$；（2）$0.2119$。

\subsection{作业第5题：正态分布的等尾概率}\label{ex:w11-h5}
设 $X,Y$ 独立同分布于 $N(10,4)$。求数 $x$，使得 $P(X+Y>x)=P(X>15)$。
\basis{原笔记第3章第2节独立正态变量的线性组合，第2章正态分布的标准化。}
\paragraph{思路}两个标准正态尾概率相等时，相应标准化端点相等。
\begin{proof}[解答]
独立性给出 $X+Y\sim N(20,8)$，故所给等式等价于
\[
1-\Phi\left(\frac{x-20}{2\sqrt2}\right)
=1-\Phi\left(\frac{15-10}{2}\right).
\]
标准正态密度处处严格为正，$\Phi$ 严格递增，所以
\[
\frac{x-20}{2\sqrt2}=\frac52,
\qquad x=20+5\sqrt2.
\]
\end{proof}
\paragraph{答案}$x=20+5\sqrt2$。

\subsection{作业第6题：逐次选数的条件分布}\label{ex:w11-h6}
从数集 $\{1,2,3,4,5\}$ 中随机选一个数 $X$，再从比 $X$ 小的数集（即集合 $\{1,2,\ldots,X\}$）中随机取第二个数，记为 $Y$。
（a）求 $X,Y$ 的联合分布列；（b）在 $Y=i$ 的条件下求 $X$ 的条件分布列，$i=1,\ldots,5$；（c）$X,Y$ 是否独立？为什么？

题目文字“比 $X$ 小”与括号中的集合不一致。以下按题目明确给出的集合 $\{1,2,\ldots,X\}$，即允许 $Y=X$ 解答；否则 $X=1$ 时第二次选取无法进行。
\basis{原笔记第3章第1节边缘分布列，第3节“离散型条件分布列与条件分布函数”。}
\paragraph{思路}先写 $X$ 的均匀分布，再写给定 $X$ 后 $Y$ 的均匀分布，最后由联合分布反求条件分布。
\begin{proof}[解答]
\begin{enumerate}[label=(\alph*)]
\item 对整数 $1\leq j\leq k\leq5$，
\[
P(X=k,Y=j)=P(X=k)P(Y=j\mid X=k)=\frac1{5k}.
\]
其他取值对概率为 $0$，所以联合分布表为
\[
\begin{array}{c|ccccc}
X\backslash Y&1&2&3&4&5\\\hline
1&1/5&0&0&0&0\\
2&1/10&1/10&0&0&0\\
3&1/15&1/15&1/15&0&0\\
4&1/20&1/20&1/20&1/20&0\\
5&1/25&1/25&1/25&1/25&1/25
\end{array}
\]
\item 对 $i=1,\ldots,5$，
\[
P(Y=i)=\frac15\sum_{r=i}^5\frac1r>0.
\]
因此
\[
P(X=k\mid Y=i)=
\begin{cases}
\dfrac{1/k}{\sum_{r=i}^5 1/r},& k=i,i+1,\ldots,5,\\
0,&\text{其他}.
\end{cases}
\]
若将各个条件分布按 $k=1,\ldots,5$ 排列，可写成
\[
\begin{array}{c|ccccc}
i&k=1&k=2&k=3&k=4&k=5\\\hline
1&60/137&30/137&20/137&15/137&12/137\\
2&0&30/77&20/77&15/77&12/77\\
3&0&0&20/47&15/47&12/47\\
4&0&0&0&5/9&4/9\\
5&0&0&0&0&1
\end{array}
\]
\item $P(X=1)=1/5$、$P(Y=5)=1/25$，但 $P(X=1,Y=5)=0$。联合概率不等于边缘概率的乘积，故不独立。
\end{enumerate}
\end{proof}
\paragraph{答案}联合分布为 $1/(5k)$（$1\leq j\leq k\leq5$）；条件分布如表；$X,Y$ 不独立。

\subsection{作业第7题：三角形区域上的条件密度}\label{ex:w11-h7}
设 $(X,Y)$ 的联合密度为 $f(x,y)=3y$，$0<x<y<1$，区域外为 $0$。
（1）已知 $X=x$（$0<x<1$）时，求 $f_{Y\mid X}(y\mid x)$；（2）求 $P(\frac13<Y<\frac34\mid X=\frac12)$。
\basis{原笔记第3章第3节“由联合密度求条件密度”。}
\paragraph{思路}条件值确定后，仍须把积分区间与支持条件 $y>x$ 取交。
\begin{proof}[解答]
\begin{enumerate}
\item 对 $0<x<1$，
\[
f_X(x)=\int_x^1 3y\,dy=\frac32(1-x^2)>0.
\]
所以
\[
f_{Y\mid X}(y\mid x)=\begin{cases}\dfrac{2y}{1-x^2},&x<y<1,\\0,&\text{其他}.\end{cases}
\]
\item 条件 $X=1/2$ 下 $Y>1/2$，故
\begin{align*}
P\left(\frac13<Y<\frac34\,\middle|\,X=\frac12\right)
&=\int_{1/2}^{3/4}\frac{2y}{1-1/4}\,dy\\
&=\frac43\left[ y^2\right]_{1/2}^{3/4}
=\frac43\left(\frac9{16}-\frac14\right)=\frac5{12}.
\end{align*}
\end{enumerate}
\end{proof}
\paragraph{答案}条件密度如上；条件概率为 $5/12$。

\subsection{作业第8题：对称三角形区域上的条件密度}\label{ex:w11-h8}
设二维连续型随机变量 $(X,Y)$ 的联合密度为
\[
f(x,y)=\frac32x,\qquad |y|<x,\quad0<x<1,
\]
区域外为 $0$。（1）求 $f_{X\mid Y}(x\mid y)$ 和 $f_{Y\mid X}(y\mid x)$；（2）求
\[
P\left(\frac13<Y<\frac34\,\middle|\,X=\frac12\right),\qquad
P\left(X>\frac23\,\middle|\,Y=-\frac12\right).
\]
\basis{原笔记第3章第1节边缘密度公式，第3节“由联合密度求条件密度”。}
\paragraph{思路}固定 $x$ 时 $y\in(-x,x)$；固定 $y$ 时 $x\in(|y|,1)$。
\begin{proof}[解答]
\begin{enumerate}
\item 两个边缘密度分别为
\[
f_X(x)=\int_{-x}^x\frac32x\,dy=3x^2,\quad0<x<1,
\]
\[
f_Y(y)=\int_{|y|}^1\frac32x\,dx=\frac34(1-y^2),\quad |y|<1,
\]
区间外为 $0$。因此，对 $|y|<1$，
\[
f_{X\mid Y}(x\mid y)=\begin{cases}\dfrac{2x}{1-y^2},&|y|<x<1,\\0,&\text{其他},\end{cases}
\]
而对 $0<x<1$，
\[
f_{Y\mid X}(y\mid x)=\begin{cases}\dfrac1{2x},&-x<y<x,\\0,&\text{其他}.\end{cases}
\]
\item 当 $X=1/2$ 时，条件支持为 $(-1/2,1/2)$，故
\[
P\left(\frac13<Y<\frac34\,\middle|\,X=\frac12\right)
=\int_{1/3}^{1/2}1\,dy=\frac16.
\]
当 $Y=-1/2$ 时，$x>1/2$，于是
\[
P\left(X>\frac23\,\middle|\,Y=-\frac12\right)
=\int_{2/3}^1\frac{2x}{1-1/4}dx
=\frac43\left(1-\frac49\right)=\frac{20}{27}.
\]
\end{enumerate}
\end{proof}
\paragraph{答案}条件密度如上；两个条件概率为 $1/6$、$20/27$。
% !TEX root = 概率论期中复习.tex
\section{第十二周：随机变量函数与次序统计量}

\subsection{教材第90页第1题}\label{ex:w12-b90-1}
设随机变量 $\xi$ 和 $\eta$ 独立且都服从均匀分布 $U(0,1)$。求 $\xi+\eta$ 的分布密度。
\basis{原笔记第3章第2节“独立随机变量求和的卷积公式”，例“两个独立均匀变量之和”。}
\paragraph{思路}卷积积分等于区间 $(0,1)$ 与 $(z-1,z)$ 的交集长度。
\begin{proof}[解答]
令 $S=\xi+\eta$。两变量密度均为 $I_{\{0<x<1\}}$，所以
\[
f_S(s)=\int_{\mathbb R}I_{\{0<x<1\}}I_{\{0<s-x<1\}}dx.
\]
当 $0<s\leq1$ 时，积分区间为 $(0,s)$，积分为 $s$；当 $1<s<2$ 时，积分区间为 $(s-1,1)$，积分为 $2-s$。其他 $s$ 对应的积分区间为空。因此
\[
f_S(s)=\begin{cases}s,&0<s\leq1,\\2-s,&1<s<2,\\0,&\text{其他}.\end{cases}
\]
\end{proof}
\paragraph{答案}$\xi+\eta$ 服从上述三角形密度。

\subsection{教材第90页第3题}\label{ex:w12-b90-3}
设随机变量 $\xi_1,\xi_2$ 独立且分别服从泊松分布 $P(\lambda_1)$ 和 $P(\lambda_2)$。求：（1）$\xi_1+\xi_2$ 的分布列；（2）在给定 $\xi_1+\xi_2=n$ 时 $\xi_1$ 的条件分布列。
\basis{原笔记第3章第2节泊松分布的可加性，第3节例“给定独立泊松变量的和”。}
\paragraph{思路}和的分布用离散卷积；条件分布用“联合概率除以条件事件概率”。
\begin{proof}[解答]
记 $S=\xi_1+\xi_2$，$\lambda_1,\lambda_2>0$。
\begin{enumerate}
\item 对非负整数 $n$，
\begin{align*}
P(S=n)
&=\sum_{k=0}^nP(\xi_1=k)P(\xi_2=n-k)\\
&=e^{-(\lambda_1+\lambda_2)}\sum_{k=0}^n
\frac{\lambda_1^k\lambda_2^{n-k}}{k!(n-k)!}\\
&=\frac{e^{-(\lambda_1+\lambda_2)}}{n!}
\sum_{k=0}^n\binom nk\lambda_1^k\lambda_2^{n-k}\\
&=e^{-(\lambda_1+\lambda_2)}\frac{(\lambda_1+\lambda_2)^n}{n!}.
\end{align*}
所以 $S\sim P(\lambda_1+\lambda_2)$。
\item 对 $k=0,\ldots,n$，
\begin{align*}
P(\xi_1=k\mid S=n)
&=\frac{P(\xi_1=k,\xi_2=n-k)}{P(S=n)}\\
&=\frac{e^{-(\lambda_1+\lambda_2)}
\lambda_1^k\lambda_2^{n-k}/[k!(n-k)!]}
{e^{-(\lambda_1+\lambda_2)}(\lambda_1+\lambda_2)^n/n!}\\
&=\binom nk\left(\frac{\lambda_1}{\lambda_1+\lambda_2}\right)^k
\left(\frac{\lambda_2}{\lambda_1+\lambda_2}\right)^{n-k}.
\end{align*}
其他 $k$ 的条件概率为 $0$。当 $n=0$ 时，该条件下 $\xi_1=0$ 的概率为 $1$。
\end{enumerate}
\end{proof}
\paragraph{答案}$S\sim P(\lambda_1+\lambda_2)$，$\xi_1\mid(S=n)\sim B(n,\lambda_1/(\lambda_1+\lambda_2))$。

\subsection{教材第90页第4题}\label{ex:w12-b90-4}
设随机变量 $\xi_1,\xi_2$ 独立且均服从指数分布 $\operatorname{Exp}(\lambda)$。证明 $\xi_1+\xi_2$ 与 $\xi_1/\xi_2$ 独立。
\basis{原笔记第3章第5节“二维一一变换的密度公式”，例“独立 Gamma 变量的和与比例”，第2节密度分解判据。}
\paragraph{思路}将和与比值一起作为新坐标，计算联合密度并分解。
\begin{proof}
记 $U=\xi_1+\xi_2$、$V=\xi_1/\xi_2$。在 $x>0,y>0$ 上，变换
\[
u=x+y,\quad v=\frac xy
\]
是一一变换，其逆变换和逆雅可比绝对值为
\[
x=\frac{uv}{1+v},\qquad y=\frac{u}{1+v},\qquad
\left|\frac{\partial(x,y)}{\partial(u,v)}\right|=\frac{u}{(1+v)^2}.
\]
原区域恰好变为 $u>0,v>0$，所以
\[
f_{U,V}(u,v)=\lambda^2e^{-\lambda u}\frac{u}{(1+v)^2},\qquad u,v>0,
\]
区域外为 $0$。两因子各自满足
\[
\int_0^\infty\lambda^2ue^{-\lambda u}du=1,\qquad
\int_0^\infty\frac1{(1+v)^2}dv=1.
\]
因此边缘密度分别为
\[
f_U(u)=\lambda^2ue^{-\lambda u}I_{\{u>0\}},\qquad
f_V(v)=\frac{I_{\{v>0\}}}{(1+v)^2},
\]
且 $f_{U,V}=f_Uf_V$，从而 $U,V$ 独立。
\end{proof}
\paragraph{答案}$\xi_1+\xi_2$ 与 $\xi_1/\xi_2$ 独立。

\subsection{教材第90页第6题}\label{ex:w12-b90-6}
设随机变量 $\xi_1,\ldots,\xi_n$ 独立且分别服从指数分布
$\operatorname{Exp}(\lambda_1),\ldots,\operatorname{Exp}(\lambda_n)$。求
$\eta_n=\xi_1\wedge\cdots\wedge\xi_n$ 的分布。
\basis{原笔记第3章第2节独立性的定义，第4节最小值分布的求法。}
\paragraph{思路}最小值超过 $t$，等价于每个变量都超过 $t$。
\begin{proof}[解答]
令 $\Lambda=\sum_{i=1}^n\lambda_i$。当 $t\geq0$ 时，
\begin{align*}
P(\eta_n>t)
&=P(\xi_1>t,\ldots,\xi_n>t)\\
&=\prod_{i=1}^nP(\xi_i>t)
=\prod_{i=1}^ne^{-\lambda_it}=e^{-\Lambda t}.
\end{align*}
每个 $\xi_i$ 以概率 $1$ 为正，所以 $t<0$ 时 $F_{\eta_n}(t)=0$。因此
\[
F_{\eta_n}(t)=\begin{cases}0,&t<0,\\1-e^{-\Lambda t},&t\geq0,\end{cases}
\qquad
f_{\eta_n}(t)=\Lambda e^{-\Lambda t}I_{\{t>0\}}.
\]
\end{proof}
\paragraph{答案}$\eta_n\sim\operatorname{Exp}(\lambda_1+\cdots+\lambda_n)$。

\subsection{作业第2题：Gamma--泊松条件分布}\label{ex:w12-h2}
设某天内空气中的水分含量 $W\sim\Gamma(t,\beta)$，其密度为
\[
f_W(w)=\frac{(\beta w)^{t-1}}{\Gamma(t)}\beta e^{-\beta w},\qquad w>0,
\]
其中 $t,\beta>0$。已知 $W=w>0$ 的条件下，当天内交通事故数量 $N$ 服从参数为 $w$ 的泊松分布。证明：已知 $N=n$ 的条件下，$W$ 的条件分布是 $\Gamma(t+n,\beta+1)$。Gamma 分布第二参数采用率参数。
\basis{原笔记第3章第3节“离散与连续变量的条件分布乘法公式”，第2章 Gamma 密度与 Gamma 积分。}
\paragraph{思路}用混合型贝叶斯公式，先将条件泊松概率与先验密度相乘，再归一化。
\begin{proof}
对非负整数 $n$ 和 $w>0$，
\[
P(N=n\mid W=w)=e^{-w}\frac{w^n}{n!}.
\]
于是离散变量与连续变量的联合概率密度表达式为
\begin{align*}
P(N=n\mid W=w)f_W(w)
&=\frac{\beta^t}{n!\Gamma(t)}w^{t+n-1}e^{-(\beta+1)w}.
\end{align*}
利用 $a,b>0$ 时的 Gamma 积分
\[
\int_0^\infty w^{a-1}e^{-bw}dw=\frac{\Gamma(a)}{b^a},
\]
得到
\[
P(N=n)=\frac{\beta^t\Gamma(t+n)}{n!\Gamma(t)(\beta+1)^{t+n}}>0.
\]
所以相应条件密度是
\begin{align*}
f_{W\mid N}(w\mid n)
&=\frac{P(N=n\mid W=w)f_W(w)}{P(N=n)}\\
&=\frac{(\beta+1)^{t+n}}{\Gamma(t+n)}
w^{t+n-1}e^{-(\beta+1)w},\qquad w>0,
\end{align*}
在 $w\leq0$ 处取 $0$。这正是形状参数 $t+n$、率参数 $\beta+1$ 的 Gamma 密度。
\end{proof}
\paragraph{答案}$W\mid(N=n)\sim\Gamma(t+n,\beta+1)$。

\subsection{作业第3题：差、积、商的密度}\label{ex:w12-h3}
设 $(X,Y)$ 的联合密度为 $f(x,y)=3y$，$0<x<y<1$，区域外为 $0$。
（1）求 $Y-X$ 的密度；（2）求 $XY$ 的密度；（3）求 $X/Y$ 的密度。
\basis{原笔记第3章第5节“二维一一变换的密度公式”及例“引入辅助变量求和的密度”。}
\paragraph{思路}每次保留 $Y$ 为辅助坐标，写逆变换、雅可比和新区域，再对辅助坐标积分。
\begin{proof}[解答]
\begin{enumerate}
\item 令 $D=Y-X$，保留 $Y$，则逆变换为 $x=y-d$，逆雅可比绝对值为 $1$。原区域等价于 $0<d<y<1$。因此
\[
f_{D,Y}(d,y)=3y,\qquad0<d<y<1,
\]
区域外为 $0$。对 $0<d<1$，
\[
f_D(d)=\int_d^13y\,dy=\frac32(1-d^2),
\]
其余处为 $0$。
\item 令 $Z=XY$，逆变换为 $x=z/y$，逆雅可比绝对值为 $1/y$。条件 $0<x<y<1$ 等价于 $0<z<1,\sqrt z<y<1$，故
\[
f_{Z,Y}(z,y)=3y\frac1y=3,\qquad0<z<1,\ \sqrt z<y<1.
\]
于是
\[
f_Z(z)=\int_{\sqrt z}^13\,dy=3(1-\sqrt z),\qquad0<z<1,
\]
其余处为 $0$。
\item 令 $Q=X/Y$，逆变换为 $x=qy$，逆雅可比绝对值为 $y$。原区域等价于 $0<q<1,0<y<1$。所以
\[
f_{Q,Y}(q,y)=3y^2,\qquad0<q<1,\ 0<y<1,
\]
\[
f_Q(q)=\int_0^13y^2dy=1,\qquad0<q<1,
\]
其余处为 $0$。因此 $Q\sim U(0,1)$。
\end{enumerate}
\end{proof}
\paragraph{答案}
\[
f_{Y-X}(z)=\frac32(1-z^2)I_{\{0<z<1\}},\quad
f_{XY}(z)=3(1-\sqrt z)I_{\{0<z<1\}},\quad
f_{X/Y}(z)=I_{\{0<z<1\}}.
\]

\subsection{作业第4题：指数样本的极差}\label{ex:w12-h4}
设 $X_1,\ldots,X_n$ 独立同分布于参数为 $\lambda$ 的指数分布。求极差 $R=X_{(n)}-X_{(1)}$ 的密度函数。
\basis{原笔记第3章第4节“两个次序统计量的联合密度”及例“样本极差的分布”。}
\paragraph{思路}将原笔记的极差分布函数公式代入指数分布，再对 $r$ 求导。
\begin{proof}[解答]
先设 $n\geq2$。共同分布函数与密度为
\[
F(x)=1-e^{-\lambda x},\qquad f(x)=\lambda e^{-\lambda x},\qquad x>0.
\]
原笔记的极差分布公式给出，对 $r\geq0$，
\begin{align*}
F_R(r)
&=n\int_0^\infty[F(x+r)-F(x)]^{n-1}f(x)dx\\
&=n\int_0^\infty
\bigl[e^{-\lambda x}(1-e^{-\lambda r})\bigr]^{n-1}
\lambda e^{-\lambda x}dx\\
&=n\lambda(1-e^{-\lambda r})^{n-1}\int_0^\infty e^{-n\lambda x}dx\\
&=(1-e^{-\lambda r})^{n-1}.
\end{align*}
$r<0$ 时分布函数为 $0$。对 $r>0$ 求导，得到
\[
f_R(r)=(n-1)\lambda e^{-\lambda r}(1-e^{-\lambda r})^{n-2},
\]
其余处可取 $0$。分布函数在 $0$ 处无跳跃，且在无穷处趋于 $1$，所以该密度没有漏掉点质量。

若 $n=1$，则 $R=X_1-X_1=0$ 几乎必然，是集中于 $0$ 的退化分布，不存在通常意义下的概率密度。
\end{proof}
\paragraph{答案}$n\geq2$ 时，$f_R(r)=(n-1)\lambda e^{-\lambda r}(1-e^{-\lambda r})^{n-2}I_{\{r>0\}}$；$n=1$ 时 $R=0$ 几乎必然。

\subsection{作业第5题：五台发动机组成的机器寿命}\label{ex:w12-h5}
某机器里有五台发动机，寿命独立同分布，密度为 $f(t)=te^{-t}$，$t>0$，其他处为 $0$。只要五台发动机中至少有三台正常工作，机器就能正常工作。求机器正常工作时间的密度。提示：机器寿命可以表示为某个次序统计量。
\basis{原笔记第3章第4节“次序统计量”及“单个次序统计量的密度”。}
\paragraph{思路}机器在第三台发动机失效时停止工作，寿命为五个发动机寿命的第三次序统计量。
\begin{proof}[解答]
记五台发动机寿命为 $T_1,\ldots,T_5$，按递增顺序排列为 $T_{(1)}<\cdots<T_{(5)}$；等值情况概率为 $0$。
前两台失效后仍有三台工作；第三台失效后只剩两台工作。因此机器寿命为 $T=T_{(3)}$。

单台发动机的分布函数为
\[
F(t)=\int_0^t se^{-s}ds=1-(1+t)e^{-t},\qquad t>0.
\]
单个次序统计量密度公式给出
\begin{align*}
f_T(t)
&=\frac{5!}{(3-1)!(5-3)!}[F(t)]^{3-1}[1-F(t)]^{5-3}f(t)\\
&=30[1-(1+t)e^{-t}]^2[(1+t)e^{-t}]^2te^{-t}\\
&=30t(1+t)^2e^{-3t}[1-(1+t)e^{-t}]^2,\qquad t>0.
\end{align*}
$t\leq0$ 时密度取 $0$。
\end{proof}
\paragraph{答案}机器寿命为 $T_{(3)}$，密度为
\[
f_T(t)=30t(1+t)^2e^{-3t}[1-(1+t)e^{-t}]^2I_{\{t>0\}}.
\]

\subsection{作业第6题：平方和与比值的联合分布}\label{ex:w12-h6}
设 $X,Y$ 独立同分布于标准正态分布 $N(0,1)$。令 $U=X^2+Y^2$、$V=X/Y$。求 $(U,V)$ 的联合密度和边缘密度，并判断 $U,V$ 是否相互独立。
\basis{原笔记第3章第5节“二维一一变换的密度公式”及例“标准正态变量的极坐标变换”，第2节独立性判据。}
\paragraph{思路}变换在整个平面上不是一一的，按 $Y>0$ 与 $Y<0$ 分成两个分支，分别换元后相加。
\begin{proof}[解答]
$P(Y=0)=0$，所以 $V$ 在 $Y=0$ 上任意定义不改变分布。对 $u>0,v\in\R$，逆变换有两个分支：
\[
y_\pm=\pm\sqrt{\frac{u}{1+v^2}},\qquad x_\pm=vy_\pm.
\]
两个分支分别对应 $y>0$ 与 $y<0$，各自是一一变换。正向雅可比为
\[
\frac{\partial(u,v)}{\partial(x,y)}
=\begin{vmatrix}2x&2y\\1/y&-x/y^2\end{vmatrix}
=-2\frac{x^2+y^2}{y^2}=-2(1+v^2),
\]
所以每一分支的逆雅可比绝对值为 $1/[2(1+v^2)]$。

原联合密度为 $(2\pi)^{-1}e^{-(x^2+y^2)/2}$。两个分支对变换后密度的贡献相同，故
\begin{align*}
f_{U,V}(u,v)
&=2\cdot\frac1{2\pi}e^{-u/2}\frac1{2(1+v^2)}\\
&=\frac{e^{-u/2}}{2\pi(1+v^2)},\qquad u>0,\ v\in\R.
\end{align*}
$u\leq0$ 时密度取 $0$。分别积分得到
\[
f_U(u)=\frac12e^{-u/2}I_{\{u>0\}},\qquad
f_V(v)=\frac1{\pi(1+v^2)},\quad v\in\R,
\]
其中使用了 $\int_{\mathbb R}(1+v^2)^{-1}dv=\pi$ 与 $\int_0^\infty e^{-u/2}du=2$。
联合密度等于上述两个边缘密度的乘积，故 $U,V$ 相互独立。
\end{proof}
\paragraph{答案}$U\sim\operatorname{Exp}(1/2)=\chi^2(2)$，$V$ 服从标准柯西分布；联合密度如上，二者独立。

\subsection{作业第7题：不同指数变量的最先发生概率}\label{ex:w12-h7}
设 $X_1,\ldots,X_n$ 相互独立，且 $X_i\sim\operatorname{Exp}(\lambda_i)$，$1\leq i\leq n$。证明
\[
P(X_i=\min\{X_1,\ldots,X_n\})
=\frac{\lambda_i}{\lambda_1+\cdots+\lambda_n},\qquad1\leq i\leq n.
\]
提示：可以考虑 $n$ 重积分，也可考虑其他方法。
\basis{原笔记第3章第2节独立随机变量的联合密度，第4节最小值分布的求法。}
\paragraph{思路}固定第 $i$ 个变量取值 $t$，其余变量均须大于 $t$；按这个区域逐次积分。
\begin{proof}
先设 $n\geq2$。独立连续变量取相同值的概率为 $0$，所以 $X_i$ 为最小值与 $X_j>X_i$（所有 $j\ne i$）仅相差零概率事件。记 $\Lambda=\sum_{j=1}^n\lambda_j$。由联合密度的乘积形式，
\begin{align*}
P(X_i=\min_jX_j)
&=\int_0^\infty\lambda_i e^{-\lambda_i t}
\prod_{j\ne i}\left(\int_t^\infty\lambda_je^{-\lambda_jx_j}dx_j\right)dt\\
&=\int_0^\infty\lambda_i e^{-\lambda_i t}
\prod_{j\ne i}e^{-\lambda_jt}dt\\
&=\lambda_i\int_0^\infty e^{-\Lambda t}dt
=\frac{\lambda_i}{\Lambda}.
\end{align*}
若 $n=1$，所求概率为 $1$，也等于 $\lambda_1/\lambda_1$。
\end{proof}
\paragraph{答案}第 $i$ 个变量最先发生的概率为 $\lambda_i/(\lambda_1+\cdots+\lambda_n)$。

\clearpage\part{检索与答案速查}
\section{作业答案速查}
\subsection*{第1周}
\begin{enumerate}
\item \hyperref[ex:w01-t1]{教材第9--10页第1题}：依次为$A\cup B\cup C$，$AB\cup AC\cup BC$，$A^cB^cC^c$，$(AB\cup AC\cup BC)^c$，$(ABC)^c$，$ABC^c\cup AB^cC\cup A^cBC$。
\item \hyperref[ex:w01-t2]{教材第10页第2题}：$13/30$。
\item \hyperref[ex:w01-t4]{教材第10页第4题}：有放回$169/512$；无放回$3/8$。
\item \hyperref[ex:w01-t5]{教材第10页第5题}：$p_1=[m(m-1)+n(n-1)]/[M(M-1)]$，$p_2=(m^2+n^2)/M^2$，且$p_2-p_1=2mn/[M^2(M-1)]>0$。
\item \hyperref[ex:w01-t6]{教材第10页第6题}：$1/2$。
\item \hyperref[ex:w01-t7]{教材第10页第7题}：按原题“小于”理解，$0<a<10/9$，不存在正的最小边长。
\item \hyperref[ex:w01-h2]{作业第2题：可列德摩根律}：两式均成立。
\item \hyperref[ex:w01-h3]{作业第3题：恰有两个事件发生}：两边均表示恰有两个事件发生。
\end{enumerate}
\subsection*{第2周}
\begin{enumerate}
\item \hyperref[ex:w02-t1]{教材第19页第1题}：不存在可数无穷等可能且每个结果有正概率的模型。
\item \hyperref[ex:w02-t2]{教材第19页第2题(1)}：$\mathscr E$是$E$上的$\sigma$代数，补运算必须相对于$E$进行。
\item \hyperref[ex:w02-t4]{教材第19页第4题}：所给下界成立；结合非负性还可写成$P(AB)\geq\max\{0,P(A)+P(B)-1\}$。
\item \hyperref[ex:w02-t5]{教材第19页第5题}：$P(A\cap B)\in[1/12,1/3]$。
\item \hyperref[ex:w02-h2]{作业第2题：偏离独立乘积的界}：$1/4$为可达到的界，例如$A=B$且$P(A)=1/2$时取等号。
\item \hyperref[ex:w02-h3]{作业第3题：开区间的可列表示}：该可列并表示成立。
\item \hyperref[ex:w02-h4]{作业第4题：无限重复试验的事件极限}：单减；极限为永远不成功；至少成功一次的概率为1。
\end{enumerate}
\subsection*{第3周}
\begin{enumerate}
\item \hyperref[ex:w03-t3]{教材第19页第3题}：两项均成立，且$\mu(\Omega)=\sum_{k=1}^n\mu_k(\Omega_k)$。
\item \hyperref[ex:w03-t6]{教材第19页第6题}：可列个满概率事件的交仍为满概率事件。
\item \hyperref[ex:w03-t241]{教材第24页第1题}：结论成立，且题设自动保证$P(A)>0$。
\item \hyperref[ex:w03-t242]{教材第24页第2题}：(1)一定成立；(2)不一定，以上反例中$P(A)=3/10<P(B)=7/10$。
\item \hyperref[ex:w03-t254]{教材第25页第4题}：$528/5915\approx0.08926$。
\item \hyperref[ex:w03-h2]{作业第2题：一排夫妻均不相邻}：$\displaystyle\sum_{r=0}^n(-1)^r\binom nr2^r(2n-r)!/(2n)!$。
\item \hyperref[ex:w03-h3]{作业第3题：已知补事件条件概率}：$0.73$。
\item \hyperref[ex:w03-h4]{作业第4题：玻璃器皿连续三次未摔坏}：$0.028$。
\item \hyperref[ex:w03-h5]{作业第5题：两种天气的递推}：$\bigl[1+(2p-1)^{n-1}\bigr]/2$。
\end{enumerate}
\subsection*{第4周}
\begin{enumerate}
\item \hyperref[ex:w04-t25-3]{教材第25页第3题}：$\P(D)=(1+4p)/5$；$\P(K\mid D)=5p/(1+4p)$。
\item \hyperref[ex:w04-t25-5]{教材第25页第5题}：两种信号的条件误码率均为$0.05$时，答案为$38/39$；只给总误码率则条件不足。
\item \hyperref[ex:w04-t25-6]{教材第25页第6题}：$\displaystyle\frac{\sum_{i=1}^ki^{n+1}}{k\sum_{i=1}^ki^n}$。
\item \hyperref[ex:w04-t32-1]{教材第32页第1题}：两个结论均成立；事件与自身独立当且仅当其概率为0或1。
\item \hyperref[ex:w04-t32-2]{教材第32页第2题}：3孩时独立；4孩时不独立。
\item \hyperref[ex:w04-t33-6]{教材第33页第6题}：$p_n=\frac12+(c-\frac12)(2p-1)^{n-1}$，极限为$1/2$。
\item \hyperref[ex:w04-h3]{作业第3题：两个并事件的独立性}：$A_1\cup A_2$与$A_3\cup A_4$相互独立。
\item \hyperref[ex:w04-h4]{作业第4题：失败无穷多次}：$\P(A)=1$。
\end{enumerate}
\subsection*{第5周}
\begin{enumerate}
\item \hyperref[ex:w05-t33-4]{教材第33页第4题}：$2/5$。
\item \hyperref[ex:w05-t44-5]{教材第44页第5题}：$\P(a\leq\xi\leq b)\geq1-\alpha-\beta$。
\item \hyperref[ex:w05-h2]{作业第2题：区间投点的重复试验}：(1) $1-(2/3)^k$、$1-(1/2)^k$；(2) $1-(1/4)^k$；(3) 不独立，可改取$B'=[1/6,2/3]$；(4) $C,D$独立当且仅当$A,B$独立；(5) $1/3$。
\item \hyperref[ex:w05-h3]{作业第3题：示性函数的乘积与并}：$I_{AB}=I_AI_B$，$I_{A\cup B}=I_A+I_B-I_AI_B$，均为逐点恒等式。
\item \hyperref[ex:w05-h4]{作业第4题：随机变量与常数取最小值}：$Y=X\wedge a$是原概率空间上的随机变量。
\end{enumerate}
\subsection*{第6周}
\begin{enumerate}
\item \hyperref[ex:w06-h1]{作业第1题：取出号码的中间值}：$X$在2、3、4处的概率分别为$3/10,2/5,3/10$，分布函数如上。
\item \hyperref[ex:w06-h2]{作业第2题：随机选盒后取球}：(1) $1/6,11/18,2/9$；(2) 分布函数如上。
\item \hyperref[ex:w06-h3]{作业第3题：分段分布函数的跳跃与连续部分}：(1) $1/10,1/5,7/60,1/12$，四点总概率为$1/2$；(2) $1/10,2/5,0$。
\item \hyperref[ex:w06-h4]{作业第4题：由单点概率确定分布函数参数}：$A=1/2$，$B=1/10$；所求概率为$1/4$。
\item \hyperref[ex:w06-h5]{作业第5题（思考题，选做）：分布函数的一致收敛}：$\displaystyle\sup_{t\in\R}|F_n(t)-F(t)|\to0$。
\end{enumerate}
\subsection*{第7周}
\begin{enumerate}
\item \hyperref[ex:w07-b52-1]{教材第52页第1题}：已通过的路口数满足 $\P(X=k)=2^{-(k+1)}$，$k\geq0$；包含当前路口的计数服从参数 $1/2$ 的几何分布。
\item \hyperref[ex:w07-b52-2]{教材第52页第2题}：$0<p<1$ 时，$\P(\xi=k)=p^k(1-p)+(1-p)^kp$，$k\geq1$；$p=0,1$ 时游程无限长。
\item \hyperref[ex:w07-b53-3]{教材第53页第3题}：$\P(X=k)=(1-p)^{k-1}p$（$1\leq k\leq4$），$\P(X=5)=(1-p)^4$。
\item \hyperref[ex:w07-b53-4]{教材第53页第4题}：$\eta\sim P(\lambda p)$，即 $\P(\eta=k)=e^{-\lambda p}(\lambda p)^k/k!$。
\item \hyperref[ex:w07-b53-5]{教材第53页第5题}：奇数概率为 $(1-e^{-2\lambda})/2$，偶数概率为 $(1+e^{-2\lambda})/2$。
\item \hyperref[ex:w07-b58-3]{教材第58页第3题}：$a=1/2$，$\P(X<3/2)=7/8$。
\item \hyperref[ex:w07-h2]{作业第2题：三角形密度的区间概率}：$59/72$。
\item \hyperref[ex:w07-h3]{作业第3题：储油量}：至少 $61$ 千升。
\item \hyperref[ex:w07-h4]{作业第4题：由归一化求密度参数}：$C=1/4$，所求概率为 $1-\tfrac32e^{-1/2}$。
\end{enumerate}
\subsection*{第8周}
\begin{enumerate}
\item \hyperref[ex:w08-b57-2]{教材第57页第2题}：$pf+qg$ 是密度；其分布函数为 $pF_f+qF_g$，其中 $F_f,F_g$ 分别是 $f,g$ 对应的分布函数。
\item \hyperref[ex:w08-b58-4]{教材第58页第4题}：$F(t)=\exp(-e^{-t})$，$t\in\R$。
\item \hyperref[ex:w08-h2]{作业第2题：由三角形密度求分布函数}：\[
F(t)=\begin{cases}0,&t<0,\\t^2/2,&0\leq t\leq1,\\1-(2-t)^2/2,&1<t<2,\\1,&t\geq2.\end{cases}
\]
\item \hyperref[ex:w08-h3]{作业第3题：对称密度的分布函数}：$F(c-t)+F(c+t)=1$ 对全部 $t\in\R$ 成立；特别地 $F(c)=1/2$。
\item \hyperref[ex:w08-h4]{作业第4题：先到列车的方向}：两问均为 $2/3$。
\item \hyperref[ex:w08-h5]{作业第5题：修理时间}：第一问 $0.368$；第二问字面“恰好10小时”为 $0.000$，若指10小时内修好则为 $0.393$。
\item \hyperref[ex:w08-h6]{作业第6题：旧车的剩余里程}：指数分布下为 $0.368$；保留题面“10公里”并统一单位后，均匀分布下为 $0.500$；若原意为“10千里”，则为 $0.333$。
\item \hyperref[ex:w08-h7]{作业第7题：正态分布的中心}：$\mu=3$。
\end{enumerate}
\subsection*{第9周}
\begin{enumerate}
\item \hyperref[ex:w09-b58-4]{教材第58页第4题（再次布置）}：$F(t)=\exp(-e^{-t})$，$t\in\R$。
\item \hyperref[ex:w09-b58-5]{教材第58页第5题}：$X\sim N(a,\sigma^2)$，密度如上。
\item \hyperref[ex:w09-b62-1]{教材第62页第1题}：$1/3$。
\item \hyperref[ex:w09-h2]{作业第2题：正态分布的区间概率}：依次为 $0.7977$、$0.6827$、$0.8413$。
\item \hyperref[ex:w09-h3]{作业第3题：由分位点确定收入分布}：$\mu=250000$、$\sigma\approx103782.1553$；两问概率分别为 $0.3150$、$0.1363$。
\item \hyperref[ex:w09-h4]{作业第4题：正态分位点}：$c\approx14.56$。
\item \hyperref[ex:w09-h5]{作业第5题：线段分割的长短比}：$2/5$。
\item \hyperref[ex:w09-h6]{作业第6题：由均匀分布构造指数分布}：$Y\sim\operatorname{Exp}(\lambda)$，密度为 $\lambda e^{-\lambda y}$（$y>0$），其他处取 $0$。
\item \hyperref[ex:w09-h7]{作业第7题：指数变量的尺度变换}：$Y\sim\operatorname{Exp}(\lambda/c)$。
\item \hyperref[ex:w09-h8]{作业第8题：均匀变量的绝对值}：$Y\sim U(0,1/2)$。
\item \hyperref[ex:w09-h9]{作业第9题：指数变量的余弦}：\[
f_Y(y)=\begin{cases}
\dfrac{e^{-\arccos y}+e^{-2\pi+\arccos y}}
{(1-e^{-2\pi})\sqrt{1-y^2}},&-1<y<1,\\[6pt]
0,&\text{其他。}
\end{cases}
\]
\item \hyperref[ex:w09-h10]{作业第10题及思考题：概率积分变换}：$F(X)\sim U(0,1)$；即使 $F$ 有平坦区间，结论仍成立。这里保留了题设中 $X$ 连续型的条件。
\end{enumerate}
\subsection*{第10周}
\begin{enumerate}
\item \hyperref[ex:w10-b72-1]{教材第72页第1题}：联合分布为 $p(i,j)=1/(nj)$（$1\leq i\leq j\leq n$）；两个边缘分布如上。
\item \hyperref[ex:w10-b72-2]{教材第72页第2题}：$P(\xi=\eta)=0$。
\item \hyperref[ex:w10-b72-4]{教材第72页第4题}：$a=2$；$f_\xi(x)=2e^{-2x}I_{\{x>0\}}$；后两问为 $(1-e^{-4})(1-e^{-1})$、$(1-e^{-2})^2$。
\item \hyperref[ex:w10-b72-5]{教材第72页第5题}：$f_\xi(x)=3(1-x)^2I_{\{0<x<1\}}$，$f_\eta(y)=(3y-\frac32y^2)I_{\{0<y<1\}}$。
\item \hyperref[ex:w10-b72-6]{教材第72页第6题}：不是二维概率分布函数。
\item \hyperref[ex:w10-b81-4]{教材第81页第4题}：$f_\xi(x)=2xI_{\{0<x<1\}}$、$f_\eta(y)=2yI_{\{0<y<1\}}$，且二者独立。
\item \hyperref[ex:w10-h2]{作业第2题：由事件独立性确定联合分布}：$a=0.4,b=0.1$；$X$ 的分布为 $(0.8,0.2)$，$Y$ 的分布为 $(0.5,0.5)$；相互独立。
\item \hyperref[ex:w10-h3]{作业第3题：三项分布的边缘分布}：$X_1\sim B(n,p_1)$，$X_2\sim B(n,p_2)$，分布列如上。
\item \hyperref[ex:w10-h4]{作业第4题：双曲线之间的联合密度}：$A=3/2$；边缘密度如上；两概率为 $5/16$、$19/32$。
\end{enumerate}
\subsection*{第11周}
\begin{enumerate}
\item \hyperref[ex:w11-b80-1]{教材第80页第1题}：$f_{\xi\mid\eta}(x\mid y)=\sin(x+y)/(\sin y+\cos y)$，$0<x<\pi/2$；区间外为 $0$。
\item \hyperref[ex:w11-b80-2]{教材第80页第2题}：联合分布列如上，支持为 $n\geq1,0\leq k\leq n$ 的整数对。
\item \hyperref[ex:w11-b81-3]{教材第81页第3题}：联合密度如上。
\item \hyperref[ex:w11-b81-5]{教材第81页第5题}：$f_\xi(x)=4x(1-x^2)I_{\{0<x<1\}}$，$f_\eta(y)=4y^3I_{\{0<y<1\}}$；不独立。
\item \hyperref[ex:w11-h2]{作业第2题：中间次序与两个变量之和}：（1）$1/3$；（2）
\[
f_S(s)=\begin{cases}\frac23s^3,&0<s\leq1,\\-\frac23s^3+4s-\frac83,&1<s<2,\\0,&\text{其他}.\end{cases}
\]
\item \hyperref[ex:w11-h3]{作业第3题：首次非零结果与等待次数}：$N\sim G(1-p_0)$，$P(X=j)=p_j/(1-p_0)$，且 $N,X$ 独立。
\item \hyperref[ex:w11-h4]{作业第4题：保龄球得分}：（1）$0.6554$；（2）$0.2119$。
\item \hyperref[ex:w11-h5]{作业第5题：正态分布的等尾概率}：$x=20+5\sqrt2$。
\item \hyperref[ex:w11-h6]{作业第6题：逐次选数的条件分布}：联合分布为 $1/(5k)$（$1\leq j\leq k\leq5$）；条件分布如表；$X,Y$ 不独立。
\item \hyperref[ex:w11-h7]{作业第7题：三角形区域上的条件密度}：条件密度如上；条件概率为 $5/12$。
\item \hyperref[ex:w11-h8]{作业第8题：对称三角形区域上的条件密度}：条件密度如上；两个条件概率为 $1/6$、$20/27$。
\end{enumerate}
\subsection*{第12周}
\begin{enumerate}
\item \hyperref[ex:w12-b90-1]{教材第90页第1题}：$\xi+\eta$ 服从上述三角形密度。
\item \hyperref[ex:w12-b90-3]{教材第90页第3题}：$S\sim P(\lambda_1+\lambda_2)$，$\xi_1\mid(S=n)\sim B(n,\lambda_1/(\lambda_1+\lambda_2))$。
\item \hyperref[ex:w12-b90-4]{教材第90页第4题}：$\xi_1+\xi_2$ 与 $\xi_1/\xi_2$ 独立。
\item \hyperref[ex:w12-b90-6]{教材第90页第6题}：$\eta_n\sim\operatorname{Exp}(\lambda_1+\cdots+\lambda_n)$。
\item \hyperref[ex:w12-h2]{作业第2题：Gamma--泊松条件分布}：$W\mid(N=n)\sim\Gamma(t+n,\beta+1)$。
\item \hyperref[ex:w12-h3]{作业第3题：差、积、商的密度}：\[
f_{Y-X}(z)=\frac32(1-z^2)I_{\{0<z<1\}},\quad
f_{XY}(z)=3(1-\sqrt z)I_{\{0<z<1\}},\quad
f_{X/Y}(z)=I_{\{0<z<1\}}.
\]
\item \hyperref[ex:w12-h4]{作业第4题：指数样本的极差}：$n\geq2$ 时，$f_R(r)=(n-1)\lambda e^{-\lambda r}(1-e^{-\lambda r})^{n-2}I_{\{r>0\}}$；$n=1$ 时 $R=0$ 几乎必然。
\item \hyperref[ex:w12-h5]{作业第5题：五台发动机组成的机器寿命}：机器寿命为 $T_{(3)}$，密度为
\[
f_T(t)=30t(1+t)^2e^{-3t}[1-(1+t)e^{-t}]^2I_{\{t>0\}}.
\]
\item \hyperref[ex:w12-h6]{作业第6题：平方和与比值的联合分布}：$U\sim\operatorname{Exp}(1/2)=\chi^2(2)$，$V$ 服从标准柯西分布；联合密度如上，二者独立。
\item \hyperref[ex:w12-h7]{作业第7题：不同指数变量的最先发生概率}：第 $i$ 个变量最先发生的概率为 $\lambda_i/(\lambda_1+\cdots+\lambda_n)$。
\end{enumerate}

\end{document}
